% 由 scripts/gen-ai-usage-tex.py 根据 docs/exports/ai-usage-inventory.tsv 生成。
% 勿手改；更新清单后重新运行本脚本。

\subsubsection*{序号1}
\noindent 路径：\code{os/components/wateros-abi/abi-impl/impl-dummy/src/lib.rs}\\
\noindent 行范围：1-22\\[0.4em]

\begin{rustcode}
#![no_std]
//! 本模块代码由AI完成
//! ABI 占位实现 crate：用于工作区依赖解析与编译连通性，不包含真实平台语义。
//!
//! 对外符号仅为 Cargo 模板级占位，后续可由具体 `impl-*` 包替换或删除。
...
\end{rustcode}

\subsubsection*{序号2}
\noindent 路径：\code{os/components/wateros-abi/abi-impl/impl-dummy/src/lib.rs}\\
\noindent 行范围：10\\[0.4em]

\begin{rustcode}
...
// 本方法代码由AI完成
#[inline]
pub fn add(left : u64, right : u64) -> u64 { left + right }

#[cfg(test)]
...
\end{rustcode}

\subsubsection*{序号3}
\noindent 路径：\code{os/components/wateros-abi/abi-impl/impl-dummy/src/lib.rs}\\
\noindent 行范围：18-21\\[0.4em]

\begin{rustcode}
...
    #[test]
    fn it_works() {
        let result = add(2, 2);
        assert_eq!(result, 4);
    }
...
\end{rustcode}

\subsubsection*{序号4}
\noindent 路径：\code{os/components/wateros-abi/abi-impl/impl-linux-generic64/src/lib.rs}\\
\noindent 行范围：13\\[0.4em]

\begin{rustcode}
...
/// RISC-V 64 与 LoongArch64 当前都通过 `wateros-abi` 的架构 feature 选择这张表。
// 本结构代码由AI完成
pub struct LinuxGeneric64;

impl SyscallNumberTable for LinuxGeneric64 {
...
\end{rustcode}

\subsubsection*{序号5}
\noindent 路径：\code{os/components/wateros-abi/abi-impl/impl-linux-generic64/src/lib.rs}\\
\noindent 行范围：208-383\\[0.4em]

\begin{rustcode}
...
    fn dispatched_syscall_numbers_are_unique() {
        // 本变量代码由AI完成
        let nums = [
            LinuxGeneric64::READ,
            LinuxGeneric64::READV,
...
\end{rustcode}

\subsubsection*{序号6}
\noindent 路径：\code{os/components/wateros-abi/abi-impl/impl-linux-generic64/src/lib.rs}\\
\noindent 行范围：210-373\\[0.4em]

\begin{rustcode}
...
        let nums = [
            LinuxGeneric64::READ,
            LinuxGeneric64::READV,
            LinuxGeneric64::WRITE,
            LinuxGeneric64::WRITEV,
...
\end{rustcode}

\subsubsection*{序号7}
\noindent 路径：\code{os/components/wateros-driver/driver-block/block-impl/impl-block-cache/src/manager.rs}\\
\noindent 行范围：1-44\\[0.4em]

\begin{rustcode}
//! 本模块代码由AI完成
//! 块缓存注册辅助：从 [`wateros_base_config`] 读取默认容量并包装 [`CachingBlockDevice`]。

extern crate alloc;

...
\end{rustcode}

\subsubsection*{序号8}
\noindent 路径：\code{os/components/wateros-driver/driver-block/block-impl/impl-block-cache/src/manager.rs}\\
\noindent 行范围：16\\[0.4em]

\begin{rustcode}
...
/// 块设备写穿缓存的管理入口（v1：包装与默认配置；不向下转型已注册句柄）。
// 本结构代码由AI完成
pub struct BlockCacheManager;

impl BlockCacheManager {
...
\end{rustcode}

\subsubsection*{序号9}
\noindent 路径：\code{os/components/wateros-driver/driver-block/block-impl/impl-dummy/src/lib.rs}\\
\noindent 行范围：1-21\\[0.4em]

\begin{rustcode}
//! 本模块代码由AI完成
//! 块子系统占位实现：无硬件绑定，保留最小 crate 边界以便 feature 组合。
//!
//! 后续若需「无块设备」语义，应在此集中说明而非散落在调用方。

...
\end{rustcode}

\subsubsection*{序号10}
\noindent 路径：\code{os/components/wateros-driver/driver-character/character-impl/impl-dummy/src/lib.rs}\\
\noindent 行范围：1-18\\[0.4em]

\begin{rustcode}
//! 本模块代码由AI完成
//! 字符设备占位实现：无硬件、无 DTB 绑定。

#![no_std]

...
\end{rustcode}

\subsubsection*{序号11}
\noindent 路径：\code{os/components/wateros-driver/driver-character/character-impl/impl-null-stub/src/lib.rs}\\
\noindent 行范围：1-40\\[0.4em]

\begin{rustcode}
//! 本模块代码由AI完成
//! `/dev/null` 虚拟字符设备：读 EOF，写丢弃。

#![no_std]

...
\end{rustcode}

\subsubsection*{序号12}
\noindent 路径：\code{os/components/wateros-driver/driver-character/character-impl/impl-null-stub/src/lib.rs}\\
\noindent 行范围：19\\[0.4em]

\begin{rustcode}
...
// 本结构代码由AI完成
#[derive(Debug, Clone, Copy, Default)]
pub struct NullCharacterDevice;

impl CharacterDevice for NullCharacterDevice {
...
\end{rustcode}

\subsubsection*{序号13}
\noindent 路径：\code{os/components/wateros-driver/driver-character/character-impl/impl-rtc-stub/src/lib.rs}\\
\noindent 行范围：1-56\\[0.4em]

\begin{rustcode}
//! 本模块代码由AI完成
//! 软件 RTC 字符设备 stub：`hwclock` 等通过 `/dev/misc/rtc` + `ioctl` 访问；时间语义由 syscall 层处理。

#![no_std]

...
\end{rustcode}

\subsubsection*{序号14}
\noindent 路径：\code{os/components/wateros-driver/driver-character/character-impl/impl-rtc-stub/src/lib.rs}\\
\noindent 行范围：34\\[0.4em]

\begin{rustcode}
...
// 本结构代码由AI完成
#[derive(Debug, Clone, Copy, Default)]
pub struct RtcCharacterDevice;

impl CharacterDevice for RtcCharacterDevice {
...
\end{rustcode}

\subsubsection*{序号15}
\noindent 路径：\code{os/components/wateros-driver/driver-impl/impl-dummy/src/lib.rs}\\
\noindent 行范围：1-18\\[0.4em]

\begin{rustcode}
//! 本模块代码由AI完成
//! 平台驱动占位实现：不解析 DTB、不注册设备，用于无硬件目标的构建与依赖占位。

#![no_std]

...
\end{rustcode}

\subsubsection*{序号16}
\noindent 路径：\code{os/components/wateros-driver/driver-network/network-impl/impl-dummy/src/lib.rs}\\
\noindent 行范围：1-51\\[0.4em]

\begin{rustcode}
//! 本模块代码由AI完成
//! 网络设备占位实现：无硬件交互，用作缺省回退与编译占位。
//!
//! **当前行为**：无真实帧收发；**后续替换点**：virtio-net 等实现 crate。

...
\end{rustcode}

\subsubsection*{序号17}
\noindent 路径：\code{os/components/wateros-driver/driver-network/network-impl/impl-dummy/src/lib.rs}\\
\noindent 行范围：16-18\\[0.4em]

\begin{rustcode}
...
// 本结构代码由AI完成
pub struct DummyNetworkDevice {
    mac: [u8; 6],
}

...
\end{rustcode}

\subsubsection*{序号18}
\noindent 路径：\code{os/components/wateros-fs/fs-devfs/devfs-impl/impl-kernel/src/lib.rs}\\
\noindent 行范围：1-322\\[0.4em]

\begin{rustcode}
#![no_std]
//! 本模块代码由AI完成

//! 内核 devfs 实现：枚举块/字符设备、维护路径绑定，并可选合并 DTB 占位节点。
//!
...
\end{rustcode}

\subsubsection*{序号19}
\noindent 路径：\code{os/components/wateros-fs/fs-devfs/devfs-impl/impl-kernel/src/lib.rs}\\
\noindent 行范围：23-32\\[0.4em]

\begin{rustcode}
...
struct DevFsImpl {
    // 当前 refresh 后的节点快照（块/字符/占位）。
    nodes: Vec<DevNode>,
    // 路径 → 块设备句柄；refresh 会清空后按驱动枚举重建。
    block_bindings: Vec<(String, SharedBlockDevice)>,
...
\end{rustcode}

\subsubsection*{序号20}
\noindent 路径：\code{os/components/wateros-fs/fs-devfs/devfs-impl/impl-kernel/src/lib.rs}\\
\noindent 行范围：36\\[0.4em]

\begin{rustcode}
...
/// 零大小 [`DevFsManager`] 句柄；实际状态在静态 `DEVFS` 中。
// 本结构代码由AI完成
pub struct KernelDevFsManager;

// 本变量代码由AI完成
...
\end{rustcode}

\subsubsection*{序号21}
\noindent 路径：\code{os/components/wateros-fs/fs-devfs/devfs-impl/impl-kernel/src/lib.rs}\\
\noindent 行范围：39-44\\[0.4em]

\begin{rustcode}
...
static DEVFS: Mutex<DevFsImpl> = Mutex::new(DevFsImpl {
    nodes: Vec::new(),
    block_bindings: Vec::new(),
    character_bindings: Vec::new(),
    dt_unsupported_paths: Vec::new(),
...
\end{rustcode}

\subsubsection*{序号22}
\noindent 路径：\code{os/components/wateros-fs/fs-devfs/devfs-impl/impl-kernel/src/lib.rs}\\
\noindent 行范围：48-51\\[0.4em]

\begin{rustcode}
...
// 本方法代码由AI完成
fn linux_vd_disk_path(idx: usize) -> String {
    let letter = (b'a' + (idx as u8).min(25)) as char;
    format!("/dev/vd{}", letter)
}
...
\end{rustcode}

\subsubsection*{序号23}
\noindent 路径：\code{os/components/wateros-fs/fs-devfs/devfs-impl/impl-kernel/src/lib.rs}\\
\noindent 行范围：55-64\\[0.4em]

\begin{rustcode}
...
fn push_block_alias(inner: &mut DevFsImpl, path: String, dev: SharedBlockDevice) {
    if inner.block_bindings.iter().any(|(p, _)| p == &path) {
        return;
    }
    inner.nodes.push(api_v0::DevNode {
...
\end{rustcode}

\subsubsection*{序号24}
\noindent 路径：\code{os/components/wateros-fs/fs-devfs/devfs-impl/impl-kernel/src/lib.rs}\\
\noindent 行范围：68-77\\[0.4em]

\begin{rustcode}
...
fn push_char_alias(inner: &mut DevFsImpl, path: String, dev: SharedCharacterDevice) {
    if inner.character_bindings.iter().any(|(p, _)| p == &path) {
        return;
    }
    inner.nodes.push(api_v0::DevNode {
...
\end{rustcode}

\subsubsection*{序号25}
\noindent 路径：\code{os/components/wateros-fs/fs-devfs/devfs-impl/impl-kernel/src/lib.rs}\\
\noindent 行范围：81-152\\[0.4em]

\begin{rustcode}
...
    fn refresh(&mut self) {
        let block_snapshot: alloc::vec::Vec<_> = (0..block_device_count())
            .filter_map(|idx| block_device_at(idx).map(|dev| (idx, dev)))
            .collect();
        let char_snapshot: alloc::vec::Vec<_> = (0..character_device_count())
...
\end{rustcode}

\subsubsection*{序号26}
\noindent 路径：\code{os/components/wateros-fs/fs-devfs/devfs-impl/impl-kernel/src/lib.rs}\\
\noindent 行范围：155-157\\[0.4em]

\begin{rustcode}
...
// 本方法代码由AI完成
    fn set_dt_unsupported_paths(&mut self, paths: Vec<String>) {
        DEVFS.lock().dt_unsupported_paths = paths;
    }

...
\end{rustcode}

\subsubsection*{序号27}
\noindent 路径：\code{os/components/wateros-fs/fs-devfs/devfs-impl/impl-kernel/src/lib.rs}\\
\noindent 行范围：160-162\\[0.4em]

\begin{rustcode}
...
// 本方法代码由AI完成
    fn list_nodes(&self) -> Vec<DevNode> {
        DEVFS.lock().nodes.clone()
    }

...
\end{rustcode}

\subsubsection*{序号28}
\noindent 路径：\code{os/components/wateros-fs/fs-devfs/devfs-impl/impl-kernel/src/lib.rs}\\
\noindent 行范围：165-185\\[0.4em]

\begin{rustcode}
...
    fn register_block_device(
        &mut self,
        path: &str,
        device: SharedBlockDevice,
    ) -> fs_api_v0::FsResult<()> {
...
\end{rustcode}

\subsubsection*{序号29}
\noindent 路径：\code{os/components/wateros-fs/fs-devfs/devfs-impl/impl-kernel/src/lib.rs}\\
\noindent 行范围：188-196\\[0.4em]

\begin{rustcode}
...
    fn lookup_block_device(&self, path: &str) -> fs_api_v0::FsResult<SharedBlockDevice> {
        DEVFS
            .lock()
            .block_bindings
            .iter()
...
\end{rustcode}

\subsubsection*{序号30}
\noindent 路径：\code{os/components/wateros-fs/fs-devfs/devfs-impl/impl-kernel/src/lib.rs}\\
\noindent 行范围：199-219\\[0.4em]

\begin{rustcode}
...
    fn register_character_device(
        &mut self,
        path: &str,
        device: SharedCharacterDevice,
    ) -> fs_api_v0::FsResult<()> {
...
\end{rustcode}

\subsubsection*{序号31}
\noindent 路径：\code{os/components/wateros-fs/fs-devfs/devfs-impl/impl-kernel/src/lib.rs}\\
\noindent 行范围：222-230\\[0.4em]

\begin{rustcode}
...
    fn lookup_character_device(&self, path: &str) -> fs_api_v0::FsResult<SharedCharacterDevice> {
        DEVFS
            .lock()
            .character_bindings
            .iter()
...
\end{rustcode}

\subsubsection*{序号32}
\noindent 路径：\code{os/components/wateros-fs/fs-devfs/devfs-impl/impl-kernel/src/lib.rs}\\
\noindent 行范围：233-245\\[0.4em]

\begin{rustcode}
...
    fn default_root_block_path(&self) -> Option<String> {
        let inner = DEVFS.lock();
        inner
            .nodes
            .iter()
...
\end{rustcode}

\subsubsection*{序号33}
\noindent 路径：\code{os/components/wateros-fs/fs-devfs/devfs-impl/impl-kernel/src/lib.rs}\\
\noindent 行范围：250-254\\[0.4em]

\begin{rustcode}
...
pub fn refresh() -> usize {
    let mut m = KernelDevFsManager;
    m.refresh();
    m.list_nodes().len()
}
...
\end{rustcode}

\subsubsection*{序号34}
\noindent 路径：\code{os/components/wateros-fs/fs-devfs/devfs-impl/impl-kernel/src/lib.rs}\\
\noindent 行范围：258-261\\[0.4em]

\begin{rustcode}
...
// 本方法代码由AI完成
pub fn set_dt_unsupported_paths(paths: Vec<String>) {
    let mut m = KernelDevFsManager;
    m.set_dt_unsupported_paths(paths);
}
...
\end{rustcode}

\subsubsection*{序号35}
\noindent 路径：\code{os/components/wateros-fs/fs-devfs/devfs-impl/impl-kernel/src/lib.rs}\\
\noindent 行范围：266-269\\[0.4em]

\begin{rustcode}
...
// 本方法代码由AI完成
pub fn list_nodes() -> Vec<DevNode> {
    let m = KernelDevFsManager;
    m.list_nodes()
}
...
\end{rustcode}

\subsubsection*{序号36}
\noindent 路径：\code{os/components/wateros-fs/fs-devfs/devfs-impl/impl-kernel/src/lib.rs}\\
\noindent 行范围：274-277\\[0.4em]

\begin{rustcode}
...
// 本方法代码由AI完成
pub fn lookup_block_device(path: &str) -> fs_api_v0::FsResult<SharedBlockDevice> {
    let m = KernelDevFsManager;
    m.lookup_block_device(path)
}
...
\end{rustcode}

\subsubsection*{序号37}
\noindent 路径：\code{os/components/wateros-fs/fs-devfs/devfs-impl/impl-kernel/src/lib.rs}\\
\noindent 行范围：282-285\\[0.4em]

\begin{rustcode}
...
// 本方法代码由AI完成
pub fn lookup_character_device(path: &str) -> fs_api_v0::FsResult<SharedCharacterDevice> {
    let m = KernelDevFsManager;
    m.lookup_character_device(path)
}
...
\end{rustcode}

\subsubsection*{序号38}
\noindent 路径：\code{os/components/wateros-fs/fs-devfs/devfs-impl/impl-kernel/src/lib.rs}\\
\noindent 行范围：290-293\\[0.4em]

\begin{rustcode}
...
// 本方法代码由AI完成
pub fn default_root_block_path() -> Option<String> {
    let m = KernelDevFsManager;
    m.default_root_block_path()
}
...
\end{rustcode}

\subsubsection*{序号39}
\noindent 路径：\code{os/components/wateros-fs/fs-devfs/devfs-impl/impl-kernel/src/lib.rs}\\
\noindent 行范围：297\\[0.4em]

\begin{rustcode}
...
/// devfs 的 [`FsImpl`] 注册项；仅列能力，不参与块卷挂载。
// 本结构代码由AI完成
pub struct KernelDevFsImpl;

/// 全局 devfs impl 实例，供聚合层 `registered_fs_impls()` 引用。
...
\end{rustcode}

\subsubsection*{序号40}
\noindent 路径：\code{os/components/wateros-fs/fs-devfs/devfs-impl/impl-kernel/src/lib.rs}\\
\noindent 行范围：301\\[0.4em]

\begin{rustcode}
...
/// 全局 devfs impl 实例，供聚合层 `registered_fs_impls()` 引用。
// 本变量代码由AI完成
pub static IMPL: KernelDevFsImpl = KernelDevFsImpl;

// 本变量代码由AI完成
...
\end{rustcode}

\subsubsection*{序号41}
\noindent 路径：\code{os/components/wateros-fs/fs-devfs/devfs-impl/impl-kernel/src/lib.rs}\\
\noindent 行范围：304-305\\[0.4em]

\begin{rustcode}
...

// 本变量代码由AI完成
const SUPPORTED: &[FsCapability] =
    &[FsCapability::new(FsKind::DevFs, FsAccessMode::ReadOnly)];

...
\end{rustcode}

\subsubsection*{序号42}
\noindent 路径：\code{os/components/wateros-fs/fs-devfs/devfs-impl/impl-kernel/src/lib.rs}\\
\noindent 行范围：319-321\\[0.4em]

\begin{rustcode}
...
// 本方法代码由AI完成
    fn mount_ro(&self, _device: SharedBlockDevice) -> FsResult<SharedFs> {
        Err(FsError::Unsupported)
    }
}
\end{rustcode}

\subsubsection*{序号43}
\noindent 路径：\code{os/components/wateros-fs/fs-impl/impl-devfs/src/lib.rs}\\
\noindent 行范围：1-123\\[0.4em]

\begin{rustcode}
#![no_std]
//! 本模块代码由AI完成

//! 用户态/测试向的简化 devfs 视图：枚举块设备为 Linux 风格 `/dev/vda*` 与兼容 `/dev/vblk{n}`。
//!
...
\end{rustcode}

\subsubsection*{序号44}
\noindent 路径：\code{os/components/wateros-fs/fs-impl/impl-devfs/src/lib.rs}\\
\noindent 行范围：17-22\\[0.4em]

\begin{rustcode}
...
pub enum DevNodeType {
    /// 块设备索引节点。
    Block,
    /// 预留；当前刷新逻辑不生成此变体。
    Unsupported,
...
\end{rustcode}

\subsubsection*{序号45}
\noindent 路径：\code{os/components/wateros-fs/fs-impl/impl-devfs/src/lib.rs}\\
\noindent 行范围：27-34\\[0.4em]

\begin{rustcode}
...
pub struct DevNode {
    /// 逻辑路径。
    pub path: String,
    /// 节点类型。
    pub node_type: DevNodeType,
...
\end{rustcode}

\subsubsection*{序号46}
\noindent 路径：\code{os/components/wateros-fs/fs-impl/impl-devfs/src/lib.rs}\\
\noindent 行范围：38\\[0.4em]

\begin{rustcode}
...
// 与内核 devfs 不同：仅缓存枚举快照，无动态 register 表；单 Mutex 保护 bring-up 阶段并发。
// 本变量代码由AI完成
static DEV_NODES: Mutex<Vec<DevNode>> = Mutex::new(Vec::new());

// Linux 风格磁盘名：索引 0 → `/dev/vda`。
...
\end{rustcode}

\subsubsection*{序号47}
\noindent 路径：\code{os/components/wateros-fs/fs-impl/impl-devfs/src/lib.rs}\\
\noindent 行范围：42-45\\[0.4em]

\begin{rustcode}
...
// 本方法代码由AI完成
fn linux_vd_disk_path(index: usize) -> String {
    let letter = (b'a' + (index as u8).min(25)) as char;
    format!("/dev/vd{}", letter)
}
...
\end{rustcode}

\subsubsection*{序号48}
\noindent 路径：\code{os/components/wateros-fs/fs-impl/impl-devfs/src/lib.rs}\\
\noindent 行范围：48-57\\[0.4em]

\begin{rustcode}
...
fn push_node(nodes: &mut Vec<DevNode>, path: String, index: usize) {
    if nodes.iter().any(|n| n.path == path) {
        return;
    }
    nodes.push(DevNode {
...
\end{rustcode}

\subsubsection*{序号49}
\noindent 路径：\code{os/components/wateros-fs/fs-impl/impl-devfs/src/lib.rs}\\
\noindent 行范围：61-76\\[0.4em]

\begin{rustcode}
...
pub fn refresh() -> usize {
    let count = block_device_count();
    let mut nodes = DEV_NODES.lock();
    nodes.clear();
    for idx in 0..count {
...
\end{rustcode}

\subsubsection*{序号50}
\noindent 路径：\code{os/components/wateros-fs/fs-impl/impl-devfs/src/lib.rs}\\
\noindent 行范围：87-90\\[0.4em]

\begin{rustcode}
...
// 本方法代码由AI完成
pub fn lookup_block_device(path: &str) -> FsResult<SharedBlockDevice> {
    let idx = parse_block_index(path).ok_or(FsError::NotFound)?;
    block_device_at(idx).ok_or(FsError::NotFound)
}
...
\end{rustcode}

\subsubsection*{序号51}
\noindent 路径：\code{os/components/wateros-fs/fs-impl/impl-devfs/src/lib.rs}\\
\noindent 行范围：95-101\\[0.4em]

\begin{rustcode}
...
pub fn default_root_block_path() -> Option<String> {
    if block_device_count() == 0 {
        None
    } else {
        Some(linux_vd_disk_path(0))
...
\end{rustcode}

\subsubsection*{序号52}
\noindent 路径：\code{os/components/wateros-fs/fs-impl/impl-devfs/src/lib.rs}\\
\noindent 行范围：105-123\\[0.4em]

\begin{rustcode}
...
fn parse_block_index(path: &str) -> Option<usize> {
    if let Some(suffix) = path.strip_prefix("/dev/vblk") {
        return suffix.parse::<usize>().ok();
    }
    if let Some(rest) = path.strip_prefix("/dev/vd") {
...
\end{rustcode}

\subsubsection*{序号53}
\noindent 路径：\code{os/components/wateros-fs/fs-impl/impl-ext4-rs/src/lib.rs}\\
\noindent 行范围：1-829\\[0.4em]

\begin{rustcode}
#![no_std]
//! 本模块代码由AI完成

//! 基于 `ext4_rs` crate 的 ext4 实现。
//!
...
\end{rustcode}

\subsubsection*{序号54}
\noindent 路径：\code{os/components/wateros-fs/fs-impl/impl-ext4-rs/src/lib.rs}\\
\noindent 行范围：27\\[0.4em]

\begin{rustcode}
...
/// ext2/3/4 共用的 superblock magic（Linux 布局 `s_magic = 0xEF53`）。
// 本变量代码由AI完成
const EXT4_SUPER_MAGIC : u16 = 0xEF53;
/// 主 superblock 起始字节偏移（卷头 1024 字节之后）。
// 本变量代码由AI完成
...
\end{rustcode}

\subsubsection*{序号55}
\noindent 路径：\code{os/components/wateros-fs/fs-impl/impl-ext4-rs/src/lib.rs}\\
\noindent 行范围：30\\[0.4em]

\begin{rustcode}
...
/// 主 superblock 起始字节偏移（卷头 1024 字节之后）。
// 本变量代码由AI完成
const SUPERBLOCK_OFFSET : u64 = 1024;
/// `s_magic` 在 1024 字节 superblock 内的偏移。
// 本变量代码由AI完成
...
\end{rustcode}

\subsubsection*{序号56}
\noindent 路径：\code{os/components/wateros-fs/fs-impl/impl-ext4-rs/src/lib.rs}\\
\noindent 行范围：33\\[0.4em]

\begin{rustcode}
...
/// `s_magic` 在 1024 字节 superblock 内的偏移。
// 本变量代码由AI完成
const MAGIC_OFFSET_IN_SB : usize = 0x38;
/// ext4 根目录 inode 号（固定为 2）。
// 本变量代码由AI完成
...
\end{rustcode}

\subsubsection*{序号57}
\noindent 路径：\code{os/components/wateros-fs/fs-impl/impl-ext4-rs/src/lib.rs}\\
\noindent 行范围：36\\[0.4em]

\begin{rustcode}
...
/// ext4 根目录 inode 号（固定为 2）。
// 本变量代码由AI完成
const ROOT_INODE : u32 = 2;
/// 普通文件 mode 前缀（`S_IFREG`）。
// 本变量代码由AI完成
...
\end{rustcode}

\subsubsection*{序号58}
\noindent 路径：\code{os/components/wateros-fs/fs-impl/impl-ext4-rs/src/lib.rs}\\
\noindent 行范围：39\\[0.4em]

\begin{rustcode}
...
/// 普通文件 mode 前缀（`S_IFREG`）。
// 本变量代码由AI完成
const S_IFREG : u16 = 0o100000;
/// 目录 mode 前缀（`S_IFDIR`）。
// 本变量代码由AI完成
...
\end{rustcode}

\subsubsection*{序号59}
\noindent 路径：\code{os/components/wateros-fs/fs-impl/impl-ext4-rs/src/lib.rs}\\
\noindent 行范围：42\\[0.4em]

\begin{rustcode}
...
/// 目录 mode 前缀（`S_IFDIR`）。
// 本变量代码由AI完成
const S_IFDIR : u16 = 0o040000;

// 读取 superblock magic，判定是否为 ext2/3/4 卷。
...
\end{rustcode}

\subsubsection*{序号60}
\noindent 路径：\code{os/components/wateros-fs/fs-impl/impl-ext4-rs/src/lib.rs}\\
\noindent 行范围：46-53\\[0.4em]

\begin{rustcode}
...
fn probe_ext4_magic(device : &SharedBlockDevice) -> FsResult<bool> {
    let mut buf = [0u8; 2];
    device.lock()
          .read_bytes(SUPERBLOCK_OFFSET + MAGIC_OFFSET_IN_SB as u64,
                      &mut buf)
...
\end{rustcode}

\subsubsection*{序号61}
\noindent 路径：\code{os/components/wateros-fs/fs-impl/impl-ext4-rs/src/lib.rs}\\
\noindent 行范围：57-59\\[0.4em]

\begin{rustcode}
...
// 本结构代码由AI完成
struct BlockDevAdapter {
    device : SharedBlockDevice,
}

...
\end{rustcode}

\subsubsection*{序号62}
\noindent 路径：\code{os/components/wateros-fs/fs-impl/impl-ext4-rs/src/lib.rs}\\
\noindent 行范围：63-69\\[0.4em]

\begin{rustcode}
...
    fn read_offset(&self, offset : usize) -> Vec<u8> {
        let mut out = vec![0u8; ext4_rs::BLOCK_SIZE];
        let _ = self.device
                    .lock()
                    .read_bytes(offset as u64, &mut out);
...
\end{rustcode}

\subsubsection*{序号63}
\noindent 路径：\code{os/components/wateros-fs/fs-impl/impl-ext4-rs/src/lib.rs}\\
\noindent 行范围：72-74\\[0.4em]

\begin{rustcode}
...
// 本方法代码由AI完成
    fn write_offset(&self, offset : usize, data : &[u8]) {
        let _ = block_write_bytes(&self.device, offset as u64, data);
    }
}
...
\end{rustcode}

\subsubsection*{序号64}
\noindent 路径：\code{os/components/wateros-fs/fs-impl/impl-ext4-rs/src/lib.rs}\\
\noindent 行范围：79-115\\[0.4em]

\begin{rustcode}
...
fn block_write_bytes(dev : &SharedBlockDevice,
                     start_byte : u64,
                     src : &[u8])
                     -> Result<(), DriverError> {
    if src.is_empty() {
...
\end{rustcode}

\subsubsection*{序号65}
\noindent 路径：\code{os/components/wateros-fs/fs-impl/impl-ext4-rs/src/lib.rs}\\
\noindent 行范围：119-132\\[0.4em]

\begin{rustcode}
...
fn write_partial_block(bdev : &mut dyn driver_block_api_v0::BlockDevice,
                       block : usize,
                       offset : usize,
                       data : &[u8])
                       -> Result<(), DriverError> {
...
\end{rustcode}

\subsubsection*{序号66}
\noindent 路径：\code{os/components/wateros-fs/fs-impl/impl-ext4-rs/src/lib.rs}\\
\noindent 行范围：136-146\\[0.4em]

\begin{rustcode}
...
fn map_ext4_rs(err : Ext4Error) -> FsError {
    match err.error() {
        Errno::ENOENT => FsError::NotFound,
        Errno::EEXIST => FsError::Exists,
        Errno::ENOTDIR | Errno::EISDIR => FsError::NotAFile,
...
\end{rustcode}

\subsubsection*{序号67}
\noindent 路径：\code{os/components/wateros-fs/fs-impl/impl-ext4-rs/src/lib.rs}\\
\noindent 行范围：150-160\\[0.4em]

\begin{rustcode}
...
fn map_node_type(kind : InodeFileType) -> FsNodeType {
    if kind == InodeFileType::S_IFDIR {
        FsNodeType::Directory
    } else if kind == InodeFileType::S_IFLNK {
        FsNodeType::Symlink
...
\end{rustcode}

\subsubsection*{序号68}
\noindent 路径：\code{os/components/wateros-fs/fs-impl/impl-ext4-rs/src/lib.rs}\\
\noindent 行范围：164-176\\[0.4em]

\begin{rustcode}
...
fn split_parent_and_name(path : &str) -> FsResult<(&str, &str)> {
    let p = path.trim_end_matches('/');
    if p.is_empty() || p == "/" {
        return Err(FsError::InvalidPath);
    }
...
\end{rustcode}

\subsubsection*{序号69}
\noindent 路径：\code{os/components/wateros-fs/fs-impl/impl-ext4-rs/src/lib.rs}\\
\noindent 行范围：180-187\\[0.4em]

\begin{rustcode}
...
fn ensure_dir_inode(fs : &Ext4, inode : u32) -> FsResult<()> {
    let meta = metadata_for_inode(fs, inode)?;
    if meta.node_type == FsNodeType::Directory {
        Ok(())
    } else {
...
\end{rustcode}

\subsubsection*{序号70}
\noindent 路径：\code{os/components/wateros-fs/fs-impl/impl-ext4-rs/src/lib.rs}\\
\noindent 行范围：191-216\\[0.4em]

\begin{rustcode}
...
fn lookup_inode(fs : &Ext4, path : &str) -> FsResult<u32> {
    let p = path.trim_end_matches('/');
    if p.is_empty() || p == "/" {
        return Ok(ROOT_INODE);
    }
...
\end{rustcode}

\subsubsection*{序号71}
\noindent 路径：\code{os/components/wateros-fs/fs-impl/impl-ext4-rs/src/lib.rs}\\
\noindent 行范围：220-233\\[0.4em]

\begin{rustcode}
...
fn metadata_for_inode(fs : &Ext4, inode : u32) -> FsResult<FsMetadata> {
    let inode_ref = fs.get_inode_ref(inode);
    Ok(FsMetadata { node_type : map_node_type(inode_ref.inode
                                                       .file_type()),
                    size : inode_ref.inode
...
\end{rustcode}

\subsubsection*{序号72}
\noindent 路径：\code{os/components/wateros-fs/fs-impl/impl-ext4-rs/src/lib.rs}\\
\noindent 行范围：237-254\\[0.4em]

\begin{rustcode}
...
fn walk_ext4_rs_tree(fs : &Ext4RsFs, path : &str) {
    let Ok(entries) = ReadOnlyFs::read_dir(fs, path) else {
        return;
    };
    for entry in entries {
...
\end{rustcode}

\subsubsection*{序号73}
\noindent 路径：\code{os/components/wateros-fs/fs-impl/impl-ext4-rs/src/lib.rs}\\
\noindent 行范围：258-271\\[0.4em]

\begin{rustcode}
...
fn create_regular(fs : &mut Ext4, path : &str, mode : u16) -> FsResult<u32> {
    let (parent_path, name) = split_parent_and_name(path)?;
    let parent = lookup_inode(fs, parent_path)?;
    ensure_dir_inode(fs, parent)?;
    fs.fuse_mknod(parent as u64,
...
\end{rustcode}

\subsubsection*{序号74}
\noindent 路径：\code{os/components/wateros-fs/fs-impl/impl-ext4-rs/src/lib.rs}\\
\noindent 行范围：275-290\\[0.4em]

\begin{rustcode}
...
fn create_directory(fs : &mut Ext4, path : &str, mode : u16) -> FsResult<()> {
    let (parent_path, name) = split_parent_and_name(path)?;
    let parent = lookup_inode(fs, parent_path)?;
    ensure_dir_inode(fs, parent)?;
    match fs.fuse_lookup(parent as u64, name) {
...
\end{rustcode}

\subsubsection*{序号75}
\noindent 路径：\code{os/components/wateros-fs/fs-impl/impl-ext4-rs/src/lib.rs}\\
\noindent 行范围：294-296\\[0.4em]

\begin{rustcode}
...
// 本结构代码由AI完成
pub struct Ext4RsFs {
    fs : Option<Ext4>,
}

...
\end{rustcode}

\subsubsection*{序号76}
\noindent 路径：\code{os/components/wateros-fs/fs-impl/impl-ext4-rs/src/lib.rs}\\
\noindent 行范围：305-309\\[0.4em]

\begin{rustcode}
...
    fn fs(&self) -> FsResult<&Ext4> {
        self.fs
            .as_ref()
            .ok_or(FsError::NotMounted)
    }
...
\end{rustcode}

\subsubsection*{序号77}
\noindent 路径：\code{os/components/wateros-fs/fs-impl/impl-ext4-rs/src/lib.rs}\\
\noindent 行范围：313-317\\[0.4em]

\begin{rustcode}
...
    fn fs_mut(&mut self) -> FsResult<&mut Ext4> {
        self.fs
            .as_mut()
            .ok_or(FsError::NotMounted)
    }
...
\end{rustcode}

\subsubsection*{序号78}
\noindent 路径：\code{os/components/wateros-fs/fs-impl/impl-ext4-rs/src/lib.rs}\\
\noindent 行范围：322-326\\[0.4em]

\begin{rustcode}
...
    fn mount(&mut self, device : SharedBlockDevice) -> FsResult<()> {
        let dev : Arc<dyn Ext4RsBlockDevice> = Arc::new(BlockDevAdapter { device });
        self.fs = Some(Ext4::open(dev));
        Ok(())
    }
...
\end{rustcode}

\subsubsection*{序号79}
\noindent 路径：\code{os/components/wateros-fs/fs-impl/impl-ext4-rs/src/lib.rs}\\
\noindent 行范围：332-338\\[0.4em]

\begin{rustcode}
...
    fn exists(&self, path : &str) -> FsResult<bool> {
        match lookup_inode(self.fs()?, path) {
            Ok(_) => Ok(true),
            Err(FsError::NotFound) => Ok(false),
            Err(err) => Err(err),
...
\end{rustcode}

\subsubsection*{序号80}
\noindent 路径：\code{os/components/wateros-fs/fs-impl/impl-ext4-rs/src/lib.rs}\\
\noindent 行范围：341-344\\[0.4em]

\begin{rustcode}
...
// 本方法代码由AI完成
    fn metadata(&self, path : &str) -> FsResult<FsMetadata> {
        let inode = lookup_inode(self.fs()?, path)?;
        metadata_for_inode(self.fs()?, inode)
    }
...
\end{rustcode}

\subsubsection*{序号81}
\noindent 路径：\code{os/components/wateros-fs/fs-impl/impl-ext4-rs/src/lib.rs}\\
\noindent 行范围：347-358\\[0.4em]

\begin{rustcode}
...
    fn read(&self, path : &str) -> FsResult<Vec<u8>> {
        let meta = ReadOnlyFs::metadata(self, path)?;
        if meta.node_type != FsNodeType::File {
            return Err(FsError::NotAFile);
        }
...
\end{rustcode}

\subsubsection*{序号82}
\noindent 路径：\code{os/components/wateros-fs/fs-impl/impl-ext4-rs/src/lib.rs}\\
\noindent 行范围：361-380\\[0.4em]

\begin{rustcode}
...
    fn read_range(&self, path : &str, offset : u64, buf : &mut [u8]) -> FsResult<usize> {
        if buf.is_empty() {
            return Ok(0);
        }
        let fs = self.fs()?;
...
\end{rustcode}

\subsubsection*{序号83}
\noindent 路径：\code{os/components/wateros-fs/fs-impl/impl-ext4-rs/src/lib.rs}\\
\noindent 行范围：383-388\\[0.4em]

\begin{rustcode}
...
    fn read_prefix(&self, path : &str, len : usize) -> FsResult<Vec<u8>> {
        let mut buf = vec![0u8; len];
        let n = ReadOnlyFs::read_range(self, path, 0, &mut buf)?;
        buf.truncate(n);
        Ok(buf)
...
\end{rustcode}

\subsubsection*{序号84}
\noindent 路径：\code{os/components/wateros-fs/fs-impl/impl-ext4-rs/src/lib.rs}\\
\noindent 行范围：391-416\\[0.4em]

\begin{rustcode}
...
    fn read_dir(&self, path : &str) -> FsResult<Vec<FsDirEntry>> {
        let fs = self.fs()?;
        let inode = lookup_inode(fs, path)?;
        let attr = fs.fuse_getattr(inode as u64)
                     .map_err(map_ext4_rs)?;
...
\end{rustcode}

\subsubsection*{序号85}
\noindent 路径：\code{os/components/wateros-fs/fs-impl/impl-ext4-rs/src/lib.rs}\\
\noindent 行范围：419-433\\[0.4em]

\begin{rustcode}
...
    fn read_symlink(&self, path : &str) -> FsResult<Vec<u8>> {
        let fs = self.fs()?;
        let inode = lookup_inode(fs, path)?;
        let meta = metadata_for_inode(fs, inode)?;
        if meta.node_type != FsNodeType::Symlink {
...
\end{rustcode}

\subsubsection*{序号86}
\noindent 路径：\code{os/components/wateros-fs/fs-impl/impl-ext4-rs/src/lib.rs}\\
\noindent 行范围：436-440\\[0.4em]

\begin{rustcode}
...
    fn boot_dump_all_paths(&self) {
        if self.fs.is_some() {
            walk_ext4_rs_tree(self, "/");
        }
    }
...
\end{rustcode}

\subsubsection*{序号87}
\noindent 路径：\code{os/components/wateros-fs/fs-impl/impl-ext4-rs/src/lib.rs}\\
\noindent 行范围：451-458\\[0.4em]

\begin{rustcode}
...
    fn write_regular_file_at_root(&mut self, name : &str, data : &[u8]) -> FsResult<()> {
        if name.is_empty() || name.contains('/') {
            return Err(FsError::InvalidPath);
        }
        let mut path = String::from("/");
...
\end{rustcode}

\subsubsection*{序号88}
\noindent 路径：\code{os/components/wateros-fs/fs-impl/impl-ext4-rs/src/lib.rs}\\
\noindent 行范围：461-483\\[0.4em]

\begin{rustcode}
...
    fn write_regular_file(&mut self, path : &str, data : &[u8]) -> FsResult<()> {
        let fs = self.fs_mut()?;
        let inode = match lookup_inode(fs, path) {
            Ok(inode) => {
                let meta = metadata_for_inode(fs, inode)?;
...
\end{rustcode}

\subsubsection*{序号89}
\noindent 路径：\code{os/components/wateros-fs/fs-impl/impl-ext4-rs/src/lib.rs}\\
\noindent 行范围：486-527\\[0.4em]

\begin{rustcode}
...
    fn unlink(&mut self, path : &str) -> FsResult<()> {
        let fs = self.fs_mut()?;
        let (parent_path, name) = split_parent_and_name(path)?;
        let parent = lookup_inode(fs, parent_path)?;
        ensure_dir_inode(fs, parent)?;
...
\end{rustcode}

\subsubsection*{序号90}
\noindent 路径：\code{os/components/wateros-fs/fs-impl/impl-ext4-rs/src/lib.rs}\\
\noindent 行范围：530-540\\[0.4em]

\begin{rustcode}
...
    fn rmdir(&mut self, path : &str) -> FsResult<()> {
        if !ReadOnlyFs::read_dir(self, path)?.is_empty() {
            return Err(FsError::Exists);
        }
        let fs = self.fs_mut()?;
...
\end{rustcode}

\subsubsection*{序号91}
\noindent 路径：\code{os/components/wateros-fs/fs-impl/impl-ext4-rs/src/lib.rs}\\
\noindent 行范围：543-563\\[0.4em]

\begin{rustcode}
...
    fn write_range(&mut self, path : &str, offset : u64, data : &[u8]) -> FsResult<usize> {
        if data.is_empty() {
            return Ok(0);
        }
        let fs = self.fs_mut()?;
...
\end{rustcode}

\subsubsection*{序号92}
\noindent 路径：\code{os/components/wateros-fs/fs-impl/impl-ext4-rs/src/lib.rs}\\
\noindent 行范围：566-586\\[0.4em]

\begin{rustcode}
...
    fn truncate(&mut self, path : &str, len : u64) -> FsResult<()> {
        let fs = self.fs_mut()?;
        let inode = lookup_inode(fs, path)?;
        let meta = metadata_for_inode(fs, inode)?;
        if meta.node_type != FsNodeType::File {
...
\end{rustcode}

\subsubsection*{序号93}
\noindent 路径：\code{os/components/wateros-fs/fs-impl/impl-ext4-rs/src/lib.rs}\\
\noindent 行范围：589-593\\[0.4em]

\begin{rustcode}
...
    fn mkdir(&mut self, path : &str, mode : u32) -> FsResult<()> {
        let fs = self.fs_mut()?;
        let mode = S_IFDIR | (mode as u16 & 0o7777);
        create_directory(fs, path, mode)
    }
...
\end{rustcode}

\subsubsection*{序号94}
\noindent 路径：\code{os/components/wateros-fs/fs-impl/impl-ext4-rs/src/lib.rs}\\
\noindent 行范围：596-605\\[0.4em]

\begin{rustcode}
...
    fn chmod(&mut self, path : &str, mode : u32) -> FsResult<()> {
        let fs = self.fs_mut()?;
        let inode = lookup_inode(fs, path)?;
        let meta = metadata_for_inode(fs, inode)?;
        let mut inode_ref = fs.get_inode_ref(inode);
...
\end{rustcode}

\subsubsection*{序号95}
\noindent 路径：\code{os/components/wateros-fs/fs-impl/impl-ext4-rs/src/lib.rs}\\
\noindent 行范围：608-625\\[0.4em]

\begin{rustcode}
...
    fn chown(&mut self, path : &str, uid : Option<u32>, gid : Option<u32>) -> FsResult<()> {
        if uid.is_none() && gid.is_none() {
            return ReadOnlyFs::metadata(self, path).map(|_| ());
        }
        let fs = self.fs_mut()?;
...
\end{rustcode}

\subsubsection*{序号96}
\noindent 路径：\code{os/components/wateros-fs/fs-impl/impl-ext4-rs/src/lib.rs}\\
\noindent 行范围：628-655\\[0.4em]

\begin{rustcode}
...
    fn hardlink(&mut self, existing_path : &str, new_path : &str) -> FsResult<()> {
        let fs = self.fs_mut()?;
        let existing = lookup_inode(fs, existing_path)?;
        let meta = metadata_for_inode(fs, existing)?;
        if meta.node_type == FsNodeType::Directory {
...
\end{rustcode}

\subsubsection*{序号97}
\noindent 路径：\code{os/components/wateros-fs/fs-impl/impl-ext4-rs/src/lib.rs}\\
\noindent 行范围：658-693\\[0.4em]

\begin{rustcode}
...
    fn rename(&mut self, old_path : &str, new_path : &str) -> FsResult<()> {
        let (old_parent_path, old_name) = split_parent_and_name(old_path)?;
        let (new_parent_path, new_name) = split_parent_and_name(new_path)?;
        if old_parent_path != new_parent_path {
            return Err(FsError::Unsupported);
...
\end{rustcode}

\subsubsection*{序号98}
\noindent 路径：\code{os/components/wateros-fs/fs-impl/impl-ext4-rs/src/lib.rs}\\
\noindent 行范围：696-710\\[0.4em]

\begin{rustcode}
...
    fn symlink(&mut self, link_path : &str, target : &str) -> FsResult<()> {
        let fs = self.fs_mut()?;
        let (parent_path, name) = split_parent_and_name(link_path)?;
        let parent = lookup_inode(fs, parent_path)?;
        ensure_dir_inode(fs, parent)?;
...
\end{rustcode}

\subsubsection*{序号99}
\noindent 路径：\code{os/components/wateros-fs/fs-impl/impl-ext4-rs/src/lib.rs}\\
\noindent 行范围：713-726\\[0.4em]

\begin{rustcode}
...
    fn mknod(&mut self, path : &str, mode : u32, rdev : u32) -> FsResult<()> {
        let fs = self.fs_mut()?;
        let (parent_path, name) = split_parent_and_name(path)?;
        let parent = lookup_inode(fs, parent_path)?;
        ensure_dir_inode(fs, parent)?;
...
\end{rustcode}

\subsubsection*{序号100}
\noindent 路径：\code{os/components/wateros-fs/fs-impl/impl-ext4-rs/src/lib.rs}\\
\noindent 行范围：755-769\\[0.4em]

\begin{rustcode}
...
fn write_all(fs : &Ext4, inode : u32, offset : u64, data : &[u8]) -> FsResult<()> {
    let mut done = 0usize;
    while done < data.len() {
        let n = fs.write_at(inode,
                            usize::try_from(offset + done as u64).map_err(|_| FsError::Io)?,
...
\end{rustcode}

\subsubsection*{序号101}
\noindent 路径：\code{os/components/wateros-fs/fs-impl/impl-ext4-rs/src/lib.rs}\\
\noindent 行范围：773-784\\[0.4em]

\begin{rustcode}
...
fn zero_extend_file(fs : &Ext4, inode : u32, old_size : u64, new_size : u64) -> FsResult<()> {
    let zeroes = [0u8; ext4_rs::BLOCK_SIZE];
    let mut offset = old_size;
    while offset < new_size {
        let len =
...
\end{rustcode}

\subsubsection*{序号102}
\noindent 路径：\code{os/components/wateros-fs/fs-impl/impl-ext4-rs/src/lib.rs}\\
\noindent 行范围：788\\[0.4em]

\begin{rustcode}
...
/// ext4_rs 的 [`FsImpl`] 注册类型。
// 本结构代码由AI完成
pub struct Ext4RsImpl;

/// 全局 ext4-rs impl 实例。
...
\end{rustcode}

\subsubsection*{序号103}
\noindent 路径：\code{os/components/wateros-fs/fs-impl/impl-ext4-rs/src/lib.rs}\\
\noindent 行范围：792\\[0.4em]

\begin{rustcode}
...
/// 全局 ext4-rs impl 实例。
// 本变量代码由AI完成
pub static IMPL : Ext4RsImpl = Ext4RsImpl;

// 本变量代码由AI完成
...
\end{rustcode}

\subsubsection*{序号104}
\noindent 路径：\code{os/components/wateros-fs/fs-impl/impl-ext4-rs/src/lib.rs}\\
\noindent 行范围：795-796\\[0.4em]

\begin{rustcode}
...

// 本变量代码由AI完成
const SUPPORTED : &[FsCapability] = &[FsCapability::new(FsKind::Ext4, FsAccessMode::ReadOnly),
                                      FsCapability::new(FsKind::Ext4, FsAccessMode::ReadWrite)];

...
\end{rustcode}

\subsubsection*{序号105}
\noindent 路径：\code{os/components/wateros-fs/fs-impl/impl-ext4-rs/src/lib.rs}\\
\noindent 行范围：806-812\\[0.4em]

\begin{rustcode}
...
    fn probe(&self, device : &SharedBlockDevice) -> FsResult<Option<FsKind>> {
        if probe_ext4_magic(device)? {
            Ok(Some(FsKind::Ext4))
        } else {
            Ok(None)
...
\end{rustcode}

\subsubsection*{序号106}
\noindent 路径：\code{os/components/wateros-fs/fs-impl/impl-ext4-rs/src/lib.rs}\\
\noindent 行范围：815-820\\[0.4em]

\begin{rustcode}
...
    fn mount_ro(&self, device : SharedBlockDevice) -> FsResult<SharedFs> {
        log::info!("[fs::ext4-rs] mount_ro begin");
        let mut fs = Ext4RsFs::new();
        ReadOnlyFs::mount(&mut fs, device)?;
        Ok(Arc::new(Mutex::new(LocalFs::new(Box::new(fs)))))
...
\end{rustcode}

\subsubsection*{序号107}
\noindent 路径：\code{os/components/wateros-fs/fs-impl/impl-ext4-rs/src/lib.rs}\\
\noindent 行范围：823-828\\[0.4em]

\begin{rustcode}
...
    fn mount_rw(&self, device : SharedBlockDevice) -> FsResult<SharedRwFs> {
        log::info!("[fs::ext4-rs] mount_rw begin");
        let mut fs = Ext4RsFs::new();
        ReadWriteFs::mount_rw(&mut fs, device)?;
        Ok(Arc::new(Mutex::new(LocalRwFs::new(Box::new(fs)))))
...
\end{rustcode}

\subsubsection*{序号108}
\noindent 路径：\code{os/components/wateros-fs/fs-impl/impl-ext4/src/boot_inspect.rs}\\
\noindent 行范围：1-77\\[0.4em]

\begin{rustcode}
//! 启动期 boot 树遍历时的轻量内容探测：`.sh` 文本预览与可执行 ELF 头字段（非完整 loader）。
//! 本模块代码由AI完成

use alloc::string::String;

...
\end{rustcode}

\subsubsection*{序号109}
\noindent 路径：\code{os/components/wateros-fs/fs-impl/impl-ext4/src/boot_inspect.rs}\\
\noindent 行范围：10-25\\[0.4em]

\begin{rustcode}
...
fn parse_elf64_le_header(data: &[u8]) -> Option<(u8, u8, u16, u64)> {
    if data.len() < 0x40 {
        return None;
    }
    if &data[0..4] != b"\x7fELF" {
...
\end{rustcode}

\subsubsection*{序号110}
\noindent 路径：\code{os/components/wateros-fs/fs-impl/impl-ext4/src/boot_inspect.rs}\\
\noindent 行范围：28-30\\[0.4em]

\begin{rustcode}
...
// 本方法代码由AI完成
fn name_ends_with_sh(name: &[u8]) -> bool {
    name.len() >= 12 && name[name.len() - 12..] == *b"_testcode.sh"
}

...
\end{rustcode}

\subsubsection*{序号111}
\noindent 路径：\code{os/components/wateros-fs/fs-impl/impl-ext4/src/boot_inspect.rs}\\
\noindent 行范围：33-35\\[0.4em]

\begin{rustcode}
...
// 本方法代码由AI完成
fn is_unix_executable(mode: u16) -> bool {
    mode & (0o111 as u16) != 0
}

...
\end{rustcode}

\subsubsection*{序号112}
\noindent 路径：\code{os/components/wateros-fs/fs-impl/impl-ext4/src/boot_inspect.rs}\\
\noindent 行范围：39-51\\[0.4em]

\begin{rustcode}
...
pub fn log_sh_file(path_display: &str, content: &[u8]) {
    let preview = String::from_utf8_lossy(content);
    let max_chars = 512usize;
    let shown = preview.chars().take(max_chars).collect::<String>();
    let truncated = preview.chars().count() > max_chars;
...
\end{rustcode}

\subsubsection*{序号113}
\noindent 路径：\code{os/components/wateros-fs/fs-impl/impl-ext4/src/boot_inspect.rs}\\
\noindent 行范围：55-71\\[0.4em]

\begin{rustcode}
...
pub fn log_elf_exec_if_applicable(path_display: &str, prefix: &[u8], mode: u16) {
    if !is_unix_executable(mode) {
        return;
    }
    let Some((class, data, machine, entry)) = parse_elf64_le_header(prefix) else {
...
\end{rustcode}

\subsubsection*{序号114}
\noindent 路径：\code{os/components/wateros-fs/fs-impl/impl-ext4/src/boot_inspect.rs}\\
\noindent 行范围：75-77\\[0.4em]

\begin{rustcode}
...
/// 供目录项判断：是否应按 shell 脚本路径打印正文。
// 本方法代码由AI完成
pub fn should_dump_sh_script(name: &[u8]) -> bool {
    name_ends_with_sh(name)
}
\end{rustcode}

\subsubsection*{序号115}
\noindent 路径：\code{os/components/wateros-fs/fs-impl/impl-ext4/src/lib.rs}\\
\noindent 行范围：1-97\\[0.4em]

\begin{rustcode}
#![no_std]
//! 本模块代码由AI完成

//! 单一 ext4 文件系统实现：RO 与 RW 路径均基于 `ext4plus`，
//! 通过 [`api_v0::FsImpl`] 向聚合层注册一条能力（[`api_v0::FsKind::Ext4`] 同时支持 RO 与 RW）。
...
\end{rustcode}

\subsubsection*{序号116}
\noindent 路径：\code{os/components/wateros-fs/fs-impl/impl-ext4/src/lib.rs}\\
\noindent 行范围：31\\[0.4em]

\begin{rustcode}
...
/// ext4 superblock 中标识 ext2/3/4 的 magic（与 Linux 布局一致：`s_magic` 固定为 0xEF53）。
// 本变量代码由AI完成
const EXT4_SUPER_MAGIC : u16 = 0xEF53;
/// 主 superblock 起始字节偏移（卷头 1024 字节之后）。
// 本变量代码由AI完成
...
\end{rustcode}

\subsubsection*{序号117}
\noindent 路径：\code{os/components/wateros-fs/fs-impl/impl-ext4/src/lib.rs}\\
\noindent 行范围：34\\[0.4em]

\begin{rustcode}
...
/// 主 superblock 起始字节偏移（卷头 1024 字节之后）。
// 本变量代码由AI完成
const SUPERBLOCK_OFFSET : u64 = 1024;
/// `s_magic` 在 1024 字节 superblock 内的字节偏移（见内核 `ext4_super_block` 布局）。
// 本变量代码由AI完成
...
\end{rustcode}

\subsubsection*{序号118}
\noindent 路径：\code{os/components/wateros-fs/fs-impl/impl-ext4/src/lib.rs}\\
\noindent 行范围：37\\[0.4em]

\begin{rustcode}
...
/// `s_magic` 在 1024 字节 superblock 内的字节偏移（见内核 `ext4_super_block` 布局）。
// 本变量代码由AI完成
const MAGIC_OFFSET_IN_SB : usize = 0x38;

/// 通过读取 superblock magic 判定卷是否为 ext2/3/4（轻量探测，不校验完整 checksum）。
...
\end{rustcode}

\subsubsection*{序号119}
\noindent 路径：\code{os/components/wateros-fs/fs-impl/impl-ext4/src/lib.rs}\\
\noindent 行范围：41-50\\[0.4em]

\begin{rustcode}
...
fn probe_ext4_magic(device : &SharedBlockDevice) -> FsResult<bool> {
    let mut buf = [0u8; 2];
    let r = device.lock()
                  .read_bytes(SUPERBLOCK_OFFSET + MAGIC_OFFSET_IN_SB as u64,
                              &mut buf);
...
\end{rustcode}

\subsubsection*{序号120}
\noindent 路径：\code{os/components/wateros-fs/fs-impl/impl-ext4/src/lib.rs}\\
\noindent 行范围：54\\[0.4em]

\begin{rustcode}
...
/// 共有 `FsImpl` 入口的具体类型；实例为 [`IMPL`]。
// 本结构代码由AI完成
pub struct Ext4FsImpl;

/// ext4 impl 的 ' static 注册项。聚合层应在 `registered_fs_impls()` 中放入 `&IMPL`。
...
\end{rustcode}

\subsubsection*{序号121}
\noindent 路径：\code{os/components/wateros-fs/fs-impl/impl-ext4/src/lib.rs}\\
\noindent 行范围：58\\[0.4em]

\begin{rustcode}
...
/// ext4 impl 的 ' static 注册项。聚合层应在 `registered_fs_impls()` 中放入 `&IMPL`。
// 本变量代码由AI完成
pub static IMPL : Ext4FsImpl = Ext4FsImpl;

// 本变量代码由AI完成
...
\end{rustcode}

\subsubsection*{序号122}
\noindent 路径：\code{os/components/wateros-fs/fs-impl/impl-ext4/src/lib.rs}\\
\noindent 行范围：61-62\\[0.4em]

\begin{rustcode}
...

// 本变量代码由AI完成
const SUPPORTED : &[FsCapability] = &[FsCapability::new(FsKind::Ext4, FsAccessMode::ReadOnly),
                                      FsCapability::new(FsKind::Ext4, FsAccessMode::ReadWrite)];

...
\end{rustcode}

\subsubsection*{序号123}
\noindent 路径：\code{os/components/wateros-fs/fs-impl/impl-ext4/src/lib.rs}\\
\noindent 行范围：72-78\\[0.4em]

\begin{rustcode}
...
    fn probe(&self, device : &SharedBlockDevice) -> FsResult<Option<FsKind>> {
        if probe_ext4_magic(device)? {
            Ok(Some(FsKind::Ext4))
        } else {
            Ok(None)
...
\end{rustcode}

\subsubsection*{序号124}
\noindent 路径：\code{os/components/wateros-fs/fs-impl/impl-ext4/src/lib.rs}\\
\noindent 行范围：81-87\\[0.4em]

\begin{rustcode}
...
    fn mount_ro(&self, device : SharedBlockDevice) -> FsResult<SharedFs> {
        logging::info!("[fs::ext4] mount_ro begin");
        let mut fs = Ext4Fs::new();
        ReadOnlyFs::mount(&mut fs, device)?;
        let shared : SharedFs = Arc::new(Mutex::new(LocalFs::new(Box::new(fs))));
...
\end{rustcode}

\subsubsection*{序号125}
\noindent 路径：\code{os/components/wateros-fs/fs-impl/impl-ext4/src/lib.rs}\\
\noindent 行范围：90-96\\[0.4em]

\begin{rustcode}
...
    fn mount_rw(&self, device : SharedBlockDevice) -> FsResult<SharedRwFs> {
        logging::trace!("[fs::ext4] mount_rw begin");
        let mut fs = Ext4FsRw::new();
        ReadWriteFs::mount_rw(&mut fs, device)?;
        let shared : SharedRwFs = Arc::new(Mutex::new(LocalRwFs::new(Box::new(fs))));
...
\end{rustcode}

\subsubsection*{序号126}
\noindent 路径：\code{os/components/wateros-fs/fs-impl/impl-ext4/src/ro.rs}\\
\noindent 行范围：1-291\\[0.4em]

\begin{rustcode}
//! 只读路径（基于 `ext4plus`）：实现 [`api_v0::ReadOnlyFs`] 与启动期目录树打印。
//!
//! 大块文件读路径刻意按 `driver_block_api_v0::BLOCK_SIZE` 分片，规避部分 VirtIO 小扇区组合下的一次性整读问题（见本文件中 `ReadOnlyFs::read` 实现内注释）。
//! 本模块代码由AI完成

...
\end{rustcode}

\subsubsection*{序号127}
\noindent 路径：\code{os/components/wateros-fs/fs-impl/impl-ext4/src/ro.rs}\\
\noindent 行范围：21-23\\[0.4em]

\begin{rustcode}
...
// 本结构代码由AI完成
struct BlockDeviceReader {
    device: SharedBlockDevice,
}

...
\end{rustcode}

\subsubsection*{序号128}
\noindent 路径：\code{os/components/wateros-fs/fs-impl/impl-ext4/src/ro.rs}\\
\noindent 行范围：27-36\\[0.4em]

\begin{rustcode}
...
    fn read(
        &self,
        start_byte: u64,
        dst: &mut [u8],
    ) -> Result<(), Box<dyn core::error::Error + Send + Sync + 'static>> {
...
\end{rustcode}

\subsubsection*{序号129}
\noindent 路径：\code{os/components/wateros-fs/fs-impl/impl-ext4/src/ro.rs}\\
\noindent 行范围：41\\[0.4em]

\begin{rustcode}
...
#[derive(Debug)]
// 本结构代码由AI完成
struct BlockIoError(driver_block_api_v0::DriverError);

impl core::fmt::Display for BlockIoError {
...
\end{rustcode}

\subsubsection*{序号130}
\noindent 路径：\code{os/components/wateros-fs/fs-impl/impl-ext4/src/ro.rs}\\
\noindent 行范围：45-47\\[0.4em]

\begin{rustcode}
...
// 本方法代码由AI完成
    fn fmt(&self, f: &mut core::fmt::Formatter<'_>) -> core::fmt::Result {
        write!(f, "block io error: {:?}", self.0)
    }
}
...
\end{rustcode}

\subsubsection*{序号131}
\noindent 路径：\code{os/components/wateros-fs/fs-impl/impl-ext4/src/ro.rs}\\
\noindent 行范围：54-56\\[0.4em]

\begin{rustcode}
...
// 本结构代码由AI完成
pub struct Ext4Fs {
    fs: Option<Ext4>,
}

...
\end{rustcode}

\subsubsection*{序号132}
\noindent 路径：\code{os/components/wateros-fs/fs-impl/impl-ext4/src/ro.rs}\\
\noindent 行范围：69-73\\[0.4em]

\begin{rustcode}
...
    fn mount(&mut self, device: SharedBlockDevice) -> FsResult<()> {
        let fs = Ext4::load(Box::new(BlockDeviceReader { device })).map_err(map_ext4_plus)?;
        self.fs = Some(fs);
        Ok(())
    }
...
\end{rustcode}

\subsubsection*{序号133}
\noindent 路径：\code{os/components/wateros-fs/fs-impl/impl-ext4/src/ro.rs}\\
\noindent 行范围：79-82\\[0.4em]

\begin{rustcode}
...
// 本方法代码由AI完成
    fn exists(&self, path: &str) -> FsResult<bool> {
        let path = Path::try_from(path).map_err(|_| FsError::InvalidPath)?;
        self.fs()?.exists(path).map_err(map_ext4_plus)
    }
...
\end{rustcode}

\subsubsection*{序号134}
\noindent 路径：\code{os/components/wateros-fs/fs-impl/impl-ext4/src/ro.rs}\\
\noindent 行范围：85-93\\[0.4em]

\begin{rustcode}
...
    fn metadata(&self, path: &str) -> FsResult<FsMetadata> {
        let fs = self.fs()?;
        let path = Path::try_from(path).map_err(|_| FsError::InvalidPath)?;
        let inode = fs
            .path_to_inode(path, FollowSymlinks::All)
...
\end{rustcode}

\subsubsection*{序号135}
\noindent 路径：\code{os/components/wateros-fs/fs-impl/impl-ext4/src/ro.rs}\\
\noindent 行范围：96-124\\[0.4em]

\begin{rustcode}
...
    fn read_dir(&self, path: &str) -> FsResult<Vec<FsDirEntry>> {
        let fs = self.fs()?;
        let pathv = Path::try_from(path).map_err(|_| FsError::InvalidPath)?;
        let rd = fs.read_dir(pathv).map_err(map_ext4_plus)?;
        let mut out = Vec::new();
...
\end{rustcode}

\subsubsection*{序号136}
\noindent 路径：\code{os/components/wateros-fs/fs-impl/impl-ext4/src/ro.rs}\\
\noindent 行范围：127-156\\[0.4em]

\begin{rustcode}
...
    fn read(&self, path: &str) -> FsResult<Vec<u8>> {
        // 不用 `Ext4::read` 单次 `read_inode_file`：在大量目录遍历 + 大块整读时，
        // 曾与内核侧 VirtIO 512B 扇区路径组合出现首读 ELF 头损坏；按块设备逻辑块
        // 粒度循环 `File::read_bytes` 组装整文件更稳。
        let fs = self.fs()?;
...
\end{rustcode}

\subsubsection*{序号137}
\noindent 路径：\code{os/components/wateros-fs/fs-impl/impl-ext4/src/ro.rs}\\
\noindent 行范围：159-190\\[0.4em]

\begin{rustcode}
...
    fn read_range(&self, path: &str, offset: u64, buf: &mut [u8]) -> FsResult<usize> {
        if buf.is_empty() {
            return Ok(0);
        }
        let fs = self.fs()?;
...
\end{rustcode}

\subsubsection*{序号138}
\noindent 路径：\code{os/components/wateros-fs/fs-impl/impl-ext4/src/ro.rs}\\
\noindent 行范围：193-198\\[0.4em]

\begin{rustcode}
...
    fn read_prefix(&self, path: &str, len: usize) -> FsResult<Vec<u8>> {
        let mut buf = vec![0u8; len];
        let n = self.read_range(path, 0, &mut buf)?;
        buf.truncate(n);
        Ok(buf)
...
\end{rustcode}

\subsubsection*{序号139}
\noindent 路径：\code{os/components/wateros-fs/fs-impl/impl-ext4/src/ro.rs}\\
\noindent 行范围：201-204\\[0.4em]

\begin{rustcode}
...
// 本方法代码由AI完成
    fn boot_dump_all_paths(&self) {
        let Some(ext4) = self.fs.as_ref() else { return };
        walk_ext4_tree(ext4, Path::ROOT);
    }
...
\end{rustcode}

\subsubsection*{序号140}
\noindent 路径：\code{os/components/wateros-fs/fs-impl/impl-ext4/src/ro.rs}\\
\noindent 行范围：209-231\\[0.4em]

\begin{rustcode}
...
fn read_regular_file_capped(ext4: &Ext4, path: Path<'_>, cap: usize) -> Option<Vec<u8>> {
    let meta = ext4.metadata(path).ok()?;
    if !meta.file_type().is_regular_file() {
        return None;
    }
...
\end{rustcode}

\subsubsection*{序号141}
\noindent 路径：\code{os/components/wateros-fs/fs-impl/impl-ext4/src/ro.rs}\\
\noindent 行范围：235-268\\[0.4em]

\begin{rustcode}
...
fn walk_ext4_tree(ext4: &Ext4, dir: Path<'_>) {
    let Ok(rd) = ext4.read_dir(dir) else { return };
    for item in rd {
        let Ok(ent) = item else { continue };
        let name = ent.file_name();
...
\end{rustcode}

\subsubsection*{序号142}
\noindent 路径：\code{os/components/wateros-fs/fs-impl/impl-ext4/src/ro.rs}\\
\noindent 行范围：271-291\\[0.4em]

\begin{rustcode}
...
fn map_metadata(meta: &Metadata, inode: u64) -> FsMetadata {
    let node_type = if meta.is_dir() {
        FsNodeType::Directory
    } else if meta.is_symlink() {
        FsNodeType::Symlink
...
\end{rustcode}

\subsubsection*{序号143}
\noindent 路径：\code{os/components/wateros-fs/fs-impl/impl-ext4/src/rw.rs}\\
\noindent 行范围：1-971\\[0.4em]

\begin{rustcode}
//! 读写路径（基于 `ext4plus`，beta；写路径无完整 journal，仅用于 bring-up 与小文件测试）。
//!
//! I/O 边界：块读写适配器将驱动错误装箱为 `ext4plus` 期望的 `Error` trait object；按块读改写见本模块中的 `block_write_bytes`。
//! 本模块代码由AI完成

...
\end{rustcode}

\subsubsection*{序号144}
\noindent 路径：\code{os/components/wateros-fs/fs-impl/impl-ext4/src/rw.rs}\\
\noindent 行范围：26\\[0.4em]

\begin{rustcode}
...

// 本变量代码由AI完成
static EXT4_SMALL_READ_CACHE: Mutex<SmallReadCache> = Mutex::new(SmallReadCache::new());

// 本结构代码由AI完成
...
\end{rustcode}

\subsubsection*{序号145}
\noindent 路径：\code{os/components/wateros-fs/fs-impl/impl-ext4/src/rw.rs}\\
\noindent 行范围：29-34\\[0.4em]

\begin{rustcode}
...
struct SmallReadCache {
    valid: bool,
    dev_id: usize,
    block: u64,
    data: [u8; BLOCK_SIZE],
...
\end{rustcode}

\subsubsection*{序号146}
\noindent 路径：\code{os/components/wateros-fs/fs-impl/impl-ext4/src/rw.rs}\\
\noindent 行范围：48-50\\[0.4em]

\begin{rustcode}
...
// 本方法代码由AI完成
fn shared_block_dev_id(dev: &SharedBlockDevice) -> usize {
    alloc::sync::Arc::as_ptr(dev) as *const () as usize
}

...
\end{rustcode}

\subsubsection*{序号147}
\noindent 路径：\code{os/components/wateros-fs/fs-impl/impl-ext4/src/rw.rs}\\
\noindent 行范围：53-93\\[0.4em]

\begin{rustcode}
...
fn read_with_small_cache(dev: &SharedBlockDevice,
                         start_byte: u64,
                         dst: &mut [u8])
                         -> Result<bool, DriverError> {
    if dst.is_empty() {
...
\end{rustcode}

\subsubsection*{序号148}
\noindent 路径：\code{os/components/wateros-fs/fs-impl/impl-ext4/src/rw.rs}\\
\noindent 行范围：96-119\\[0.4em]

\begin{rustcode}
...
fn invalidate_small_read_cache(dev: &SharedBlockDevice, start_byte: u64, len: usize) {
    if len == 0 {
        return;
    }
    let Ok(start) = usize::try_from(start_byte) else {
...
\end{rustcode}

\subsubsection*{序号149}
\noindent 路径：\code{os/components/wateros-fs/fs-impl/impl-ext4/src/rw.rs}\\
\noindent 行范围：123-125\\[0.4em]

\begin{rustcode}
...
// 本结构代码由AI完成
struct BlockDevRw {
    device: SharedBlockDevice,
}

...
\end{rustcode}

\subsubsection*{序号150}
\noindent 路径：\code{os/components/wateros-fs/fs-impl/impl-ext4/src/rw.rs}\\
\noindent 行范围：129-137\\[0.4em]

\begin{rustcode}
...
    fn read(
        &self,
        start_byte: u64,
        dst: &mut [u8],
    ) -> Result<(), Box<dyn Error + Send + Sync + 'static>> {
...
\end{rustcode}

\subsubsection*{序号151}
\noindent 路径：\code{os/components/wateros-fs/fs-impl/impl-ext4/src/rw.rs}\\
\noindent 行范围：142-148\\[0.4em]

\begin{rustcode}
...
    fn write(
        &self,
        start_byte: u64,
        src: &[u8],
    ) -> Result<(), Box<dyn Error + Send + Sync + 'static>> {
...
\end{rustcode}

\subsubsection*{序号152}
\noindent 路径：\code{os/components/wateros-fs/fs-impl/impl-ext4/src/rw.rs}\\
\noindent 行范围：152\\[0.4em]

\begin{rustcode}
...

// 本方法代码由AI完成
fn driver_boxed(e: DriverError) -> Box<dyn Error + Send + Sync + 'static> { Box::new(DriverErr(e)) }

#[derive(Debug)]
...
\end{rustcode}

\subsubsection*{序号153}
\noindent 路径：\code{os/components/wateros-fs/fs-impl/impl-ext4/src/rw.rs}\\
\noindent 行范围：156\\[0.4em]

\begin{rustcode}
...
#[derive(Debug)]
// 本结构代码由AI完成
struct DriverErr(DriverError);

impl core::fmt::Display for DriverErr {
...
\end{rustcode}

\subsubsection*{序号154}
\noindent 路径：\code{os/components/wateros-fs/fs-impl/impl-ext4/src/rw.rs}\\
\noindent 行范围：160-162\\[0.4em]

\begin{rustcode}
...
// 本方法代码由AI完成
    fn fmt(&self, f: &mut core::fmt::Formatter<'_>) -> core::fmt::Result {
        write!(f, "{:?}", self.0)
    }
}
...
\end{rustcode}

\subsubsection*{序号155}
\noindent 路径：\code{os/components/wateros-fs/fs-impl/impl-ext4/src/rw.rs}\\
\noindent 行范围：169-212\\[0.4em]

\begin{rustcode}
...
fn block_write_bytes(
    dev: &SharedBlockDevice,
    start_byte: u64,
    src: &[u8],
) -> Result<(), DriverError> {
...
\end{rustcode}

\subsubsection*{序号156}
\noindent 路径：\code{os/components/wateros-fs/fs-impl/impl-ext4/src/rw.rs}\\
\noindent 行范围：215-230\\[0.4em]

\begin{rustcode}
...
fn write_partial_block(
    bdev: &mut dyn driver_block_api_v0::BlockDevice,
    block: usize,
    offset: usize,
    data: &[u8],
...
\end{rustcode}

\subsubsection*{序号157}
\noindent 路径：\code{os/components/wateros-fs/fs-impl/impl-ext4/src/rw.rs}\\
\noindent 行范围：234-251\\[0.4em]

\begin{rustcode}
...
pub(crate) fn map_ext4_plus(err: Ext4Error) -> FsError {
    match err {
        Ext4Error::NotFound => FsError::NotFound,
        Ext4Error::NotAbsolute | Ext4Error::MalformedPath | Ext4Error::PathTooLong => {
            FsError::InvalidPath
...
\end{rustcode}

\subsubsection*{序号158}
\noindent 路径：\code{os/components/wateros-fs/fs-impl/impl-ext4/src/rw.rs}\\
\noindent 行范围：255-257\\[0.4em]

\begin{rustcode}
...
// 本结构代码由AI完成
pub struct Ext4FsRw {
    fs: Option<Ext4>,
}

...
\end{rustcode}

\subsubsection*{序号159}
\noindent 路径：\code{os/components/wateros-fs/fs-impl/impl-ext4/src/rw.rs}\\
\noindent 行范围：270-276\\[0.4em]

\begin{rustcode}
...
    fn mount_rw(&mut self, device: SharedBlockDevice) -> FsResult<()> {
        let reader = Box::new(BlockDevRw { device: device.clone() });
        let writer = Box::new(BlockDevRw { device });
        let fs = Ext4::load_with_writer(reader, Some(writer)).map_err(map_ext4_plus)?;
        self.fs = Some(fs);
...
\end{rustcode}

\subsubsection*{序号160}
\noindent 路径：\code{os/components/wateros-fs/fs-impl/impl-ext4/src/rw.rs}\\
\noindent 行范围：282-289\\[0.4em]

\begin{rustcode}
...
    fn write_regular_file_at_root(&mut self, name: &str, data: &[u8]) -> FsResult<()> {
        if name.is_empty() || name.contains('/') {
            return Err(FsError::InvalidPath);
        }
        let mut path = alloc::string::String::from("/");
...
\end{rustcode}

\subsubsection*{序号161}
\noindent 路径：\code{os/components/wateros-fs/fs-impl/impl-ext4/src/rw.rs}\\
\noindent 行范围：292-340\\[0.4em]

\begin{rustcode}
...
    fn write_regular_file(&mut self, path: &str, data: &[u8]) -> FsResult<()> {
        let (parent, name) = split_parent_and_name(path)?;
        let fs = self.fs()?;
        let name = DirEntryName::try_from(name).map_err(|_| FsError::InvalidPath)?;
        let parent_path = Path::try_from(parent).map_err(|_| FsError::InvalidPath)?;
...
\end{rustcode}

\subsubsection*{序号162}
\noindent 路径：\code{os/components/wateros-fs/fs-impl/impl-ext4/src/rw.rs}\\
\noindent 行范围：343-369\\[0.4em]

\begin{rustcode}
...
    fn write_range(&mut self, path: &str, offset: u64, data: &[u8]) -> FsResult<usize> {
        if data.is_empty() {
            return Ok(0);
        }
        let fs = self.fs()?;
...
\end{rustcode}

\subsubsection*{序号163}
\noindent 路径：\code{os/components/wateros-fs/fs-impl/impl-ext4/src/rw.rs}\\
\noindent 行范围：372-382\\[0.4em]

\begin{rustcode}
...
    fn truncate(&mut self, path: &str, len: u64) -> FsResult<()> {
        let fs = self.fs()?;
        let pathv = Path::try_from(path).map_err(|_| FsError::InvalidPath)?;
        let mut inode = fs
            .path_to_inode(pathv, FollowSymlinks::All)
...
\end{rustcode}

\subsubsection*{序号164}
\noindent 路径：\code{os/components/wateros-fs/fs-impl/impl-ext4/src/rw.rs}\\
\noindent 行范围：385-416\\[0.4em]

\begin{rustcode}
...
    fn mkdir(&mut self, path: &str, mode: u32) -> FsResult<()> {
        let (parent, name) = split_parent_and_name(path)?;
        let fs = self.fs()?;
        let name = DirEntryName::try_from(name).map_err(|_| FsError::InvalidPath)?;
        let parent_path = Path::try_from(parent).map_err(|_| FsError::InvalidPath)?;
...
\end{rustcode}

\subsubsection*{序号165}
\noindent 路径：\code{os/components/wateros-fs/fs-impl/impl-ext4/src/rw.rs}\\
\noindent 行范围：419-430\\[0.4em]

\begin{rustcode}
...
    fn chmod(&mut self, path: &str, mode: u32) -> FsResult<()> {
        let fs = self.fs()?;
        let pathv = Path::try_from(path).map_err(|_| FsError::InvalidPath)?;
        let mut inode = fs
            .path_to_inode(pathv, FollowSymlinks::All)
...
\end{rustcode}

\subsubsection*{序号166}
\noindent 路径：\code{os/components/wateros-fs/fs-impl/impl-ext4/src/rw.rs}\\
\noindent 行范围：433-450\\[0.4em]

\begin{rustcode}
...
    fn chown(&mut self, path: &str, uid: Option<u32>, gid: Option<u32>) -> FsResult<()> {
        if uid.is_none() && gid.is_none() {
            return self.metadata(path).map(|_| ());
        }
        let fs = self.fs()?;
...
\end{rustcode}

\subsubsection*{序号167}
\noindent 路径：\code{os/components/wateros-fs/fs-impl/impl-ext4/src/rw.rs}\\
\noindent 行范围：453-466\\[0.4em]

\begin{rustcode}
...
    fn setxattr(&mut self, path: &str, name: &str, value: &[u8]) -> FsResult<()> {
        if name.is_empty() || name.contains('\0') {
            return Err(FsError::InvalidPath);
        }
        let fs = self.fs()?;
...
\end{rustcode}

\subsubsection*{序号168}
\noindent 路径：\code{os/components/wateros-fs/fs-impl/impl-ext4/src/rw.rs}\\
\noindent 行范围：469-490\\[0.4em]

\begin{rustcode}
...
    fn getxattr(&self, path: &str, name: &str, buf: &mut [u8]) -> FsResult<usize> {
        if name.is_empty() || name.contains('\0') {
            return Err(FsError::InvalidPath);
        }
        let fs = self.fs()?;
...
\end{rustcode}

\subsubsection*{序号169}
\noindent 路径：\code{os/components/wateros-fs/fs-impl/impl-ext4/src/rw.rs}\\
\noindent 行范围：493-513\\[0.4em]

\begin{rustcode}
...
    fn listxattr(&self, path: &str, buf: &mut [u8]) -> FsResult<usize> {
        let fs = self.fs()?;
        let pathv = Path::try_from(path).map_err(|_| FsError::InvalidPath)?;
        let inode = fs
            .path_to_inode(pathv, FollowSymlinks::All)
...
\end{rustcode}

\subsubsection*{序号170}
\noindent 路径：\code{os/components/wateros-fs/fs-impl/impl-ext4/src/rw.rs}\\
\noindent 行范围：516-529\\[0.4em]

\begin{rustcode}
...
    fn removexattr(&mut self, path: &str, name: &str) -> FsResult<()> {
        if name.is_empty() || name.contains('\0') {
            return Err(FsError::InvalidPath);
        }
        let fs = self.fs()?;
...
\end{rustcode}

\subsubsection*{序号171}
\noindent 路径：\code{os/components/wateros-fs/fs-impl/impl-ext4/src/rw.rs}\\
\noindent 行范围：532-553\\[0.4em]

\begin{rustcode}
...
    fn unlink(&mut self, path: &str) -> FsResult<()> {
        let (parent, name) = split_parent_and_name(path)?;
        let fs = self.fs()?;
        let parent_path = Path::try_from(parent).map_err(|_| FsError::InvalidPath)?;
        let parent_inode = fs
...
\end{rustcode}

\subsubsection*{序号172}
\noindent 路径：\code{os/components/wateros-fs/fs-impl/impl-ext4/src/rw.rs}\\
\noindent 行范围：556-582\\[0.4em]

\begin{rustcode}
...
    fn rename(&mut self, old_path: &str, new_path: &str) -> FsResult<()> {
        let (old_parent, old_name) = split_parent_and_name(old_path)?;
        let (new_parent, new_name) = split_parent_and_name(new_path)?;
        if old_parent != new_parent {
            return Err(FsError::Unsupported);
...
\end{rustcode}

\subsubsection*{序号173}
\noindent 路径：\code{os/components/wateros-fs/fs-impl/impl-ext4/src/rw.rs}\\
\noindent 行范围：585-615\\[0.4em]

\begin{rustcode}
...
    fn hardlink(&mut self, existing_path: &str, new_path: &str) -> FsResult<()> {
        let fs = self.fs()?;
        let existing_pathv = Path::try_from(existing_path).map_err(|_| FsError::InvalidPath)?;
        let mut inode = fs
            .path_to_inode(existing_pathv, FollowSymlinks::All)
...
\end{rustcode}

\subsubsection*{序号174}
\noindent 路径：\code{os/components/wateros-fs/fs-impl/impl-ext4/src/rw.rs}\\
\noindent 行范围：618-645\\[0.4em]

\begin{rustcode}
...
    fn rmdir(&mut self, path: &str) -> FsResult<()> {
        let (parent, name) = split_parent_and_name(path)?;
        let fs = self.fs()?;
        let parent_path = Path::try_from(parent).map_err(|_| FsError::InvalidPath)?;
        let parent_inode = fs
...
\end{rustcode}

\subsubsection*{序号175}
\noindent 路径：\code{os/components/wateros-fs/fs-impl/impl-ext4/src/rw.rs}\\
\noindent 行范围：648-652\\[0.4em]

\begin{rustcode}
...
    fn exists(&self, path: &str) -> FsResult<bool> {
        let fs = self.fs()?;
        let pathv = Path::try_from(path).map_err(|_| FsError::InvalidPath)?;
        fs.exists(pathv).map_err(map_ext4_plus)
    }
...
\end{rustcode}

\subsubsection*{序号176}
\noindent 路径：\code{os/components/wateros-fs/fs-impl/impl-ext4/src/rw.rs}\\
\noindent 行范围：655-663\\[0.4em]

\begin{rustcode}
...
    fn metadata(&self, path: &str) -> FsResult<FsMetadata> {
        let fs = self.fs()?;
        let pathv = Path::try_from(path).map_err(|_| FsError::InvalidPath)?;
        let inode = fs
            .path_to_inode(pathv, FollowSymlinks::All)
...
\end{rustcode}

\subsubsection*{序号177}
\noindent 路径：\code{os/components/wateros-fs/fs-impl/impl-ext4/src/rw.rs}\\
\noindent 行范围：666-695\\[0.4em]

\begin{rustcode}
...
    fn read(&self, path: &str) -> FsResult<Vec<u8>> {
        let fs = self.fs()?;
        let pathv = Path::try_from(path).map_err(|_| FsError::InvalidPath)?;
        let meta = fs.metadata(pathv).map_err(map_ext4_plus)?;
        if !meta.file_type().is_regular_file() {
...
\end{rustcode}

\subsubsection*{序号178}
\noindent 路径：\code{os/components/wateros-fs/fs-impl/impl-ext4/src/rw.rs}\\
\noindent 行范围：698-732\\[0.4em]

\begin{rustcode}
...
    fn read_range(&self, path: &str, offset: u64, buf: &mut [u8]) -> FsResult<usize> {
        if buf.is_empty() {
            return Ok(0);
        }
        let fs = self.fs()?;
...
\end{rustcode}

\subsubsection*{序号179}
\noindent 路径：\code{os/components/wateros-fs/fs-impl/impl-ext4/src/rw.rs}\\
\noindent 行范围：735-763\\[0.4em]

\begin{rustcode}
...
    fn read_dir(&self, path: &str) -> FsResult<Vec<FsDirEntry>> {
        let fs = self.fs()?;
        let pathv = Path::try_from(path).map_err(|_| FsError::InvalidPath)?;
        let mut rd = fs.read_dir(pathv).map_err(map_ext4_plus)?;
        let mut out = Vec::new();
...
\end{rustcode}

\subsubsection*{序号180}
\noindent 路径：\code{os/components/wateros-fs/fs-impl/impl-ext4/src/rw.rs}\\
\noindent 行范围：767-786\\[0.4em]

\begin{rustcode}
...
fn map_rw_metadata(meta: &Metadata, inode: u64) -> FsMetadata {
    let node_type = if meta.is_dir() {
        FsNodeType::Directory
    } else if meta.is_symlink() {
        FsNodeType::Symlink
...
\end{rustcode}

\subsubsection*{序号181}
\noindent 路径：\code{os/components/wateros-fs/fs-impl/impl-ext4/src/rw.rs}\\
\noindent 行范围：790-792\\[0.4em]

\begin{rustcode}
...
// 本方法代码由AI完成
fn linux_mkdir_mode_to_inode_mode(mode: u32) -> InodeMode {
    InodeMode::S_IFDIR | linux_mode_perm_to_inode_mode(mode)
}

...
\end{rustcode}

\subsubsection*{序号182}
\noindent 路径：\code{os/components/wateros-fs/fs-impl/impl-ext4/src/rw.rs}\\
\noindent 行范围：796-804\\[0.4em]

\begin{rustcode}
...
fn inode_type_bits(mode: InodeMode) -> InodeMode {
    mode & (InodeMode::S_IFIFO
        | InodeMode::S_IFCHR
        | InodeMode::S_IFDIR
        | InodeMode::S_IFBLK
...
\end{rustcode}

\subsubsection*{序号183}
\noindent 路径：\code{os/components/wateros-fs/fs-impl/impl-ext4/src/rw.rs}\\
\noindent 行范围：808-848\\[0.4em]

\begin{rustcode}
...
fn linux_mode_perm_to_inode_mode(mode: u32) -> InodeMode {
    let perm = mode & 0o7777;
    let mut m = InodeMode::empty();
    if perm & 0o4000 != 0 {
        m |= InodeMode::S_ISUID;
...
\end{rustcode}

\subsubsection*{序号184}
\noindent 路径：\code{os/components/wateros-fs/fs-impl/impl-ext4/src/rw.rs}\\
\noindent 行范围：852-863\\[0.4em]

\begin{rustcode}
...
fn split_parent_and_name(path: &str) -> FsResult<(&str, &str)> {
    let p = path.trim_end_matches('/');
    if p.is_empty() || p == "/" {
        return Err(FsError::InvalidPath);
    }
...
\end{rustcode}

\subsubsection*{序号185}
\noindent 路径：\code{os/components/wateros-fs/fs-impl/impl-ext4/src/selftest.rs}\\
\noindent 行范围：1-132\\[0.4em]

\begin{rustcode}
//! ext4 自检：RO 与 RW 根卷路径。
//! 本模块代码由AI完成

use alloc::{string::String, vec};
use api_v0::{FsNodeType, FsResult, ReadOnlyFs, ReadWriteFs, SharedFs, SharedRwFs};
...
\end{rustcode}

\subsubsection*{序号186}
\noindent 路径：\code{os/components/wateros-fs/fs-impl/impl-ext4/src/selftest.rs}\\
\noindent 行范围：10-28\\[0.4em]

\begin{rustcode}
...
pub fn ro_self_test(fs: SharedFs) -> FsResult<()> {
    let fs = fs.lock();
    match fs.metadata("/") {
        Ok(meta) => {
            logging::info!(
...
\end{rustcode}

\subsubsection*{序号187}
\noindent 路径：\code{os/components/wateros-fs/fs-impl/impl-ext4/src/selftest.rs}\\
\noindent 行范围：32-50\\[0.4em]

\begin{rustcode}
...
pub fn rw_self_test(rw: SharedRwFs) -> FsResult<()> {
    let fs = rw.lock();
    match fs.metadata("/") {
        Ok(meta) => {
            logging::info!(
...
\end{rustcode}

\subsubsection*{序号188}
\noindent 路径：\code{os/components/wateros-fs/fs-impl/impl-ext4/src/selftest.rs}\\
\noindent 行范围：54-81\\[0.4em]

\begin{rustcode}
...
fn ro_range_smoke_readonly(fs: &dyn ReadOnlyFs) -> FsResult<()> {
    let candidates = [
        "/elf/000_hello_world.elf",
        "/src/bin/000_hello_world.rs",
    ];
...
\end{rustcode}

\subsubsection*{序号189}
\noindent 路径：\code{os/components/wateros-fs/fs-impl/impl-ext4/src/selftest.rs}\\
\noindent 行范围：84-112\\[0.4em]

\begin{rustcode}
...
fn rw_range_smoke(fs: &dyn ReadWriteFs) -> FsResult<()> {
    let candidates = [
        "/elf/000_hello_world.elf",
        "/src/bin/000_hello_world.rs",
        "/glibc/basic/text.txt",
...
\end{rustcode}

\subsubsection*{序号190}
\noindent 路径：\code{os/components/wateros-fs/fs-impl/impl-ext4/src/selftest.rs}\\
\noindent 行范围：116-132\\[0.4em]

\begin{rustcode}
...
pub fn rw_mkdir_verify(rw: SharedRwFs, dir_name: &str) -> FsResult<()> {
    if dir_name.is_empty() || dir_name.contains('/') {
        return Err(api_v0::FsError::InvalidPath);
    }
    let mut path = String::from("/");
...
\end{rustcode}

\subsubsection*{序号191}
\noindent 路径：\code{os/components/wateros-fs/fs-procfs/procfs-impl/impl-kernel/src/lib.rs}\\
\noindent 行范围：1-614\\[0.4em]

\begin{rustcode}
#![no_std]
//! 本模块代码由AI完成

//! 内核 procfs：从 task/cred/mm 与 VFS 回调生成 `/proc` 内容。

...
\end{rustcode}

\subsubsection*{序号192}
\noindent 路径：\code{os/components/wateros-fs/fs-procfs/procfs-impl/impl-kernel/src/lib.rs}\\
\noindent 行范围：23\\[0.4em]

\begin{rustcode}
...

// 本变量代码由AI完成
static ARGV_LOOKUP: Mutex<Option<TaskArgvLookup>> = Mutex::new(None);
// 本变量代码由AI完成
static EXE_LOOKUP: Mutex<Option<TaskExeLookup>> = Mutex::new(None);
...
\end{rustcode}

\subsubsection*{序号193}
\noindent 路径：\code{os/components/wateros-fs/fs-procfs/procfs-impl/impl-kernel/src/lib.rs}\\
\noindent 行范围：25\\[0.4em]

\begin{rustcode}
...
static ARGV_LOOKUP: Mutex<Option<TaskArgvLookup>> = Mutex::new(None);
// 本变量代码由AI完成
static EXE_LOOKUP: Mutex<Option<TaskExeLookup>> = Mutex::new(None);
// 本变量代码由AI完成
static MOUNT_LOOKUP: Mutex<Option<MountListLookup>> = Mutex::new(None);
...
\end{rustcode}

\subsubsection*{序号194}
\noindent 路径：\code{os/components/wateros-fs/fs-procfs/procfs-impl/impl-kernel/src/lib.rs}\\
\noindent 行范围：27\\[0.4em]

\begin{rustcode}
...
static EXE_LOOKUP: Mutex<Option<TaskExeLookup>> = Mutex::new(None);
// 本变量代码由AI完成
static MOUNT_LOOKUP: Mutex<Option<MountListLookup>> = Mutex::new(None);

/// 注册按 leader task id 查询 argv 的回调（VFS 层在 init 时注入）。
...
\end{rustcode}

\subsubsection*{序号195}
\noindent 路径：\code{os/components/wateros-fs/fs-procfs/procfs-impl/impl-kernel/src/lib.rs}\\
\noindent 行范围：31-33\\[0.4em]

\begin{rustcode}
...
// 本方法代码由AI完成
pub fn register_task_argv_lookup(f: TaskArgvLookup) {
    *ARGV_LOOKUP.lock() = Some(f);
}

...
\end{rustcode}

\subsubsection*{序号196}
\noindent 路径：\code{os/components/wateros-fs/fs-procfs/procfs-impl/impl-kernel/src/lib.rs}\\
\noindent 行范围：37-39\\[0.4em]

\begin{rustcode}
...
// 本方法代码由AI完成
pub fn register_task_exe_lookup(f: TaskExeLookup) {
    *EXE_LOOKUP.lock() = Some(f);
}

...
\end{rustcode}

\subsubsection*{序号197}
\noindent 路径：\code{os/components/wateros-fs/fs-procfs/procfs-impl/impl-kernel/src/lib.rs}\\
\noindent 行范围：43-45\\[0.4em]

\begin{rustcode}
...
// 本方法代码由AI完成
pub fn register_mount_list_lookup(f: MountListLookup) {
    *MOUNT_LOOKUP.lock() = Some(f);
}

...
\end{rustcode}

\subsubsection*{序号198}
\noindent 路径：\code{os/components/wateros-fs/fs-procfs/procfs-impl/impl-kernel/src/lib.rs}\\
\noindent 行范围：49-52\\[0.4em]

\begin{rustcode}
...
// 本方法代码由AI完成
fn argv_for(leader: TaskId) -> Option<Vec<String>> {
    let lookup = *ARGV_LOOKUP.lock();
    lookup.and_then(|f| f(leader))
}
...
\end{rustcode}

\subsubsection*{序号199}
\noindent 路径：\code{os/components/wateros-fs/fs-procfs/procfs-impl/impl-kernel/src/lib.rs}\\
\noindent 行范围：56-59\\[0.4em]

\begin{rustcode}
...
// 本方法代码由AI完成
fn exe_for(leader: TaskId) -> Option<String> {
    let lookup = *EXE_LOOKUP.lock();
    lookup.and_then(|f| f(leader))
}
...
\end{rustcode}

\subsubsection*{序号200}
\noindent 路径：\code{os/components/wateros-fs/fs-procfs/procfs-impl/impl-kernel/src/lib.rs}\\
\noindent 行范围：63-66\\[0.4em]

\begin{rustcode}
...
// 本方法代码由AI完成
fn mount_lines() -> Vec<ProcMountLine> {
    let lookup = *MOUNT_LOOKUP.lock();
    lookup.map(|f| f()).unwrap_or_default()
}
...
\end{rustcode}

\subsubsection*{序号201}
\noindent 路径：\code{os/components/wateros-fs/fs-procfs/procfs-impl/impl-kernel/src/lib.rs}\\
\noindent 行范围：88-104\\[0.4em]

\begin{rustcode}
...
fn proc_inode(node: ProcNode) -> u64 {
    match node {
        ProcNode::Root => 1,
        ProcNode::Meminfo => 2,
        ProcNode::Cpuinfo => 7,
...
\end{rustcode}

\subsubsection*{序号202}
\noindent 路径：\code{os/components/wateros-fs/fs-procfs/procfs-impl/impl-kernel/src/lib.rs}\\
\noindent 行范围：108-141\\[0.4em]

\begin{rustcode}
...
fn normalize_rel(path: &str) -> String {
    use alloc::borrow::Cow;

    let abs: Cow<'_, str> = if path.is_empty() {
        Cow::Borrowed("/")
...
\end{rustcode}

\subsubsection*{序号203}
\noindent 路径：\code{os/components/wateros-fs/fs-procfs/procfs-impl/impl-kernel/src/lib.rs}\\
\noindent 行范围：145-151\\[0.4em]

\begin{rustcode}
...
fn parse_pid(name: &str) -> Option<ProcessId> {
    if name == "self" {
        Some(task::current_process_task_snapshot()?.pid)
    } else {
        Some(ProcessId::from_raw(name.parse::<usize>().ok()?))
...
\end{rustcode}

\subsubsection*{序号204}
\noindent 路径：\code{os/components/wateros-fs/fs-procfs/procfs-impl/impl-kernel/src/lib.rs}\\
\noindent 行范围：155-178\\[0.4em]

\begin{rustcode}
...
fn parse_node(path: &str) -> Option<ProcNode> {
    let p = normalize_rel(path);
    if p == "/" {
        return Some(ProcNode::Root);
    }
...
\end{rustcode}

\subsubsection*{序号205}
\noindent 路径：\code{os/components/wateros-fs/fs-procfs/procfs-impl/impl-kernel/src/lib.rs}\\
\noindent 行范围：182-184\\[0.4em]

\begin{rustcode}
...
// 本方法代码由AI完成
fn process_visible(pid: ProcessId) -> bool {
    task::process_snapshot(pid).is_some()
}

...
\end{rustcode}

\subsubsection*{序号206}
\noindent 路径：\code{os/components/wateros-fs/fs-procfs/procfs-impl/impl-kernel/src/lib.rs}\\
\noindent 行范围：188-199\\[0.4em]

\begin{rustcode}
...
fn comm_for(pid: ProcessId) -> String {
    let leader = task::leader_task_for_process(pid).unwrap_or(0);
    if let Some(argv) = argv_for(leader) {
        if let Some(arg0) = argv.first() {
            return basename(arg0);
...
\end{rustcode}

\subsubsection*{序号207}
\noindent 路径：\code{os/components/wateros-fs/fs-procfs/procfs-impl/impl-kernel/src/lib.rs}\\
\noindent 行范围：202-204\\[0.4em]

\begin{rustcode}
...
// 本方法代码由AI完成
fn format_cpuinfo() -> Vec<u8> {
    b"processor\t: 0\ncpu family\t: 0\nmodel name\t: WaterOS QEMU virt\n".to_vec()
}

...
\end{rustcode}

\subsubsection*{序号208}
\noindent 路径：\code{os/components/wateros-fs/fs-procfs/procfs-impl/impl-kernel/src/lib.rs}\\
\noindent 行范围：207-224\\[0.4em]

\begin{rustcode}
...
fn format_cgroups() -> Vec<u8> {
    alloc::format!(
        "#subsys_name\thierarchy\tnum_cgroups\tenabled\n\
memory\t1\t1\t1\n\
cpuset\t2\t1\t1\n\
...
\end{rustcode}

\subsubsection*{序号209}
\noindent 路径：\code{os/components/wateros-fs/fs-procfs/procfs-impl/impl-kernel/src/lib.rs}\\
\noindent 行范围：228-230\\[0.4em]

\begin{rustcode}
...
// 本方法代码由AI完成
fn basename(path: &str) -> String {
    path.rsplit('/').next().unwrap_or(path).to_string()
}

...
\end{rustcode}

\subsubsection*{序号210}
\noindent 路径：\code{os/components/wateros-fs/fs-procfs/procfs-impl/impl-kernel/src/lib.rs}\\
\noindent 行范围：234-244\\[0.4em]

\begin{rustcode}
...
fn state_char(process: ProcessState, leader_state: Option<TaskState>) -> char {
    match process {
        ProcessState::Exited(_) | ProcessState::Exiting(_) => 'Z',
        ProcessState::Stopped { .. } => 'T',
        ProcessState::Running => match leader_state {
...
\end{rustcode}

\subsubsection*{序号211}
\noindent 路径：\code{os/components/wateros-fs/fs-procfs/procfs-impl/impl-kernel/src/lib.rs}\\
\noindent 行范围：247-279\\[0.4em]

\begin{rustcode}
...
fn format_stat(pid: ProcessId) -> FsResult<Vec<u8>> {
    let process = task::process_snapshot(pid).ok_or(FsError::NotFound)?;
    let leader = process.leader_task_id;
    // Linux utime/stime are in USER_HZ jiffies; tick_count tracks scheduler ticks (~100/s).
    let jiffies = task::task_snapshot(leader)
...
\end{rustcode}

\subsubsection*{序号212}
\noindent 路径：\code{os/components/wateros-fs/fs-procfs/procfs-impl/impl-kernel/src/lib.rs}\\
\noindent 行范围：282-317\\[0.4em]

\begin{rustcode}
...
fn format_status(pid: ProcessId) -> FsResult<Vec<u8>> {
    let process = task::process_snapshot(pid).ok_or(FsError::NotFound)?;
    let leader = process.leader_task_id;
    let cred = cred::credentials_for(leader);
    let comm = comm_for(pid);
...
\end{rustcode}

\subsubsection*{序号213}
\noindent 路径：\code{os/components/wateros-fs/fs-procfs/procfs-impl/impl-kernel/src/lib.rs}\\
\noindent 行范围：322-326\\[0.4em]

\begin{rustcode}
...
struct ProcMemoryKb {
    size_kb : usize,
    rss_kb : usize,
    private_dirty_kb : usize,
}
...
\end{rustcode}

\subsubsection*{序号214}
\noindent 路径：\code{os/components/wateros-fs/fs-procfs/procfs-impl/impl-kernel/src/lib.rs}\\
\noindent 行范围：329-336\\[0.4em]

\begin{rustcode}
...
fn process_memory_kb(pid : ProcessId) -> FsResult<ProcMemoryKb> {
    let process = task::process_snapshot(pid).ok_or(FsError::NotFound)?;
    let thread_extra = process.task_count.saturating_sub(1) * 64;
    let heap_like = 4096usize.saturating_add(thread_extra);
    Ok(ProcMemoryKb { size_kb : heap_like,
...
\end{rustcode}

\subsubsection*{序号215}
\noindent 路径：\code{os/components/wateros-fs/fs-procfs/procfs-impl/impl-kernel/src/lib.rs}\\
\noindent 行范围：339-348\\[0.4em]

\begin{rustcode}
...
fn format_maps(pid : ProcessId) -> FsResult<Vec<u8>> {
    let mem = process_memory_kb(pid)?;
    let heap_start = 0x1000_0000usize;
    let heap_end = heap_start.saturating_add(mem.size_kb.saturating_mul(1024));
    Ok(format!(
...
\end{rustcode}

\subsubsection*{序号216}
\noindent 路径：\code{os/components/wateros-fs/fs-procfs/procfs-impl/impl-kernel/src/lib.rs}\\
\noindent 行范围：351-365\\[0.4em]

\begin{rustcode}
...
fn format_smaps(pid : ProcessId) -> FsResult<Vec<u8>> {
    let mem = process_memory_kb(pid)?;
    let start = 0x1000_0000usize;
    let end = start.saturating_add(mem.size_kb.saturating_mul(1024));
    let line = format!(
...
\end{rustcode}

\subsubsection*{序号217}
\noindent 路径：\code{os/components/wateros-fs/fs-procfs/procfs-impl/impl-kernel/src/lib.rs}\\
\noindent 行范围：368-387\\[0.4em]

\begin{rustcode}
...
fn format_cmdline(pid: ProcessId) -> FsResult<Vec<u8>> {
    let process = task::process_snapshot(pid).ok_or(FsError::NotFound)?;
    let leader = process.leader_task_id;
    let mut out = Vec::new();
    if let Some(argv) = argv_for(leader) {
...
\end{rustcode}

\subsubsection*{序号218}
\noindent 路径：\code{os/components/wateros-fs/fs-procfs/procfs-impl/impl-kernel/src/lib.rs}\\
\noindent 行范围：390-399\\[0.4em]

\begin{rustcode}
...
fn format_meminfo() -> Vec<u8> {
    let stats = mm_frame_alloctor::frame_mem_stats();
    format!(
        "MemTotal:\t{}\tkB\nMemFree:\t{}\tkB\nMemAvailable:\t{}\tkB\nBuffers:\t0\tkB\nCached:\t0\tkB\n",
        stats.total_bytes() / 1024,
...
\end{rustcode}

\subsubsection*{序号219}
\noindent 路径：\code{os/components/wateros-fs/fs-procfs/procfs-impl/impl-kernel/src/lib.rs}\\
\noindent 行范围：402-412\\[0.4em]

\begin{rustcode}
...
fn format_mounts() -> Vec<u8> {
    let mut out = Vec::new();
    for line in mount_lines() {
        let row = format!(
            "{} {} {} rw,relatime 0 0\n",
...
\end{rustcode}

\subsubsection*{序号220}
\noindent 路径：\code{os/components/wateros-fs/fs-procfs/procfs-impl/impl-kernel/src/lib.rs}\\
\noindent 行范围：416\\[0.4em]

\begin{rustcode}
...
/// 内核 procfs 只读视图（零大小；无实例状态）。
// 本结构代码由AI完成
pub struct KernelProcFs;

/// 返回全局 procfs 视图句柄。
...
\end{rustcode}

\subsubsection*{序号221}
\noindent 路径：\code{os/components/wateros-fs/fs-procfs/procfs-impl/impl-kernel/src/lib.rs}\\
\noindent 行范围：426-440\\[0.4em]

\begin{rustcode}
...
    fn exists(&self, rel_path: &str) -> FsResult<bool> {
        let Some(node) = parse_node(rel_path) else {
            return Ok(false);
        };
        Ok(match node {
...
\end{rustcode}

\subsubsection*{序号222}
\noindent 路径：\code{os/components/wateros-fs/fs-procfs/procfs-impl/impl-kernel/src/lib.rs}\\
\noindent 行范围：443-488\\[0.4em]

\begin{rustcode}
...
    fn metadata(&self, rel_path: &str) -> FsResult<FsMetadata> {
        let node = parse_node(rel_path).ok_or(FsError::NotFound)?;
        match node {
            ProcNode::Root | ProcNode::PidDir(_) => Ok(FsMetadata {
                node_type: FsNodeType::Directory,
...
\end{rustcode}

\subsubsection*{序号223}
\noindent 路径：\code{os/components/wateros-fs/fs-procfs/procfs-impl/impl-kernel/src/lib.rs}\\
\noindent 行范围：491-507\\[0.4em]

\begin{rustcode}
...
    fn read(&self, rel_path: &str) -> FsResult<Vec<u8>> {
        let node = parse_node(rel_path).ok_or(FsError::NotFound)?;
        match node {
            ProcNode::Root | ProcNode::PidDir(_) => Err(FsError::NotAFile),
            ProcNode::Meminfo => Ok(format_meminfo()),
...
\end{rustcode}

\subsubsection*{序号224}
\noindent 路径：\code{os/components/wateros-fs/fs-procfs/procfs-impl/impl-kernel/src/lib.rs}\\
\noindent 行范围：510-573\\[0.4em]

\begin{rustcode}
...
    fn read_dir(&self, rel_path: &str) -> FsResult<Vec<FsDirEntry>> {
        let node = parse_node(rel_path).ok_or(FsError::NotFound)?;
        match node {
            ProcNode::Root => {
                let mut entries = vec![
...
\end{rustcode}

\subsubsection*{序号225}
\noindent 路径：\code{os/components/wateros-fs/fs-procfs/procfs-impl/impl-kernel/src/lib.rs}\\
\noindent 行范围：578\\[0.4em]

\begin{rustcode}
...
/// procfs 的 [`FsImpl`] 注册项；仅列能力，不参与块卷挂载。
// 本结构代码由AI完成
pub struct KernelProcFsImpl;

/// 全局 procfs impl 实例。
...
\end{rustcode}

\subsubsection*{序号226}
\noindent 路径：\code{os/components/wateros-fs/fs-procfs/procfs-impl/impl-kernel/src/lib.rs}\\
\noindent 行范围：582\\[0.4em]

\begin{rustcode}
...
/// 全局 procfs impl 实例。
// 本变量代码由AI完成
pub static IMPL: KernelProcFsImpl = KernelProcFsImpl;

// 本变量代码由AI完成
...
\end{rustcode}

\subsubsection*{序号227}
\noindent 路径：\code{os/components/wateros-fs/fs-procfs/procfs-impl/impl-kernel/src/lib.rs}\\
\noindent 行范围：585-586\\[0.4em]

\begin{rustcode}
...

// 本变量代码由AI完成
const SUPPORTED: &[FsCapability] =
    &[FsCapability::new(FsKind::Other("procfs"), FsAccessMode::ReadOnly)];

...
\end{rustcode}

\subsubsection*{序号228}
\noindent 路径：\code{os/components/wateros-fs/fs-procfs/procfs-impl/impl-kernel/src/lib.rs}\\
\noindent 行范围：600-605\\[0.4em]

\begin{rustcode}
...
    fn mount_ro(
        &self,
        _device: driver_block_api_v0::SharedBlockDevice,
    ) -> fs_api_v0::FsResult<fs_api_v0::SharedFs> {
        Err(FsError::Unsupported)
...
\end{rustcode}

\subsubsection*{序号229}
\noindent 路径：\code{os/components/wateros-fs/fs-procfs/procfs-impl/impl-kernel/src/lib.rs}\\
\noindent 行范围：610-614\\[0.4em]

\begin{rustcode}
...
pub fn test() {
    let v = view();
    let _ = v.read_dir("/");
    logging::info!("[fs::procfs] self_test ok");
}
\end{rustcode}

\subsubsection*{序号230}
\noindent 路径：\code{os/components/wateros-fs/fs-rootfs/rootfs-impl/impl-kernel/src/lib.rs}\\
\noindent 行范围：1-208\\[0.4em]

\begin{rustcode}
#![no_std]
//! 本模块代码由AI完成

//! 内核 rootfs 实现：静态槽位保存当前根 [`SharedFs`]、根块路径与聚合层注入的 [`FsImpl`]。
//!
...
\end{rustcode}

\subsubsection*{序号231}
\noindent 路径：\code{os/components/wateros-fs/fs-rootfs/rootfs-impl/impl-kernel/src/lib.rs}\\
\noindent 行范围：17\\[0.4em]

\begin{rustcode}
...

// 本变量代码由AI完成
static MOUNT_GENERATION: AtomicU64 = AtomicU64::new(0);

/// 零大小 [`RootFsManager`] 句柄；读写 `static ROOT_FS` 等状态。
...
\end{rustcode}

\subsubsection*{序号232}
\noindent 路径：\code{os/components/wateros-fs/fs-rootfs/rootfs-impl/impl-kernel/src/lib.rs}\\
\noindent 行范围：21\\[0.4em]

\begin{rustcode}
...
/// 零大小 [`RootFsManager`] 句柄；读写 `static ROOT_FS` 等状态。
// 本结构代码由AI完成
pub struct KernelRootFsManager;

// 本变量代码由AI完成
...
\end{rustcode}

\subsubsection*{序号233}
\noindent 路径：\code{os/components/wateros-fs/fs-rootfs/rootfs-impl/impl-kernel/src/lib.rs}\\
\noindent 行范围：24\\[0.4em]

\begin{rustcode}
...

// 本变量代码由AI完成
static ROOT_FS: Mutex<Option<fs_api_v0::SharedFs>> = Mutex::new(None);
// 本变量代码由AI完成
static ROOT_RW_FS: Mutex<Option<fs_api_v0::SharedRwFs>> = Mutex::new(None);
...
\end{rustcode}

\subsubsection*{序号234}
\noindent 路径：\code{os/components/wateros-fs/fs-rootfs/rootfs-impl/impl-kernel/src/lib.rs}\\
\noindent 行范围：26\\[0.4em]

\begin{rustcode}
...
static ROOT_FS: Mutex<Option<fs_api_v0::SharedFs>> = Mutex::new(None);
// 本变量代码由AI完成
static ROOT_RW_FS: Mutex<Option<fs_api_v0::SharedRwFs>> = Mutex::new(None);
// 本变量代码由AI完成
static ROOT_DEV_PATH: Mutex<Option<String>> = Mutex::new(None);
...
\end{rustcode}

\subsubsection*{序号235}
\noindent 路径：\code{os/components/wateros-fs/fs-rootfs/rootfs-impl/impl-kernel/src/lib.rs}\\
\noindent 行范围：28\\[0.4em]

\begin{rustcode}
...
static ROOT_RW_FS: Mutex<Option<fs_api_v0::SharedRwFs>> = Mutex::new(None);
// 本变量代码由AI完成
static ROOT_DEV_PATH: Mutex<Option<String>> = Mutex::new(None);
// 生命周期：'static 引用指向注册表中的 impl，由链接期/聚合 init 保证长于内核运行期。
/// 由聚合层在启动期注入的「当前根 FS impl」。聚合层根据 `probe` 选定后调用 [`set_active_fs_impl`]。
...
\end{rustcode}

\subsubsection*{序号236}
\noindent 路径：\code{os/components/wateros-fs/fs-rootfs/rootfs-impl/impl-kernel/src/lib.rs}\\
\noindent 行范围：32\\[0.4em]

\begin{rustcode}
...
/// 由聚合层在启动期注入的「当前根 FS impl」。聚合层根据 `probe` 选定后调用 [`set_active_fs_impl`]。
// 本变量代码由AI完成
static ACTIVE_FS_IMPL: Mutex<Option<&'static dyn FsImpl>> = Mutex::new(None);

impl api_v0::RootFsManager for KernelRootFsManager {
...
\end{rustcode}

\subsubsection*{序号237}
\noindent 路径：\code{os/components/wateros-fs/fs-rootfs/rootfs-impl/impl-kernel/src/lib.rs}\\
\noindent 行范围：36-38\\[0.4em]

\begin{rustcode}
...
// 本方法代码由AI完成
    fn set_root_fs(&mut self, fs: fs_api_v0::SharedFs) {
        *ROOT_FS.lock() = Some(fs);
    }

...
\end{rustcode}

\subsubsection*{序号238}
\noindent 路径：\code{os/components/wateros-fs/fs-rootfs/rootfs-impl/impl-kernel/src/lib.rs}\\
\noindent 行范围：45-49\\[0.4em]

\begin{rustcode}
...
    fn clear_root_fs(&mut self) {
        *ROOT_FS.lock() = None;
        *ROOT_RW_FS.lock() = None;
        *ROOT_DEV_PATH.lock() = None;
    }
...
\end{rustcode}

\subsubsection*{序号239}
\noindent 路径：\code{os/components/wateros-fs/fs-rootfs/rootfs-impl/impl-kernel/src/lib.rs}\\
\noindent 行范围：52-63\\[0.4em]

\begin{rustcode}
...
    fn mount_root_from_block_path(&mut self, path: &str) -> fs_api_v0::FsResult<()> {
        // devfs 解析路径 → 活动 FsImpl RO 挂载 → 记录路径，三步任一失败即向上返回。
        let device = devfs::active_impl::lookup_block_device(path)?;
        let imp = ACTIVE_FS_IMPL
            .lock()
...
\end{rustcode}

\subsubsection*{序号240}
\noindent 路径：\code{os/components/wateros-fs/fs-rootfs/rootfs-impl/impl-kernel/src/lib.rs}\\
\noindent 行范围：71-73\\[0.4em]

\begin{rustcode}
...
// 本方法代码由AI完成
pub fn set_active_fs_impl(imp: &'static dyn FsImpl) {
    *ACTIVE_FS_IMPL.lock() = Some(imp);
}

...
\end{rustcode}

\subsubsection*{序号241}
\noindent 路径：\code{os/components/wateros-fs/fs-rootfs/rootfs-impl/impl-kernel/src/lib.rs}\\
\noindent 行范围：81-88\\[0.4em]

\begin{rustcode}
...
pub fn mount_default_root() -> fs_api_v0::FsResult<()> {
    let Some(path) = devfs::active_impl::default_root_block_path() else {
        return Err(fs_api_v0::FsError::NotMounted);
    };
    logging::info!("[fs::rootfs] mount default root RO from {}", path);
...
\end{rustcode}

\subsubsection*{序号242}
\noindent 路径：\code{os/components/wateros-fs/fs-rootfs/rootfs-impl/impl-kernel/src/lib.rs}\\
\noindent 行范围：92-97\\[0.4em]

\begin{rustcode}
...
pub fn mount_default_root_rw() -> fs_api_v0::FsResult<()> {
    let Some(path) = devfs::active_impl::default_root_block_path() else {
        return Err(fs_api_v0::FsError::NotMounted);
    };
    mount_root_rw_from_block_path(path.as_str())
...
\end{rustcode}

\subsubsection*{序号243}
\noindent 路径：\code{os/components/wateros-fs/fs-rootfs/rootfs-impl/impl-kernel/src/lib.rs}\\
\noindent 行范围：101-112\\[0.4em]

\begin{rustcode}
...
pub fn mount_root_rw_from_block_path(path: &str) -> fs_api_v0::FsResult<()> {
    let device = devfs::active_impl::lookup_block_device(path)?;
    let imp = ACTIVE_FS_IMPL
        .lock()
        .ok_or(fs_api_v0::FsError::Unsupported)?;
...
\end{rustcode}

\subsubsection*{序号244}
\noindent 路径：\code{os/components/wateros-fs/fs-rootfs/rootfs-impl/impl-kernel/src/lib.rs}\\
\noindent 行范围：117-120\\[0.4em]

\begin{rustcode}
...
// 本方法代码由AI完成
pub fn root_fs() -> Option<fs_api_v0::SharedFs> {
    let mgr = KernelRootFsManager;
    mgr.root_fs()
}
...
\end{rustcode}

\subsubsection*{序号245}
\noindent 路径：\code{os/components/wateros-fs/fs-rootfs/rootfs-impl/impl-kernel/src/lib.rs}\\
\noindent 行范围：131-134\\[0.4em]

\begin{rustcode}
...
// 本方法代码由AI完成
pub fn current_root_device_path() -> Option<String> {
    let mgr = KernelRootFsManager;
    mgr.current_root_device_path()
}
...
\end{rustcode}

\subsubsection*{序号246}
\noindent 路径：\code{os/components/wateros-fs/fs-rootfs/rootfs-impl/impl-kernel/src/lib.rs}\\
\noindent 行范围：150-180\\[0.4em]

\begin{rustcode}
...
pub fn mount_aux_ro_from_block_path(path: &str) -> fs_api_v0::FsResult<fs_api_v0::SharedFs> {
    let device = devfs::active_impl::lookup_block_device(path)?;
    if let Some(root_path) = current_root_device_path() {
        if let Ok(root_dev) = devfs::active_impl::lookup_block_device(root_path.as_str()) {
            if Arc::ptr_eq(&device, &root_dev) {
...
\end{rustcode}

\subsubsection*{序号247}
\noindent 路径：\code{os/components/wateros-fs/fs-rootfs/rootfs-impl/impl-kernel/src/lib.rs}\\
\noindent 行范围：184-207\\[0.4em]

\begin{rustcode}
...
pub fn mount_aux_rw_from_block_path(path: &str) -> fs_api_v0::FsResult<fs_api_v0::SharedRwFs> {
    let device = devfs::active_impl::lookup_block_device(path)?;
    if let Some(root_path) = current_root_device_path() {
        if let Ok(root_dev) = devfs::active_impl::lookup_block_device(root_path.as_str()) {
            if Arc::ptr_eq(&device, &root_dev) {
...
\end{rustcode}

\subsubsection*{序号248}
\noindent 路径：\code{os/components/wateros-ipc/ipc-futex/futex-impl/impl-task/src/hub.rs}\\
\noindent 行范围：1-184\\[0.4em]

\begin{rustcode}
//! 基于 `ipc-waitqueue` 的全局 futex 枢纽。
//! 本模块代码由AI完成

extern crate alloc;

...
\end{rustcode}

\subsubsection*{序号249}
\noindent 路径：\code{os/components/wateros-ipc/ipc-futex/futex-impl/impl-task/src/hub.rs}\\
\noindent 行范围：16-19\\[0.4em]

\begin{rustcode}
...
// 本结构代码由AI完成
struct FutexTables {
    queues: BTreeMap<FutexKey, WaitQueue>,
    robust: BTreeMap<TaskId, RobustState>,
}
...
\end{rustcode}

\subsubsection*{序号250}
\noindent 路径：\code{os/components/wateros-ipc/ipc-futex/futex-impl/impl-task/src/hub.rs}\\
\noindent 行范围：23-25\\[0.4em]

\begin{rustcode}
...
// 本结构代码由AI完成
pub struct FutexHub {
    inner: Mutex<FutexTables>,
}

...
\end{rustcode}

\subsubsection*{序号251}
\noindent 路径：\code{os/components/wateros-ipc/ipc-futex/futex-impl/impl-task/src/hub.rs}\\
\noindent 行范围：35-38\\[0.4em]

\begin{rustcode}
...
// 本方法代码由AI完成
    fn with_tables<R>(&self, f: impl FnOnce(&mut FutexTables) -> R) -> R {
        let mut guard = self.inner.lock();
        f(&mut guard)
    }
...
\end{rustcode}

\subsubsection*{序号252}
\noindent 路径：\code{os/components/wateros-ipc/ipc-futex/futex-impl/impl-task/src/hub.rs}\\
\noindent 行范围：41-46\\[0.4em]

\begin{rustcode}
...
    fn get_queue(tables: &mut FutexTables, key: FutexKey) -> WaitQueue {
        *tables
            .queues
            .entry(key)
            .or_insert_with(WaitQueue::new)
...
\end{rustcode}

\subsubsection*{序号253}
\noindent 路径：\code{os/components/wateros-ipc/ipc-futex/futex-impl/impl-task/src/hub.rs}\\
\noindent 行范围：49-56\\[0.4em]

\begin{rustcode}
...
    fn cleanup_empty_queue(tables: &mut FutexTables, key: FutexKey) {
        let Some(wq) = tables.queues.get(&key).copied() else {
            return;
        };
        if wq.try_release_empty() {
...
\end{rustcode}

\subsubsection*{序号254}
\noindent 路径：\code{os/components/wateros-ipc/ipc-futex/futex-impl/impl-task/src/hub.rs}\\
\noindent 行范围：60-96\\[0.4em]

\begin{rustcode}
...
    pub fn wait_while(
        &self,
        key: FutexKey,
        timeout: Option<TaskTick>,
        mut condition: impl FnMut() -> bool,
...
\end{rustcode}

\subsubsection*{序号255}
\noindent 路径：\code{os/components/wateros-ipc/ipc-futex/futex-impl/impl-task/src/hub.rs}\\
\noindent 行范围：101-115\\[0.4em]

\begin{rustcode}
...
    pub fn requeue(&self,
                   from_key : FutexKey,
                   to_key : FutexKey,
                   wake_count : u32,
                   requeue_count : u32)
...
\end{rustcode}

\subsubsection*{序号256}
\noindent 路径：\code{os/components/wateros-ipc/ipc-futex/futex-impl/impl-task/src/hub.rs}\\
\noindent 行范围：119-124\\[0.4em]

\begin{rustcode}
...
static GLOBAL_HUB: FutexHub = FutexHub {
    inner: Mutex::new(FutexTables {
        queues: BTreeMap::new(),
        robust: BTreeMap::new(),
    }),
...
\end{rustcode}

\subsubsection*{序号257}
\noindent 路径：\code{os/components/wateros-ipc/ipc-futex/futex-impl/impl-task/src/hub.rs}\\
\noindent 行范围：128-143\\[0.4em]

\begin{rustcode}
...
    fn wake(&self, key: FutexKey, max_wake: u32) -> FutexResult<usize> {
        let wq = self.with_tables(|tables| tables.queues.get(&key).copied());
        let Some(wq) = wq else {
            return Ok(0);
        };
...
\end{rustcode}

\subsubsection*{序号258}
\noindent 路径：\code{os/components/wateros-ipc/ipc-futex/futex-impl/impl-task/src/hub.rs}\\
\noindent 行范围：146-151\\[0.4em]

\begin{rustcode}
...
    fn wake_all(&self, key: FutexKey) -> FutexResult<usize> {
        let wq = self.with_tables(|tables| tables.queues.get(&key).copied());
        let woken = wq.map(|wq| wq.wake_all()).unwrap_or(0);
        self.with_tables(|tables| Self::cleanup_empty_queue(tables, key));
        Ok(woken)
...
\end{rustcode}

\subsubsection*{序号259}
\noindent 路径：\code{os/components/wateros-ipc/ipc-futex/futex-impl/impl-task/src/hub.rs}\\
\noindent 行范围：154-165\\[0.4em]

\begin{rustcode}
...
    fn set_robust_list(&self, task: TaskId, head: usize, len: usize) -> FutexResult<()> {
        if len != ROBUST_LIST_HEAD_SIZE {
            return Err(FutexError::Invalid);
        }
        self.with_tables(|tables| {
...
\end{rustcode}

\subsubsection*{序号260}
\noindent 路径：\code{os/components/wateros-ipc/ipc-futex/futex-impl/impl-task/src/hub.rs}\\
\noindent 行范围：168-176\\[0.4em]

\begin{rustcode}
...
    fn get_robust_list(&self, task: TaskId) -> FutexResult<(usize, usize)> {
        Ok(self.with_tables(|tables| {
            tables
                .robust
                .get(&task)
...
\end{rustcode}

\subsubsection*{序号261}
\noindent 路径：\code{os/components/wateros-ipc/ipc-futex/futex-impl/impl-task/src/hub.rs}\\
\noindent 行范围：179-183\\[0.4em]

\begin{rustcode}
...
    fn drop_robust_list(&self, task: TaskId) {
        self.with_tables(|tables| {
            tables.robust.remove(&task);
        });
    }
...
\end{rustcode}

\subsubsection*{序号262}
\noindent 路径：\code{os/components/wateros-ipc/ipc-futex/futex-impl/impl-task/src/lib.rs}\\
\noindent 行范围：1-42\\[0.4em]

\begin{rustcode}
#![no_std]
//! Futex task 实现：全局 `FutexHub` + `ipc-waitqueue` 阻塞/唤醒 + robust 侧表。
//! 本模块代码由AI完成

extern crate alloc;
...
\end{rustcode}

\subsubsection*{序号263}
\noindent 路径：\code{os/components/wateros-ipc/ipc-futex/futex-impl/impl-task/src/lib.rs}\\
\noindent 行范围：14-32\\[0.4em]

\begin{rustcode}
...
pub fn test() {
    use api_v0::KernelFutexOps;

    api_v0::test();
    let hub = FutexHub::global();
...
\end{rustcode}

\subsubsection*{序号264}
\noindent 路径：\code{os/components/wateros-ipc/ipc-pipe/pipe-impl/impl-ringbuf/src/endpoint.rs}\\
\noindent 行范围：1-208\\[0.4em]

\begin{rustcode}
//! fd 表 pipe 端点实现。
//! 本模块代码由AI完成

extern crate alloc;

...
\end{rustcode}

\subsubsection*{序号265}
\noindent 路径：\code{os/components/wateros-ipc/ipc-pipe/pipe-impl/impl-ringbuf/src/endpoint.rs}\\
\noindent 行范围：15-19\\[0.4em]

\begin{rustcode}
...
pub struct PipeEndpoint {
    pipe: Arc<Pipe>,
    kind: PipeEndpointKind,
    nonblocking: Cell<bool>,
}
...
\end{rustcode}

\subsubsection*{序号266}
\noindent 路径：\code{os/components/wateros-ipc/ipc-pipe/pipe-impl/impl-ringbuf/src/endpoint.rs}\\
\noindent 行范围：23-33\\[0.4em]

\begin{rustcode}
...
    fn clone(&self) -> Self {
        match self.kind {
            PipeEndpointKind::Read => self.pipe.acquire_read(),
            PipeEndpointKind::Write => self.pipe.acquire_write(),
        }
...
\end{rustcode}

\subsubsection*{序号267}
\noindent 路径：\code{os/components/wateros-ipc/ipc-pipe/pipe-impl/impl-ringbuf/src/endpoint.rs}\\
\noindent 行范围：39-41\\[0.4em]

\begin{rustcode}
...
// 本方法代码由AI完成
    pub fn pair(nonblocking: bool) -> (Self, Self) {
        PipeEndpointOps::pair(nonblocking)
    }

...
\end{rustcode}

\subsubsection*{序号268}
\noindent 路径：\code{os/components/wateros-ipc/ipc-pipe/pipe-impl/impl-ringbuf/src/endpoint.rs}\\
\noindent 行范围：57-59\\[0.4em]

\begin{rustcode}
...
// 本方法代码由AI完成
    pub fn set_nonblocking(&self, value: bool) {
        self.nonblocking.set(value);
    }

...
\end{rustcode}

\subsubsection*{序号269}
\noindent 路径：\code{os/components/wateros-ipc/ipc-pipe/pipe-impl/impl-ringbuf/src/endpoint.rs}\\
\noindent 行范围：69-71\\[0.4em]

\begin{rustcode}
...
// 本方法代码由AI完成
    pub fn set_pipe_capacity(&self, capacity: usize) -> PipeResult<usize> {
        self.pipe.set_capacity(capacity)
    }

...
\end{rustcode}

\subsubsection*{序号270}
\noindent 路径：\code{os/components/wateros-ipc/ipc-pipe/pipe-impl/impl-ringbuf/src/endpoint.rs}\\
\noindent 行范围：75-77\\[0.4em]

\begin{rustcode}
...
// 本方法代码由AI完成
    pub fn read(&self, out: &mut [u8]) -> PipeResult<usize> {
        PipeEndpointOps::read(self, out)
    }

...
\end{rustcode}

\subsubsection*{序号271}
\noindent 路径：\code{os/components/wateros-ipc/ipc-pipe/pipe-impl/impl-ringbuf/src/endpoint.rs}\\
\noindent 行范围：81-83\\[0.4em]

\begin{rustcode}
...
// 本方法代码由AI完成
    pub fn write(&self, input: &[u8]) -> PipeResult<usize> {
        PipeEndpointOps::write(self, input)
    }

...
\end{rustcode}

\subsubsection*{序号272}
\noindent 路径：\code{os/components/wateros-ipc/ipc-pipe/pipe-impl/impl-ringbuf/src/endpoint.rs}\\
\noindent 行范围：87-89\\[0.4em]

\begin{rustcode}
...
// 本方法代码由AI完成
    pub fn close(&self) {
        PipeEndpointOps::close(self);
    }
}
...
\end{rustcode}

\subsubsection*{序号273}
\noindent 路径：\code{os/components/wateros-ipc/ipc-pipe/pipe-impl/impl-ringbuf/src/endpoint.rs}\\
\noindent 行范围：95-111\\[0.4em]

\begin{rustcode}
...
    fn pair(nonblocking: bool) -> (Self, Self) {
        let pipe = Arc::new(Pipe::new());
        pipe.acquire_read();
        pipe.acquire_write();
        (
...
\end{rustcode}

\subsubsection*{序号274}
\noindent 路径：\code{os/components/wateros-ipc/ipc-pipe/pipe-impl/impl-ringbuf/src/endpoint.rs}\\
\noindent 行范围：124-133\\[0.4em]

\begin{rustcode}
...
    fn read(&self, out: &mut [u8]) -> PipeResult<usize> {
        if self.kind != PipeEndpointKind::Read {
            return Err(PipeError::BrokenPipe);
        }
        if self.nonblocking.get() {
...
\end{rustcode}

\subsubsection*{序号275}
\noindent 路径：\code{os/components/wateros-ipc/ipc-pipe/pipe-impl/impl-ringbuf/src/endpoint.rs}\\
\noindent 行范围：136-145\\[0.4em]

\begin{rustcode}
...
    fn write(&self, input: &[u8]) -> PipeResult<usize> {
        if self.kind != PipeEndpointKind::Write {
            return Err(PipeError::BrokenPipe);
        }
        if self.nonblocking.get() {
...
\end{rustcode}

\subsubsection*{序号276}
\noindent 路径：\code{os/components/wateros-ipc/ipc-pipe/pipe-impl/impl-ringbuf/src/endpoint.rs}\\
\noindent 行范围：148-153\\[0.4em]

\begin{rustcode}
...
    fn close(&self) {
        match self.kind {
            PipeEndpointKind::Read => self.pipe.release_read(),
            PipeEndpointKind::Write => self.pipe.release_write(),
        }
...
\end{rustcode}

\subsubsection*{序号277}
\noindent 路径：\code{os/components/wateros-ipc/ipc-pipe/pipe-impl/impl-ringbuf/src/endpoint.rs}\\
\noindent 行范围：157-177\\[0.4em]

\begin{rustcode}
...
    fn poll_revents(&self, events: i16) -> PipeResult<i16> {
// 本变量代码由AI完成
        const POLLIN: i16 = 0x001;
// 本变量代码由AI完成
        const POLLOUT: i16 = 0x004;
...
\end{rustcode}

\subsubsection*{序号278}
\noindent 路径：\code{os/components/wateros-ipc/ipc-pipe/pipe-impl/impl-ringbuf/src/endpoint.rs}\\
\noindent 行范围：159\\[0.4em]

\begin{rustcode}
...
    fn poll_revents(&self, events: i16) -> PipeResult<i16> {
// 本变量代码由AI完成
        const POLLIN: i16 = 0x001;
// 本变量代码由AI完成
        const POLLOUT: i16 = 0x004;
...
\end{rustcode}

\subsubsection*{序号279}
\noindent 路径：\code{os/components/wateros-ipc/ipc-pipe/pipe-impl/impl-ringbuf/src/endpoint.rs}\\
\noindent 行范围：161\\[0.4em]

\begin{rustcode}
...
        const POLLIN: i16 = 0x001;
// 本变量代码由AI完成
        const POLLOUT: i16 = 0x004;
        let raw = match self.kind {
            PipeEndpointKind::Read => self.pipe.poll_revents_read(),
...
\end{rustcode}

\subsubsection*{序号280}
\noindent 路径：\code{os/components/wateros-ipc/ipc-pipe/pipe-impl/impl-ringbuf/src/endpoint.rs}\\
\noindent 行范围：180-207\\[0.4em]

\begin{rustcode}
...
    fn poll_wait_for_ticks(
        &self,
        events: i16,
        timeout_ticks: u64,
        still_waiting: &mut dyn FnMut() -> bool,
...
\end{rustcode}

\subsubsection*{序号281}
\noindent 路径：\code{os/components/wateros-ipc/ipc-pipe/pipe-impl/impl-ringbuf/src/endpoint.rs}\\
\noindent 行范围：187\\[0.4em]

\begin{rustcode}
...
    ) -> PipeResult<()> {
// 本变量代码由AI完成
        const POLLIN: i16 = 0x001;
// 本变量代码由AI完成
        const POLLOUT: i16 = 0x004;
...
\end{rustcode}

\subsubsection*{序号282}
\noindent 路径：\code{os/components/wateros-ipc/ipc-pipe/pipe-impl/impl-ringbuf/src/endpoint.rs}\\
\noindent 行范围：189\\[0.4em]

\begin{rustcode}
...
        const POLLIN: i16 = 0x001;
// 本变量代码由AI完成
        const POLLOUT: i16 = 0x004;
        if events & POLLIN != 0 && self.kind == PipeEndpointKind::Read {
            let result = self
...
\end{rustcode}

\subsubsection*{序号283}
\noindent 路径：\code{os/components/wateros-ipc/ipc-pipe/pipe-impl/impl-ringbuf/src/kernel_pipe.rs}\\
\noindent 行范围：1-438\\[0.4em]

\begin{rustcode}
//! 固定容量 ring buffer pipe 实现。
//! 本模块代码由AI完成

extern crate alloc;

...
\end{rustcode}

\subsubsection*{序号284}
\noindent 路径：\code{os/components/wateros-ipc/ipc-pipe/pipe-impl/impl-ringbuf/src/kernel_pipe.rs}\\
\noindent 行范围：13\\[0.4em]

\begin{rustcode}
...

// 本变量代码由AI完成
const POLLIN: i16 = 0x001;
// 本变量代码由AI完成
const POLLOUT: i16 = 0x004;
...
\end{rustcode}

\subsubsection*{序号285}
\noindent 路径：\code{os/components/wateros-ipc/ipc-pipe/pipe-impl/impl-ringbuf/src/kernel_pipe.rs}\\
\noindent 行范围：15\\[0.4em]

\begin{rustcode}
...
const POLLIN: i16 = 0x001;
// 本变量代码由AI完成
const POLLOUT: i16 = 0x004;
// 本变量代码由AI完成
const POLLHUP: i16 = 0x010;
...
\end{rustcode}

\subsubsection*{序号286}
\noindent 路径：\code{os/components/wateros-ipc/ipc-pipe/pipe-impl/impl-ringbuf/src/kernel_pipe.rs}\\
\noindent 行范围：17\\[0.4em]

\begin{rustcode}
...
const POLLOUT: i16 = 0x004;
// 本变量代码由AI完成
const POLLHUP: i16 = 0x010;
// 本变量代码由AI完成
const POLLERR: i16 = 0x008;
...
\end{rustcode}

\subsubsection*{序号287}
\noindent 路径：\code{os/components/wateros-ipc/ipc-pipe/pipe-impl/impl-ringbuf/src/kernel_pipe.rs}\\
\noindent 行范围：19\\[0.4em]

\begin{rustcode}
...
const POLLHUP: i16 = 0x010;
// 本变量代码由AI完成
const POLLERR: i16 = 0x008;

// 本结构代码由AI完成
...
\end{rustcode}

\subsubsection*{序号288}
\noindent 路径：\code{os/components/wateros-ipc/ipc-pipe/pipe-impl/impl-ringbuf/src/kernel_pipe.rs}\\
\noindent 行范围：22-32\\[0.4em]

\begin{rustcode}
...
struct PipeState {
    buf: Vec<u8>,
    head: usize,
    len: usize,
    read_open: bool,
...
\end{rustcode}

\subsubsection*{序号289}
\noindent 路径：\code{os/components/wateros-ipc/ipc-pipe/pipe-impl/impl-ringbuf/src/kernel_pipe.rs}\\
\noindent 行范围：36-49\\[0.4em]

\begin{rustcode}
...
    fn with_capacity(capacity: usize) -> PipeResult<Self> {
        if capacity == 0 {
            return Err(PipeError::InvalidCapacity);
        }
        Ok(Self {
...
\end{rustcode}

\subsubsection*{序号290}
\noindent 路径：\code{os/components/wateros-ipc/ipc-pipe/pipe-impl/impl-ringbuf/src/kernel_pipe.rs}\\
\noindent 行范围：72-87\\[0.4em]

\begin{rustcode}
...
    fn read_into(&mut self, out: &mut [u8]) -> usize {
        let count = out.len().min(self.len);
        if count == 0 {
            return 0;
        }
...
\end{rustcode}

\subsubsection*{序号291}
\noindent 路径：\code{os/components/wateros-ipc/ipc-pipe/pipe-impl/impl-ringbuf/src/kernel_pipe.rs}\\
\noindent 行范围：90-105\\[0.4em]

\begin{rustcode}
...
    fn write_from(&mut self, input: &[u8]) -> usize {
        let count = input.len().min(self.free_len());
        if count == 0 {
            return 0;
        }
...
\end{rustcode}

\subsubsection*{序号292}
\noindent 路径：\code{os/components/wateros-ipc/ipc-pipe/pipe-impl/impl-ringbuf/src/kernel_pipe.rs}\\
\noindent 行范围：110-114\\[0.4em]

\begin{rustcode}
...
pub struct Pipe {
    state: Mutex<PipeState>,
    read_wait: WaitQueue,
    write_wait: WaitQueue,
}
...
\end{rustcode}

\subsubsection*{序号293}
\noindent 路径：\code{os/components/wateros-ipc/ipc-pipe/pipe-impl/impl-ringbuf/src/kernel_pipe.rs}\\
\noindent 行范围：125-127\\[0.4em]

\begin{rustcode}
...
// 本方法代码由AI完成
    pub fn with_capacity(capacity: usize) -> PipeResult<Self> {
        KernelPipe::with_capacity(capacity)
    }

...
\end{rustcode}

\subsubsection*{序号294}
\noindent 路径：\code{os/components/wateros-ipc/ipc-pipe/pipe-impl/impl-ringbuf/src/kernel_pipe.rs}\\
\noindent 行范围：137-139\\[0.4em]

\begin{rustcode}
...
// 本方法代码由AI完成
    pub fn set_capacity(&self, capacity: usize) -> PipeResult<usize> {
        KernelPipe::set_capacity(self, capacity)
    }

...
\end{rustcode}

\subsubsection*{序号295}
\noindent 路径：\code{os/components/wateros-ipc/ipc-pipe/pipe-impl/impl-ringbuf/src/kernel_pipe.rs}\\
\noindent 行范围：149-151\\[0.4em]

\begin{rustcode}
...
// 本方法代码由AI完成
    pub fn try_read(&self, out: &mut [u8]) -> PipeResult<usize> {
        KernelPipe::try_read(self, out)
    }

...
\end{rustcode}

\subsubsection*{序号296}
\noindent 路径：\code{os/components/wateros-ipc/ipc-pipe/pipe-impl/impl-ringbuf/src/kernel_pipe.rs}\\
\noindent 行范围：155-157\\[0.4em]

\begin{rustcode}
...
// 本方法代码由AI完成
    pub fn read(&self, out: &mut [u8]) -> PipeResult<usize> {
        KernelPipe::read(self, out)
    }

...
\end{rustcode}

\subsubsection*{序号297}
\noindent 路径：\code{os/components/wateros-ipc/ipc-pipe/pipe-impl/impl-ringbuf/src/kernel_pipe.rs}\\
\noindent 行范围：161-163\\[0.4em]

\begin{rustcode}
...
// 本方法代码由AI完成
    pub fn try_write(&self, input: &[u8]) -> PipeResult<usize> {
        KernelPipe::try_write(self, input)
    }

...
\end{rustcode}

\subsubsection*{序号298}
\noindent 路径：\code{os/components/wateros-ipc/ipc-pipe/pipe-impl/impl-ringbuf/src/kernel_pipe.rs}\\
\noindent 行范围：167-169\\[0.4em]

\begin{rustcode}
...
// 本方法代码由AI完成
    pub fn write(&self, input: &[u8]) -> PipeResult<usize> {
        KernelPipe::write(self, input)
    }

...
\end{rustcode}

\subsubsection*{序号299}
\noindent 路径：\code{os/components/wateros-ipc/ipc-pipe/pipe-impl/impl-ringbuf/src/kernel_pipe.rs}\\
\noindent 行范围：173-175\\[0.4em]

\begin{rustcode}
...
// 本方法代码由AI完成
    pub fn close_read(&self) {
        KernelPipe::close_read(self);
    }

...
\end{rustcode}

\subsubsection*{序号300}
\noindent 路径：\code{os/components/wateros-ipc/ipc-pipe/pipe-impl/impl-ringbuf/src/kernel_pipe.rs}\\
\noindent 行范围：179-181\\[0.4em]

\begin{rustcode}
...
// 本方法代码由AI完成
    pub fn close_write(&self) {
        KernelPipe::close_write(self);
    }

...
\end{rustcode}

\subsubsection*{序号301}
\noindent 路径：\code{os/components/wateros-ipc/ipc-pipe/pipe-impl/impl-ringbuf/src/kernel_pipe.rs}\\
\noindent 行范围：185-187\\[0.4em]

\begin{rustcode}
...
// 本方法代码由AI完成
    pub fn acquire_read(&self) {
        self.state.lock().read_refs += 1;
    }

...
\end{rustcode}

\subsubsection*{序号302}
\noindent 路径：\code{os/components/wateros-ipc/ipc-pipe/pipe-impl/impl-ringbuf/src/kernel_pipe.rs}\\
\noindent 行范围：191-193\\[0.4em]

\begin{rustcode}
...
// 本方法代码由AI完成
    pub fn acquire_write(&self) {
        self.state.lock().write_refs += 1;
    }

...
\end{rustcode}

\subsubsection*{序号303}
\noindent 路径：\code{os/components/wateros-ipc/ipc-pipe/pipe-impl/impl-ringbuf/src/kernel_pipe.rs}\\
\noindent 行范围：197-207\\[0.4em]

\begin{rustcode}
...
    pub fn release_read(&self) {
        let mut state = self.state.lock();
        if state.read_refs > 0 {
            state.read_refs -= 1;
        }
...
\end{rustcode}

\subsubsection*{序号304}
\noindent 路径：\code{os/components/wateros-ipc/ipc-pipe/pipe-impl/impl-ringbuf/src/kernel_pipe.rs}\\
\noindent 行范围：211-221\\[0.4em]

\begin{rustcode}
...
    pub fn release_write(&self) {
        let mut state = self.state.lock();
        if state.write_refs > 0 {
            state.write_refs -= 1;
        }
...
\end{rustcode}

\subsubsection*{序号305}
\noindent 路径：\code{os/components/wateros-ipc/ipc-pipe/pipe-impl/impl-ringbuf/src/kernel_pipe.rs}\\
\noindent 行范围：225-238\\[0.4em]

\begin{rustcode}
...
    pub fn poll_revents_read(&self) -> i16 {
        let state = self.state.lock();
        if !state.read_open {
            return POLLHUP;
        }
...
\end{rustcode}

\subsubsection*{序号306}
\noindent 路径：\code{os/components/wateros-ipc/ipc-pipe/pipe-impl/impl-ringbuf/src/kernel_pipe.rs}\\
\noindent 行范围：242-255\\[0.4em]

\begin{rustcode}
...
    pub fn poll_revents_write(&self) -> i16 {
        let state = self.state.lock();
        if !state.write_open {
            return POLLHUP | POLLERR;
        }
...
\end{rustcode}

\subsubsection*{序号307}
\noindent 路径：\code{os/components/wateros-ipc/ipc-pipe/pipe-impl/impl-ringbuf/src/kernel_pipe.rs}\\
\noindent 行范围：259-267\\[0.4em]

\begin{rustcode}
...
    pub fn poll_wait_read_for_ticks(
        &self,
        timeout_ticks: u64,
        still_waiting: &mut dyn FnMut() -> bool,
    ) -> TaskWaitResult {
...
\end{rustcode}

\subsubsection*{序号308}
\noindent 路径：\code{os/components/wateros-ipc/ipc-pipe/pipe-impl/impl-ringbuf/src/kernel_pipe.rs}\\
\noindent 行范围：271-279\\[0.4em]

\begin{rustcode}
...
    pub fn poll_wait_write_for_ticks(
        &self,
        timeout_ticks: u64,
        still_waiting: &mut dyn FnMut() -> bool,
    ) -> TaskWaitResult {
...
\end{rustcode}

\subsubsection*{序号309}
\noindent 路径：\code{os/components/wateros-ipc/ipc-pipe/pipe-impl/impl-ringbuf/src/kernel_pipe.rs}\\
\noindent 行范围：282-285\\[0.4em]

\begin{rustcode}
...
// 本方法代码由AI完成
    fn read_poll_blocked(&self) -> bool {
        let state = self.state.lock();
        state.is_empty() && state.write_open
    }
...
\end{rustcode}

\subsubsection*{序号310}
\noindent 路径：\code{os/components/wateros-ipc/ipc-pipe/pipe-impl/impl-ringbuf/src/kernel_pipe.rs}\\
\noindent 行范围：288-291\\[0.4em]

\begin{rustcode}
...
// 本方法代码由AI完成
    fn write_poll_blocked(&self) -> bool {
        let state = self.state.lock();
        state.is_full() && state.read_open
    }
...
\end{rustcode}

\subsubsection*{序号311}
\noindent 路径：\code{os/components/wateros-ipc/ipc-pipe/pipe-impl/impl-ringbuf/src/kernel_pipe.rs}\\
\noindent 行范围：303-309\\[0.4em]

\begin{rustcode}
...
    fn with_capacity(capacity: usize) -> PipeResult<Self> {
        Ok(Self {
            state: Mutex::new(PipeState::with_capacity(capacity)?),
            read_wait: WaitQueue::new(),
            write_wait: WaitQueue::new(),
...
\end{rustcode}

\subsubsection*{序号312}
\noindent 路径：\code{os/components/wateros-ipc/ipc-pipe/pipe-impl/impl-ringbuf/src/kernel_pipe.rs}\\
\noindent 行范围：312-314\\[0.4em]

\begin{rustcode}
...
// 本方法代码由AI完成
    fn capacity(&self) -> usize {
        self.state.lock().capacity()
    }

...
\end{rustcode}

\subsubsection*{序号313}
\noindent 路径：\code{os/components/wateros-ipc/ipc-pipe/pipe-impl/impl-ringbuf/src/kernel_pipe.rs}\\
\noindent 行范围：317-328\\[0.4em]

\begin{rustcode}
...
    fn set_capacity(&self, capacity: usize) -> PipeResult<usize> {
        if capacity == 0 {
            return Err(PipeError::InvalidCapacity);
        }
        let mut state = self.state.lock();
...
\end{rustcode}

\subsubsection*{序号314}
\noindent 路径：\code{os/components/wateros-ipc/ipc-pipe/pipe-impl/impl-ringbuf/src/kernel_pipe.rs}\\
\noindent 行范围：331-333\\[0.4em]

\begin{rustcode}
...
// 本方法代码由AI完成
    fn len(&self) -> usize {
        self.state.lock().len
    }

...
\end{rustcode}

\subsubsection*{序号315}
\noindent 路径：\code{os/components/wateros-ipc/ipc-pipe/pipe-impl/impl-ringbuf/src/kernel_pipe.rs}\\
\noindent 行范围：336-352\\[0.4em]

\begin{rustcode}
...
    fn try_read(&self, out: &mut [u8]) -> PipeResult<usize> {
        if out.is_empty() {
            return Ok(0);
        }
        let mut state = self.state.lock();
...
\end{rustcode}

\subsubsection*{序号316}
\noindent 路径：\code{os/components/wateros-ipc/ipc-pipe/pipe-impl/impl-ringbuf/src/kernel_pipe.rs}\\
\noindent 行范围：355-370\\[0.4em]

\begin{rustcode}
...
    fn read(&self, out: &mut [u8]) -> PipeResult<usize> {
        loop {
            match self.try_read(out) {
                Err(PipeError::WouldBlock) => {
                    let result = self.read_wait.wait_current_while(|| {
...
\end{rustcode}

\subsubsection*{序号317}
\noindent 路径：\code{os/components/wateros-ipc/ipc-pipe/pipe-impl/impl-ringbuf/src/kernel_pipe.rs}\\
\noindent 行范围：373-388\\[0.4em]

\begin{rustcode}
...
    fn try_write(&self, input: &[u8]) -> PipeResult<usize> {
        if input.is_empty() {
            return Ok(0);
        }
        let mut state = self.state.lock();
...
\end{rustcode}

\subsubsection*{序号318}
\noindent 路径：\code{os/components/wateros-ipc/ipc-pipe/pipe-impl/impl-ringbuf/src/kernel_pipe.rs}\\
\noindent 行范围：391-415\\[0.4em]

\begin{rustcode}
...
    fn write(&self, input: &[u8]) -> PipeResult<usize> {
        let mut written = 0usize;
        while written < input.len() {
            match self.try_write(&input[written..]) {
                Ok(0) => break,
...
\end{rustcode}

\subsubsection*{序号319}
\noindent 路径：\code{os/components/wateros-ipc/ipc-pipe/pipe-impl/impl-ringbuf/src/kernel_pipe.rs}\\
\noindent 行范围：418-421\\[0.4em]

\begin{rustcode}
...
// 本方法代码由AI完成
    fn close_read(&self) {
        self.state.lock().read_open = false;
        self.write_wait.wake_all();
    }
...
\end{rustcode}

\subsubsection*{序号320}
\noindent 路径：\code{os/components/wateros-ipc/ipc-pipe/pipe-impl/impl-ringbuf/src/kernel_pipe.rs}\\
\noindent 行范围：424-427\\[0.4em]

\begin{rustcode}
...
// 本方法代码由AI完成
    fn close_write(&self) {
        self.state.lock().write_open = false;
        self.read_wait.wake_all();
    }
...
\end{rustcode}

\subsubsection*{序号321}
\noindent 路径：\code{os/components/wateros-ipc/ipc-pipe/pipe-impl/impl-ringbuf/src/kernel_pipe.rs}\\
\noindent 行范围：430-432\\[0.4em]

\begin{rustcode}
...
// 本方法代码由AI完成
    fn poll_revents_read(&self) -> i16 {
        Pipe::poll_revents_read(self)
    }

...
\end{rustcode}

\subsubsection*{序号322}
\noindent 路径：\code{os/components/wateros-ipc/ipc-pipe/pipe-impl/impl-ringbuf/src/kernel_pipe.rs}\\
\noindent 行范围：435-437\\[0.4em]

\begin{rustcode}
...
// 本方法代码由AI完成
    fn poll_revents_write(&self) -> i16 {
        Pipe::poll_revents_write(self)
    }
}
\end{rustcode}

\subsubsection*{序号323}
\noindent 路径：\code{os/components/wateros-ipc/ipc-pipe/pipe-impl/impl-ringbuf/src/lib.rs}\\
\noindent 行范围：1-75\\[0.4em]

\begin{rustcode}
#![no_std]
//! 管道 v0 ring buffer 实现：固定容量 ring buffer + task waitqueue。
//!
//! 行为为可用的内核内部 pipe，实现 [`api_v0::KernelPipe`] 与 [`api_v0::PipeEndpointOps`]。
//! 本模块代码由AI完成
...
\end{rustcode}

\subsubsection*{序号324}
\noindent 路径：\code{os/components/wateros-ipc/ipc-pipe/pipe-impl/impl-ringbuf/src/lib.rs}\\
\noindent 行范围：15-75\\[0.4em]

\begin{rustcode}
...
pub fn test() {
    let pipe = Pipe::with_capacity(8).expect("pipe capacity should be valid");
    assert_eq!(pipe.capacity(), 8);
    assert_eq!(pipe.len(), 0);

...
\end{rustcode}

\subsubsection*{序号325}
\noindent 路径：\code{os/components/wateros-ipc/ipc-waitqueue/waitqueue-impl/impl-task/src/lib.rs}\\
\noindent 行范围：1-187\\[0.4em]

\begin{rustcode}
#![no_std]
//! IPC 视角下的任务等待队列：对 `wateros_task::WaitQueue` 的薄包装，便于 IPC
//! 模块与调度子系统解耦命名。
//!
//! 不变量：不在此类型上附加第二套等待状态；唤醒与 tick 超时语义与
...
\end{rustcode}

\subsubsection*{序号326}
\noindent 路径：\code{os/components/wateros-ipc/ipc-waitqueue/waitqueue-impl/impl-task/src/lib.rs}\\
\noindent 行范围：26-29\\[0.4em]

\begin{rustcode}
...
// 本结构代码由AI完成
pub struct WaitQueue {
    // 仅委托底层队列；不在 IPC 层缓存 waiters 或附加锁。
    inner : TaskWaitQueue,
}
...
\end{rustcode}

\subsubsection*{序号327}
\noindent 路径：\code{os/components/wateros-ipc/ipc-waitqueue/waitqueue-impl/impl-task/src/lib.rs}\\
\noindent 行范围：43-46\\[0.4em]

\begin{rustcode}
...
// 本方法代码由AI完成
    pub fn try_release_empty(&self) -> bool {
        self.inner
            .try_release_empty()
    }
...
\end{rustcode}

\subsubsection*{序号328}
\noindent 路径：\code{os/components/wateros-ipc/ipc-waitqueue/waitqueue-impl/impl-task/src/lib.rs}\\
\noindent 行范围：58-61\\[0.4em]

\begin{rustcode}
...
// 本方法代码由AI完成
    pub fn wait_current(&self) -> TaskWaitResult {
        self.inner
            .wait_current()
    }
...
\end{rustcode}

\subsubsection*{序号329}
\noindent 路径：\code{os/components/wateros-ipc/ipc-waitqueue/waitqueue-impl/impl-task/src/lib.rs}\\
\noindent 行范围：66-69\\[0.4em]

\begin{rustcode}
...
// 本方法代码由AI完成
    pub fn wait_current_for_ticks(&self, timeout_ticks : TaskTick) -> TaskWaitResult {
        self.inner
            .wait_current_for_ticks(timeout_ticks)
    }
...
\end{rustcode}

\subsubsection*{序号330}
\noindent 路径：\code{os/components/wateros-ipc/ipc-waitqueue/waitqueue-impl/impl-task/src/lib.rs}\\
\noindent 行范围：74-79\\[0.4em]

\begin{rustcode}
...
    pub fn wait_current_while(&self,
                              condition : impl FnOnce() -> bool)
                              -> TaskWaitResult {
        self.inner
            .wait_current_while(condition)
...
\end{rustcode}

\subsubsection*{序号331}
\noindent 路径：\code{os/components/wateros-ipc/ipc-waitqueue/waitqueue-impl/impl-task/src/lib.rs}\\
\noindent 行范围：84-90\\[0.4em]

\begin{rustcode}
...
    pub fn wait_current_while_for_ticks(&self,
                                        timeout_ticks : TaskTick,
                                        condition : impl FnOnce() -> bool)
                                        -> TaskWaitResult {
        self.inner
...
\end{rustcode}

\subsubsection*{序号332}
\noindent 路径：\code{os/components/wateros-ipc/ipc-waitqueue/waitqueue-impl/impl-task/src/lib.rs}\\
\noindent 行范围：95-98\\[0.4em]

\begin{rustcode}
...
// 本方法代码由AI完成
    pub fn wake_one(&self) -> Option<TaskId> {
        self.inner
            .wake_one()
    }
...
\end{rustcode}

\subsubsection*{序号333}
\noindent 路径：\code{os/components/wateros-ipc/ipc-waitqueue/waitqueue-impl/impl-task/src/lib.rs}\\
\noindent 行范围：103-106\\[0.4em]

\begin{rustcode}
...
// 本方法代码由AI完成
    pub fn wake_all(&self) -> usize {
        self.inner
            .wake_all()
    }
...
\end{rustcode}

\subsubsection*{序号334}
\noindent 路径：\code{os/components/wateros-ipc/ipc-waitqueue/waitqueue-impl/impl-task/src/lib.rs}\\
\noindent 行范围：111-118\\[0.4em]

\begin{rustcode}
...
    pub fn requeue_to(&self,
                      target : Self,
                      wake_count : usize,
                      requeue_count : usize)
                      -> usize {
...
\end{rustcode}

\subsubsection*{序号335}
\noindent 路径：\code{os/components/wateros-ipc/ipc-waitqueue/waitqueue-impl/impl-task/src/lib.rs}\\
\noindent 行范围：135-138\\[0.4em]

\begin{rustcode}
...
// 本方法代码由AI完成
    fn wait_handle(&self) -> TaskWaitHandle {
        self.inner
            .wait_handle()
    }
...
\end{rustcode}

\subsubsection*{序号336}
\noindent 路径：\code{os/components/wateros-ipc/ipc-waitqueue/waitqueue-impl/impl-task/src/lib.rs}\\
\noindent 行范围：142-145\\[0.4em]

\begin{rustcode}
...
// 本方法代码由AI完成
    fn wait_current(&self) -> TaskWaitResult {
        self.inner
            .wait_current()
    }
...
\end{rustcode}

\subsubsection*{序号337}
\noindent 路径：\code{os/components/wateros-ipc/ipc-waitqueue/waitqueue-impl/impl-task/src/lib.rs}\\
\noindent 行范围：149-152\\[0.4em]

\begin{rustcode}
...
// 本方法代码由AI完成
    fn wait_current_for_ticks(&self, timeout_ticks : TaskTick) -> TaskWaitResult {
        self.inner
            .wait_current_for_ticks(timeout_ticks)
    }
...
\end{rustcode}

\subsubsection*{序号338}
\noindent 路径：\code{os/components/wateros-ipc/ipc-waitqueue/waitqueue-impl/impl-task/src/lib.rs}\\
\noindent 行范围：156-160\\[0.4em]

\begin{rustcode}
...
    fn wait_current_while<F>(&self, condition : F) -> TaskWaitResult
        where F : FnOnce() -> bool {
        self.inner
            .wait_current_while(condition)
    }
...
\end{rustcode}

\subsubsection*{序号339}
\noindent 路径：\code{os/components/wateros-ipc/ipc-waitqueue/waitqueue-impl/impl-task/src/lib.rs}\\
\noindent 行范围：164-172\\[0.4em]

\begin{rustcode}
...
    fn wait_current_while_for_ticks<F>(&self,
                                       timeout_ticks : TaskTick,
                                       condition : F)
                                       -> TaskWaitResult
        where F : FnOnce() -> bool
...
\end{rustcode}

\subsubsection*{序号340}
\noindent 路径：\code{os/components/wateros-ipc/ipc-waitqueue/waitqueue-impl/impl-task/src/lib.rs}\\
\noindent 行范围：176-179\\[0.4em]

\begin{rustcode}
...
// 本方法代码由AI完成
    fn wake_one(&self) -> Option<TaskId> {
        self.inner
            .wake_one()
    }
...
\end{rustcode}

\subsubsection*{序号341}
\noindent 路径：\code{os/components/wateros-ipc/ipc-waitqueue/waitqueue-impl/impl-task/src/lib.rs}\\
\noindent 行范围：183-186\\[0.4em]

\begin{rustcode}
...
// 本方法代码由AI完成
    fn wake_all(&self) -> usize {
        self.inner
            .wake_all()
    }
...
\end{rustcode}

\subsubsection*{序号342}
\noindent 路径：\code{os/components/wateros-klog/klog-api/api-v0/src/action.rs}\\
\noindent 行范围：1-52\\[0.4em]

\begin{rustcode}
//! Linux `syslog(2)` / `klogctl` action 常量（与 man page 对齐）。
//! 本模块代码由AI完成

/// `SYSLOG_ACTION_CLOSE`
pub const SYSLOG_ACTION_CLOSE: i32 = 0;
...
\end{rustcode}

\subsubsection*{序号343}
\noindent 路径：\code{os/components/wateros-klog/klog-api/api-v0/src/action.rs}\\
\noindent 行范围：30-45\\[0.4em]

\begin{rustcode}
...
pub fn decode_action(raw: i32) -> Option<i32> {
    match raw {
        SYSLOG_ACTION_CLOSE
        | SYSLOG_ACTION_OPEN
        | SYSLOG_ACTION_READ
...
\end{rustcode}

\subsubsection*{序号344}
\noindent 路径：\code{os/components/wateros-klog/klog-api/api-v0/src/action.rs}\\
\noindent 行范围：50-52\\[0.4em]

\begin{rustcode}
...
// 本方法代码由AI完成
#[inline]
pub fn is_write_priority(raw: i32) -> bool {
    decode_action(raw).is_none() && raw != 0
}
\end{rustcode}

\subsubsection*{序号345}
\noindent 路径：\code{os/components/wateros-klog/klog-api/api-v0/src/error.rs}\\
\noindent 行范围：1-22\\[0.4em]

\begin{rustcode}
//! 环操作错误（内核侧；syscall 层映射为 errno 或 panic）。
//! 本模块代码由AI完成

/// `KlogStore` 操作结果。
// 本结构代码由AI完成
...
\end{rustcode}

\subsubsection*{序号346}
\noindent 路径：\code{os/components/wateros-klog/klog-api/api-v0/src/error.rs}\\
\noindent 行范围：7-12\\[0.4em]

\begin{rustcode}
...
pub enum KlogError {
    /// 无未读记录（READ 路径返回 0，非错误）。
    NoUnread,
    /// 用户/调用方缓冲过小。
    BufferTooSmall,
...
\end{rustcode}

\subsubsection*{序号347}
\noindent 路径：\code{os/components/wateros-klog/klog-api/api-v0/src/error.rs}\\
\noindent 行范围：17-22\\[0.4em]

\begin{rustcode}
...
pub struct AppendResult {
    /// 分配到的序号。
    pub seq: u64,
    /// 是否发生截断。
    pub truncated: bool,
...
\end{rustcode}

\subsubsection*{序号348}
\noindent 路径：\code{os/components/wateros-klog/klog-api/api-v0/src/flags.rs}\\
\noindent 行范围：1-34\\[0.4em]

\begin{rustcode}
//! 记录标志位（内核私有语义，不导出给用户态裸结构）。
//! 本模块代码由AI完成

/// [`KlogRecordMeta::flags`] 位域。
// 本结构代码由AI完成
...
\end{rustcode}

\subsubsection*{序号349}
\noindent 路径：\code{os/components/wateros-klog/klog-api/api-v0/src/flags.rs}\\
\noindent 行范围：7\\[0.4em]

\begin{rustcode}
...
// 本结构代码由AI完成
#[derive(Clone, Copy, Debug, PartialEq, Eq, Default)]
pub struct KlogFlags(pub u8);

impl KlogFlags {
...
\end{rustcode}

\subsubsection*{序号350}
\noindent 路径：\code{os/components/wateros-klog/klog-api/api-v0/src/lib.rs}\\
\noindent 行范围：1-17\\[0.4em]

\begin{rustcode}
#![no_std]
//! `wateros-klog` API v0：记录元数据、syslog action 常量与环存储 trait。
//!
//! 本 crate 仅定义契约，不包含全局存储或平台时间源。
//! 本模块代码由AI完成
...
\end{rustcode}

\subsubsection*{序号351}
\noindent 路径：\code{os/components/wateros-klog/klog-api/api-v0/src/meta.rs}\\
\noindent 行范围：1-84\\[0.4em]

\begin{rustcode}
//! 单条内核消息的记录头（固定布局，供内核观测；不 `copy_to_user` 裸导出）。
//! 本模块代码由AI完成

use crate::KlogFlags;

...
\end{rustcode}

\subsubsection*{序号352}
\noindent 路径：\code{os/components/wateros-klog/klog-api/api-v0/src/meta.rs}\\
\noindent 行范围：32-47\\[0.4em]

\begin{rustcode}
...
pub struct KlogRecordMeta {
    /// 提交后单调递增序号。
    pub seq: u64,
    /// 单调时钟纳秒（由聚合层填入）。
    pub ts_nsec: u64,
...
\end{rustcode}

\subsubsection*{序号353}
\noindent 路径：\code{os/components/wateros-klog/klog-api/api-v0/src/meta.rs}\\
\noindent 行范围：53-70\\[0.4em]

\begin{rustcode}
...
    pub const fn new(
        ts_nsec: u64,
        text_len: u16,
        facility: u8,
        level: u8,
...
\end{rustcode}

\subsubsection*{序号354}
\noindent 路径：\code{os/components/wateros-klog/klog-api/api-v0/src/store.rs}\\
\noindent 行范围：1-54\\[0.4em]

\begin{rustcode}
//! 内核消息环存储 trait。
//! 本模块代码由AI完成

use crate::{AppendResult, KlogError, KlogRecordMeta};

...
\end{rustcode}

\subsubsection*{序号355}
\noindent 路径：\code{os/components/wateros-klog/klog-api/api-v0/src/store.rs}\\
\noindent 行范围：8-13\\[0.4em]

\begin{rustcode}
...
pub struct KlogRecordView<'a> {
    /// 记录头。
    pub meta: KlogRecordMeta,
    /// 正文。
    pub text: &'a [u8],
...
\end{rustcode}

\subsubsection*{序号356}
\noindent 路径：\code{os/components/wateros-klog/klog-api/api-v0/src/store.rs}\\
\noindent 行范围：17-38\\[0.4em]

\begin{rustcode}
...
pub trait KlogStore {
    /// 追加一条记录；`meta` 中 `seq` 由实现写入。
    fn append(&mut self, meta: &mut KlogRecordMeta, text: &[u8]) -> AppendResult;

    /// 统计快照。
...
\end{rustcode}

\subsubsection*{序号357}
\noindent 路径：\code{os/components/wateros-klog/klog-api/api-v0/src/store.rs}\\
\noindent 行范围：43-54\\[0.4em]

\begin{rustcode}
...
pub struct KlogStats {
    /// 成功提交条数。
    pub records_committed: u64,
    /// 因槽满覆盖丢弃的条数。
    pub records_dropped: u64,
...
\end{rustcode}

\subsubsection*{序号358}
\noindent 路径：\code{os/components/wateros-klog/klog-impl/klog-ringbuf/src/lib.rs}\\
\noindent 行范围：1-347\\[0.4em]

\begin{rustcode}
#![no_std]
//! desc 槽 + 每槽变长正文（上限 `KLOG_MAX_RECORD_BYTES`）的环形 klog 实现。
//! 本模块代码由AI完成

use api_v0::{
...
\end{rustcode}

\subsubsection*{序号359}
\noindent 路径：\code{os/components/wateros-klog/klog-impl/klog-ringbuf/src/lib.rs}\\
\noindent 行范围：18-24\\[0.4em]

\begin{rustcode}
...
struct Slot {
    /// 槽是否持有有效记录（被覆盖后置 false）。
    valid: bool,
    meta: KlogRecordMeta,
    /// 单条正文上限 `KLOG_MAX_RECORD_BYTES`。
...
\end{rustcode}

\subsubsection*{序号360}
\noindent 路径：\code{os/components/wateros-klog/klog-impl/klog-ringbuf/src/lib.rs}\\
\noindent 行范围：28-37\\[0.4em]

\begin{rustcode}
...
pub struct KlogRingbufInner {
    slots: [Slot; KLOG_DESC_SLOTS],
    head: usize,
    count: usize,
    next_seq: u64,
...
\end{rustcode}

\subsubsection*{序号361}
\noindent 路径：\code{os/components/wateros-klog/klog-impl/klog-ringbuf/src/lib.rs}\\
\noindent 行范围：118-141\\[0.4em]

\begin{rustcode}
...
    pub fn iter_from<F>(&self, start_seq: u64, f: &mut F)
    where
        F: FnMut(KlogRecordView<'_>),
    {
        let mut buf = [0u64; KLOG_DESC_SLOTS];
...
\end{rustcode}

\subsubsection*{序号362}
\noindent 路径：\code{os/components/wateros-klog/klog-impl/klog-ringbuf/src/lib.rs}\\
\noindent 行范围：147-187\\[0.4em]

\begin{rustcode}
...
    fn append(&mut self, meta: &mut KlogRecordMeta, text: &[u8]) -> AppendResult {
        let mut flags = KlogFlags(meta.flags);
        let copy_len = text.len().min(KLOG_MAX_RECORD_BYTES);
        let truncated = text.len() > copy_len;
        if truncated {
...
\end{rustcode}

\subsubsection*{序号363}
\noindent 路径：\code{os/components/wateros-klog/klog-impl/klog-ringbuf/src/lib.rs}\\
\noindent 行范围：224-245\\[0.4em]

\begin{rustcode}
...
    fn peek_next_unread(&self) -> Result<KlogRecordView<'_>, KlogError> {
        let mut min_unread: Option<u64> = None;
        self.for_each_valid_seq(|seq| {
            if seq >= self.read_cursor_seq {
                min_unread = Some(match min_unread {
...
\end{rustcode}

\subsubsection*{序号364}
\noindent 路径：\code{os/components/wateros-klog/klog-impl/klog-ringbuf/src/lib.rs}\\
\noindent 行范围：262\\[0.4em]

\begin{rustcode}
...
/// 全局内核消息环。
// 本结构代码由AI完成
pub struct KlogRingbuf;

// 本变量代码由AI完成
...
\end{rustcode}

\subsubsection*{序号365}
\noindent 路径：\code{os/components/wateros-klog/klog-impl/klog-ringbuf/src/lib.rs}\\
\noindent 行范围：265\\[0.4em]

\begin{rustcode}
...

// 本变量代码由AI完成
static KLOG: Mutex<Option<KlogRingbufInner>> = Mutex::new(None);

// 本结构代码由AI完成
...
\end{rustcode}

\subsubsection*{序号366}
\noindent 路径：\code{os/components/wateros-klog/klog-impl/klog-ringbuf/src/lib.rs}\\
\noindent 行范围：268-270\\[0.4em]

\begin{rustcode}
...
// 本结构代码由AI完成
struct KlogInterruptGuard {
    state : ArchInterruptState,
}

...
\end{rustcode}

\subsubsection*{序号367}
\noindent 路径：\code{os/components/wateros-klog/klog-impl/klog-ringbuf/src/lib.rs}\\
\noindent 行范围：289-294\\[0.4em]

\begin{rustcode}
...
fn ensure_inner(guard: &mut Option<KlogRingbufInner>) -> &mut KlogRingbufInner {
    if guard.is_none() {
        *guard = Some(KlogRingbufInner::default());
    }
    guard.as_mut().unwrap()
...
\end{rustcode}

\subsubsection*{序号368}
\noindent 路径：\code{os/components/wateros-klog/klog-impl/klog-ringbuf/src/lib.rs}\\
\noindent 行范围：300-304\\[0.4em]

\begin{rustcode}
...
    pub fn init() {
        let mut guard = KLOG.lock();
        let inner = ensure_inner(&mut guard);
        inner.reset();
    }
...
\end{rustcode}

\subsubsection*{序号369}
\noindent 路径：\code{os/components/wateros-klog/klog-impl/klog-ringbuf/src/lib.rs}\\
\noindent 行范围：309-313\\[0.4em]

\begin{rustcode}
...
    pub fn with<R>(f: impl FnOnce(&mut KlogRingbufInner) -> R) -> R {
        let _irq = KlogInterruptGuard::new();
        let mut guard = KLOG.lock();
        f(ensure_inner(&mut guard))
    }
...
\end{rustcode}

\subsubsection*{序号370}
\noindent 路径：\code{os/components/wateros-klog/klog-impl/klog-ringbuf/src/lib.rs}\\
\noindent 行范围：317-324\\[0.4em]

\begin{rustcode}
...
    pub fn iter_from<F>(start_seq: u64, mut f: F)
    where
        F: FnMut(KlogRecordView<'_>),
    {
        let _irq = KlogInterruptGuard::new();
...
\end{rustcode}

\subsubsection*{序号371}
\noindent 路径：\code{os/components/wateros-klog/src/export.rs}\\
\noindent 行范围：1-34\\[0.4em]

\begin{rustcode}
//! 用户态可见的 syslog 线格式（traditional）。
//! 本模块代码由AI完成

use api_v0::KlogRecordMeta;

...
\end{rustcode}

\subsubsection*{序号372}
\noindent 路径：\code{os/components/wateros-klog/src/export.rs}\\
\noindent 行范围：11-34\\[0.4em]

\begin{rustcode}
...
pub fn format_traditional(meta: &KlogRecordMeta, text: &[u8], out: &mut [u8]) -> usize {
    let level_ch = [meta.traditional_level_char()];
    let prefix = [b'<', level_ch[0], b'>'];
    let mut written = 0usize;
    for &b in &prefix {
...
\end{rustcode}

\subsubsection*{序号373}
\noindent 路径：\code{os/components/wateros-klog/src/lib.rs}\\
\noindent 行范围：1-165\\[0.4em]

\begin{rustcode}
#![no_std]
//! 内核消息环聚合层：全局环、`klog_*!` 宏与 `sys_syslog` 内核语义。
//! 本模块代码由AI完成

pub use api_v0 as api;
...
\end{rustcode}

\subsubsection*{序号374}
\noindent 路径：\code{os/components/wateros-klog/src/lib.rs}\\
\noindent 行范围：19-21\\[0.4em]

\begin{rustcode}
...
#[inline]
pub fn init() {
    Ring::init();
}

...
\end{rustcode}

\subsubsection*{序号375}
\noindent 路径：\code{os/components/wateros-klog/src/lib.rs}\\
\noindent 行范围：32-42\\[0.4em]

\begin{rustcode}
...
pub fn record(level: u8, facility: u8, text: &[u8]) -> AppendResult {
    let mut meta = KlogRecordMeta::new(
        ts_nsec_now(),
        0,
        facility,
...
\end{rustcode}

\subsubsection*{序号376}
\noindent 路径：\code{os/components/wateros-klog/src/lib.rs}\\
\noindent 行范围：135-138\\[0.4em]

\begin{rustcode}
...
// 本结构代码由AI完成
pub struct KlogFmtBuffer {
    buf: [u8; 512],
    len: usize,
}
...
\end{rustcode}

\subsubsection*{序号377}
\noindent 路径：\code{os/components/wateros-klog/src/lib.rs}\\
\noindent 行范围：144-146\\[0.4em]

\begin{rustcode}
...
    #[inline]
    pub const fn new() -> Self {
        Self { buf: [0; 512], len: 0 }
    }

...
\end{rustcode}

\subsubsection*{序号378}
\noindent 路径：\code{os/components/wateros-klog/src/lib.rs}\\
\noindent 行范围：157-164\\[0.4em]

\begin{rustcode}
...
    fn write_str(&mut self, s: &str) -> core::fmt::Result {
        let bytes = s.as_bytes();
        let room = self.buf.len().saturating_sub(self.len);
        let n = bytes.len().min(room);
        self.buf[self.len..self.len + n].copy_from_slice(&bytes[..n]);
...
\end{rustcode}

\subsubsection*{序号379}
\noindent 路径：\code{os/components/wateros-klog/src/syscall.rs}\\
\noindent 行范围：1-113\\[0.4em]

\begin{rustcode}
//! `sys_syslog` 语义（内核缓冲侧）；用户指针由 `wateros-syscall` 负责拷贝。
//! 本模块代码由AI完成

use api_v0::{
    is_write_priority, KlogError, KlogFlags, KlogRecordMeta, KlogStore, SYSLOG_ACTION_CLEAR,
...
\end{rustcode}

\subsubsection*{序号380}
\noindent 路径：\code{os/components/wateros-klog/src/syscall.rs}\\
\noindent 行范围：16\\[0.4em]

\begin{rustcode}
...

// 本变量代码由AI完成
const KERNEL_LINE_MAX: usize = 2048;

/// 处理 `sys_syslog` action；`user_buf` 为 syscall 层提供的内核侧缓冲。
...
\end{rustcode}

\subsubsection*{序号381}
\noindent 路径：\code{os/components/wateros-klog/src/syscall.rs}\\
\noindent 行范围：20-39\\[0.4em]

\begin{rustcode}
...
pub fn dispatch_kernel(action: i32, user_buf: &mut [u8], user_len: usize) -> isize {
    if is_write_priority(action) {
        return write_priority(action, user_buf, user_len);
    }
    match action {
...
\end{rustcode}

\subsubsection*{序号382}
\noindent 路径：\code{os/components/wateros-klog/src/syscall.rs}\\
\noindent 行范围：43-56\\[0.4em]

\begin{rustcode}
...
fn read_one(buf: &mut [u8], len: usize, advance: bool) -> isize {
    let mut line = [0u8; KERNEL_LINE_MAX];
    KlogRingbuf::with(|ring| match ring.peek_next_unread() {
        Ok(view) => {
            let n = format_traditional(&view.meta, view.text, &mut line);
...
\end{rustcode}

\subsubsection*{序号383}
\noindent 路径：\code{os/components/wateros-klog/src/syscall.rs}\\
\noindent 行范围：59-87\\[0.4em]

\begin{rustcode}
...
fn read_all(buf: &mut [u8], len: usize) -> isize {
    let mut total = 0isize;
    let mut offset = 0usize;
    loop {
        if offset >= len {
...
\end{rustcode}

\subsubsection*{序号384}
\noindent 路径：\code{os/components/wateros-klog/src/syscall.rs}\\
\noindent 行范围：91-105\\[0.4em]

\begin{rustcode}
...
fn write_priority(priority: i32, msg: &[u8], _msg_len: usize) -> isize {
    let level = ((priority >> 3) & 7) as u8;
    let facility = (priority & 7) as u8;
    let flags = KlogFlags::empty().with(KlogFlags::USER);
    let mut meta = KlogRecordMeta::new(
...
\end{rustcode}

\subsubsection*{序号385}
\noindent 路径：\code{os/components/wateros-klog/src/syscall.rs}\\
\noindent 行范围：109-113\\[0.4em]

\begin{rustcode}
...
fn copy_out(buf: &mut [u8], len: usize, src: &[u8]) -> isize {
    let n = src.len().min(len);
    buf[..n].copy_from_slice(&src[..n]);
    n as isize
}
\end{rustcode}

\subsubsection*{序号386}
\noindent 路径：\code{os/components/wateros-mm/mm-impl/impl-dummy/src/lib.rs}\\
\noindent 行范围：1-54\\[0.4em]

\begin{rustcode}
//! 本模块代码由AI完成
//! `mm-impl` 桩：无 Sv39、无真实 `satp`；用于未启用 `wateros-mm` 的 `impl-sv39` 的构建。
//!
//! `kernel_mm_impl::from_elf_path` 等返回固定错误，避免链接真实 FS/页表路径。

...
\end{rustcode}

\subsubsection*{序号387}
\noindent 路径：\code{os/components/wateros-mm/mm-impl/impl-loongarch64/src/kernel_executable.rs}\\
\noindent 行范围：1-65\\[0.4em]

\begin{rustcode}
//! 本模块代码由AI完成
//! 从根卷路径装载可执行文件：ELF 直载或 shebang 脚本解析后加载解释器。

extern crate alloc;

...
\end{rustcode}

\subsubsection*{序号388}
\noindent 路径：\code{os/components/wateros-mm/mm-impl/impl-loongarch64/src/lib.rs}\\
\noindent 行范围：96-112\\[0.4em]

\begin{rustcode}
...
    pub fn fork_user_aspace(parent_aspace_ptr : usize) -> api_v0::error::MmResult<(usize, usize)> {
        use crate::pagetable::LoongArch64AddressSpace;
        use alloc::boxed::Box;
        use api_v0::address_space::AddressSpaceOps;
        use api_v0::error::MmError;
...
\end{rustcode}

\subsubsection*{序号389}
\noindent 路径：\code{os/components/wateros-mm/mm-impl/impl-loongarch64/src/lib.rs}\\
\noindent 行范围：115-127\\[0.4em]

\begin{rustcode}
...
    pub fn handle_cow_fault(parent_aspace_ptr : usize,
                            fault_addr : usize)
                            -> api_v0::error::MmResult<bool> {
        use crate::pagetable::LoongArch64AddressSpace;
        use api_v0::addr::VirtAddr;
...
\end{rustcode}

\subsubsection*{序号390}
\noindent 路径：\code{os/components/wateros-mm/mm-impl/impl-loongarch64/src/lib.rs}\\
\noindent 行范围：133-143\\[0.4em]

\begin{rustcode}
...
    pub fn drop_user_aspace(aspace_ptr : usize) {
        use crate::pagetable::LoongArch64AddressSpace;
        use alloc::boxed::Box;

        if aspace_ptr == 0 {
...
\end{rustcode}

\subsubsection*{序号391}
\noindent 路径：\code{os/components/wateros-mm/mm-impl/impl-loongarch64/src/pagetable.rs}\\
\noindent 行范围：234-241\\[0.4em]

\begin{rustcode}
...
pub(crate) struct LazyFileVma {
    pub start : VirtAddr,
    pub end : VirtAddr,
    pub perm : PagePerm,
    pub file_offset : usize,
...
\end{rustcode}

\subsubsection*{序号392}
\noindent 路径：\code{os/components/wateros-mm/mm-impl/impl-loongarch64/src/pagetable.rs}\\
\noindent 行范围：264-267\\[0.4em]

\begin{rustcode}
...
#[derive(Clone, Copy)]
pub(crate) struct SharedAnonVma {
    pub start : VirtAddr,
    pub end : VirtAddr,
}
...
\end{rustcode}

\subsubsection*{序号393}
\noindent 路径：\code{os/components/wateros-mm/mm-impl/impl-loongarch64/src/pagetable.rs}\\
\noindent 行范围：624-657\\[0.4em]

\begin{rustcode}
...
    pub fn fork_cow(&mut self) -> MmResult<LoongArch64AddressSpace> {
        log::trace!("[mm-fork] LoongArch64AddressSpace::fork begin root_ppn={}",
                    self.root.0);
        let child_root = alloc_table_frame_zeroed()?;
        // SAFETY: 刚分配并清零的帧作为子地址空间根页表。
...
\end{rustcode}

\subsubsection*{序号394}
\noindent 路径：\code{os/components/wateros-mm/mm-impl/impl-loongarch64/src/pagetable.rs}\\
\noindent 行范围：660-687\\[0.4em]

\begin{rustcode}
...
    fn handle_cow_page(&mut self, vpn : VirtPageNum) -> MmResult<bool> {
        let Some((pte, level)) = self.walk_find(vpn)? else {
            return Ok(false);
        };
        let flags = pte.flags();
...
\end{rustcode}

\subsubsection*{序号395}
\noindent 路径：\code{os/components/wateros-mm/mm-impl/impl-loongarch64/src/pagetable.rs}\\
\noindent 行范围：690-692\\[0.4em]

\begin{rustcode}
...
    // 本方法代码由AI完成
    pub fn handle_cow_fault(&mut self, fault_addr : VirtAddr) -> MmResult<bool> {
        self.handle_cow_page(fault_addr.floor_page())
    }

...
\end{rustcode}

\subsubsection*{序号396}
\noindent 路径：\code{os/components/wateros-mm/mm-impl/impl-loongarch64/src/pagetable.rs}\\
\noindent 行范围：695-741\\[0.4em]

\begin{rustcode}
...
    pub fn handle_lazy_page_fault<A>(&mut self,
                                     allocator : &mut A,
                                     fault_addr : VirtAddr,
                                     access : PageFaultAccess)
                                     -> MmResult<bool>
...
\end{rustcode}

\subsubsection*{序号397}
\noindent 路径：\code{os/components/wateros-mm/mm-impl/impl-loongarch64/src/pagetable.rs}\\
\noindent 行范围：744-773\\[0.4em]

\begin{rustcode}
...
    pub fn ensure_private_for_write(&mut self, vpn : VirtPageNum) -> MmResult<bool> {
        if self.handle_cow_page(vpn)? {
            return Ok(true);
        }
        let Some((pte, level)) = self.walk_find(vpn)? else {
...
\end{rustcode}

\subsubsection*{序号398}
\noindent 路径：\code{os/components/wateros-mm/mm-impl/impl-loongarch64/src/user_access.rs}\\
\noindent 行范围：1-129\\[0.4em]

\begin{rustcode}
//! 本模块代码由AI完成
//! LoongArch64 用户缓冲区 [`api_v0::user_access::UserMemoryOps`]。

use api_v0::addr::{VirtAddr, PAGE_SIZE};
use api_v0::address_space::AddressSpaceOps;
...
\end{rustcode}

\subsubsection*{序号399}
\noindent 路径：\code{os/components/wateros-mm/mm-impl/impl-loongarch64/src/user_aspace.rs}\\
\noindent 行范围：1-29\\[0.4em]

\begin{rustcode}
//! 本模块代码由AI完成
//! 将 [`api_v0::kernel_bringup::LoadedElf::user_aspace_ptr`] 解析为
//! [`LoongArch64AddressSpace`]， 供上层在闭包内调用 [`api_v0::brk::HeapBrk`] /
//! [`api_v0::mmap::MmapOps`] 等机制原语。
//!
...
\end{rustcode}

\subsubsection*{序号400}
\noindent 路径：\code{os/components/wateros-mm/mm-impl/impl-loongarch64/src/user_heap_mmap.rs}\\
\noindent 行范围：1-554\\[0.4em]

\begin{rustcode}
//! 本模块代码由AI完成
//! 用户堆 `brk` 与匿名/文件 `mmap`/`munmap`/`mprotect` 的 LoongArch64 实现。
//!
//! # 缺页与 trap
//!
...
\end{rustcode}

\subsubsection*{序号401}
\noindent 路径：\code{os/components/wateros-mm/mm-impl/impl-sv39/src/kernel_executable.rs}\\
\noindent 行范围：1-73\\[0.4em]

\begin{rustcode}
//! 本模块代码由AI完成
//! 从根卷路径装载可执行文件：ELF 直载或 shebang 脚本解析后加载解释器。

extern crate alloc;

...
\end{rustcode}

\subsubsection*{序号402}
\noindent 路径：\code{os/components/wateros-mm/mm-impl/impl-sv39/src/lib.rs}\\
\noindent 行范围：98-113\\[0.4em]

\begin{rustcode}
...
    pub fn fork_user_aspace(parent_aspace_ptr : usize) -> api_v0::error::MmResult<(usize, usize)> {
        use crate::pagetable::Sv39AddressSpace;
        use alloc::boxed::Box;
        use api_v0::address_space::AddressSpaceOps;
        use api_v0::error::MmError;
...
\end{rustcode}

\subsubsection*{序号403}
\noindent 路径：\code{os/components/wateros-mm/mm-impl/impl-sv39/src/lib.rs}\\
\noindent 行范围：116-128\\[0.4em]

\begin{rustcode}
...
    pub fn handle_cow_fault(parent_aspace_ptr : usize,
                            fault_addr : usize)
                            -> api_v0::error::MmResult<bool> {
        use crate::pagetable::Sv39AddressSpace;
        use api_v0::addr::VirtAddr;
...
\end{rustcode}

\subsubsection*{序号404}
\noindent 路径：\code{os/components/wateros-mm/mm-impl/impl-sv39/src/lib.rs}\\
\noindent 行范围：134-144\\[0.4em]

\begin{rustcode}
...
    pub fn drop_user_aspace(aspace_ptr : usize) {
        use crate::pagetable::Sv39AddressSpace;
        use alloc::boxed::Box;

        if aspace_ptr == 0 {
...
\end{rustcode}

\subsubsection*{序号405}
\noindent 路径：\code{os/components/wateros-mm/mm-impl/impl-sv39/src/pagetable.rs}\\
\noindent 行范围：252-259\\[0.4em]

\begin{rustcode}
...
pub(crate) struct LazyFileVma {
    pub start : VirtAddr,
    pub end : VirtAddr,
    pub perm : PagePerm,
    pub file_offset : usize,
...
\end{rustcode}

\subsubsection*{序号406}
\noindent 路径：\code{os/components/wateros-mm/mm-impl/impl-sv39/src/pagetable.rs}\\
\noindent 行范围：285-288\\[0.4em]

\begin{rustcode}
...
#[derive(Clone, Copy)]
pub(crate) struct SharedAnonVma {
    pub start : VirtAddr,
    pub end : VirtAddr,
}
...
\end{rustcode}

\subsubsection*{序号407}
\noindent 路径：\code{os/components/wateros-mm/mm-impl/impl-sv39/src/pagetable.rs}\\
\noindent 行范围：690-725\\[0.4em]

\begin{rustcode}
...
    pub fn fork_cow(&mut self) -> MmResult<Sv39AddressSpace> {
        log::trace!("[mm-fork] Sv39AddressSpace::fork begin root_ppn={}",
                    self.root.0);
        let child_root = alloc_table_frame_zeroed()?;
        // SAFETY: 刚分配并清零的帧作为子地址空间根页表。
...
\end{rustcode}

\subsubsection*{序号408}
\noindent 路径：\code{os/components/wateros-mm/mm-impl/impl-sv39/src/pagetable.rs}\\
\noindent 行范围：745-775\\[0.4em]

\begin{rustcode}
...
    fn handle_cow_page(&mut self, vpn : VirtPageNum) -> MmResult<bool> {
        let Some((pte, level)) = self.walk_find(vpn)? else {
            return Ok(false);
        };
        let flags = pte.flags();
...
\end{rustcode}

\subsubsection*{序号409}
\noindent 路径：\code{os/components/wateros-mm/mm-impl/impl-sv39/src/pagetable.rs}\\
\noindent 行范围：778-780\\[0.4em]

\begin{rustcode}
...
    // 本方法代码由AI完成
    pub fn handle_cow_fault(&mut self, fault_addr : VirtAddr) -> MmResult<bool> {
        self.handle_cow_page(fault_addr.floor_page())
    }

...
\end{rustcode}

\subsubsection*{序号410}
\noindent 路径：\code{os/components/wateros-mm/mm-impl/impl-sv39/src/pagetable.rs}\\
\noindent 行范围：783-832\\[0.4em]

\begin{rustcode}
...
    pub fn handle_lazy_page_fault<A>(&mut self,
                                     allocator : &mut A,
                                     fault_addr : VirtAddr,
                                     access : PageFaultAccess)
                                     -> MmResult<bool>
...
\end{rustcode}

\subsubsection*{序号411}
\noindent 路径：\code{os/components/wateros-mm/mm-impl/impl-sv39/src/pagetable.rs}\\
\noindent 行范围：835-865\\[0.4em]

\begin{rustcode}
...
    pub fn ensure_private_for_write(&mut self, vpn : VirtPageNum) -> MmResult<bool> {
        if self.handle_cow_page(vpn)? {
            return Ok(true);
        }
        let Some((pte, level)) = self.walk_find(vpn)? else {
...
\end{rustcode}

\subsubsection*{序号412}
\noindent 路径：\code{os/components/wateros-mm/mm-impl/impl-sv39/src/user_access.rs}\\
\noindent 行范围：1-170\\[0.4em]

\begin{rustcode}
//! 本模块代码由AI完成
//! Sv39 用户缓冲区 [`api_v0::user_access::UserMemoryOps`]。
//!
//! - **读/写**（`copy_from_user` / `copy_to_user`）：软件 walk 用户页表 + 内核恒等访问 PA，
//!   避免在 syscall 处理中临时切换 `satp`（ioctl 成功路径依赖 caller-saved 寄存器保真）。
...
\end{rustcode}

\subsubsection*{序号413}
\noindent 路径：\code{os/components/wateros-mm/mm-impl/impl-sv39/src/user_aspace.rs}\\
\noindent 行范围：1-27\\[0.4em]

\begin{rustcode}
//! 本模块代码由AI完成
//! 将 [`api_v0::kernel_bringup::LoadedElf::user_aspace_ptr`] 解析为 [`Sv39AddressSpace`]，
//! 供上层在闭包内调用 [`api_v0::brk::HeapBrk`] / [`api_v0::mmap::MmapOps`] 等机制原语。
//!
//! **Safety**：`handle` 须来自 bring-up 泄漏的用户地址空间，且与当前任务安装的 `satp` 一致。
...
\end{rustcode}

\subsubsection*{序号414}
\noindent 路径：\code{os/components/wateros-mm/mm-impl/impl-sv39/src/user_heap_mmap.rs}\\
\noindent 行范围：1-555\\[0.4em]

\begin{rustcode}
//! 本模块代码由AI完成
//! 用户堆 `brk` 与匿名/文件 `mmap`/`munmap`/`mprotect` 的 Sv39 实现。
//!
//! # 缺页与 trap
//!
...
\end{rustcode}

\subsubsection*{序号415}
\noindent 路径：\code{os/components/wateros-platform/platform-arch/arch-impl/impl-dummy/src/lib.rs}\\
\noindent 行范围：1-20\\[0.4em]

\begin{rustcode}
#![no_std]

//! 本模块代码由AI完成
//! **占位**架构实现：用于未启用真实 ISA profile 的构建或依赖占位；非生产路径。
//!
...
\end{rustcode}

\subsubsection*{序号416}
\noindent 路径：\code{os/components/wateros-platform/platform-impl/impl-dummy/src/lib.rs}\\
\noindent 行范围：1-109\\[0.4em]

\begin{rustcode}
#![no_std]

//! 本模块代码由AI完成
//! **占位**平台实现：启动参数与时间频率 trait 的桩，用于未绑定真实环境的构建。
//!
...
\end{rustcode}

\subsubsection*{序号417}
\noindent 路径：\code{os/components/wateros-platform/platform-impl/impl-qemu-loongarch64-virt/src/lib.rs}\\
\noindent 行范围：20-24\\[0.4em]

\begin{rustcode}
...
    pub struct QEMULoongArch64VirtBootArgs {
        arg0: usize,
        arg1: usize,
        arg2: usize,
    }
...
\end{rustcode}

\subsubsection*{序号418}
\noindent 路径：\code{os/components/wateros-platform/platform-impl/impl-qemu-loongarch64-virt/src/lib.rs}\\
\noindent 行范围：53-60\\[0.4em]

\begin{rustcode}
...
    pub struct QEMULoongArch64VirtBootContext {
        /// 固件 `a0`。
        pub arg0: usize,
        /// 固件 `a1`。
        pub arg1: usize,
...
\end{rustcode}

\subsubsection*{序号419}
\noindent 路径：\code{os/components/wateros-platform/platform-impl/impl-qemu-loongarch64-virt/src/lib.rs}\\
\noindent 行范围：84\\[0.4em]

\begin{rustcode}
...
    /// QEMU virt 上 StableCounter 常用频率（Hz）；与 arch `rdtime.d` 刻度一致。
    // 本结构代码由AI完成
    pub struct QEMULoongArch64VirtTime;

    impl PlatformTime for QEMULoongArch64VirtTime {
...
\end{rustcode}

\subsubsection*{序号420}
\noindent 路径：\code{os/components/wateros-platform/platform-impl/impl-qemu-riscv64-opensbi/src/console.rs}\\
\noindent 行范围：1-31\\[0.4em]

\begin{rustcode}
//! 本模块代码由AI完成
//! OpenSBI 控制台后端。

use api_v0::console::{PlatformConsoleError, PlatformConsoleResult};
#[allow(unused)]
...
\end{rustcode}

\subsubsection*{序号421}
\noindent 路径：\code{os/components/wateros-platform/platform-impl/impl-qemu-riscv64-opensbi/src/lib.rs}\\
\noindent 行范围：22-27\\[0.4em]

\begin{rustcode}
...
    pub struct QEMURiscv64OpenSBIBootArgs {
        /// 固件 `a0`（通常为 hart id）。
        _arg0 : usize,
        /// 固件 `a1`（通常为 DTB 物理地址）。
        _arg1 : usize,
...
\end{rustcode}

\subsubsection*{序号422}
\noindent 路径：\code{os/components/wateros-platform/platform-impl/impl-qemu-riscv64-opensbi/src/lib.rs}\\
\noindent 行范围：46-51\\[0.4em]

\begin{rustcode}
...
    pub struct QEMURiscv64OpenSBIBootContext {
        /// 启动 hart id。
        _hart_id : base::cpu::CPUHartID,
        /// 设备树 blob 物理地址。
        _dtb_pa : base::boot::DTBPA,
...
\end{rustcode}

\subsubsection*{序号423}
\noindent 路径：\code{os/components/wateros-platform/platform-impl/impl-qemu-riscv64-opensbi/src/lib.rs}\\
\noindent 行范围：76\\[0.4em]

\begin{rustcode}
...
    /// QEMU virt 上常用的 timebase 频率（Hz）；引导期可由 DTB 覆盖，此处为回退常量。
    // 本结构代码由AI完成
    pub struct QEMURiscv64OpenSBITime;

    impl PlatformTime for QEMURiscv64OpenSBITime {
...
\end{rustcode}

\subsubsection*{序号424}
\noindent 路径：\code{os/components/wateros-platform/platform-impl/impl-qemu-riscv64-opensbi/src/reset.rs}\\
\noindent 行范围：1-90\\[0.4em]

\begin{rustcode}
//! 本模块代码由AI完成
//! OpenSBI system reset 后端。

use api_v0::reset::{
    PlatformResetError, PlatformResetReason, PlatformResetResult, PlatformResetType,
...
\end{rustcode}

\subsubsection*{序号425}
\noindent 路径：\code{os/components/wateros-platform/platform-impl/impl-qemu-riscv64-opensbi/src/reset.rs}\\
\noindent 行范围：11-18\\[0.4em]

\begin{rustcode}
...
pub enum OpenSBIResetType {
    /// 关机。
    Shutdown,
    /// 冷重启。
    ColdReboot,
...
\end{rustcode}

\subsubsection*{序号426}
\noindent 路径：\code{os/components/wateros-platform/platform-impl/impl-qemu-riscv64-opensbi/src/reset.rs}\\
\noindent 行范围：22-27\\[0.4em]

\begin{rustcode}
...
pub enum OpenSBIResetReason {
    /// 无特定原因。
    NoReason,
    /// 系统故障。
    SystemFailure,
...
\end{rustcode}

\subsubsection*{序号427}
\noindent 路径：\code{os/components/wateros-platform/platform-impl/impl-qemu-riscv64-opensbi/src/timer.rs}\\
\noindent 行范围：1-17\\[0.4em]

\begin{rustcode}
//! 本模块代码由AI完成
//! OpenSBI timer 后端：将绝对 tick deadline 交给 SBI `set_timer`。

use api_v0::timer::{
    PlatformDeadlineTimerError, PlatformDeadlineTimerResult, PlatformTimerDeadline,
...
\end{rustcode}

\subsubsection*{序号428}
\noindent 路径：\code{os/components/wateros-syscall/syscall-impl/impl-kernel/src/epoll_fd.rs}\\
\noindent 行范围：1-289\\[0.4em]

\begin{rustcode}
//! epoll fd → [`EpollInstance`] 映射表（[`VfsIoHandle`] 无法向下转型）。

//! 本模块代码由AI完成
extern crate alloc;

...
\end{rustcode}

\subsubsection*{序号429}
\noindent 路径：\code{os/components/wateros-syscall/syscall-impl/impl-kernel/src/epoll_fd.rs}\\
\noindent 行范围：20-23\\[0.4em]

\begin{rustcode}
...
// 本结构代码由AI完成
pub(crate) struct EpollInterest {
    pub events: u32,
    pub data: u64,
}
...
\end{rustcode}

\subsubsection*{序号430}
\noindent 路径：\code{os/components/wateros-syscall/syscall-impl/impl-kernel/src/epoll_fd.rs}\\
\noindent 行范围：28-31\\[0.4em]

\begin{rustcode}
...
// 本结构代码由AI完成
pub(crate) struct EpollInstance {
    pub interests: BTreeMap<usize, EpollInterest>,
    inode: u64,
}
...
\end{rustcode}

\subsubsection*{序号431}
\noindent 路径：\code{os/components/wateros-syscall/syscall-impl/impl-kernel/src/epoll_fd.rs}\\
\noindent 行范围：250-268\\[0.4em]

\begin{rustcode}
...
pub(crate) fn epoll_to_poll_events(events: u32) -> i16 {
    let mut poll = 0i16;
    if events & EPOLLIN != 0 {
        poll |= POLLIN;
    }
...
\end{rustcode}

\subsubsection*{序号432}
\noindent 路径：\code{os/components/wateros-syscall/syscall-impl/impl-kernel/src/epoll_fd.rs}\\
\noindent 行范围：271-289\\[0.4em]

\begin{rustcode}
...
pub(crate) fn poll_to_epoll_events(revents: i16) -> u32 {
    let mut events = 0u32;
    if revents & POLLIN != 0 {
        events |= EPOLLIN;
    }
...
\end{rustcode}

\subsubsection*{序号433}
\noindent 路径：\code{os/components/wateros-syscall/syscall-impl/impl-kernel/src/fallible_buf.rs}\\
\noindent 行范围：1-54\\[0.4em]

\begin{rustcode}
//! Syscall 内核缓冲的可失败分配，避免 OOM 时触发全局 `alloc_error_handler`。

//! 本模块代码由AI完成
extern crate alloc;

...
\end{rustcode}

\subsubsection*{序号434}
\noindent 路径：\code{os/components/wateros-syscall/syscall-impl/impl-kernel/src/fallible_buf.rs}\\
\noindent 行范围：12\\[0.4em]

\begin{rustcode}
...
/// 与 `read`/`write`/`pread` 族 syscall 对齐的上限。
// 本变量代码由AI完成
pub const SYSCALL_IO_MAX : usize = 4 * 1024 * 1024;

/// 与 `sendto`/`recvfrom`/`sendmsg` 对齐的上限。
...
\end{rustcode}

\subsubsection*{序号435}
\noindent 路径：\code{os/components/wateros-syscall/syscall-impl/impl-kernel/src/fallible_buf.rs}\\
\noindent 行范围：16\\[0.4em]

\begin{rustcode}
...
/// 与 `sendto`/`recvfrom`/`sendmsg` 对齐的上限。
// 本变量代码由AI完成
pub const SYSCALL_SOCK_IO_MAX : usize = 64 * 1024;

/// `sched_getaffinity` / `sched_setaffinity` 用户缓冲上界。
...
\end{rustcode}

\subsubsection*{序号436}
\noindent 路径：\code{os/components/wateros-syscall/syscall-impl/impl-kernel/src/fallible_buf.rs}\\
\noindent 行范围：26-38\\[0.4em]

\begin{rustcode}
...
pub fn try_kbuf(len : usize, max : usize) -> Result<Vec<u8>, ErrNo> {
    if len > max {
        return Err(ErrNo::EINVAL);
    }
    if len == 0 {
...
\end{rustcode}

\subsubsection*{序号437}
\noindent 路径：\code{os/components/wateros-syscall/syscall-impl/impl-kernel/src/lib.rs}\\
\noindent 行范围：1-979\\[0.4em]

\begin{rustcode}
#![no_std]
//! 内核 syscall 实现：各 `sys_*` 具体语义，并实现 `wateros-syscall-api-v0` 的分发 trait。
//! 本模块代码由AI完成

extern crate alloc;
...
\end{rustcode}

\subsubsection*{序号438}
\noindent 路径：\code{os/components/wateros-syscall/syscall-impl/impl-kernel/src/lib.rs}\\
\noindent 行范围：79\\[0.4em]

\begin{rustcode}
...
/// 聚合 crate feature 选中的内核 syscall 分发器。
// 本结构代码由AI完成
pub struct KernelSyscallDispatcher;

const SYS_STATX: usize = 291;
...
\end{rustcode}

\subsubsection*{序号439}
\noindent 路径：\code{os/components/wateros-syscall/syscall-impl/impl-kernel/src/lib.rs}\\
\noindent 行范围：93-141\\[0.4em]

\begin{rustcode}
...
fn dispatch_syscall_aliases(syscall_nr: usize, args: SyscallArgs) -> isize {
    if syscall_nr == SYS_FSTATAT {
        return sys::sys_fstatat(args).0;
    }
    if syscall_nr == SYS_STATX {
...
\end{rustcode}

\subsubsection*{序号440}
\noindent 路径：\code{os/components/wateros-syscall/syscall-impl/impl-kernel/src/lib.rs}\\
\noindent 行范围：144-977\\[0.4em]

\begin{rustcode}
...
impl api_v0::SyscallDispatcher for KernelSyscallDispatcher {
    type NumberTable = ActiveSyscallNumberTable;

    #[inline]
    fn dispatch_yield(_args: SyscallArgs) -> isize {
...
\end{rustcode}

\subsubsection*{序号441}
\noindent 路径：\code{os/components/wateros-syscall/syscall-impl/impl-kernel/src/linux_stat.rs}\\
\noindent 行范围：1-212\\[0.4em]

\begin{rustcode}
//! Linux riscv64 `struct stat` 布局（与 musl/glibc 用户态一致）。

//! 本模块代码由AI完成
use vfs::api::VfsMetadata;
use vfs::api::VfsNodeType;
...
\end{rustcode}

\subsubsection*{序号442}
\noindent 路径：\code{os/components/wateros-syscall/syscall-impl/impl-kernel/src/linux_stat.rs}\\
\noindent 行范围：11-32\\[0.4em]

\begin{rustcode}
...
pub struct LinuxStat {
    pub st_dev: u64,
    pub st_ino: u64,
    pub st_mode: u32,
    pub st_nlink: u32,
...
\end{rustcode}

\subsubsection*{序号443}
\noindent 路径：\code{os/components/wateros-syscall/syscall-impl/impl-kernel/src/linux_stat.rs}\\
\noindent 行范围：49-74\\[0.4em]

\begin{rustcode}
...
pub struct LinuxStatx {
    pub stx_mask: u32,
    pub stx_blksize: u32,
    pub stx_attributes: u64,
    pub stx_nlink: u32,
...
\end{rustcode}

\subsubsection*{序号444}
\noindent 路径：\code{os/components/wateros-syscall/syscall-impl/impl-kernel/src/linux_stat.rs}\\
\noindent 行范围：127-150\\[0.4em]

\begin{rustcode}
...
pub(crate) fn fill_linux_stat(meta: &VfsMetadata, size: u64) -> LinuxStat {
    let mode = linux_mode(meta);
    LinuxStat {
        st_dev: linux_dev(meta.device_major, meta.device_minor),
        st_ino: meta.inode,
...
\end{rustcode}

\subsubsection*{序号445}
\noindent 路径：\code{os/components/wateros-syscall/syscall-impl/impl-kernel/src/linux_stat.rs}\\
\noindent 行范围：153-172\\[0.4em]

\begin{rustcode}
...
pub(crate) fn fill_linux_statx(meta: &VfsMetadata, size: u64, _requested_mask: u32) -> LinuxStatx {
    let mode = linux_mode(meta);
    LinuxStatx {
        stx_mask: STATX_SUPPORTED,
        stx_blksize: 4096,
...
\end{rustcode}

\subsubsection*{序号446}
\noindent 路径：\code{os/components/wateros-syscall/syscall-impl/impl-kernel/src/mm_util.rs}\\
\noindent 行范围：1-83\\[0.4em]

\begin{rustcode}
//! 系统调用层与 [`wateros-mm`] API 之间的错误与标志转换。
//! 本模块代码由AI完成

use abi::errno::ErrNo;

...
\end{rustcode}

\subsubsection*{序号447}
\noindent 路径：\code{os/components/wateros-syscall/syscall-impl/impl-kernel/src/mm_util.rs}\\
\noindent 行范围：12-20\\[0.4em]

\begin{rustcode}
...
pub(crate) fn mm_err_to_errno(e: mm::api::error::MmError) -> ErrNo {
    use mm::api::error::MmError;
    match e {
        MmError::OutOfMemory | MmError::FrameAlloc(_) => ErrNo::ENOMEM,
        MmError::InvalidAddress | MmError::AlreadyMapped | MmError::NotMapped => ErrNo::EINVAL,
...
\end{rustcode}

\subsubsection*{序号448}
\noindent 路径：\code{os/components/wateros-syscall/syscall-impl/impl-kernel/src/mm_util.rs}\\
\noindent 行范围：23-36\\[0.4em]

\begin{rustcode}
...
pub(crate) fn linux_mmap_prot_to_perm(prot: i32) -> mm::api::perm::PagePerm {
    use mm::api::perm::PagePerm;
    let mut p = PagePerm::U;
    if prot & 1 != 0 {
        p |= PagePerm::R;
...
\end{rustcode}

\subsubsection*{序号449}
\noindent 路径：\code{os/components/wateros-syscall/syscall-impl/impl-kernel/src/mm_util.rs}\\
\noindent 行范围：39-59\\[0.4em]

\begin{rustcode}
...
pub(crate) fn linux_mmap_flags_to_map_flags(flags: u32) -> mm::api::flags::MapFlags {
    use mm::api::flags::MapFlags;
    const MAP_SHARED: u32 = 0x01;
    const MAP_PRIVATE: u32 = 0x02;
    const MAP_FIXED: u32 = 0x10;
...
\end{rustcode}

\subsubsection*{序号450}
\noindent 路径：\code{os/components/wateros-syscall/syscall-impl/impl-kernel/src/poll_engine.rs}\\
\noindent 行范围：1-697\\[0.4em]

\begin{rustcode}
//! 共享 `poll` / `ppoll` / `pselect6` / `select` 就绪扫描与阻塞等待。

//! 本模块代码由AI完成
extern crate alloc;

...
\end{rustcode}

\subsubsection*{序号451}
\noindent 路径：\code{os/components/wateros-syscall/syscall-impl/impl-kernel/src/poll_engine.rs}\\
\noindent 行范围：17\\[0.4em]

\begin{rustcode}
...

// 本变量代码由AI完成
pub(crate) const POLLIN: i16 = 0x001;
// 本变量代码由AI完成
pub(crate) const POLLOUT: i16 = 0x004;
...
\end{rustcode}

\subsubsection*{序号452}
\noindent 路径：\code{os/components/wateros-syscall/syscall-impl/impl-kernel/src/poll_engine.rs}\\
\noindent 行范围：19\\[0.4em]

\begin{rustcode}
...
pub(crate) const POLLIN: i16 = 0x001;
// 本变量代码由AI完成
pub(crate) const POLLOUT: i16 = 0x004;
pub(crate) const POLLPRI: i16 = 0x002;
// 本变量代码由AI完成
...
\end{rustcode}

\subsubsection*{序号453}
\noindent 路径：\code{os/components/wateros-syscall/syscall-impl/impl-kernel/src/poll_engine.rs}\\
\noindent 行范围：22\\[0.4em]

\begin{rustcode}
...
pub(crate) const POLLPRI: i16 = 0x002;
// 本变量代码由AI完成
pub(crate) const POLLERR: i16 = 0x008;
// 本变量代码由AI完成
pub(crate) const POLLHUP: i16 = 0x010;
...
\end{rustcode}

\subsubsection*{序号454}
\noindent 路径：\code{os/components/wateros-syscall/syscall-impl/impl-kernel/src/poll_engine.rs}\\
\noindent 行范围：24\\[0.4em]

\begin{rustcode}
...
pub(crate) const POLLERR: i16 = 0x008;
// 本变量代码由AI完成
pub(crate) const POLLHUP: i16 = 0x010;
// 本变量代码由AI完成
pub(crate) const POLLNVAL: i16 = 0x020;
...
\end{rustcode}

\subsubsection*{序号455}
\noindent 路径：\code{os/components/wateros-syscall/syscall-impl/impl-kernel/src/poll_engine.rs}\\
\noindent 行范围：26\\[0.4em]

\begin{rustcode}
...
pub(crate) const POLLHUP: i16 = 0x010;
// 本变量代码由AI完成
pub(crate) const POLLNVAL: i16 = 0x020;

// 本变量代码由AI完成
...
\end{rustcode}

\subsubsection*{序号456}
\noindent 路径：\code{os/components/wateros-syscall/syscall-impl/impl-kernel/src/poll_engine.rs}\\
\noindent 行范围：29\\[0.4em]

\begin{rustcode}
...

// 本变量代码由AI完成
pub(crate) const FD_SETSIZE: usize = 1024;
const FD_SET_WORDS: usize = FD_SETSIZE / 64;
const SOCKET_READY_YIELD_SPINS: usize = 4;
...
\end{rustcode}

\subsubsection*{序号457}
\noindent 路径：\code{os/components/wateros-syscall/syscall-impl/impl-kernel/src/poll_engine.rs}\\
\noindent 行范围：36-40\\[0.4em]

\begin{rustcode}
...
pub(crate) struct PollFd {
    pub fd: i32,
    pub events: i16,
    pub revents: i16,
}
...
\end{rustcode}

\subsubsection*{序号458}
\noindent 路径：\code{os/components/wateros-syscall/syscall-impl/impl-kernel/src/poll_engine.rs}\\
\noindent 行范围：45-48\\[0.4em]

\begin{rustcode}
...
// 本结构代码由AI完成
pub(crate) struct UserTimespec {
    pub sec: isize,
    pub nsec: isize,
}
...
\end{rustcode}

\subsubsection*{序号459}
\noindent 路径：\code{os/components/wateros-syscall/syscall-impl/impl-kernel/src/poll_engine.rs}\\
\noindent 行范围：233-241\\[0.4em]

\begin{rustcode}
...
pub(crate) fn poll_revents_fd(fd: usize, events: i16) -> i16 {
    if socket_fd::lookup(fd).is_some() {
        return poll_socket_revents(fd, events);
    }
    match vfs::fd::with_current_io(fd, |handle| handle.poll_revents(events)) {
...
\end{rustcode}

\subsubsection*{序号460}
\noindent 路径：\code{os/components/wateros-syscall/syscall-impl/impl-kernel/src/poll_engine.rs}\\
\noindent 行范围：359-378\\[0.4em]

\begin{rustcode}
...
pub(crate) fn do_poll_with_deadline(
    fds_ptr: usize,
    nfds: usize,
    deadline: PollDeadline,
) -> UserRet {
...
\end{rustcode}

\subsubsection*{序号461}
\noindent 路径：\code{os/components/wateros-syscall/syscall-impl/impl-kernel/src/poll_engine.rs}\\
\noindent 行范围：667-697\\[0.4em]

\begin{rustcode}
...
pub(crate) fn do_pselect_with_deadline(
    nfds: usize,
    readfds_ptr: usize,
    writefds_ptr: usize,
    exceptfds_ptr: usize,
...
\end{rustcode}

\subsubsection*{序号462}
\noindent 路径：\code{os/components/wateros-syscall/syscall-impl/impl-kernel/src/socket_block.rs}\\
\noindent 行范围：1-17\\[0.4em]

\begin{rustcode}
//! 阻塞套接字等待：真阻塞 + 可中断（`EINTR`），非阻塞立即 `EAGAIN`。

//! 本模块代码由AI完成
use abi::errno::ErrNo;

...
\end{rustcode}

\subsubsection*{序号463}
\noindent 路径：\code{os/components/wateros-syscall/syscall-impl/impl-kernel/src/socket_block.rs}\\
\noindent 行范围：8-17\\[0.4em]

\begin{rustcode}
...
pub(crate) fn socket_blocking_tick(nonblocking: bool, task_id: usize) -> Result<(), ErrNo> {
    if nonblocking {
        return Err(ErrNo::EAGAIN);
    }
    if ipc::signal::with_registry(|registry| registry.has_deliverable(task_id).unwrap_or(false)) {
...
\end{rustcode}

\subsubsection*{序号464}
\noindent 路径：\code{os/components/wateros-syscall/syscall-impl/impl-kernel/src/socket_fd.rs}\\
\noindent 行范围：1-187\\[0.4em]

\begin{rustcode}
//! Socket fd → network [`SocketRef`] 映射表。
//!
//! 因 [`VfsIoHandle`] 不支持向下转型，每个 socket fd 的共享 socket 状态在此独立维护。

//! 本模块代码由AI完成
...
\end{rustcode}

\subsubsection*{序号465}
\noindent 路径：\code{os/components/wateros-syscall/syscall-impl/impl-kernel/src/stat_times.rs}\\
\noindent 行范围：1-75\\[0.4em]

\begin{rustcode}
//! VFS 元数据尚无 atime/mtime 字段前，syscall 层临时覆盖时间戳。

//! 本模块代码由AI完成
extern crate alloc;

...
\end{rustcode}

\subsubsection*{序号466}
\noindent 路径：\code{os/components/wateros-syscall/syscall-impl/impl-kernel/src/stat_times.rs}\\
\noindent 行范围：15-18\\[0.4em]

\begin{rustcode}
...
// 本结构代码由AI完成
pub(crate) struct StatTime {
    pub sec: i64,
    pub nsec: i64,
}
...
\end{rustcode}

\subsubsection*{序号467}
\noindent 路径：\code{os/components/wateros-syscall/syscall-impl/impl-kernel/src/sys/accept.rs}\\
\noindent 行范围：1-181\\[0.4em]

\begin{rustcode}
//! `accept4(2)`：接受 TCP 连接并返回新 fd。

//! 本模块代码由AI完成
use abi::errno::ErrNo;
use abi::syscall_args::SyscallArgs;
...
\end{rustcode}

\subsubsection*{序号468}
\noindent 路径：\code{os/components/wateros-syscall/syscall-impl/impl-kernel/src/sys/accept.rs}\\
\noindent 行范围：23-28\\[0.4em]

\begin{rustcode}
...
struct SockAddrIn {
    sin_family: u16,
    sin_port: u16,
    sin_addr: [u8; 4],
    sin_zero: [u8; 8],
...
\end{rustcode}

\subsubsection*{序号469}
\noindent 路径：\code{os/components/wateros-syscall/syscall-impl/impl-kernel/src/sys/accept.rs}\\
\noindent 行范围：31-37\\[0.4em]

\begin{rustcode}
...
pub(crate) fn sys_accept4(args: SyscallArgs) -> UserRet {
    let fd = args.arg(0);
    let addr_ptr = args.arg(1);
    let addrlen_ptr = args.arg(2);
    let flags = args.arg(3);
...
\end{rustcode}

\subsubsection*{序号470}
\noindent 路径：\code{os/components/wateros-syscall/syscall-impl/impl-kernel/src/sys/accept.rs}\\
\noindent 行范围：40-45\\[0.4em]

\begin{rustcode}
...
pub(crate) fn sys_accept(args: SyscallArgs) -> UserRet {
    let fd = args.arg(0);
    let addr_ptr = args.arg(1);
    let addrlen_ptr = args.arg(2);
    accept_inner(fd, addr_ptr, addrlen_ptr, 0)
...
\end{rustcode}

\subsubsection*{序号471}
\noindent 路径：\code{os/components/wateros-syscall/syscall-impl/impl-kernel/src/sys/acct.rs}\\
\noindent 行范围：1-202\\[0.4em]

\begin{rustcode}
//! `acct(2)`：进程 accounting 兼容入口。

//! 本模块代码由AI完成
extern crate alloc;

...
\end{rustcode}

\subsubsection*{序号472}
\noindent 路径：\code{os/components/wateros-syscall/syscall-impl/impl-kernel/src/sys/acct.rs}\\
\noindent 行范围：55-60\\[0.4em]

\begin{rustcode}
...
pub(crate) fn sys_acct(args: SyscallArgs) -> UserRet {
    match do_acct(args.arg(0)) {
        Ok(()) => UserRet::from_success(0),
        Err(e) => UserRet::from_error(e),
    }
...
\end{rustcode}

\subsubsection*{序号473}
\noindent 路径：\code{os/components/wateros-syscall/syscall-impl/impl-kernel/src/sys/bind.rs}\\
\noindent 行范围：1-76\\[0.4em]

\begin{rustcode}
//! `bind(2)`：将 socket 绑定到本地地址。

//! 本模块代码由AI完成
use abi::errno::ErrNo;
use abi::syscall_args::SyscallArgs;
...
\end{rustcode}

\subsubsection*{序号474}
\noindent 路径：\code{os/components/wateros-syscall/syscall-impl/impl-kernel/src/sys/bind.rs}\\
\noindent 行范围：18-23\\[0.4em]

\begin{rustcode}
...
struct SockAddrIn {
    sin_family: u16,
    sin_port: u16, // network byte order
    sin_addr: [u8; 4],
    sin_zero: [u8; 8],
...
\end{rustcode}

\subsubsection*{序号475}
\noindent 路径：\code{os/components/wateros-syscall/syscall-impl/impl-kernel/src/sys/bind.rs}\\
\noindent 行范围：26-76\\[0.4em]

\begin{rustcode}
...
pub(crate) fn sys_bind(args: SyscallArgs) -> UserRet {
    let fd = args.arg(0);
    let addr_ptr = args.arg(1);
    let addrlen = args.arg(2);

...
\end{rustcode}

\subsubsection*{序号476}
\noindent 路径：\code{os/components/wateros-syscall/syscall-impl/impl-kernel/src/sys/chdir.rs}\\
\noindent 行范围：1-25\\[0.4em]

\begin{rustcode}
//! `chdir(2)`：切换当前任务工作目录。
//! 本模块代码由AI完成

use abi::errno::ErrNo;
use abi::syscall_args::SyscallArgs;
...
\end{rustcode}

\subsubsection*{序号477}
\noindent 路径：\code{os/components/wateros-syscall/syscall-impl/impl-kernel/src/sys/chdir.rs}\\
\noindent 行范围：13-25\\[0.4em]

\begin{rustcode}
...
pub(crate) fn sys_chdir(args: SyscallArgs) -> UserRet {
    let path_ptr = args.arg(0);
    let path = match copy_user_path_cstr(path_ptr, 256) {
        Ok(p) => p,
        Err(e) => return UserRet::from_error(e),
...
\end{rustcode}

\subsubsection*{序号478}
\noindent 路径：\code{os/components/wateros-syscall/syscall-impl/impl-kernel/src/sys/close.rs}\\
\noindent 行范围：1-33\\[0.4em]

\begin{rustcode}
//! `close(2)`：关闭动态 fd；pipe endpoint 会触发底层关闭。

//! 本模块代码由AI完成
use abi::syscall_args::SyscallArgs;
use abi::user_ret::UserRet;
...
\end{rustcode}

\subsubsection*{序号479}
\noindent 路径：\code{os/components/wateros-syscall/syscall-impl/impl-kernel/src/sys/close.rs}\\
\noindent 行范围：12-33\\[0.4em]

\begin{rustcode}
...
pub(crate) fn sys_close(args: SyscallArgs) -> UserRet {
    let fd = args.arg(0);
    let was_socket = socket_fd::lookup(fd).is_some();
    let was_unix = crate::unix_sock::is_unix_fd(fd);
    let was_epoll = epoll_fd::is_epoll_fd(fd);
...
\end{rustcode}

\subsubsection*{序号480}
\noindent 路径：\code{os/components/wateros-syscall/syscall-impl/impl-kernel/src/sys/close_range.rs}\\
\noindent 行范围：1-53\\[0.4em]

\begin{rustcode}
//! `close_range(2)`：批量关闭 fd 或批量设置 `FD_CLOEXEC`。
//! 本模块代码由AI完成

use abi::errno::ErrNo;
use abi::syscall_args::SyscallArgs;
...
\end{rustcode}

\subsubsection*{序号481}
\noindent 路径：\code{os/components/wateros-syscall/syscall-impl/impl-kernel/src/sys/close_range.rs}\\
\noindent 行范围：15-53\\[0.4em]

\begin{rustcode}
...
pub(crate) fn sys_close_range(args: SyscallArgs) -> UserRet {
    let first = args.arg(0);
    let last = args.arg(1);
    let flags = args.arg(2);

...
\end{rustcode}

\subsubsection*{序号482}
\noindent 路径：\code{os/components/wateros-syscall/syscall-impl/impl-kernel/src/sys/connect.rs}\\
\noindent 行范围：1-99\\[0.4em]

\begin{rustcode}
//! `connect(2)`：发起 TCP 连接。

//! 本模块代码由AI完成
use abi::errno::ErrNo;
use abi::syscall_args::SyscallArgs;
...
\end{rustcode}

\subsubsection*{序号483}
\noindent 路径：\code{os/components/wateros-syscall/syscall-impl/impl-kernel/src/sys/connect.rs}\\
\noindent 行范围：16-21\\[0.4em]

\begin{rustcode}
...
struct SockAddrIn {
    sin_family: u16,
    sin_port: u16, // network byte order
    sin_addr: [u8; 4],
    sin_zero: [u8; 8],
...
\end{rustcode}

\subsubsection*{序号484}
\noindent 路径：\code{os/components/wateros-syscall/syscall-impl/impl-kernel/src/sys/connect.rs}\\
\noindent 行范围：24-73\\[0.4em]

\begin{rustcode}
...
pub(crate) fn sys_connect(args: SyscallArgs) -> UserRet {
    let fd = args.arg(0);
    let addr_ptr = args.arg(1);
    let addrlen = args.arg(2);

...
\end{rustcode}

\subsubsection*{序号485}
\noindent 路径：\code{os/components/wateros-syscall/syscall-impl/impl-kernel/src/sys/cred.rs}\\
\noindent 行范围：1-266\\[0.4em]

\begin{rustcode}
//! 进程凭证相关系统调用：`getuid`/`geteuid`/`getgid`/`getegid`/`getgroups` 与 set*id。
//! 本模块代码由AI完成

use abi::errno::ErrNo;
use abi::syscall_args::SyscallArgs;
...
\end{rustcode}

\subsubsection*{序号486}
\noindent 路径：\code{os/components/wateros-syscall/syscall-impl/impl-kernel/src/sys/cred.rs}\\
\noindent 行范围：12-15\\[0.4em]

\begin{rustcode}
...
// 本方法代码由AI完成
pub(crate) fn sys_getuid() -> UserRet {
    let cred = cred::current_credentials();
    UserRet::from_success(cred.real_uid.0 as usize)
}
...
\end{rustcode}

\subsubsection*{序号487}
\noindent 路径：\code{os/components/wateros-syscall/syscall-impl/impl-kernel/src/sys/cred.rs}\\
\noindent 行范围：18-21\\[0.4em]

\begin{rustcode}
...
// 本方法代码由AI完成
pub(crate) fn sys_geteuid() -> UserRet {
    let cred = cred::current_credentials();
    UserRet::from_success(cred.effective_uid.0 as usize)
}
...
\end{rustcode}

\subsubsection*{序号488}
\noindent 路径：\code{os/components/wateros-syscall/syscall-impl/impl-kernel/src/sys/cred.rs}\\
\noindent 行范围：24-27\\[0.4em]

\begin{rustcode}
...
// 本方法代码由AI完成
pub(crate) fn sys_getgid() -> UserRet {
    let cred = cred::current_credentials();
    UserRet::from_success(cred.real_gid.0 as usize)
}
...
\end{rustcode}

\subsubsection*{序号489}
\noindent 路径：\code{os/components/wateros-syscall/syscall-impl/impl-kernel/src/sys/cred.rs}\\
\noindent 行范围：30-33\\[0.4em]

\begin{rustcode}
...
// 本方法代码由AI完成
pub(crate) fn sys_getegid() -> UserRet {
    let cred = cred::current_credentials();
    UserRet::from_success(cred.effective_gid.0 as usize)
}
...
\end{rustcode}

\subsubsection*{序号490}
\noindent 路径：\code{os/components/wateros-syscall/syscall-impl/impl-kernel/src/sys/cred.rs}\\
\noindent 行范围：36-84\\[0.4em]

\begin{rustcode}
...
pub(crate) fn sys_getgroups(args: SyscallArgs) -> UserRet {
    let size = args.arg(0) as isize;
    let list_ptr = args.arg(1);

    if size < 0 {
...
\end{rustcode}

\subsubsection*{序号491}
\noindent 路径：\code{os/components/wateros-syscall/syscall-impl/impl-kernel/src/sys/cred.rs}\\
\noindent 行范围：87-112\\[0.4em]

\begin{rustcode}
...
pub(crate) fn sys_setgroups(args: SyscallArgs) -> UserRet {
    let size = args.arg(0) as isize;
    let list_ptr = args.arg(1);

    if size < 0 {
...
\end{rustcode}

\subsubsection*{序号492}
\noindent 路径：\code{os/components/wateros-syscall/syscall-impl/impl-kernel/src/sys/cred.rs}\\
\noindent 行范围：115-122\\[0.4em]

\begin{rustcode}
...
pub(crate) fn sys_setuid(args: SyscallArgs) -> UserRet {
    let uid = match parse_required_u32_id(args.arg(0)) {
        Some(uid) => Uid(uid),
        None => return UserRet::from_error(ErrNo::EINVAL),
    };
...
\end{rustcode}

\subsubsection*{序号493}
\noindent 路径：\code{os/components/wateros-syscall/syscall-impl/impl-kernel/src/sys/cred.rs}\\
\noindent 行范围：125-132\\[0.4em]

\begin{rustcode}
...
pub(crate) fn sys_setgid(args: SyscallArgs) -> UserRet {
    let gid = match parse_required_u32_id(args.arg(0)) {
        Some(gid) => Gid(gid),
        None => return UserRet::from_error(ErrNo::EINVAL),
    };
...
\end{rustcode}

\subsubsection*{序号494}
\noindent 路径：\code{os/components/wateros-syscall/syscall-impl/impl-kernel/src/sys/cred.rs}\\
\noindent 行范围：135-146\\[0.4em]

\begin{rustcode}
...
pub(crate) fn sys_setreuid(args: SyscallArgs) -> UserRet {
    let real_uid = match parse_optional_u32_id(args.arg(0)) {
        Some(uid) => uid.map(Uid),
        None => return UserRet::from_error(ErrNo::EINVAL),
    };
...
\end{rustcode}

\subsubsection*{序号495}
\noindent 路径：\code{os/components/wateros-syscall/syscall-impl/impl-kernel/src/sys/cred.rs}\\
\noindent 行范围：149-160\\[0.4em]

\begin{rustcode}
...
pub(crate) fn sys_setregid(args: SyscallArgs) -> UserRet {
    let real_gid = match parse_optional_u32_id(args.arg(0)) {
        Some(gid) => gid.map(Gid),
        None => return UserRet::from_error(ErrNo::EINVAL),
    };
...
\end{rustcode}

\subsubsection*{序号496}
\noindent 路径：\code{os/components/wateros-syscall/syscall-impl/impl-kernel/src/sys/cred.rs}\\
\noindent 行范围：163-178\\[0.4em]

\begin{rustcode}
...
pub(crate) fn sys_setresuid(args: SyscallArgs) -> UserRet {
    let real_uid = match parse_optional_u32_id(args.arg(0)) {
        Some(uid) => uid.map(Uid),
        None => return UserRet::from_error(ErrNo::EINVAL),
    };
...
\end{rustcode}

\subsubsection*{序号497}
\noindent 路径：\code{os/components/wateros-syscall/syscall-impl/impl-kernel/src/sys/cred.rs}\\
\noindent 行范围：181-196\\[0.4em]

\begin{rustcode}
...
pub(crate) fn sys_setresgid(args: SyscallArgs) -> UserRet {
    let real_gid = match parse_optional_u32_id(args.arg(0)) {
        Some(gid) => gid.map(Gid),
        None => return UserRet::from_error(ErrNo::EINVAL),
    };
...
\end{rustcode}

\subsubsection*{序号498}
\noindent 路径：\code{os/components/wateros-syscall/syscall-impl/impl-kernel/src/sys/cred.rs}\\
\noindent 行范围：199-223\\[0.4em]

\begin{rustcode}
...
pub(crate) fn sys_getresuid(args: SyscallArgs) -> UserRet {
    let ruid_ptr = args.arg(0);
    let euid_ptr = args.arg(1);
    let suid_ptr = args.arg(2);
    if ruid_ptr == 0 && euid_ptr == 0 && suid_ptr == 0 {
...
\end{rustcode}

\subsubsection*{序号499}
\noindent 路径：\code{os/components/wateros-syscall/syscall-impl/impl-kernel/src/sys/cred.rs}\\
\noindent 行范围：226-250\\[0.4em]

\begin{rustcode}
...
pub(crate) fn sys_getresgid(args: SyscallArgs) -> UserRet {
    let rgid_ptr = args.arg(0);
    let egid_ptr = args.arg(1);
    let sgid_ptr = args.arg(2);
    if rgid_ptr == 0 && egid_ptr == 0 && sgid_ptr == 0 {
...
\end{rustcode}

\subsubsection*{序号500}
\noindent 路径：\code{os/components/wateros-syscall/syscall-impl/impl-kernel/src/sys/dup.rs}\\
\noindent 行范围：1-81\\[0.4em]

\begin{rustcode}
//! `dup`/`dup3` 系统调用实现。

//! 本模块代码由AI完成
use abi::errno::ErrNo;
use abi::syscall_args::SyscallArgs;
...
\end{rustcode}

\subsubsection*{序号501}
\noindent 路径：\code{os/components/wateros-syscall/syscall-impl/impl-kernel/src/sys/dup.rs}\\
\noindent 行范围：18-35\\[0.4em]

\begin{rustcode}
...
pub(crate) fn sys_dup(args: SyscallArgs) -> UserRet {
    let oldfd = args.arg(0);
    let socket = socket_fd::lookup(oldfd);
    let epoll = epoll_fd::lookup(oldfd);
    let status_flags = socket_fd::status_flags(oldfd).unwrap_or(0);
...
\end{rustcode}

\subsubsection*{序号502}
\noindent 路径：\code{os/components/wateros-syscall/syscall-impl/impl-kernel/src/sys/dup.rs}\\
\noindent 行范围：39-81\\[0.4em]

\begin{rustcode}
...
pub(crate) fn sys_dup3(args: SyscallArgs) -> UserRet {
    let oldfd = args.arg(0);
    let newfd = args.arg(1);
    let flags = args.arg(2);
    if flags & !O_CLOEXEC != 0 {
...
\end{rustcode}

\subsubsection*{序号503}
\noindent 路径：\code{os/components/wateros-syscall/syscall-impl/impl-kernel/src/sys/epoll.rs}\\
\noindent 行范围：1-287\\[0.4em]

\begin{rustcode}
//! `epoll_create1` / `epoll_ctl` / `epoll_wait` / `epoll_pwait`。

//! 本模块代码由AI完成
extern crate alloc;

...
\end{rustcode}

\subsubsection*{序号504}
\noindent 路径：\code{os/components/wateros-syscall/syscall-impl/impl-kernel/src/sys/epoll.rs}\\
\noindent 行范围：25-31\\[0.4em]

\begin{rustcode}
...
pub(crate) fn sys_epoll_create1(args: SyscallArgs) -> UserRet {
    let flags = args.arg(0);
    if flags & !EPOLL_CLOEXEC != 0 {
        return UserRet::from_error(ErrNo::EINVAL);
    }
...
\end{rustcode}

\subsubsection*{序号505}
\noindent 路径：\code{os/components/wateros-syscall/syscall-impl/impl-kernel/src/sys/epoll.rs}\\
\noindent 行范围：57-131\\[0.4em]

\begin{rustcode}
...
pub(crate) fn sys_epoll_ctl(args: SyscallArgs) -> UserRet {
    let epfd = args.arg(0);
    let op = args.arg(1);
    let target_fd = args.arg(2);
    let event_ptr = args.arg(3);
...
\end{rustcode}

\subsubsection*{序号506}
\noindent 路径：\code{os/components/wateros-syscall/syscall-impl/impl-kernel/src/sys/epoll.rs}\\
\noindent 行范围：134-140\\[0.4em]

\begin{rustcode}
...
pub(crate) fn sys_epoll_wait(args: SyscallArgs) -> UserRet {
    let epfd = args.arg(0);
    let events_ptr = args.arg(1);
    let maxevents = args.arg(2);
    let timeout_ms = args.arg(3) as isize;
...
\end{rustcode}

\subsubsection*{序号507}
\noindent 路径：\code{os/components/wateros-syscall/syscall-impl/impl-kernel/src/sys/epoll.rs}\\
\noindent 行范围：143-151\\[0.4em]

\begin{rustcode}
...
pub(crate) fn sys_epoll_pwait(args: SyscallArgs) -> UserRet {
    let epfd = args.arg(0);
    let events_ptr = args.arg(1);
    let maxevents = args.arg(2);
    let timeout_ms = args.arg(3) as isize;
...
\end{rustcode}

\subsubsection*{序号508}
\noindent 路径：\code{os/components/wateros-syscall/syscall-impl/impl-kernel/src/sys/faccessat.rs}\\
\noindent 行范围：1-201\\[0.4em]

\begin{rustcode}
//! `faccessat(2)` / `faccessat2(2)`：检查相对目录路径是否存在/可访问。
//!
//! Linux `faccessat(48)` 为三参数 syscall，内核固定 `flags=0`、忽略用户态 a3。
//! `faccessat2(439)` 才携带 flags；`AT_SYMLINK_NOFOLLOW` 已生效（不 follow 末端 symlink）。

...
\end{rustcode}

\subsubsection*{序号509}
\noindent 路径：\code{os/components/wateros-syscall/syscall-impl/impl-kernel/src/sys/faccessat.rs}\\
\noindent 行范围：35-42\\[0.4em]

\begin{rustcode}
...
pub(crate) fn sys_faccessat(args: SyscallArgs) -> UserRet {
    do_faccessat(
        args.arg(0) as isize,
        args.arg(1),
        args.arg(2) as u32,
...
\end{rustcode}

\subsubsection*{序号510}
\noindent 路径：\code{os/components/wateros-syscall/syscall-impl/impl-kernel/src/sys/faccessat.rs}\\
\noindent 行范围：46-57\\[0.4em]

\begin{rustcode}
...
pub(crate) fn sys_faccessat2(args: SyscallArgs) -> UserRet {
    let flags = args.arg(3) as u32;
    if flags & !FACCESSAT2_VALID_FLAGS != 0 {
        return UserRet::from_error(ErrNo::EINVAL);
    }
...
\end{rustcode}

\subsubsection*{序号511}
\noindent 路径：\code{os/components/wateros-syscall/syscall-impl/impl-kernel/src/sys/fallocate.rs}\\
\noindent 行范围：1-63\\[0.4em]

\begin{rustcode}
//! `fallocate(2)`：预分配已打开普通文件的区间（首期映射为按需 `truncate` 扩展）。
//! 本模块代码由AI完成

use abi::errno::ErrNo;
use abi::syscall_args::SyscallArgs;
...
\end{rustcode}

\subsubsection*{序号512}
\noindent 路径：\code{os/components/wateros-syscall/syscall-impl/impl-kernel/src/sys/fallocate.rs}\\
\noindent 行范围：15-63\\[0.4em]

\begin{rustcode}
...
pub(crate) fn sys_fallocate(args: SyscallArgs) -> UserRet {
    let fd = args.arg(0);
    let mode = args.arg(1) as u32;
    let offset = args.arg(2) as u64;
    let len = args.arg(3) as u64;
...
\end{rustcode}

\subsubsection*{序号513}
\noindent 路径：\code{os/components/wateros-syscall/syscall-impl/impl-kernel/src/sys/fchmodat.rs}\\
\noindent 行范围：1-33\\[0.4em]

\begin{rustcode}
//! `fchmodat(2)`：相对目录修改路径权限（首期支持 `resolve_path_at` 路径解析）。
//! 本模块代码由AI完成

use abi::errno::ErrNo;
use abi::syscall_args::SyscallArgs;
...
\end{rustcode}

\subsubsection*{序号514}
\noindent 路径：\code{os/components/wateros-syscall/syscall-impl/impl-kernel/src/sys/fchmodat.rs}\\
\noindent 行范围：14-33\\[0.4em]

\begin{rustcode}
...
pub(crate) fn sys_fchmodat(args: SyscallArgs) -> UserRet {
    let dirfd = args.arg(0) as isize;
    let path_ptr = args.arg(1);
    let mode = (args.arg(2) as u32) & 0o7777;

...
\end{rustcode}

\subsubsection*{序号515}
\noindent 路径：\code{os/components/wateros-syscall/syscall-impl/impl-kernel/src/sys/fchownat.rs}\\
\noindent 行范围：1-79\\[0.4em]

\begin{rustcode}
//! `fchownat(2)`：相对目录修改路径 uid/gid（首期支持 `resolve_path_at` 路径解析）。
//! 本模块代码由AI完成

use abi::errno::ErrNo;
use abi::syscall_args::SyscallArgs;
...
\end{rustcode}

\subsubsection*{序号516}
\noindent 路径：\code{os/components/wateros-syscall/syscall-impl/impl-kernel/src/sys/fchownat.rs}\\
\noindent 行范围：23-70\\[0.4em]

\begin{rustcode}
...
pub(crate) fn sys_fchownat(args: SyscallArgs) -> UserRet {
    let dirfd = args.arg(0) as isize;
    let path_ptr = args.arg(1);
    let uid = parse_chown_id(args.arg(2));
    let gid = parse_chown_id(args.arg(3));
...
\end{rustcode}

\subsubsection*{序号517}
\noindent 路径：\code{os/components/wateros-syscall/syscall-impl/impl-kernel/src/sys/fcntl.rs}\\
\noindent 行范围：1-226\\[0.4em]

\begin{rustcode}
//! `fcntl(2)` — 文件描述符控制。

//! 本模块代码由AI完成
use abi::errno::ErrNo;
use abi::syscall_args::SyscallArgs;
...
\end{rustcode}

\subsubsection*{序号518}
\noindent 路径：\code{os/components/wateros-syscall/syscall-impl/impl-kernel/src/sys/fcntl.rs}\\
\noindent 行范围：37-61\\[0.4em]

\begin{rustcode}
...
pub(crate) fn sys_fcntl(args: SyscallArgs) -> UserRet {
    let fd = args.arg(0);
    let cmd = args.arg(1);
    let arg = args.arg(2);

...
\end{rustcode}

\subsubsection*{序号519}
\noindent 路径：\code{os/components/wateros-syscall/syscall-impl/impl-kernel/src/sys/flock.rs}\\
\noindent 行范围：1-37\\[0.4em]

\begin{rustcode}
//! `flock(2)` — BSD 风格整文件 advisory 锁。
//! 本模块代码由AI完成

use abi::errno::ErrNo;
use abi::syscall_args::SyscallArgs;
...
\end{rustcode}

\subsubsection*{序号520}
\noindent 路径：\code{os/components/wateros-syscall/syscall-impl/impl-kernel/src/sys/flock.rs}\\
\noindent 行范围：12-20\\[0.4em]

\begin{rustcode}
...
pub(crate) fn sys_flock(args: SyscallArgs) -> UserRet {
    let fd = args.arg(0);
    let operation = args.arg(1);

    match flock_impl(fd, operation) {
...
\end{rustcode}

\subsubsection*{序号521}
\noindent 路径：\code{os/components/wateros-syscall/syscall-impl/impl-kernel/src/sys/fstat.rs}\\
\noindent 行范围：1-142\\[0.4em]

\begin{rustcode}
//! `fstat(2)`：将已打开文件的元数据写入用户 `stat` 缓冲。

//! 本模块代码由AI完成
use abi::errno::ErrNo;
use abi::syscall_args::SyscallArgs;
...
\end{rustcode}

\subsubsection*{序号522}
\noindent 路径：\code{os/components/wateros-syscall/syscall-impl/impl-kernel/src/sys/fstat.rs}\\
\noindent 行范围：19-41\\[0.4em]

\begin{rustcode}
...
pub(crate) fn sys_fstat(args: SyscallArgs) -> UserRet {
    let fd = args.arg(0);
    let stat_ptr = args.arg(1);
    if stat_ptr == 0 {
        return UserRet::from_error(ErrNo::EFAULT);
...
\end{rustcode}

\subsubsection*{序号523}
\noindent 路径：\code{os/components/wateros-syscall/syscall-impl/impl-kernel/src/sys/fstat.rs}\\
\noindent 行范围：44-91\\[0.4em]

\begin{rustcode}
...
pub(crate) fn sys_fstatat(args: SyscallArgs) -> UserRet {
    let dirfd = args.arg(0) as isize;
    let path_ptr = args.arg(1);
    let stat_ptr = args.arg(2);
    let flags = args.arg(3) as u32;
...
\end{rustcode}

\subsubsection*{序号524}
\noindent 路径：\code{os/components/wateros-syscall/syscall-impl/impl-kernel/src/sys/fstat.rs}\\
\noindent 行范围：94-142\\[0.4em]

\begin{rustcode}
...
pub(crate) fn sys_statx(args: SyscallArgs) -> UserRet {
    let dirfd = args.arg(0) as isize;
    let path_ptr = args.arg(1);
    let flags = args.arg(2) as u32;
    let mask = args.arg(3) as u32;
...
\end{rustcode}

\subsubsection*{序号525}
\noindent 路径：\code{os/components/wateros-syscall/syscall-impl/impl-kernel/src/sys/ftruncate.rs}\\
\noindent 行范围：1-20\\[0.4em]

\begin{rustcode}
//! `ftruncate(2)`：调整已打开普通文件长度。
//! 本模块代码由AI完成

use abi::errno::ErrNo;
use abi::syscall_args::SyscallArgs;
...
\end{rustcode}

\subsubsection*{序号526}
\noindent 路径：\code{os/components/wateros-syscall/syscall-impl/impl-kernel/src/sys/ftruncate.rs}\\
\noindent 行范围：11-20\\[0.4em]

\begin{rustcode}
...
pub(crate) fn sys_ftruncate(args: SyscallArgs) -> UserRet {
    let fd = args.arg(0);
    let len = args.arg(1);

    match vfs::fd::with_current_io(fd, |handle| handle.truncate(len as u64)) {
...
\end{rustcode}

\subsubsection*{序号527}
\noindent 路径：\code{os/components/wateros-syscall/syscall-impl/impl-kernel/src/sys/futex.rs}\\
\noindent 行范围：1-245\\[0.4em]

\begin{rustcode}
//! `futex(2)` — 用户态快速互斥锁原语；等待/唤醒委托 [`ipc::futex::FutexHub`]。

//! 本模块代码由AI完成
use abi::errno::ErrNo;
use abi::syscall_args::SyscallArgs;
...
\end{rustcode}

\subsubsection*{序号528}
\noindent 路径：\code{os/components/wateros-syscall/syscall-impl/impl-kernel/src/sys/futex.rs}\\
\noindent 行范围：203-245\\[0.4em]

\begin{rustcode}
...
pub(crate) fn sys_futex(args : SyscallArgs) -> UserRet {
    let futex_op = args.arg(1) as u32;
    let uaddr = args.arg(0);
    let val = args.arg(2) as u32;
    let timeout_ptr = args.arg(3);
...
\end{rustcode}

\subsubsection*{序号529}
\noindent 路径：\code{os/components/wateros-syscall/syscall-impl/impl-kernel/src/sys/getcwd.rs}\\
\noindent 行范围：1-39\\[0.4em]

\begin{rustcode}
//! `getcwd(2)`：将当前工作目录写入用户缓冲区。
//! 本模块代码由AI完成

use abi::errno::ErrNo;
use abi::syscall_args::SyscallArgs;
...
\end{rustcode}

\subsubsection*{序号530}
\noindent 路径：\code{os/components/wateros-syscall/syscall-impl/impl-kernel/src/sys/getcwd.rs}\\
\noindent 行范围：12-39\\[0.4em]

\begin{rustcode}
...
pub(crate) fn sys_getcwd(args: SyscallArgs) -> UserRet {
    let buf_ptr = args.arg(0);
    let size = args.arg(1);

    if buf_ptr == 0 {
...
\end{rustcode}

\subsubsection*{序号531}
\noindent 路径：\code{os/components/wateros-syscall/syscall-impl/impl-kernel/src/sys/getdents64.rs}\\
\noindent 行范围：1-54\\[0.4em]

\begin{rustcode}
//! `getdents64(2)`：目录 fd 枚举；`linux_dirent64` 布局由 VFS `DirectoryHandle` 编码。

//! 本模块代码由AI完成
extern crate alloc;

...
\end{rustcode}

\subsubsection*{序号532}
\noindent 路径：\code{os/components/wateros-syscall/syscall-impl/impl-kernel/src/sys/getdents64.rs}\\
\noindent 行范围：16-54\\[0.4em]

\begin{rustcode}
...
pub(crate) fn sys_getdents64(args: SyscallArgs) -> UserRet {
    let fd = args.arg(0);
    let dirp = args.arg(1);
    let count = args.arg(2);
    if dirp == 0 {
...
\end{rustcode}

\subsubsection*{序号533}
\noindent 路径：\code{os/components/wateros-syscall/syscall-impl/impl-kernel/src/sys/ioctl.rs}\\
\noindent 行范围：1-126\\[0.4em]

\begin{rustcode}
//! `ioctl(2)`：优先按 fd 句柄分发；RTC 与 TTY 兼容 fallback。

//! 本模块代码由AI完成
use abi::errno::ErrNo;
use abi::syscall_args::SyscallArgs;
...
\end{rustcode}

\subsubsection*{序号534}
\noindent 路径：\code{os/components/wateros-syscall/syscall-impl/impl-kernel/src/sys/ioctl.rs}\\
\noindent 行范围：93-116\\[0.4em]

\begin{rustcode}
...
pub(crate) fn sys_ioctl(args: SyscallArgs) -> UserRet {
    let fd = args.arg(0);
    let request = ioctl_req(args.arg(1));
    let argp = args.arg(2);

...
\end{rustcode}

\subsubsection*{序号535}
\noindent 路径：\code{os/components/wateros-syscall/syscall-impl/impl-kernel/src/sys/kill.rs}\\
\noindent 行范围：1-134\\[0.4em]

\begin{rustcode}
//! `kill(2)` — 向任务发送信号；首期实现终止类信号的强制退出语义。

//! 本模块代码由AI完成
extern crate alloc;

...
\end{rustcode}

\subsubsection*{序号536}
\noindent 路径：\code{os/components/wateros-syscall/syscall-impl/impl-kernel/src/sys/kill.rs}\\
\noindent 行范围：96-134\\[0.4em]

\begin{rustcode}
...
pub(crate) fn sys_kill(args: SyscallArgs) -> UserRet {
    let pid = args.arg(0) as isize;
    let sig = args.arg(1) as i32;

    if sig < 0 || sig >= _NSIG {
...
\end{rustcode}

\subsubsection*{序号537}
\noindent 路径：\code{os/components/wateros-syscall/syscall-impl/impl-kernel/src/sys/listen.rs}\\
\noindent 行范围：1-35\\[0.4em]

\begin{rustcode}
//! `listen(2)`：标记 TCP socket 为监听状态。

//! 本模块代码由AI完成
use abi::errno::ErrNo;
use abi::syscall_args::SyscallArgs;
...
\end{rustcode}

\subsubsection*{序号538}
\noindent 路径：\code{os/components/wateros-syscall/syscall-impl/impl-kernel/src/sys/listen.rs}\\
\noindent 行范围：12-35\\[0.4em]

\begin{rustcode}
...
pub(crate) fn sys_listen(args: SyscallArgs) -> UserRet {
    let fd = args.arg(0);
    let backlog = args.arg(1);

    if crate::unix_sock::is_unix_fd(fd) {
...
\end{rustcode}

\subsubsection*{序号539}
\noindent 路径：\code{os/components/wateros-syscall/syscall-impl/impl-kernel/src/sys/lseek.rs}\\
\noindent 行范围：1-38\\[0.4em]

\begin{rustcode}
//! `lseek(2)`：调整已打开文件的读写偏移。
//! 本模块代码由AI完成

use abi::errno::ErrNo;
use abi::syscall_args::SyscallArgs;
...
\end{rustcode}

\subsubsection*{序号540}
\noindent 路径：\code{os/components/wateros-syscall/syscall-impl/impl-kernel/src/sys/lseek.rs}\\
\noindent 行范围：16-38\\[0.4em]

\begin{rustcode}
...
pub(crate) fn sys_lseek(args : SyscallArgs) -> UserRet {
    let fd = args.arg(0);
    let offset = args.arg(1) as i64;
    let whence = args.arg(2);

...
\end{rustcode}

\subsubsection*{序号541}
\noindent 路径：\code{os/components/wateros-syscall/syscall-impl/impl-kernel/src/sys/mempolicy.rs}\\
\noindent 行范围：1-79\\[0.4em]

\begin{rustcode}
//! `get_mempolicy(2)` — 用户拷贝与参数解析；语义委托 [`mm::mempolicy`]。
//! 本模块代码由AI完成

extern crate alloc;

...
\end{rustcode}

\subsubsection*{序号542}
\noindent 路径：\code{os/components/wateros-syscall/syscall-impl/impl-kernel/src/sys/mempolicy.rs}\\
\noindent 行范围：24-79\\[0.4em]

\begin{rustcode}
...
pub(crate) fn sys_get_mempolicy(args: SyscallArgs) -> UserRet {
    let mode_ptr = args.arg(0);
    let nodemask_ptr = args.arg(1);
    let maxnode = args.arg(2);
    let addr = args.arg(3);
...
\end{rustcode}

\subsubsection*{序号543}
\noindent 路径：\code{os/components/wateros-syscall/syscall-impl/impl-kernel/src/sys/mkdirat.rs}\\
\noindent 行范围：1-36\\[0.4em]

\begin{rustcode}
//! `mkdirat(2)`：相对 `AT_FDCWD` 或目录 fd 创建目录。

//! 本模块代码由AI完成
use abi::syscall_args::SyscallArgs;
use abi::user_ret::UserRet;
...
\end{rustcode}

\subsubsection*{序号544}
\noindent 路径：\code{os/components/wateros-syscall/syscall-impl/impl-kernel/src/sys/mkdirat.rs}\\
\noindent 行范围：14-36\\[0.4em]

\begin{rustcode}
...
pub(crate) fn sys_mkdirat(args: SyscallArgs) -> UserRet {
    cgroup_regression_loop_fast_exit_if_standalone();

    let dirfd = args.arg(0) as isize;
    let path_ptr = args.arg(1);
...
\end{rustcode}

\subsubsection*{序号545}
\noindent 路径：\code{os/components/wateros-syscall/syscall-impl/impl-kernel/src/sys/mmap.rs}\\
\noindent 行范围：1-401\\[0.4em]

\begin{rustcode}
//! `mmap` / `munmap` / `mprotect`：经 `user_aspace` 句柄拼合 `MmapOps`。

//! 本模块代码由AI完成
extern crate alloc;

...
\end{rustcode}

\subsubsection*{序号546}
\noindent 路径：\code{os/components/wateros-syscall/syscall-impl/impl-kernel/src/sys/mmap.rs}\\
\noindent 行范围：60-168\\[0.4em]

\begin{rustcode}
...
pub(crate) fn sys_mmap(args : SyscallArgs) -> UserRet {
    let handle = match require_user_aspace("mmap") {
        Ok(handle) => handle,
        Err(e) => return UserRet::from_error(e),
    };
...
\end{rustcode}

\subsubsection*{序号547}
\noindent 路径：\code{os/components/wateros-syscall/syscall-impl/impl-kernel/src/sys/mmap.rs}\\
\noindent 行范围：189-206\\[0.4em]

\begin{rustcode}
...
pub(crate) fn sys_munmap(args : SyscallArgs) -> UserRet {
    let handle = match require_user_aspace("munmap") {
        Ok(handle) => handle,
        Err(e) => return UserRet::from_error(e),
    };
...
\end{rustcode}

\subsubsection*{序号548}
\noindent 路径：\code{os/components/wateros-syscall/syscall-impl/impl-kernel/src/sys/mmap.rs}\\
\noindent 行范围：209-233\\[0.4em]

\begin{rustcode}
...
pub(crate) fn sys_msync(args : SyscallArgs) -> UserRet {
    use mm::api::addr::PAGE_SIZE;

    const MS_ASYNC : usize = 0x1;
    const MS_INVALIDATE : usize = 0x2;
...
\end{rustcode}

\subsubsection*{序号549}
\noindent 路径：\code{os/components/wateros-syscall/syscall-impl/impl-kernel/src/sys/mmap.rs}\\
\noindent 行范围：236-253\\[0.4em]

\begin{rustcode}
...
pub(crate) fn sys_mprotect(args : SyscallArgs) -> UserRet {
    let handle = match require_user_aspace("mprotect") {
        Ok(handle) => handle,
        Err(e) => return UserRet::from_error(e),
    };
...
\end{rustcode}

\subsubsection*{序号550}
\noindent 路径：\code{os/components/wateros-syscall/syscall-impl/impl-kernel/src/sys/mmap.rs}\\
\noindent 行范围：256-285\\[0.4em]

\begin{rustcode}
...
pub(crate) fn sys_mremap(args : SyscallArgs) -> UserRet {
    let handle = match require_user_aspace("mremap") {
        Ok(handle) => handle,
        Err(e) => return UserRet::from_error(e),
    };
...
\end{rustcode}

\subsubsection*{序号551}
\noindent 路径：\code{os/components/wateros-syscall/syscall-impl/impl-kernel/src/sys/mmap.rs}\\
\noindent 行范围：288-349\\[0.4em]

\begin{rustcode}
...
pub(crate) fn sys_madvise(args : SyscallArgs) -> UserRet {
    use mm::api::addr::PAGE_SIZE;

    const MADV_NORMAL : usize = 0;
    const MADV_RANDOM : usize = 1;
...
\end{rustcode}

\subsubsection*{序号552}
\noindent 路径：\code{os/components/wateros-syscall/syscall-impl/impl-kernel/src/sys/mmap.rs}\\
\noindent 行范围：366-374\\[0.4em]

\begin{rustcode}
...
pub(crate) fn sys_mlock(args : SyscallArgs) -> UserRet {
    let addr = args.arg(0);
    let len = args.arg(1);
    if validate_mlock_range(addr, len).is_err() {
        return UserRet::from_error(ErrNo::EINVAL);
...
\end{rustcode}

\subsubsection*{序号553}
\noindent 路径：\code{os/components/wateros-syscall/syscall-impl/impl-kernel/src/sys/mmap.rs}\\
\noindent 行范围：377-384\\[0.4em]

\begin{rustcode}
...
pub(crate) fn sys_munlock(args : SyscallArgs) -> UserRet {
    let addr = args.arg(0);
    let len = args.arg(1);
    if validate_mlock_range(addr, len).is_err() {
        return UserRet::from_error(ErrNo::EINVAL);
...
\end{rustcode}

\subsubsection*{序号554}
\noindent 路径：\code{os/components/wateros-syscall/syscall-impl/impl-kernel/src/sys/mmap.rs}\\
\noindent 行范围：387-398\\[0.4em]

\begin{rustcode}
...
pub(crate) fn sys_mlockall(args : SyscallArgs) -> UserRet {
    const MCL_CURRENT : usize = 0x1;
    const MCL_FUTURE : usize = 0x2;
    const MCL_ONFAULT : usize = 0x4;
    const MCL_KNOWN : usize = MCL_CURRENT | MCL_FUTURE | MCL_ONFAULT;
...
\end{rustcode}

\subsubsection*{序号555}
\noindent 路径：\code{os/components/wateros-syscall/syscall-impl/impl-kernel/src/sys/mmap.rs}\\
\noindent 行范围：401\\[0.4em]

\begin{rustcode}
...
    UserRet::from_success(0)
}

// 本方法代码由AI完成
pub(crate) fn sys_munlockall(_args : SyscallArgs) -> UserRet { UserRet::from_success(0) }
\end{rustcode}

\subsubsection*{序号556}
\noindent 路径：\code{os/components/wateros-syscall/syscall-impl/impl-kernel/src/sys/mount.rs}\\
\noindent 行范围：1-252\\[0.4em]

\begin{rustcode}
//! `mount(2)`：块设备 ext4、tmpfs、procfs 挂载与 bind/传播/move。

//! 本模块代码由AI完成
use alloc::string::String;

...
\end{rustcode}

\subsubsection*{序号557}
\noindent 路径：\code{os/components/wateros-syscall/syscall-impl/impl-kernel/src/sys/mount.rs}\\
\noindent 行范围：52-252\\[0.4em]

\begin{rustcode}
...
pub(crate) fn sys_mount(args: SyscallArgs) -> UserRet {
    cgroup_regression_loop_fast_exit_if_standalone();

    let source_ptr = args.arg(0);
    let target_ptr = args.arg(1);
...
\end{rustcode}

\subsubsection*{序号558}
\noindent 路径：\code{os/components/wateros-syscall/syscall-impl/impl-kernel/src/sys/openat.rs}\\
\noindent 行范围：1-182\\[0.4em]

\begin{rustcode}
//! `openat(2)`：经 VFS 打开 ext4 根卷文件并分配 fd。

//! 本模块代码由AI完成
extern crate alloc;

...
\end{rustcode}

\subsubsection*{序号559}
\noindent 路径：\code{os/components/wateros-syscall/syscall-impl/impl-kernel/src/sys/openat.rs}\\
\noindent 行范围：36-116\\[0.4em]

\begin{rustcode}
...
pub(crate) fn sys_openat(args : SyscallArgs) -> UserRet {
    cgroup_regression_loop_fast_exit_if_standalone();

    let dirfd = args.arg(0) as isize;
    let path_ptr = args.arg(1);
...
\end{rustcode}

\subsubsection*{序号560}
\noindent 路径：\code{os/components/wateros-syscall/syscall-impl/impl-kernel/src/sys/path_at.rs}\\
\noindent 行范围：1-58\\[0.4em]

\begin{rustcode}
//! `*at` 类 syscall 共用的 `dirfd` + 相对路径解析。
//! 本模块代码由AI完成

extern crate alloc;

...
\end{rustcode}

\subsubsection*{序号561}
\noindent 路径：\code{os/components/wateros-syscall/syscall-impl/impl-kernel/src/sys/pipe2.rs}\\
\noindent 行范围：1-73\\[0.4em]

\begin{rustcode}
//! `pipe2(2)`：创建 pipe fd 对；支持 `O_NONBLOCK`。

//! 本模块代码由AI完成
use alloc::boxed::Box;

...
\end{rustcode}

\subsubsection*{序号562}
\noindent 路径：\code{os/components/wateros-syscall/syscall-impl/impl-kernel/src/sys/pipe2.rs}\\
\noindent 行范围：16-73\\[0.4em]

\begin{rustcode}
...
pub(crate) fn sys_pipe2(args: SyscallArgs) -> UserRet {
    let pipefd_ptr = args.arg(0);
    let flags = args.arg(1);
    if pipefd_ptr == 0 {
        return UserRet::from_error(ErrNo::EFAULT);
...
\end{rustcode}

\subsubsection*{序号563}
\noindent 路径：\code{os/components/wateros-syscall/syscall-impl/impl-kernel/src/sys/poll.rs}\\
\noindent 行范围：1-19\\[0.4em]

\begin{rustcode}
//! `poll(2)`（号 271）：委托共享 [`poll_engine`]。

//! 本模块代码由AI完成
use abi::syscall_args::SyscallArgs;
use abi::user_ret::UserRet;
...
\end{rustcode}

\subsubsection*{序号564}
\noindent 路径：\code{os/components/wateros-syscall/syscall-impl/impl-kernel/src/sys/poll.rs}\\
\noindent 行范围：10-19\\[0.4em]

\begin{rustcode}
...
pub(crate) fn sys_poll(args: SyscallArgs) -> UserRet {
    let fds_ptr = args.arg(0);
    let nfds = args.arg(1);
    let timeout_ms = args.arg(2) as isize;
    let deadline = match PollDeadline::from_poll_millis(timeout_ms) {
...
\end{rustcode}

\subsubsection*{序号565}
\noindent 路径：\code{os/components/wateros-syscall/syscall-impl/impl-kernel/src/sys/poll_multiplex.rs}\\
\noindent 行范围：1-66\\[0.4em]

\begin{rustcode}
//! `ppoll`(73) / `pselect6`(72) / `select`(23)。

//! 本模块代码由AI完成
use abi::syscall_args::SyscallArgs;
use abi::user_ret::UserRet;
...
\end{rustcode}

\subsubsection*{序号566}
\noindent 路径：\code{os/components/wateros-syscall/syscall-impl/impl-kernel/src/sys/poll_multiplex.rs}\\
\noindent 行范围：12-29\\[0.4em]

\begin{rustcode}
...
pub(crate) fn sys_ppoll(args: SyscallArgs) -> UserRet {
    let fds_ptr = args.arg(0);
    let nfds = args.arg(1);
    let timeout_ptr = args.arg(2);
    let sigmask_ptr = args.arg(3);
...
\end{rustcode}

\subsubsection*{序号567}
\noindent 路径：\code{os/components/wateros-syscall/syscall-impl/impl-kernel/src/sys/poll_multiplex.rs}\\
\noindent 行范围：32-51\\[0.4em]

\begin{rustcode}
...
pub(crate) fn sys_pselect6(args: SyscallArgs) -> UserRet {
    let nfds = args.arg(0);
    let readfds = args.arg(1);
    let writefds = args.arg(2);
    let exceptfds = args.arg(3);
...
\end{rustcode}

\subsubsection*{序号568}
\noindent 路径：\code{os/components/wateros-syscall/syscall-impl/impl-kernel/src/sys/poll_multiplex.rs}\\
\noindent 行范围：54-66\\[0.4em]

\begin{rustcode}
...
pub(crate) fn sys_select(args: SyscallArgs) -> UserRet {
    let nfds = args.arg(0);
    let readfds = args.arg(1);
    let writefds = args.arg(2);
    let exceptfds = args.arg(3);
...
\end{rustcode}

\subsubsection*{序号569}
\noindent 路径：\code{os/components/wateros-syscall/syscall-impl/impl-kernel/src/sys/posix_at_io.rs}\\
\noindent 行范围：1-274\\[0.4em]

\begin{rustcode}
//! `pread64` / `pwrite64` / `preadv` / `pwritev`：按绝对文件偏移读写，不改变 fd 当前偏移。

//! 本模块代码由AI完成
extern crate alloc;

...
\end{rustcode}

\subsubsection*{序号570}
\noindent 路径：\code{os/components/wateros-syscall/syscall-impl/impl-kernel/src/sys/posix_at_io.rs}\\
\noindent 行范围：127-160\\[0.4em]

\begin{rustcode}
...
pub(crate) fn sys_pread64(args : SyscallArgs) -> UserRet {
    let fd = args.arg(0);
    let ptr = args.arg(1);
    let len = args.arg(2);
    if len == 0 {
...
\end{rustcode}

\subsubsection*{序号571}
\noindent 路径：\code{os/components/wateros-syscall/syscall-impl/impl-kernel/src/sys/posix_at_io.rs}\\
\noindent 行范围：163-193\\[0.4em]

\begin{rustcode}
...
pub(crate) fn sys_pwrite64(args : SyscallArgs) -> UserRet {
    let fd = args.arg(0);
    let ptr = args.arg(1);
    let len = args.arg(2);
    if len == 0 {
...
\end{rustcode}

\subsubsection*{序号572}
\noindent 路径：\code{os/components/wateros-syscall/syscall-impl/impl-kernel/src/sys/posix_at_io.rs}\\
\noindent 行范围：196-249\\[0.4em]

\begin{rustcode}
...
pub(crate) fn sys_preadv(args : SyscallArgs) -> UserRet {
    let fd = args.arg(0);
    let iov_ptr = args.arg(1);
    let iovcnt = args.arg(2);
    let want = match total_iov_len(iov_ptr, iovcnt) {
...
\end{rustcode}

\subsubsection*{序号573}
\noindent 路径：\code{os/components/wateros-syscall/syscall-impl/impl-kernel/src/sys/posix_at_io.rs}\\
\noindent 行范围：252-274\\[0.4em]

\begin{rustcode}
...
pub(crate) fn sys_pwritev(args : SyscallArgs) -> UserRet {
    let fd = args.arg(0);
    let iov_ptr = args.arg(1);
    let iovcnt = args.arg(2);
    let offset = match offset_from_arg(args.arg(3)) {
...
\end{rustcode}

\subsubsection*{序号574}
\noindent 路径：\code{os/components/wateros-syscall/syscall-impl/impl-kernel/src/sys/priority.rs}\\
\noindent 行范围：1-165\\[0.4em]

\begin{rustcode}
//! `setpriority(2)` / `getpriority(2)`：仅维护 per-process nice 变量，不参与调度。

//! 本模块代码由AI完成
use abi::errno::ErrNo;
use abi::syscall_args::SyscallArgs;
...
\end{rustcode}

\subsubsection*{序号575}
\noindent 路径：\code{os/components/wateros-syscall/syscall-impl/impl-kernel/src/sys/priority.rs}\\
\noindent 行范围：89-127\\[0.4em]

\begin{rustcode}
...
pub(crate) fn sys_setpriority(args: SyscallArgs) -> UserRet {
    let which = args.arg(0) as i32;
    let who = args.arg(1) as i32;
    let prio = clamp_nice(args.arg(2) as i32);

...
\end{rustcode}

\subsubsection*{序号576}
\noindent 路径：\code{os/components/wateros-syscall/syscall-impl/impl-kernel/src/sys/priority.rs}\\
\noindent 行范围：130-165\\[0.4em]

\begin{rustcode}
...
pub(crate) fn sys_getpriority(args: SyscallArgs) -> UserRet {
    let which = args.arg(0) as i32;
    let who = args.arg(1) as i32;

    if !matches!(which, PRIO_PROCESS | PRIO_PGRP | PRIO_USER) {
...
\end{rustcode}

\subsubsection*{序号577}
\noindent 路径：\code{os/components/wateros-syscall/syscall-impl/impl-kernel/src/sys/read.rs}\\
\noindent 行范围：1-268\\[0.4em]

\begin{rustcode}
//! `read(2)`：支持 pipe 读端；stdin 暂未接真实输入。

//! 本模块代码由AI完成
extern crate alloc;

...
\end{rustcode}

\subsubsection*{序号578}
\noindent 路径：\code{os/components/wateros-syscall/syscall-impl/impl-kernel/src/sys/read.rs}\\
\noindent 行范围：37-67\\[0.4em]

\begin{rustcode}
...
pub(crate) fn sys_read(args : SyscallArgs) -> UserRet {
    let fd = args.arg(0);
    let ptr = args.arg(1);
    let len = args.arg(2);
    if len == 0 {
...
\end{rustcode}

\subsubsection*{序号579}
\noindent 路径：\code{os/components/wateros-syscall/syscall-impl/impl-kernel/src/sys/read.rs}\\
\noindent 行范围：70-189\\[0.4em]

\begin{rustcode}
...
pub(crate) fn sys_readv(args : SyscallArgs) -> UserRet {
    // iozone 调试：最早期 trace，在参数解包前
    let fd = args.arg(0);
    let iov_ptr = args.arg(1);
    let iovcnt = args.arg(2);
...
\end{rustcode}

\subsubsection*{序号580}
\noindent 路径：\code{os/components/wateros-syscall/syscall-impl/impl-kernel/src/sys/recvfrom.rs}\\
\noindent 行范围：1-161\\[0.4em]

\begin{rustcode}
//! `recvfrom(2)`：接收 TCP/UDP 数据。

//! 本模块代码由AI完成
use abi::errno::ErrNo;
use abi::syscall_args::SyscallArgs;
...
\end{rustcode}

\subsubsection*{序号581}
\noindent 路径：\code{os/components/wateros-syscall/syscall-impl/impl-kernel/src/sys/recvfrom.rs}\\
\noindent 行范围：16-21\\[0.4em]

\begin{rustcode}
...
struct SockAddrIn {
    sin_family: u16,
    sin_port: u16,
    sin_addr: [u8; 4],
    sin_zero: [u8; 8],
...
\end{rustcode}

\subsubsection*{序号582}
\noindent 路径：\code{os/components/wateros-syscall/syscall-impl/impl-kernel/src/sys/recvfrom.rs}\\
\noindent 行范围：24-112\\[0.4em]

\begin{rustcode}
...
pub(crate) fn sys_recvfrom(args: SyscallArgs) -> UserRet {
    let fd = args.arg(0);
    let buf_ptr = args.arg(1);
    let len = args.arg(2);
    let _flags = args.arg(3);
...
\end{rustcode}

\subsubsection*{序号583}
\noindent 路径：\code{os/components/wateros-syscall/syscall-impl/impl-kernel/src/sys/renameat2.rs}\\
\noindent 行范围：1-64\\[0.4em]

\begin{rustcode}
//! `renameat2(2)`：bring-up 同父目录 rename（文件与目录）；非 journal 原子语义。
//! 本模块代码由AI完成

use abi::errno::ErrNo;
use abi::syscall_args::SyscallArgs;
...
\end{rustcode}

\subsubsection*{序号584}
\noindent 路径：\code{os/components/wateros-syscall/syscall-impl/impl-kernel/src/sys/renameat2.rs}\\
\noindent 行范围：17-64\\[0.4em]

\begin{rustcode}
...
pub(crate) fn sys_renameat2(args: SyscallArgs) -> UserRet {
    let old_dirfd = args.arg(0) as isize;
    let old_path_ptr = args.arg(1);
    let new_dirfd = args.arg(2) as isize;
    let new_path_ptr = args.arg(3);
...
\end{rustcode}

\subsubsection*{序号585}
\noindent 路径：\code{os/components/wateros-syscall/syscall-impl/impl-kernel/src/sys/robust.rs}\\
\noindent 行范围：1-191\\[0.4em]

\begin{rustcode}
//! `set_robust_list` / `get_robust_list` 与线程退出时的 robust futex 深清理。

//! 本模块代码由AI完成
use abi::errno::ErrNo;
use abi::syscall_args::SyscallArgs;
...
\end{rustcode}

\subsubsection*{序号586}
\noindent 路径：\code{os/components/wateros-syscall/syscall-impl/impl-kernel/src/sys/robust.rs}\\
\noindent 行范围：18-27\\[0.4em]

\begin{rustcode}
...
pub(crate) fn futex_error_to_errno(error: FutexError) -> ErrNo {
    match error {
        FutexError::Again => ErrNo::EAGAIN,
        FutexError::Fault => ErrNo::EFAULT,
        FutexError::Invalid => ErrNo::EINVAL,
...
\end{rustcode}

\subsubsection*{序号587}
\noindent 路径：\code{os/components/wateros-syscall/syscall-impl/impl-kernel/src/sys/robust.rs}\\
\noindent 行范围：138-157\\[0.4em]

\begin{rustcode}
...
pub(crate) fn sys_set_robust_list(args: SyscallArgs) -> UserRet {
    let head = args.arg(0);
    let len = args.arg(1);
    if len != ROBUST_LIST_HEAD_SIZE {
        return UserRet::from_error(ErrNo::EINVAL);
...
\end{rustcode}

\subsubsection*{序号588}
\noindent 路径：\code{os/components/wateros-syscall/syscall-impl/impl-kernel/src/sys/robust.rs}\\
\noindent 行范围：160-191\\[0.4em]

\begin{rustcode}
...
pub(crate) fn sys_get_robust_list(args: SyscallArgs) -> UserRet {
    let pid = args.arg(0) as isize;
    let head_out = args.arg(1);
    let len_out = args.arg(2);

...
\end{rustcode}

\subsubsection*{序号589}
\noindent 路径：\code{os/components/wateros-syscall/syscall-impl/impl-kernel/src/sys/rtc.rs}\\
\noindent 行范围：1-96\\[0.4em]

\begin{rustcode}
//! `/dev/misc/rtc` 等软件 RTC 的 `ioctl(2)` 处理。

//! 本模块代码由AI完成
use abi::errno::ErrNo;
use abi::user_ret::UserRet;
...
\end{rustcode}

\subsubsection*{序号590}
\noindent 路径：\code{os/components/wateros-syscall/syscall-impl/impl-kernel/src/sys/rtc.rs}\\
\noindent 行范围：60-96\\[0.4em]

\begin{rustcode}
...
pub(crate) fn sys_rtc_ioctl(request: u32, argp: usize) -> UserRet {
    match request {
        RTC_RD_TIME => {
            if argp == 0 {
                return UserRet::from_error(ErrNo::EFAULT);
...
\end{rustcode}

\subsubsection*{序号591}
\noindent 路径：\code{os/components/wateros-syscall/syscall-impl/impl-kernel/src/sys/sched.rs}\\
\noindent 行范围：1-342\\[0.4em]

\begin{rustcode}
//! `sched_*` 系统调用：参数解析与用户拷贝；语义委托 [`task`] 调度原语。
//! 本模块代码由AI完成

extern crate alloc;

...
\end{rustcode}

\subsubsection*{序号592}
\noindent 路径：\code{os/components/wateros-syscall/syscall-impl/impl-kernel/src/sys/sched.rs}\\
\noindent 行范围：19-21\\[0.4em]

\begin{rustcode}
...
#[derive(Clone, Copy, Default)]
struct UserSchedParam {
    sched_priority: i32,
}

...
\end{rustcode}

\subsubsection*{序号593}
\noindent 路径：\code{os/components/wateros-syscall/syscall-impl/impl-kernel/src/sys/sched.rs}\\
\noindent 行范围：26-37\\[0.4em]

\begin{rustcode}
...
struct UserSchedAttr {
    size: u32,
    sched_policy: u32,
    sched_flags: u64,
    sched_nice: i32,
...
\end{rustcode}

\subsubsection*{序号594}
\noindent 路径：\code{os/components/wateros-syscall/syscall-impl/impl-kernel/src/sys/sched.rs}\\
\noindent 行范围：61-82\\[0.4em]

\begin{rustcode}
...
pub(crate) fn sys_sched_setparam(args: SyscallArgs) -> UserRet {
    let pid = args.arg(0) as isize;
    let param_ptr = args.arg(1);
    if param_ptr == 0 {
        return UserRet::from_error(ErrNo::EFAULT);
...
\end{rustcode}

\subsubsection*{序号595}
\noindent 路径：\code{os/components/wateros-syscall/syscall-impl/impl-kernel/src/sys/sched.rs}\\
\noindent 行范围：86-112\\[0.4em]

\begin{rustcode}
...
pub(crate) fn sys_sched_setscheduler(args: SyscallArgs) -> UserRet {
    let pid = args.arg(0) as isize;
    let policy_raw = args.arg(1) as isize;
    let param_ptr = args.arg(2);
    if param_ptr == 0 {
...
\end{rustcode}

\subsubsection*{序号596}
\noindent 路径：\code{os/components/wateros-syscall/syscall-impl/impl-kernel/src/sys/sched.rs}\\
\noindent 行范围：116-126\\[0.4em]

\begin{rustcode}
...
pub(crate) fn sys_sched_getscheduler(args: SyscallArgs) -> UserRet {
    let pid = args.arg(0) as isize;
    let task_id = match task::resolve_sched_pid(pid) {
        Ok(id) => id,
        Err(e) => return UserRet::from_error(sched_err_to_errno(e)),
...
\end{rustcode}

\subsubsection*{序号597}
\noindent 路径：\code{os/components/wateros-syscall/syscall-impl/impl-kernel/src/sys/sched.rs}\\
\noindent 行范围：130-151\\[0.4em]

\begin{rustcode}
...
pub(crate) fn sys_sched_getparam(args: SyscallArgs) -> UserRet {
    let pid = args.arg(0) as isize;
    let param_ptr = args.arg(1);
    if param_ptr == 0 {
        return UserRet::from_error(ErrNo::EFAULT);
...
\end{rustcode}

\subsubsection*{序号598}
\noindent 路径：\code{os/components/wateros-syscall/syscall-impl/impl-kernel/src/sys/sched.rs}\\
\noindent 行范围：155-183\\[0.4em]

\begin{rustcode}
...
pub(crate) fn sys_sched_setaffinity(args: SyscallArgs) -> UserRet {
    let pid = args.arg(0) as isize;
    let cpusetsize = args.arg(1);
    let mask_ptr = args.arg(2);
    if mask_ptr == 0 {
...
\end{rustcode}

\subsubsection*{序号599}
\noindent 路径：\code{os/components/wateros-syscall/syscall-impl/impl-kernel/src/sys/sched.rs}\\
\noindent 行范围：187-217\\[0.4em]

\begin{rustcode}
...
pub(crate) fn sys_sched_getaffinity(args: SyscallArgs) -> UserRet {
    let pid = args.arg(0) as isize;
    let cpusetsize = args.arg(1);
    let mask_ptr = args.arg(2);
    if mask_ptr == 0 {
...
\end{rustcode}

\subsubsection*{序号600}
\noindent 路径：\code{os/components/wateros-syscall/syscall-impl/impl-kernel/src/sys/sched.rs}\\
\noindent 行范围：221-267\\[0.4em]

\begin{rustcode}
...
pub(crate) fn sys_sched_setattr(args: SyscallArgs) -> UserRet {
    let pid = args.arg(0) as isize;
    let attr_ptr = args.arg(1);
    let flags = args.arg(2);
    if attr_ptr == 0 {
...
\end{rustcode}

\subsubsection*{序号601}
\noindent 路径：\code{os/components/wateros-syscall/syscall-impl/impl-kernel/src/sys/sched.rs}\\
\noindent 行范围：271-314\\[0.4em]

\begin{rustcode}
...
pub(crate) fn sys_sched_getattr(args: SyscallArgs) -> UserRet {
    let pid = args.arg(0) as isize;
    let attr_ptr = args.arg(1);
    let user_size = args.arg(2);
    let flags = args.arg(3);
...
\end{rustcode}

\subsubsection*{序号602}
\noindent 路径：\code{os/components/wateros-syscall/syscall-impl/impl-kernel/src/sys/sched.rs}\\
\noindent 行范围：318-328\\[0.4em]

\begin{rustcode}
...
pub(crate) fn sys_sched_get_priority_max(args: SyscallArgs) -> UserRet {
    let policy = match policy_from_arg(args.arg(0) as isize) {
        Ok(value) => value,
        Err(e) => return UserRet::from_error(e),
    };
...
\end{rustcode}

\subsubsection*{序号603}
\noindent 路径：\code{os/components/wateros-syscall/syscall-impl/impl-kernel/src/sys/sched.rs}\\
\noindent 行范围：332-342\\[0.4em]

\begin{rustcode}
...
pub(crate) fn sys_sched_get_priority_min(args: SyscallArgs) -> UserRet {
    let policy = match policy_from_arg(args.arg(0) as isize) {
        Ok(value) => value,
        Err(e) => return UserRet::from_error(e),
    };
...
\end{rustcode}

\subsubsection*{序号604}
\noindent 路径：\code{os/components/wateros-syscall/syscall-impl/impl-kernel/src/sys/sendfile.rs}\\
\noindent 行范围：1-97\\[0.4em]

\begin{rustcode}
//! `sendfile(2)`：在已打开 fd 间拷贝数据（内核缓冲，最小语义）。

//! 本模块代码由AI完成
extern crate alloc;

...
\end{rustcode}

\subsubsection*{序号605}
\noindent 路径：\code{os/components/wateros-syscall/syscall-impl/impl-kernel/src/sys/sendfile.rs}\\
\noindent 行范围：18-97\\[0.4em]

\begin{rustcode}
...
pub(crate) fn sys_sendfile(args: SyscallArgs) -> UserRet {
    let out_fd = args.arg(0);
    let in_fd = args.arg(1);
    let offset_ptr = args.arg(2);
    let count = args.arg(3);
...
\end{rustcode}

\subsubsection*{序号606}
\noindent 路径：\code{os/components/wateros-syscall/syscall-impl/impl-kernel/src/sys/sendmsg.rs}\\
\noindent 行范围：1-311\\[0.4em]

\begin{rustcode}
//! `sendmsg(2)` / `recvmsg(2)` — 将 scatter/gather I/O 转换为内部 send/recv 调用。

//! 本模块代码由AI完成
use abi::errno::ErrNo;
use abi::syscall_args::SyscallArgs;
...
\end{rustcode}

\subsubsection*{序号607}
\noindent 路径：\code{os/components/wateros-syscall/syscall-impl/impl-kernel/src/sys/sendmsg.rs}\\
\noindent 行范围：37-42\\[0.4em]

\begin{rustcode}
...
struct SockAddrIn {
    sin_family: u16,
    sin_port: u16,
    sin_addr: [u8; 4],
    sin_zero: [u8; 8],
...
\end{rustcode}

\subsubsection*{序号608}
\noindent 路径：\code{os/components/wateros-syscall/syscall-impl/impl-kernel/src/sys/sendmsg.rs}\\
\noindent 行范围：49-129\\[0.4em]

\begin{rustcode}
...
pub(crate) fn sys_sendmsg(args: SyscallArgs) -> UserRet {
    let fd = args.arg(0);
    let msg_ptr = args.arg(1);
    let _flags = args.arg(2);

...
\end{rustcode}

\subsubsection*{序号609}
\noindent 路径：\code{os/components/wateros-syscall/syscall-impl/impl-kernel/src/sys/sendmsg.rs}\\
\noindent 行范围：132-237\\[0.4em]

\begin{rustcode}
...
pub(crate) fn sys_recvmsg(args: SyscallArgs) -> UserRet {
    let fd = args.arg(0);
    let msg_ptr = args.arg(1);
    let flags = args.arg(2);

...
\end{rustcode}

\subsubsection*{序号610}
\noindent 路径：\code{os/components/wateros-syscall/syscall-impl/impl-kernel/src/sys/sendto.rs}\\
\noindent 行范围：1-174\\[0.4em]

\begin{rustcode}
//! `sendto(2)`：UDP 发送数据报到指定地址。

//! 本模块代码由AI完成
use abi::errno::ErrNo;
use abi::syscall_args::SyscallArgs;
...
\end{rustcode}

\subsubsection*{序号611}
\noindent 路径：\code{os/components/wateros-syscall/syscall-impl/impl-kernel/src/sys/sendto.rs}\\
\noindent 行范围：18-23\\[0.4em]

\begin{rustcode}
...
struct SockAddrIn {
    sin_family: u16,
    sin_port: u16,
    sin_addr: [u8; 4],
    sin_zero: [u8; 8],
...
\end{rustcode}

\subsubsection*{序号612}
\noindent 路径：\code{os/components/wateros-syscall/syscall-impl/impl-kernel/src/sys/sendto.rs}\\
\noindent 行范围：26-98\\[0.4em]

\begin{rustcode}
...
pub(crate) fn sys_sendto(args: SyscallArgs) -> UserRet {
    let fd = args.arg(0);
    let buf_ptr = args.arg(1);
    let len = args.arg(2);
    let _flags = args.arg(3);
...
\end{rustcode}

\subsubsection*{序号613}
\noindent 路径：\code{os/components/wateros-syscall/syscall-impl/impl-kernel/src/sys/shm.rs}\\
\noindent 行范围：1-238\\[0.4em]

\begin{rustcode}
//! SysV 共享内存 syscall 子集（`shmget` / `shmctl` / `shmat` / `shmdt`）。

//! 本模块代码由AI完成
use abi::errno::ErrNo;
use abi::syscall_args::SyscallArgs;
...
\end{rustcode}

\subsubsection*{序号614}
\noindent 路径：\code{os/components/wateros-syscall/syscall-impl/impl-kernel/src/sys/shm.rs}\\
\noindent 行范围：15-23\\[0.4em]

\begin{rustcode}
...
pub(crate) fn sys_shmget(args: SyscallArgs) -> UserRet {
    let key = args.arg(0);
    let size = args.arg(1);
    let flags = args.arg(2);
    match ipc::shm::registry().lock().create_or_get(key, size, flags) {
...
\end{rustcode}

\subsubsection*{序号615}
\noindent 路径：\code{os/components/wateros-syscall/syscall-impl/impl-kernel/src/sys/shm.rs}\\
\noindent 行范围：26-37\\[0.4em]

\begin{rustcode}
...
pub(crate) fn sys_shmctl(args: SyscallArgs) -> UserRet {
    let shmid = args.arg(0);
    let cmd = args.arg(1);

    match cmd {
...
\end{rustcode}

\subsubsection*{序号616}
\noindent 路径：\code{os/components/wateros-syscall/syscall-impl/impl-kernel/src/sys/shm.rs}\\
\noindent 行范围：40-87\\[0.4em]

\begin{rustcode}
...
pub(crate) fn sys_shmat(args: SyscallArgs) -> UserRet {
    let Some(handle) = current_user_aspace_handle() else {
        return UserRet::from_error(ErrNo::EFAULT);
    };
    let shmid = args.arg(0);
...
\end{rustcode}

\subsubsection*{序号617}
\noindent 路径：\code{os/components/wateros-syscall/syscall-impl/impl-kernel/src/sys/shm.rs}\\
\noindent 行范围：90-101\\[0.4em]

\begin{rustcode}
...
pub(crate) fn sys_shmdt(args: SyscallArgs) -> UserRet {
    let Some(handle) = current_user_aspace_handle() else {
        return UserRet::from_error(ErrNo::EFAULT);
    };
    let base = args.arg(0);
...
\end{rustcode}

\subsubsection*{序号618}
\noindent 路径：\code{os/components/wateros-syscall/syscall-impl/impl-kernel/src/sys/shm.rs}\\
\noindent 行范围：226-234\\[0.4em]

\begin{rustcode}
...
fn shm_error_to_errno(error: ShmError) -> ErrNo {
    match error {
        ShmError::Invalid => ErrNo::EINVAL,
        ShmError::Exists => ErrNo::EEXIST,
        ShmError::NoEntry => ErrNo::ENOENT,
...
\end{rustcode}

\subsubsection*{序号619}
\noindent 路径：\code{os/components/wateros-syscall/syscall-impl/impl-kernel/src/sys/shutdown.rs}\\
\noindent 行范围：1-30\\[0.4em]

\begin{rustcode}
//! `shutdown(2)` — 极简存根。
//! 本模块代码由AI完成

use abi::errno::ErrNo;
use abi::syscall_args::SyscallArgs;
...
\end{rustcode}

\subsubsection*{序号620}
\noindent 路径：\code{os/components/wateros-syscall/syscall-impl/impl-kernel/src/sys/shutdown.rs}\\
\noindent 行范围：12-30\\[0.4em]

\begin{rustcode}
...
pub(crate) fn sys_shutdown(args: SyscallArgs) -> UserRet {
    let fd = args.arg(0);
    let how = args.arg(1);

    if how > 2 {
...
\end{rustcode}

\subsubsection*{序号621}
\noindent 路径：\code{os/components/wateros-syscall/syscall-impl/impl-kernel/src/sys/socket.rs}\\
\noindent 行范围：1-101\\[0.4em]

\begin{rustcode}
//! `socket(2)`：创建 socket 并分配 fd。

//! 本模块代码由AI完成
extern crate alloc;

...
\end{rustcode}

\subsubsection*{序号622}
\noindent 路径：\code{os/components/wateros-syscall/syscall-impl/impl-kernel/src/sys/socket.rs}\\
\noindent 行范围：25-101\\[0.4em]

\begin{rustcode}
...
pub(crate) fn sys_socket(args: SyscallArgs) -> UserRet {
    let domain = args.arg(0);
    let mut typ = args.arg(1);
    let _protocol = args.arg(2);

...
\end{rustcode}

\subsubsection*{序号623}
\noindent 路径：\code{os/components/wateros-syscall/syscall-impl/impl-kernel/src/sys/socketpair.rs}\\
\noindent 行范围：1-81\\[0.4em]

\begin{rustcode}
//! `socketpair(2)`：创建一对已连接的 AF_UNIX stream socket fd。

//! 本模块代码由AI完成
extern crate alloc;

...
\end{rustcode}

\subsubsection*{序号624}
\noindent 路径：\code{os/components/wateros-syscall/syscall-impl/impl-kernel/src/sys/socketpair.rs}\\
\noindent 行范围：23-81\\[0.4em]

\begin{rustcode}
...
pub(crate) fn sys_socketpair(args: SyscallArgs) -> UserRet {
    let domain = args.arg(0);
    let mut typ = args.arg(1);
    let protocol = args.arg(2);
    let sv_ptr = args.arg(3);
...
\end{rustcode}

\subsubsection*{序号625}
\noindent 路径：\code{os/components/wateros-syscall/syscall-impl/impl-kernel/src/sys/sockname.rs}\\
\noindent 行范围：1-116\\[0.4em]

\begin{rustcode}
//! `getsockname(2)` / `getpeername(2)` — 获取 socket 地址。

//! 本模块代码由AI完成
use abi::errno::ErrNo;
use abi::syscall_args::SyscallArgs;
...
\end{rustcode}

\subsubsection*{序号626}
\noindent 路径：\code{os/components/wateros-syscall/syscall-impl/impl-kernel/src/sys/sockname.rs}\\
\noindent 行范围：15-20\\[0.4em]

\begin{rustcode}
...
struct SockAddrIn {
    sin_family: u16,
    sin_port: u16,
    sin_addr: [u8; 4],
    sin_zero: [u8; 8],
...
\end{rustcode}

\subsubsection*{序号627}
\noindent 路径：\code{os/components/wateros-syscall/syscall-impl/impl-kernel/src/sys/sockname.rs}\\
\noindent 行范围：23-69\\[0.4em]

\begin{rustcode}
...
pub(crate) fn sys_getsockname(args: SyscallArgs) -> UserRet {
    let fd = args.arg(0);
    let addr_ptr = args.arg(1);
    let addrlen_ptr = args.arg(2);

...
\end{rustcode}

\subsubsection*{序号628}
\noindent 路径：\code{os/components/wateros-syscall/syscall-impl/impl-kernel/src/sys/sockname.rs}\\
\noindent 行范围：72-116\\[0.4em]

\begin{rustcode}
...
pub(crate) fn sys_getpeername(args: SyscallArgs) -> UserRet {
    let fd = args.arg(0);
    let addr_ptr = args.arg(1);
    let addrlen_ptr = args.arg(2);

...
\end{rustcode}

\subsubsection*{序号629}
\noindent 路径：\code{os/components/wateros-syscall/syscall-impl/impl-kernel/src/sys/sockopt.rs}\\
\noindent 行范围：1-86\\[0.4em]

\begin{rustcode}
//! `setsockopt(2)` / `getsockopt(2)` — 极简存根。
//! 本模块代码由AI完成

use abi::errno::ErrNo;
use abi::syscall_args::SyscallArgs;
...
\end{rustcode}

\subsubsection*{序号630}
\noindent 路径：\code{os/components/wateros-syscall/syscall-impl/impl-kernel/src/sys/sockopt.rs}\\
\noindent 行范围：14-48\\[0.4em]

\begin{rustcode}
...
pub(crate) fn sys_setsockopt(args: SyscallArgs) -> UserRet {
    let fd = args.arg(0);
    let level = args.arg(1);
    let optname = args.arg(2);
    let optval = args.arg(3);
...
\end{rustcode}

\subsubsection*{序号631}
\noindent 路径：\code{os/components/wateros-syscall/syscall-impl/impl-kernel/src/sys/sockopt.rs}\\
\noindent 行范围：51-86\\[0.4em]

\begin{rustcode}
...
pub(crate) fn sys_getsockopt(args: SyscallArgs) -> UserRet {
    let fd = args.arg(0);
    let level = args.arg(1);
    let optname = args.arg(2);
    let optval = args.arg(3);
...
\end{rustcode}

\subsubsection*{序号632}
\noindent 路径：\code{os/components/wateros-syscall/syscall-impl/impl-kernel/src/sys/statfs.rs}\\
\noindent 行范围：1-86\\[0.4em]

\begin{rustcode}
//! `statfs(2)`：bring-up 最小兼容实现。
//! 本模块代码由AI完成

use abi::errno::ErrNo;
use abi::syscall_args::SyscallArgs;
...
\end{rustcode}

\subsubsection*{序号633}
\noindent 路径：\code{os/components/wateros-syscall/syscall-impl/impl-kernel/src/sys/statfs.rs}\\
\noindent 行范围：15\\[0.4em]

\begin{rustcode}
...

// 本变量代码由AI完成
const EXT4_SUPER_MAGIC: isize = 0xEF53;
const STATFS_BLOCK_SIZE: isize = 4096;
const STATFS_TOTAL_BLOCKS: isize = 1024 * 1024;
...
\end{rustcode}

\subsubsection*{序号634}
\noindent 路径：\code{os/components/wateros-syscall/syscall-impl/impl-kernel/src/sys/statfs.rs}\\
\noindent 行范围：26-39\\[0.4em]

\begin{rustcode}
...
struct LinuxStatFs {
    f_type: isize,
    f_bsize: isize,
    f_blocks: isize,
    f_bfree: isize,
...
\end{rustcode}

\subsubsection*{序号635}
\noindent 路径：\code{os/components/wateros-syscall/syscall-impl/impl-kernel/src/sys/statfs.rs}\\
\noindent 行范围：44-86\\[0.4em]

\begin{rustcode}
...
pub(crate) fn sys_statfs(args: SyscallArgs) -> UserRet {
    let path_ptr = args.arg(0);
    let buf_ptr = args.arg(1);
    if buf_ptr == 0 {
        return UserRet::from_error(ErrNo::EFAULT);
...
\end{rustcode}

\subsubsection*{序号636}
\noindent 路径：\code{os/components/wateros-syscall/syscall-impl/impl-kernel/src/sys/symlinkat.rs}\\
\noindent 行范围：1-43\\[0.4em]

\begin{rustcode}
//! `symlinkat(2)`：相对目录 fd 创建符号链接。
//! 本模块代码由AI完成

use abi::errno::ErrNo;
use abi::syscall_args::SyscallArgs;
...
\end{rustcode}

\subsubsection*{序号637}
\noindent 路径：\code{os/components/wateros-syscall/syscall-impl/impl-kernel/src/sys/symlinkat.rs}\\
\noindent 行范围：16-43\\[0.4em]

\begin{rustcode}
...
pub(crate) fn sys_symlinkat(args: SyscallArgs) -> UserRet {
    let target_ptr = args.arg(0);
    let dirfd = args.arg(1) as isize;
    let linkpath_ptr = args.arg(2);

...
\end{rustcode}

\subsubsection*{序号638}
\noindent 路径：\code{os/components/wateros-syscall/syscall-impl/impl-kernel/src/sys/sync.rs}\\
\noindent 行范围：1-43\\[0.4em]

\begin{rustcode}
//! `sync(2)` / `fsync(2)` / `fdatasync(2)`：将脏数据刷回后端。
//! 本模块代码由AI完成

use abi::errno::ErrNo;
use abi::syscall_args::SyscallArgs;
...
\end{rustcode}

\subsubsection*{序号639}
\noindent 路径：\code{os/components/wateros-syscall/syscall-impl/impl-kernel/src/sys/sync.rs}\\
\noindent 行范围：11-30\\[0.4em]

\begin{rustcode}
...
fn sync_fd(op : &str, fd : usize) -> UserRet {
    use vfs::api::VfsNodeType;

    match vfs::fd::with_current_io(fd, |handle| {
        let meta = handle.metadata()?;
...
\end{rustcode}

\subsubsection*{序号640}
\noindent 路径：\code{os/components/wateros-syscall/syscall-impl/impl-kernel/src/sys/sync.rs}\\
\noindent 行范围：33\\[0.4em]

\begin{rustcode}
...

// 本方法代码由AI完成
pub(crate) fn sys_fsync(args : SyscallArgs) -> UserRet { sync_fd("fsync", args.arg(0)) }

// 本方法代码由AI完成
...
\end{rustcode}

\subsubsection*{序号641}
\noindent 路径：\code{os/components/wateros-syscall/syscall-impl/impl-kernel/src/sys/sync.rs}\\
\noindent 行范围：36\\[0.4em]

\begin{rustcode}
...

// 本方法代码由AI完成
pub(crate) fn sys_fdatasync(args : SyscallArgs) -> UserRet { sync_fd("fdatasync", args.arg(0)) }

// 本方法代码由AI完成
...
\end{rustcode}

\subsubsection*{序号642}
\noindent 路径：\code{os/components/wateros-syscall/syscall-impl/impl-kernel/src/sys/sync.rs}\\
\noindent 行范围：39-43\\[0.4em]

\begin{rustcode}
...
pub(crate) fn sys_sync(_args : SyscallArgs) -> UserRet {
    // Linux sync(2) 发起系统级写回并始终返回 0；单个写回错误由后续文件操作报告。
    let _ = vfs::fd::flush_all_open_files();
    UserRet::from_success(0)
}
\end{rustcode}

\subsubsection*{序号643}
\noindent 路径：\code{os/components/wateros-syscall/syscall-impl/impl-kernel/src/sys/syslog.rs}\\
\noindent 行范围：1-59\\[0.4em]

\begin{rustcode}
//! `syslog(2)` → 内核 klog 环。

//! 本模块代码由AI完成
use abi::errno::ErrNo;
use abi::syscall_args::SyscallArgs;
...
\end{rustcode}

\subsubsection*{序号644}
\noindent 路径：\code{os/components/wateros-syscall/syscall-impl/impl-kernel/src/sys/syslog.rs}\\
\noindent 行范围：18-59\\[0.4em]

\begin{rustcode}
...
pub(crate) fn sys_syslog(args: SyscallArgs) -> UserRet {
    let action = args.arg(0) as i32;
    let buf_ptr = args.arg(1);
    let len = args.arg(2);

...
\end{rustcode}

\subsubsection*{序号645}
\noindent 路径：\code{os/components/wateros-syscall/syscall-impl/impl-kernel/src/sys/truncate.rs}\\
\noindent 行范围：1-37\\[0.4em]

\begin{rustcode}
//! `truncate(2)`：按路径调整普通文件长度。
//! 本模块代码由AI完成

use abi::errno::ErrNo;
use abi::syscall_args::SyscallArgs;
...
\end{rustcode}

\subsubsection*{序号646}
\noindent 路径：\code{os/components/wateros-syscall/syscall-impl/impl-kernel/src/sys/truncate.rs}\\
\noindent 行范围：14-37\\[0.4em]

\begin{rustcode}
...
pub(crate) fn sys_truncate(args: SyscallArgs) -> UserRet {
    let path_ptr = args.arg(0);
    let len = args.arg(1) as u64;

    if path_ptr == 0 {
...
\end{rustcode}

\subsubsection*{序号647}
\noindent 路径：\code{os/components/wateros-syscall/syscall-impl/impl-kernel/src/sys/umount2.rs}\\
\noindent 行范围：1-48\\[0.4em]

\begin{rustcode}
//! `umount2(2)`：卸载辅助卷挂载点；支持 `MNT_DETACH`。
//! 本模块代码由AI完成

use abi::errno::ErrNo;
use abi::syscall_args::SyscallArgs;
...
\end{rustcode}

\subsubsection*{序号648}
\noindent 路径：\code{os/components/wateros-syscall/syscall-impl/impl-kernel/src/sys/umount2.rs}\\
\noindent 行范围：16-48\\[0.4em]

\begin{rustcode}
...
pub(crate) fn sys_umount2(args: SyscallArgs) -> UserRet {
    let target_ptr = args.arg(0);
    let flags = args.arg(1) as u32;

    if flags != 0 && flags != MNT_DETACH {
...
\end{rustcode}

\subsubsection*{序号649}
\noindent 路径：\code{os/components/wateros-syscall/syscall-impl/impl-kernel/src/sys/unlinkat.rs}\\
\noindent 行范围：1-43\\[0.4em]

\begin{rustcode}
//! `unlinkat(2)`：相对 cwd / 目录 fd 删除文件或空目录。

//! 本模块代码由AI完成
use abi::errno::ErrNo;
use abi::syscall_args::SyscallArgs;
...
\end{rustcode}

\subsubsection*{序号650}
\noindent 路径：\code{os/components/wateros-syscall/syscall-impl/impl-kernel/src/sys/unlinkat.rs}\\
\noindent 行范围：15-43\\[0.4em]

\begin{rustcode}
...
pub(crate) fn sys_unlinkat(args : SyscallArgs) -> UserRet {
    cgroup_regression_loop_fast_exit_if_standalone();

    let dirfd = args.arg(0) as isize;
    let path_ptr = args.arg(1);
...
\end{rustcode}

\subsubsection*{序号651}
\noindent 路径：\code{os/components/wateros-syscall/syscall-impl/impl-kernel/src/sys/unshare.rs}\\
\noindent 行范围：1-26\\[0.4em]

\begin{rustcode}
//! `unshare(2)`：创建新挂载命名空间（`CLONE_NEWNS`）。
//! 本模块代码由AI完成

use abi::errno::ErrNo;
use abi::syscall_args::SyscallArgs;
...
\end{rustcode}

\subsubsection*{序号652}
\noindent 路径：\code{os/components/wateros-syscall/syscall-impl/impl-kernel/src/sys/unshare.rs}\\
\noindent 行范围：11-26\\[0.4em]

\begin{rustcode}
...
pub(crate) fn sys_unshare(args: SyscallArgs) -> UserRet {
    let flags = args.arg(0);
    if flags == 0 {
        return UserRet::from_error(ErrNo::EINVAL);
    }
...
\end{rustcode}

\subsubsection*{序号653}
\noindent 路径：\code{os/components/wateros-syscall/syscall-impl/impl-kernel/src/sys/utimensat.rs}\\
\noindent 行范围：1-144\\[0.4em]

\begin{rustcode}
//! `utimensat(2)`：bring-up 最小兼容实现，暂不持久化 atime/mtime。
//! 本模块代码由AI完成

use abi::errno::ErrNo;
use abi::syscall_args::SyscallArgs;
...
\end{rustcode}

\subsubsection*{序号654}
\noindent 路径：\code{os/components/wateros-syscall/syscall-impl/impl-kernel/src/sys/utimensat.rs}\\
\noindent 行范围：23-26\\[0.4em]

\begin{rustcode}
...
// 本结构代码由AI完成
struct UserTimespec {
    sec: isize,
    nsec: isize,
}
...
\end{rustcode}

\subsubsection*{序号655}
\noindent 路径：\code{os/components/wateros-syscall/syscall-impl/impl-kernel/src/sys/utimensat.rs}\\
\noindent 行范围：36-83\\[0.4em]

\begin{rustcode}
...
pub(crate) fn sys_utimensat(args: SyscallArgs) -> UserRet {
    let dirfd = args.arg(0) as isize;
    let path_ptr = args.arg(1);
    let times_ptr = args.arg(2);
    let flags = args.arg(3) as u32;
...
\end{rustcode}

\subsubsection*{序号656}
\noindent 路径：\code{os/components/wateros-syscall/syscall-impl/impl-kernel/src/sys/write.rs}\\
\noindent 行范围：1-225\\[0.4em]

\begin{rustcode}
//! `write(2)`：fd 1/2 走控制台；pipe 写端走 IPC。

//! 本模块代码由AI完成
extern crate alloc;

...
\end{rustcode}

\subsubsection*{序号657}
\noindent 路径：\code{os/components/wateros-syscall/syscall-impl/impl-kernel/src/sys/write.rs}\\
\noindent 行范围：34-85\\[0.4em]

\begin{rustcode}
...
pub(crate) fn sys_write(args : SyscallArgs) -> UserRet {
    let fd = args.arg(0);
    let ptr = args.arg(1);
    let len = args.arg(2);
    if len == 0 {
...
\end{rustcode}

\subsubsection*{序号658}
\noindent 路径：\code{os/components/wateros-syscall/syscall-impl/impl-kernel/src/sys/write.rs}\\
\noindent 行范围：88-143\\[0.4em]

\begin{rustcode}
...
pub(crate) fn sys_writev(args : SyscallArgs) -> UserRet {
    let fd = args.arg(0);
    let iov_ptr = args.arg(1);
    let iovcnt = args.arg(2);
    if iovcnt == 0 {
...
\end{rustcode}

\subsubsection*{序号659}
\noindent 路径：\code{os/components/wateros-syscall/syscall-impl/impl-kernel/src/sys/xattr.rs}\\
\noindent 行范围：1-378\\[0.4em]

\begin{rustcode}
//! 扩展属性 syscall：`setxattr`/`getxattr`/`listxattr`/`removexattr` 及 l*/f* 变体。
//! 本模块代码由AI完成

extern crate alloc;

...
\end{rustcode}

\subsubsection*{序号660}
\noindent 路径：\code{os/components/wateros-syscall/syscall-impl/impl-kernel/src/sys/xattr.rs}\\
\noindent 行范围：89-91\\[0.4em]

\begin{rustcode}
...
// 本方法代码由AI完成
pub(crate) fn sys_setxattr(args: SyscallArgs) -> UserRet {
    path_setxattr(args, true)
}

...
\end{rustcode}

\subsubsection*{序号661}
\noindent 路径：\code{os/components/wateros-syscall/syscall-impl/impl-kernel/src/sys/xattr.rs}\\
\noindent 行范围：94-96\\[0.4em]

\begin{rustcode}
...
// 本方法代码由AI完成
pub(crate) fn sys_lsetxattr(args: SyscallArgs) -> UserRet {
    path_setxattr(args, false)
}

...
\end{rustcode}

\subsubsection*{序号662}
\noindent 路径：\code{os/components/wateros-syscall/syscall-impl/impl-kernel/src/sys/xattr.rs}\\
\noindent 行范围：125-149\\[0.4em]

\begin{rustcode}
...
pub(crate) fn sys_fsetxattr(args: SyscallArgs) -> UserRet {
    let fd = args.arg(0);
    let name_ptr = args.arg(1);
    let value_ptr = args.arg(2);
    let size = args.arg(3);
...
\end{rustcode}

\subsubsection*{序号663}
\noindent 路径：\code{os/components/wateros-syscall/syscall-impl/impl-kernel/src/sys/xattr.rs}\\
\noindent 行范围：152-154\\[0.4em]

\begin{rustcode}
...
// 本方法代码由AI完成
pub(crate) fn sys_getxattr(args: SyscallArgs) -> UserRet {
    path_getxattr(args, true)
}

...
\end{rustcode}

\subsubsection*{序号664}
\noindent 路径：\code{os/components/wateros-syscall/syscall-impl/impl-kernel/src/sys/xattr.rs}\\
\noindent 行范围：157-159\\[0.4em]

\begin{rustcode}
...
// 本方法代码由AI完成
pub(crate) fn sys_lgetxattr(args: SyscallArgs) -> UserRet {
    path_getxattr(args, false)
}

...
\end{rustcode}

\subsubsection*{序号665}
\noindent 路径：\code{os/components/wateros-syscall/syscall-impl/impl-kernel/src/sys/xattr.rs}\\
\noindent 行范围：204-244\\[0.4em]

\begin{rustcode}
...
pub(crate) fn sys_fgetxattr(args: SyscallArgs) -> UserRet {
    let fd = args.arg(0);
    let name_ptr = args.arg(1);
    let value_ptr = args.arg(2);
    let size = args.arg(3);
...
\end{rustcode}

\subsubsection*{序号666}
\noindent 路径：\code{os/components/wateros-syscall/syscall-impl/impl-kernel/src/sys/xattr.rs}\\
\noindent 行范围：247-249\\[0.4em]

\begin{rustcode}
...
// 本方法代码由AI完成
pub(crate) fn sys_listxattr(args: SyscallArgs) -> UserRet {
    path_listxattr(args, true)
}

...
\end{rustcode}

\subsubsection*{序号667}
\noindent 路径：\code{os/components/wateros-syscall/syscall-impl/impl-kernel/src/sys/xattr.rs}\\
\noindent 行范围：252-254\\[0.4em]

\begin{rustcode}
...
// 本方法代码由AI完成
pub(crate) fn sys_llistxattr(args: SyscallArgs) -> UserRet {
    path_listxattr(args, false)
}

...
\end{rustcode}

\subsubsection*{序号668}
\noindent 路径：\code{os/components/wateros-syscall/syscall-impl/impl-kernel/src/sys/xattr.rs}\\
\noindent 行范围：294-329\\[0.4em]

\begin{rustcode}
...
pub(crate) fn sys_flistxattr(args: SyscallArgs) -> UserRet {
    let fd = args.arg(0);
    let list_ptr = args.arg(1);
    let size = args.arg(2);

...
\end{rustcode}

\subsubsection*{序号669}
\noindent 路径：\code{os/components/wateros-syscall/syscall-impl/impl-kernel/src/sys/xattr.rs}\\
\noindent 行范围：332-334\\[0.4em]

\begin{rustcode}
...
// 本方法代码由AI完成
pub(crate) fn sys_removexattr(args: SyscallArgs) -> UserRet {
    path_removexattr(args, true)
}

...
\end{rustcode}

\subsubsection*{序号670}
\noindent 路径：\code{os/components/wateros-syscall/syscall-impl/impl-kernel/src/sys/xattr.rs}\\
\noindent 行范围：337-339\\[0.4em]

\begin{rustcode}
...
// 本方法代码由AI完成
pub(crate) fn sys_lremovexattr(args: SyscallArgs) -> UserRet {
    path_removexattr(args, false)
}

...
\end{rustcode}

\subsubsection*{序号671}
\noindent 路径：\code{os/components/wateros-syscall/syscall-impl/impl-kernel/src/sys/xattr.rs}\\
\noindent 行范围：361-378\\[0.4em]

\begin{rustcode}
...
pub(crate) fn sys_fremovexattr(args: SyscallArgs) -> UserRet {
    let fd = args.arg(0);
    let name_ptr = args.arg(1);

    let path = match path_from_fd(fd) {
...
\end{rustcode}

\subsubsection*{序号672}
\noindent 路径：\code{os/components/wateros-syscall/syscall-impl/impl-kernel/src/syscall_nr_dispatch.rs}\\
\noindent 行范围：1-335\\[0.4em]

\begin{rustcode}
// Syscall 号单次 match 分发（H-3）：替代 `SyscallKind::decode` + 巨型 match。
// 本模块代码由AI完成

/// 按裸 syscall 号分发；未命中时走旁路号与 ENOSYS。
#[inline]
...
\end{rustcode}

\subsubsection*{序号673}
\noindent 路径：\code{os/components/wateros-syscall/syscall-impl/impl-kernel/src/syscall_nr_dispatch.rs}\\
\noindent 行范围：7-316\\[0.4em]

\begin{rustcode}
...
pub fn dispatch_syscall_by_nr(syscall_nr: usize, syscall_args: SyscallArgs) -> isize {
    use api_v0::SyscallDispatcher;
    match syscall_nr {
        n if n == <ActiveSyscallNumberTable as SyscallNumberTable>::YIELD.raw() =>
            KernelSyscallDispatcher::dispatch_yield(syscall_args),
...
\end{rustcode}

\subsubsection*{序号674}
\noindent 路径：\code{os/components/wateros-syscall/syscall-impl/impl-kernel/src/syscall_nr_dispatch.rs}\\
\noindent 行范围：321-335\\[0.4em]

\begin{rustcode}
...
pub fn is_restartable_syscall_nr(syscall_nr: usize) -> bool {
    use ActiveSyscallNumberTable as T;
    syscall_nr == <T as SyscallNumberTable>::READ.raw()
        || syscall_nr == <T as SyscallNumberTable>::READV.raw()
        || syscall_nr == <T as SyscallNumberTable>::WRITE.raw()
...
\end{rustcode}

\subsubsection*{序号675}
\noindent 路径：\code{os/components/wateros-syscall/syscall-impl/impl-kernel/src/unix_sock.rs}\\
\noindent 行范围：1-611\\[0.4em]

\begin{rustcode}
//! AF_UNIX 域套接字：pathname / abstract bind、stream listen/accept/connect、dgram 投递。

//! 本模块代码由AI完成
extern crate alloc;

...
\end{rustcode}

\subsubsection*{序号676}
\noindent 路径：\code{os/components/wateros-syscall/syscall-impl/impl-kernel/src/unix_sock.rs}\\
\noindent 行范围：38-41\\[0.4em]

\begin{rustcode}
...
// 本结构代码由AI完成
pub(crate) enum UnixSockType {
    Stream,
    Dgram,
}
...
\end{rustcode}

\subsubsection*{序号677}
\noindent 路径：\code{os/components/wateros-syscall/syscall-impl/impl-kernel/src/unix_sock.rs}\\
\noindent 行范围：135-160\\[0.4em]

\begin{rustcode}
...
pub(crate) fn alloc_unix_socket(
    typ: usize,
    status_flags: usize,
) -> Result<(Box<dyn VfsIoHandle>, UnixSockRef), ErrNo> {
    let sock_type = match typ {
...
\end{rustcode}

\subsubsection*{序号678}
\noindent 路径：\code{os/components/wateros-syscall/syscall-impl/impl-kernel/src/unix_sock.rs}\\
\noindent 行范围：196-227\\[0.4em]

\begin{rustcode}
...
pub(crate) fn bind(fd: usize, addr_ptr: usize, addrlen: usize) -> Result<(), ErrNo> {
    let addr = parse_sockaddr_un(addr_ptr, addrlen)?;
    let sock = lookup_current(fd)?;
    {
        let inner = sock.inner.lock();
...
\end{rustcode}

\subsubsection*{序号679}
\noindent 路径：\code{os/components/wateros-syscall/syscall-impl/impl-kernel/src/unix_sock.rs}\\
\noindent 行范围：230-242\\[0.4em]

\begin{rustcode}
...
pub(crate) fn listen(fd: usize, _backlog: usize) -> Result<(), ErrNo> {
    let sock = lookup_current(fd)?;
    let mut inner = sock.inner.lock();
    let key = inner.bound_key.clone().ok_or(ErrNo::EINVAL)?;
    if inner.sock_type != UnixSockType::Stream {
...
\end{rustcode}

\subsubsection*{序号680}
\noindent 路径：\code{os/components/wateros-syscall/syscall-impl/impl-kernel/src/unix_sock.rs}\\
\noindent 行范围：245-256\\[0.4em]

\begin{rustcode}
...
pub(crate) fn connect(fd: usize, addr_ptr: usize, addrlen: usize) -> Result<(), ErrNo> {
    let addr = parse_sockaddr_un(addr_ptr, addrlen)?;
    let sock = lookup_current(fd)?;
    let mut inner = sock.inner.lock();
    match inner.sock_type {
...
\end{rustcode}

\subsubsection*{序号681}
\noindent 路径：\code{os/components/wateros-syscall/syscall-impl/impl-kernel/src/user_copy.rs}\\
\noindent 行范围：1-182\\[0.4em]

\begin{rustcode}
//! 用户虚拟地址与内核缓冲区之间的安全拷贝。

//! 本模块代码由AI完成
extern crate alloc;

...
\end{rustcode}

\subsubsection*{序号682}
\noindent 路径：\code{os/components/wateros-syscall/syscall-impl/impl-kernel/src/user_copy.rs}\\
\noindent 行范围：18-21\\[0.4em]

\begin{rustcode}
...
// 本方法代码由AI完成
pub(crate) fn user_aspace_required() -> Result<ActiveUserMemoryOps, ErrNo> {
    let handle = current_user_aspace_handle().ok_or(ErrNo::EFAULT)?;
    Ok(ActiveUserMemoryOps::new(handle))
}
...
\end{rustcode}

\subsubsection*{序号683}
\noindent 路径：\code{os/components/wateros-syscall/syscall-impl/impl-kernel/src/user_copy.rs}\\
\noindent 行范围：24-37\\[0.4em]

\begin{rustcode}
...
pub(crate) fn copy_from_user(buf: &mut [u8], ptr: usize) -> Result<usize, ErrNo> {
    if buf.is_empty() {
        return Ok(0);
    }
    if ptr == 0 {
...
\end{rustcode}

\subsubsection*{序号684}
\noindent 路径：\code{os/components/wateros-syscall/syscall-impl/impl-kernel/src/user_copy.rs}\\
\noindent 行范围：40-63\\[0.4em]

\begin{rustcode}
...
pub(crate) fn copy_to_user(ptr: usize, buf: &[u8]) -> Result<usize, ErrNo> {
    if buf.is_empty() {
        return Ok(0);
    }
    if ptr == 0 {
...
\end{rustcode}

\subsubsection*{序号685}
\noindent 路径：\code{os/components/wateros-syscall/syscall-impl/impl-kernel/src/user_copy.rs}\\
\noindent 行范围：93-119\\[0.4em]

\begin{rustcode}
...
pub(crate) fn copy_user_path_cstr(ptr: usize, max: usize) -> Result<String, ErrNo> {
    if ptr == 0 {
        return Err(ErrNo::EFAULT);
    }
    if max == 0 {
...
\end{rustcode}

\subsubsection*{序号686}
\noindent 路径：\code{os/components/wateros-syscall/syscall-impl/impl-kernel/src/user_copy.rs}\\
\noindent 行范围：122-136\\[0.4em]

\begin{rustcode}
...
pub(crate) fn copy_to_user_struct<T: Copy>(ptr: usize, value: &T) -> Result<(), ErrNo> {
    let bytes = unsafe {
        core::slice::from_raw_parts(
            (value as *const T) as *const u8,
            core::mem::size_of::<T>(),
...
\end{rustcode}

\subsubsection*{序号687}
\noindent 路径：\code{os/components/wateros-syscall/syscall-impl/impl-kernel/src/user_copy.rs}\\
\noindent 行范围：139-166\\[0.4em]

\begin{rustcode}
...
pub(crate) fn copy_to_user_struct_in_aspace<T: Copy>(
    handle: usize,
    ptr: usize,
    value: &T,
) -> Result<(), ErrNo> {
...
\end{rustcode}

\subsubsection*{序号688}
\noindent 路径：\code{os/components/wateros-syscall/syscall-impl/impl-kernel/src/user_copy.rs}\\
\noindent 行范围：169-182\\[0.4em]

\begin{rustcode}
...
pub(crate) fn copy_from_user_struct<T: Copy>(ptr: usize) -> Result<T, ErrNo> {
    let mut value = core::mem::MaybeUninit::<T>::uninit();
    let bytes = unsafe {
        core::slice::from_raw_parts_mut(
            value.as_mut_ptr() as *mut u8,
...
\end{rustcode}

\subsubsection*{序号689}
\noindent 路径：\code{os/components/wateros-syscall/syscall-impl/impl-kernel/src/vfs_util.rs}\\
\noindent 行范围：1-70\\[0.4em]

\begin{rustcode}
//! [`VfsError`] 到 ABI [`ErrNo`] 的映射。
//! 本模块代码由AI完成

use abi::errno::ErrNo;
use vfs::api::{VfsError, VfsOpenFlags};
...
\end{rustcode}

\subsubsection*{序号690}
\noindent 路径：\code{os/components/wateros-syscall/syscall-impl/impl-kernel/src/vfs_util.rs}\\
\noindent 行范围：8-29\\[0.4em]

\begin{rustcode}
...
pub(crate) fn vfs_error_to_errno(err : VfsError) -> ErrNo {
    match err {
        VfsError::BadFd => ErrNo::EBADF,
        VfsError::Busy => ErrNo::EBUSY,
        VfsError::WouldBlock => ErrNo::EAGAIN,
...
\end{rustcode}

\subsubsection*{序号691}
\noindent 路径：\code{os/components/wateros-syscall/syscall-impl/impl-kernel/src/vfs_util.rs}\\
\noindent 行范围：33-38\\[0.4em]

\begin{rustcode}
...
pub(crate) fn vfs_io_at_error_to_errno(err : VfsError) -> ErrNo {
    match err {
        VfsError::Unsupported => ErrNo::ESPIPE,
        other => vfs_error_to_errno(other),
    }
...
\end{rustcode}

\subsubsection*{序号692}
\noindent 路径：\code{os/components/wateros-syscall/syscall-impl/impl-kernel/src/vfs_util.rs}\\
\noindent 行范围：42-70\\[0.4em]

\begin{rustcode}
...
pub(crate) fn linux_open_flags_to_vfs(flags : u32) -> VfsOpenFlags {
    const O_ACCMODE : u32 = 3;
    const O_WRONLY : u32 = 1;
    const O_RDWR : u32 = 2;
    const O_CREAT : u32 = 0o100;
...
\end{rustcode}

\subsubsection*{序号693}
\noindent 路径：\code{os/components/wateros-task/task-scheduler/scheduler-impl/impl-multi-class/src/rt_fifo_queue.rs}\\
\noindent 行范围：1-154\\[0.4em]

\begin{rustcode}
//! 本模块代码由AI完成
//! `SCHED_FIFO` 就绪队列：99 个优先级桶，每桶 FIFO，无时间片。

extern crate alloc;

...
\end{rustcode}

\subsubsection*{序号694}
\noindent 路径：\code{os/components/wateros-task/task-scheduler/scheduler-impl/impl-multi-class/src/rt_fifo_queue.rs}\\
\noindent 行范围：23-25\\[0.4em]

\begin{rustcode}
...
// 本结构代码由AI完成
pub struct RtFifoRunQueue {
    buckets : [VecDeque<TaskId>; RT_BUCKET_COUNT],
}

...
\end{rustcode}

\subsubsection*{序号695}
\noindent 路径：\code{os/components/wateros-task/task-scheduler/scheduler-impl/impl-multi-class/src/rt_rr_queue.rs}\\
\noindent 行范围：1-243\\[0.4em]

\begin{rustcode}
//! 本模块代码由AI完成
//! `SCHED_RR` 就绪队列：优先级桶 + 同优先级时间片轮转。

extern crate alloc;

...
\end{rustcode}

\subsubsection*{序号696}
\noindent 路径：\code{os/components/wateros-task/task-scheduler/scheduler-impl/impl-multi-class/src/rt_rr_queue.rs}\\
\noindent 行范围：27-32\\[0.4em]

\begin{rustcode}
...
pub enum RrTickAction {
    /// 继续运行当前 RR 任务。
    ContinueRunning,
    /// 时间片耗尽，让出给同优先级下一个 RR 任务。
    YieldToSamePriority,
...
\end{rustcode}

\subsubsection*{序号697}
\noindent 路径：\code{os/components/wateros-task/task-scheduler/scheduler-impl/impl-multi-class/src/rt_rr_queue.rs}\\
\noindent 行范围：36-40\\[0.4em]

\begin{rustcode}
...
pub struct RtRrRunQueue {
    buckets : [VecDeque<TaskId>; RT_BUCKET_COUNT],
    current : Option<(TaskId, i32)>,
    remaining_ticks : u64,
}
...
\end{rustcode}

\subsubsection*{序号698}
\noindent 路径：\code{os/components/wateros-utils/src/lib.rs}\\
\noindent 行范围：1-22\\[0.4em]

\begin{rustcode}
#![no_std]
//! 通用小工具与可复用例程的聚合 crate（当前为早期占位）。
//!
//! 后续可在此集中与内核其它组件无强耦合的纯函数或数据结构；
//! 避免反向依赖 `wateros-base` 等平台类型，以保持依赖 DAG 清晰。
...
\end{rustcode}

\subsubsection*{序号699}
\noindent 路径：\code{os/components/wateros-utils/src/lib.rs}\\
\noindent 行范围：11\\[0.4em]

\begin{rustcode}
...
// 本方法代码由AI完成
#[inline]
pub fn add(left : u64, right : u64) -> u64 { left + right }

#[cfg(test)]
...
\end{rustcode}

\subsubsection*{序号700}
\noindent 路径：\code{os/components/wateros-vfs/vfs-impl/impl-fd-session/src/char_dev_handle.rs}\\
\noindent 行范围：1-371\\[0.4em]

\begin{rustcode}
//! 字符设备 [`VfsIoHandle`]：包装 [`SharedCharacterDevice`]。
//! 本模块代码由AI完成

extern crate alloc;

...
\end{rustcode}

\subsubsection*{序号701}
\noindent 路径：\code{os/components/wateros-vfs/vfs-impl/impl-fd-session/src/char_dev_handle.rs}\\
\noindent 行范围：16\\[0.4em]

\begin{rustcode}
...
/// QEMU 无主机输入时，经若干次 poll 后注入的合成密码（供 LTP ask_password.sh 等非交互场景）。
// 本变量代码由AI完成
const HEADLESS_STUB_INPUT: &[u8] = b"password\n";
// 本变量代码由AI完成
const HEADLESS_POLL_THRESHOLD: u64 = 2;
...
\end{rustcode}

\subsubsection*{序号702}
\noindent 路径：\code{os/components/wateros-vfs/vfs-impl/impl-fd-session/src/char_dev_handle.rs}\\
\noindent 行范围：18\\[0.4em]

\begin{rustcode}
...
const HEADLESS_STUB_INPUT: &[u8] = b"password\n";
// 本变量代码由AI完成
const HEADLESS_POLL_THRESHOLD: u64 = 2;

// 本变量代码由AI完成
...
\end{rustcode}

\subsubsection*{序号703}
\noindent 路径：\code{os/components/wateros-vfs/vfs-impl/impl-fd-session/src/char_dev_handle.rs}\\
\noindent 行范围：21\\[0.4em]

\begin{rustcode}
...

// 本变量代码由AI完成
static HEADLESS_STUB_OFFSET: AtomicU64 = AtomicU64::new(0);
// 本变量代码由AI完成
static HEADLESS_SERIAL_POLLS: AtomicU64 = AtomicU64::new(0);
...
\end{rustcode}

\subsubsection*{序号704}
\noindent 路径：\code{os/components/wateros-vfs/vfs-impl/impl-fd-session/src/char_dev_handle.rs}\\
\noindent 行范围：23\\[0.4em]

\begin{rustcode}
...
static HEADLESS_STUB_OFFSET: AtomicU64 = AtomicU64::new(0);
// 本变量代码由AI完成
static HEADLESS_SERIAL_POLLS: AtomicU64 = AtomicU64::new(0);

// 本方法代码由AI完成
...
\end{rustcode}

\subsubsection*{序号705}
\noindent 路径：\code{os/components/wateros-vfs/vfs-impl/impl-fd-session/src/char_dev_handle.rs}\\
\noindent 行范围：26-33\\[0.4em]

\begin{rustcode}
...
fn map_driver_err(e: DriverError) -> VfsError {
    match e {
        DriverError::Unsupported => VfsError::Unsupported,
        DriverError::InvalidParam => VfsError::InvalidPath,
        DriverError::NotFound => VfsError::NotFound,
...
\end{rustcode}

\subsubsection*{序号706}
\noindent 路径：\code{os/components/wateros-vfs/vfs-impl/impl-fd-session/src/char_dev_handle.rs}\\
\noindent 行范围：36-49\\[0.4em]

\begin{rustcode}
...
fn char_metadata(mode: u16, inode: u64) -> VfsMetadata {
    VfsMetadata {
        node_type: VfsNodeType::Special,
        size: 0,
        mode,
...
\end{rustcode}

\subsubsection*{序号707}
\noindent 路径：\code{os/components/wateros-vfs/vfs-impl/impl-fd-session/src/char_dev_handle.rs}\\
\noindent 行范围：52-59\\[0.4em]

\begin{rustcode}
...
fn path_inode(path: &str) -> u64 {
    let mut hash = 0xcbf2_9ce4_8422_2325u64;
    for byte in path.as_bytes() {
        hash ^= u64::from(*byte);
        hash = hash.wrapping_mul(0x100_0000_01b3);
...
\end{rustcode}

\subsubsection*{序号708}
\noindent 路径：\code{os/components/wateros-vfs/vfs-impl/impl-fd-session/src/char_dev_handle.rs}\\
\noindent 行范围：63-70\\[0.4em]

\begin{rustcode}
...
pub struct CharDevHandle {
    device: SharedCharacterDevice,
    stdin_eof: bool,
    rtc: bool,
    nonblocking: Cell<bool>,
...
\end{rustcode}

\subsubsection*{序号709}
\noindent 路径：\code{os/components/wateros-vfs/vfs-impl/impl-fd-session/src/char_dev_handle.rs}\\
\noindent 行范围：74-83\\[0.4em]

\begin{rustcode}
...
    pub fn new(device: SharedCharacterDevice, stdin_eof: bool) -> Self {
        Self {
            device,
            stdin_eof,
            rtc: false,
...
\end{rustcode}

\subsubsection*{序号710}
\noindent 路径：\code{os/components/wateros-vfs/vfs-impl/impl-fd-session/src/char_dev_handle.rs}\\
\noindent 行范围：86-95\\[0.4em]

\begin{rustcode}
...
    pub fn new_rtc(device: SharedCharacterDevice) -> Self {
        Self {
            device,
            stdin_eof: false,
            rtc: true,
...
\end{rustcode}

\subsubsection*{序号711}
\noindent 路径：\code{os/components/wateros-vfs/vfs-impl/impl-fd-session/src/char_dev_handle.rs}\\
\noindent 行范围：98-117\\[0.4em]

\begin{rustcode}
...
    pub fn from_devfs_path(device: SharedCharacterDevice, path: &str) -> Self {
        if path == "/dev/null" {
            Self {
                device,
                stdin_eof: false,
...
\end{rustcode}

\subsubsection*{序号712}
\noindent 路径：\code{os/components/wateros-vfs/vfs-impl/impl-fd-session/src/char_dev_handle.rs}\\
\noindent 行范围：120-122\\[0.4em]

\begin{rustcode}
...
// 本方法代码由AI完成
    pub fn new_stdin(device: SharedCharacterDevice) -> Self {
        Self::new(device, true)
    }

...
\end{rustcode}

\subsubsection*{序号713}
\noindent 路径：\code{os/components/wateros-vfs/vfs-impl/impl-fd-session/src/char_dev_handle.rs}\\
\noindent 行范围：125-127\\[0.4em]

\begin{rustcode}
...
// 本方法代码由AI完成
    pub fn new_stdout(device: SharedCharacterDevice) -> Self {
        Self::new(device, false)
    }
}
...
\end{rustcode}

\subsubsection*{序号714}
\noindent 路径：\code{os/components/wateros-vfs/vfs-impl/impl-fd-session/src/char_dev_handle.rs}\\
\noindent 行范围：131-133\\[0.4em]

\begin{rustcode}
...
// 本方法代码由AI完成
fn headless_stub_pending() -> bool {
    HEADLESS_STUB_OFFSET.load(Ordering::Relaxed) < HEADLESS_STUB_INPUT.len() as u64
}

...
\end{rustcode}

\subsubsection*{序号715}
\noindent 路径：\code{os/components/wateros-vfs/vfs-impl/impl-fd-session/src/char_dev_handle.rs}\\
\noindent 行范围：136-144\\[0.4em]

\begin{rustcode}
...
fn arm_headless_stub_if_needed() {
    if headless_stub_pending() {
        return;
    }
    let polls = HEADLESS_SERIAL_POLLS.fetch_add(1, Ordering::Relaxed);
...
\end{rustcode}

\subsubsection*{序号716}
\noindent 路径：\code{os/components/wateros-vfs/vfs-impl/impl-fd-session/src/char_dev_handle.rs}\\
\noindent 行范围：147-156\\[0.4em]

\begin{rustcode}
...
fn read_headless_stub(buf: &mut [u8]) -> usize {
    let offset = HEADLESS_STUB_OFFSET.load(Ordering::Relaxed) as usize;
    if offset >= HEADLESS_STUB_INPUT.len() {
        return 0;
    }
...
\end{rustcode}

\subsubsection*{序号717}
\noindent 路径：\code{os/components/wateros-vfs/vfs-impl/impl-fd-session/src/char_dev_handle.rs}\\
\noindent 行范围：159-173\\[0.4em]

\begin{rustcode}
...
fn try_read_serial_tty(device: &SharedCharacterDevice, buf: &mut [u8]) -> VfsResult<usize> {
    let mut guard = device.lock();
    match guard.read(buf) {
        Ok(n) if n > 0 => Ok(n),
        Ok(_) | Err(DriverError::Unsupported) => {
...
\end{rustcode}

\subsubsection*{序号718}
\noindent 路径：\code{os/components/wateros-vfs/vfs-impl/impl-fd-session/src/char_dev_handle.rs}\\
\noindent 行范围：176-203\\[0.4em]

\begin{rustcode}
...
fn serial_poll_revents(device: &SharedCharacterDevice, events: i16) -> VfsResult<i16> {
// 本变量代码由AI完成
    const POLLIN: i16 = 0x001;
// 本变量代码由AI完成
    const POLLOUT: i16 = 0x004;
...
\end{rustcode}

\subsubsection*{序号719}
\noindent 路径：\code{os/components/wateros-vfs/vfs-impl/impl-fd-session/src/char_dev_handle.rs}\\
\noindent 行范围：178\\[0.4em]

\begin{rustcode}
...
fn serial_poll_revents(device: &SharedCharacterDevice, events: i16) -> VfsResult<i16> {
// 本变量代码由AI完成
    const POLLIN: i16 = 0x001;
// 本变量代码由AI完成
    const POLLOUT: i16 = 0x004;
...
\end{rustcode}

\subsubsection*{序号720}
\noindent 路径：\code{os/components/wateros-vfs/vfs-impl/impl-fd-session/src/char_dev_handle.rs}\\
\noindent 行范围：180\\[0.4em]

\begin{rustcode}
...
    const POLLIN: i16 = 0x001;
// 本变量代码由AI完成
    const POLLOUT: i16 = 0x004;
    let mut guard = device.lock();
    let mut revents = guard.poll_revents(events).map_err(map_driver_err)?;
...
\end{rustcode}

\subsubsection*{序号721}
\noindent 路径：\code{os/components/wateros-vfs/vfs-impl/impl-fd-session/src/char_dev_handle.rs}\\
\noindent 行范围：207-237\\[0.4em]

\begin{rustcode}
...
    fn read(&mut self, buf: &mut [u8]) -> VfsResult<usize> {
        if buf.is_empty() {
            return Ok(0);
        }
        if self.rtc {
...
\end{rustcode}

\subsubsection*{序号722}
\noindent 路径：\code{os/components/wateros-vfs/vfs-impl/impl-fd-session/src/char_dev_handle.rs}\\
\noindent 行范围：231\\[0.4em]

\begin{rustcode}
...
                Err(VfsError::WouldBlock) => {
// 本变量代码由AI完成
                    const POLLIN: i16 = 0x001;
                    self.poll_wait_for_ticks(POLLIN, 1, &mut || true)?;
                }
...
\end{rustcode}

\subsubsection*{序号723}
\noindent 路径：\code{os/components/wateros-vfs/vfs-impl/impl-fd-session/src/char_dev_handle.rs}\\
\noindent 行范围：240-243\\[0.4em]

\begin{rustcode}
...
// 本方法代码由AI完成
    fn write(&mut self, buf: &[u8]) -> VfsResult<usize> {
        let mut guard = self.device.lock();
        guard.write(buf).map_err(map_driver_err)
    }
...
\end{rustcode}

\subsubsection*{序号724}
\noindent 路径：\code{os/components/wateros-vfs/vfs-impl/impl-fd-session/src/char_dev_handle.rs}\\
\noindent 行范围：246-266\\[0.4em]

\begin{rustcode}
...
    fn poll_revents(&mut self, events: i16) -> VfsResult<i16> {
        if self.rtc {
            let mut guard = self.device.lock();
            return guard.poll_revents(events).map_err(map_driver_err);
        }
...
\end{rustcode}

\subsubsection*{序号725}
\noindent 路径：\code{os/components/wateros-vfs/vfs-impl/impl-fd-session/src/char_dev_handle.rs}\\
\noindent 行范围：253\\[0.4em]

\begin{rustcode}
...
        if self.stdin_eof {
// 本变量代码由AI完成
            const POLLIN: i16 = 0x001;
// 本变量代码由AI完成
            const POLLOUT: i16 = 0x004;
...
\end{rustcode}

\subsubsection*{序号726}
\noindent 路径：\code{os/components/wateros-vfs/vfs-impl/impl-fd-session/src/char_dev_handle.rs}\\
\noindent 行范围：255\\[0.4em]

\begin{rustcode}
...
            const POLLIN: i16 = 0x001;
// 本变量代码由AI完成
            const POLLOUT: i16 = 0x004;
            let mut revents = 0i16;
            if events & POLLIN != 0 {
...
\end{rustcode}

\subsubsection*{序号727}
\noindent 路径：\code{os/components/wateros-vfs/vfs-impl/impl-fd-session/src/char_dev_handle.rs}\\
\noindent 行范围：269-292\\[0.4em]

\begin{rustcode}
...
    fn poll_wait_for_ticks(
        &mut self,
        events: i16,
        timeout_ticks: u64,
        still_waiting: &mut dyn FnMut() -> bool,
...
\end{rustcode}

\subsubsection*{序号728}
\noindent 路径：\code{os/components/wateros-vfs/vfs-impl/impl-fd-session/src/char_dev_handle.rs}\\
\noindent 行范围：279\\[0.4em]

\begin{rustcode}
...
        }
// 本变量代码由AI完成
        const POLLIN: i16 = 0x001;
        if events & POLLIN == 0 {
            return Ok(());
...
\end{rustcode}

\subsubsection*{序号729}
\noindent 路径：\code{os/components/wateros-vfs/vfs-impl/impl-fd-session/src/char_dev_handle.rs}\\
\noindent 行范围：295-301\\[0.4em]

\begin{rustcode}
...
    fn ioctl(&mut self, request: usize, arg: usize) -> VfsResult<isize> {
        if !self.rtc {
            return Err(VfsError::Unsupported);
        }
        let mut guard = self.device.lock();
...
\end{rustcode}

\subsubsection*{序号730}
\noindent 路径：\code{os/components/wateros-vfs/vfs-impl/impl-fd-session/src/char_dev_handle.rs}\\
\noindent 行范围：314-316\\[0.4em]

\begin{rustcode}
...
// 本方法代码由AI完成
    fn metadata(&self) -> VfsResult<VfsMetadata> {
        Ok(char_metadata(self.mode, self.inode))
    }

...
\end{rustcode}

\subsubsection*{序号731}
\noindent 路径：\code{os/components/wateros-vfs/vfs-impl/impl-fd-session/src/char_dev_handle.rs}\\
\noindent 行范围：319-328\\[0.4em]

\begin{rustcode}
...
    fn duplicate(&self) -> VfsResult<Box<dyn VfsIoHandle>> {
        Ok(Box::new(Self {
            device: self.device.clone(),
            stdin_eof: self.stdin_eof,
            rtc: self.rtc,
...
\end{rustcode}

\subsubsection*{序号732}
\noindent 路径：\code{os/components/wateros-vfs/vfs-impl/impl-fd-session/src/char_dev_handle.rs}\\
\noindent 行范围：331-337\\[0.4em]

\begin{rustcode}
...
    fn open_status_flags(&self) -> u32 {
        if self.nonblocking.get() {
            0o0004000
        } else {
            0
...
\end{rustcode}

\subsubsection*{序号733}
\noindent 路径：\code{os/components/wateros-vfs/vfs-impl/impl-fd-session/src/char_dev_handle.rs}\\
\noindent 行范围：340-343\\[0.4em]

\begin{rustcode}
...
// 本方法代码由AI完成
    fn set_open_status_flags(&mut self, flags: u32) -> VfsResult<()> {
        self.nonblocking.set(flags & 0o0004000 != 0);
        Ok(())
    }
...
\end{rustcode}

\subsubsection*{序号734}
\noindent 路径：\code{os/components/wateros-vfs/vfs-impl/impl-fd-session/src/char_dev_handle.rs}\\
\noindent 行范围：353-365\\[0.4em]

\begin{rustcode}
...
fn mode_for_devfs_path(path: &str) -> u16 {
    if path == "/dev/null" {
        0o20666
    } else if path == "/dev/random" || path == "/dev/urandom" {
        0o20666
...
\end{rustcode}

\subsubsection*{序号735}
\noindent 路径：\code{os/components/wateros-vfs/vfs-impl/impl-fd-session/src/char_dev_handle.rs}\\
\noindent 行范围：369-371\\[0.4em]

\begin{rustcode}
...
/// 未打开 fd 时按 devfs 路径返回字符设备元数据（`fstatat` / `faccessat`）。
// 本方法代码由AI完成
pub fn metadata_for_devfs_path(path: &str) -> VfsMetadata {
    char_metadata(mode_for_devfs_path(path), path_inode(path))
}
\end{rustcode}

\subsubsection*{序号736}
\noindent 路径：\code{os/components/wateros-vfs/vfs-impl/impl-fd-session/src/cwd.rs}\\
\noindent 行范围：1-174\\[0.4em]

\begin{rustcode}
//! 以 [`task::TaskId`] 为 key 的 per-task 工作目录字符串。
//! 本模块代码由AI完成

extern crate alloc;

...
\end{rustcode}

\subsubsection*{序号737}
\noindent 路径：\code{os/components/wateros-vfs/vfs-impl/impl-fd-session/src/cwd.rs}\\
\noindent 行范围：15-21\\[0.4em]

\begin{rustcode}
...
pub struct PerTaskCwdRegistry {
    cwd_tables: BTreeMap<task::TaskId, String>,
    exe_paths: BTreeMap<task::TaskId, String>,
    argv_vectors: BTreeMap<task::TaskId, Vec<String>>,
    owners: BTreeMap<task::TaskId, task::TaskId>,
...
\end{rustcode}

\subsubsection*{序号738}
\noindent 路径：\code{os/components/wateros-vfs/vfs-impl/impl-fd-session/src/cwd.rs}\\
\noindent 行范围：36-40\\[0.4em]

\begin{rustcode}
...
    fn ensure_owner(&mut self, task_id: task::TaskId) -> task::TaskId {
        self.owners.entry(task_id).or_insert(task_id);
        self.ref_counts.entry(task_id).or_insert(1);
        self.effective_owner(task_id)
    }
...
\end{rustcode}

\subsubsection*{序号739}
\noindent 路径：\code{os/components/wateros-vfs/vfs-impl/impl-fd-session/src/cwd.rs}\\
\noindent 行范围：44-49\\[0.4em]

\begin{rustcode}
...
    fn effective_owner(&self, task_id: task::TaskId) -> task::TaskId {
        self.owners
            .get(&task_id)
            .copied()
            .unwrap_or(task_id)
...
\end{rustcode}

\subsubsection*{序号740}
\noindent 路径：\code{os/components/wateros-vfs/vfs-impl/impl-fd-session/src/cwd.rs}\\
\noindent 行范围：52-61\\[0.4em]

\begin{rustcode}
...
    fn release_owner(&mut self, task_id: task::TaskId) -> Option<task::TaskId> {
        let owner = self.owners.remove(&task_id)?;
        if let Some(count) = self.ref_counts.get_mut(&owner) {
            *count = count.saturating_sub(1);
            if *count == 0 {
...
\end{rustcode}

\subsubsection*{序号741}
\noindent 路径：\code{os/components/wateros-vfs/vfs-impl/impl-fd-session/src/cwd.rs}\\
\noindent 行范围：65-68\\[0.4em]

\begin{rustcode}
...
// 本方法代码由AI完成
    pub fn init_task_cwd(&mut self, task_id: task::TaskId) {
        let owner = self.ensure_owner(task_id);
        self.cwd_tables.insert(owner, String::from("/"));
    }
...
\end{rustcode}

\subsubsection*{序号742}
\noindent 路径：\code{os/components/wateros-vfs/vfs-impl/impl-fd-session/src/cwd.rs}\\
\noindent 行范围：73-79\\[0.4em]

\begin{rustcode}
...
    pub fn get_cwd(&self, task_id: task::TaskId) -> &str {
        let owner = self.effective_owner(task_id);
        self.cwd_tables
            .get(&owner)
            .map(String::as_str)
...
\end{rustcode}

\subsubsection*{序号743}
\noindent 路径：\code{os/components/wateros-vfs/vfs-impl/impl-fd-session/src/cwd.rs}\\
\noindent 行范围：83-86\\[0.4em]

\begin{rustcode}
...
// 本方法代码由AI完成
    pub fn ensure_task_cwd(&mut self, task_id: task::TaskId) {
        let owner = self.ensure_owner(task_id);
        self.cwd_tables.entry(owner).or_insert_with(|| String::from("/"));
    }
...
\end{rustcode}

\subsubsection*{序号744}
\noindent 路径：\code{os/components/wateros-vfs/vfs-impl/impl-fd-session/src/cwd.rs}\\
\noindent 行范围：89-93\\[0.4em]

\begin{rustcode}
...
    pub fn get_cwd_mut(&mut self, task_id: task::TaskId) -> &mut String {
        self.ensure_task_cwd(task_id);
        let owner = self.effective_owner(task_id);
        self.cwd_tables.get_mut(&owner).expect("init_task_cwd")
    }
...
\end{rustcode}

\subsubsection*{序号745}
\noindent 路径：\code{os/components/wateros-vfs/vfs-impl/impl-fd-session/src/cwd.rs}\\
\noindent 行范围：97-111\\[0.4em]

\begin{rustcode}
...
    pub fn drop_task(&mut self, task_id: task::TaskId) {
        let Some(owner) = self.release_owner(task_id) else {
            return;
        };
        if self.ref_counts.get(&owner).copied().unwrap_or(0) == 0 {
...
\end{rustcode}

\subsubsection*{序号746}
\noindent 路径：\code{os/components/wateros-vfs/vfs-impl/impl-fd-session/src/cwd.rs}\\
\noindent 行范围：114-133\\[0.4em]

\begin{rustcode}
...
    pub fn copy_cwd_from_parent(&mut self, child: task::TaskId, parent: task::TaskId) {
        let parent_cwd = self.get_cwd(parent).to_string();
        let parent_exe = self.get_exe_path(parent).map(ToString::to_string);
        let parent_argv = self.get_argv(parent).map(|v| v.to_vec());
        if self.owners.contains_key(&child) {
...
\end{rustcode}

\subsubsection*{序号747}
\noindent 路径：\code{os/components/wateros-vfs/vfs-impl/impl-fd-session/src/cwd.rs}\\
\noindent 行范围：136-145\\[0.4em]

\begin{rustcode}
...
    pub fn share_cwd_from_parent(&mut self, child: task::TaskId, parent: task::TaskId) {
        self.ensure_task_cwd(parent);
        if self.owners.contains_key(&child) {
            self.drop_task(child);
        }
...
\end{rustcode}

\subsubsection*{序号748}
\noindent 路径：\code{os/components/wateros-vfs/vfs-impl/impl-fd-session/src/cwd.rs}\\
\noindent 行范围：148-151\\[0.4em]

\begin{rustcode}
...
// 本方法代码由AI完成
    pub fn set_exe_path(&mut self, task_id: task::TaskId, exe_path: &str) {
        let owner = self.ensure_owner(task_id);
        self.exe_paths.insert(owner, String::from(exe_path));
    }
...
\end{rustcode}

\subsubsection*{序号749}
\noindent 路径：\code{os/components/wateros-vfs/vfs-impl/impl-fd-session/src/cwd.rs}\\
\noindent 行范围：154-159\\[0.4em]

\begin{rustcode}
...
    pub fn get_exe_path(&self, task_id: task::TaskId) -> Option<&str> {
        let owner = self.effective_owner(task_id);
        self.exe_paths
            .get(&owner)
            .map(String::as_str)
...
\end{rustcode}

\subsubsection*{序号750}
\noindent 路径：\code{os/components/wateros-vfs/vfs-impl/impl-fd-session/src/cwd.rs}\\
\noindent 行范围：162-165\\[0.4em]

\begin{rustcode}
...
// 本方法代码由AI完成
    pub fn set_argv(&mut self, task_id: task::TaskId, argv: Vec<String>) {
        let owner = self.ensure_owner(task_id);
        self.argv_vectors.insert(owner, argv);
    }
...
\end{rustcode}

\subsubsection*{序号751}
\noindent 路径：\code{os/components/wateros-vfs/vfs-impl/impl-fd-session/src/cwd.rs}\\
\noindent 行范围：168-173\\[0.4em]

\begin{rustcode}
...
    pub fn get_argv(&self, task_id: task::TaskId) -> Option<&[String]> {
        let owner = self.effective_owner(task_id);
        self.argv_vectors
            .get(&owner)
            .map(Vec::as_slice)
...
\end{rustcode}

\subsubsection*{序号752}
\noindent 路径：\code{os/components/wateros-vfs/vfs-impl/impl-fd-session/src/file_lock.rs}\\
\noindent 行范围：1-578\\[0.4em]

\begin{rustcode}
//! Per-inode POSIX 记录锁与 `flock(2)` 整文件锁状态。
//! 本模块代码由AI完成

extern crate alloc;

...
\end{rustcode}

\subsubsection*{序号753}
\noindent 路径：\code{os/components/wateros-vfs/vfs-impl/impl-fd-session/src/file_lock.rs}\\
\noindent 行范围：27-33\\[0.4em]

\begin{rustcode}
...
pub struct Flock {
    pub l_type: i16,
    pub l_whence: i16,
    pub l_start: i64,
    pub l_len: i64,
...
\end{rustcode}

\subsubsection*{序号754}
\noindent 路径：\code{os/components/wateros-vfs/vfs-impl/impl-fd-session/src/file_lock.rs}\\
\noindent 行范围：43-46\\[0.4em]

\begin{rustcode}
...
// 本结构代码由AI完成
pub struct InodeKey {
    pub mount_id: u64,
    pub inode: u64,
}
...
\end{rustcode}

\subsubsection*{序号755}
\noindent 路径：\code{os/components/wateros-vfs/vfs-impl/impl-fd-session/src/file_lock.rs}\\
\noindent 行范围：56-61\\[0.4em]

\begin{rustcode}
...
struct PosixLock {
    pid: ProcessId,
    typ: LockType,
    start: u64,
    len: u64,
...
\end{rustcode}

\subsubsection*{序号756}
\noindent 路径：\code{os/components/wateros-vfs/vfs-impl/impl-fd-session/src/file_lock.rs}\\
\noindent 行范围：64-67\\[0.4em]

\begin{rustcode}
...
// 本结构代码由AI完成
struct FlockState {
    shared_holders: Vec<ProcessId>,
    exclusive: Option<ProcessId>,
}
...
\end{rustcode}

\subsubsection*{序号757}
\noindent 路径：\code{os/components/wateros-vfs/vfs-impl/impl-fd-session/src/file_lock.rs}\\
\noindent 行范围：71-76\\[0.4em]

\begin{rustcode}
...
    fn new() -> Self {
        Self {
            shared_holders: Vec::new(),
            exclusive: None,
        }
...
\end{rustcode}

\subsubsection*{序号758}
\noindent 路径：\code{os/components/wateros-vfs/vfs-impl/impl-fd-session/src/file_lock.rs}\\
\noindent 行范围：79-84\\[0.4em]

\begin{rustcode}
...
    fn clear_pid(&mut self, pid: ProcessId) {
        if self.exclusive == Some(pid) {
            self.exclusive = None;
        }
        self.shared_holders.retain(|p| *p != pid);
...
\end{rustcode}

\subsubsection*{序号759}
\noindent 路径：\code{os/components/wateros-vfs/vfs-impl/impl-fd-session/src/file_lock.rs}\\
\noindent 行范围：88-91\\[0.4em]

\begin{rustcode}
...
// 本结构代码由AI完成
struct InodeLockData {
    posix: Vec<PosixLock>,
    flock: FlockState,
}
...
\end{rustcode}

\subsubsection*{序号760}
\noindent 路径：\code{os/components/wateros-vfs/vfs-impl/impl-fd-session/src/file_lock.rs}\\
\noindent 行范围：94-97\\[0.4em]

\begin{rustcode}
...
// 本结构代码由AI完成
struct InodeLocks {
    data: Mutex<InodeLockData>,
    wait: WaitQueue,
}
...
\end{rustcode}

\subsubsection*{序号761}
\noindent 路径：\code{os/components/wateros-vfs/vfs-impl/impl-fd-session/src/file_lock.rs}\\
\noindent 行范围：101-109\\[0.4em]

\begin{rustcode}
...
    fn new() -> Arc<Self> {
        Arc::new(Self {
            data: Mutex::new(InodeLockData {
                posix: Vec::new(),
                flock: FlockState::new(),
...
\end{rustcode}

\subsubsection*{序号762}
\noindent 路径：\code{os/components/wateros-vfs/vfs-impl/impl-fd-session/src/file_lock.rs}\\
\noindent 行范围：113\\[0.4em]

\begin{rustcode}
...

// 本变量代码由AI完成
static LOCK_TABLE: Mutex<BTreeMap<InodeKey, Arc<InodeLocks>>> = Mutex::new(BTreeMap::new());

// 本方法代码由AI完成
...
\end{rustcode}

\subsubsection*{序号763}
\noindent 路径：\code{os/components/wateros-vfs/vfs-impl/impl-fd-session/src/file_lock.rs}\\
\noindent 行范围：116-122\\[0.4em]

\begin{rustcode}
...
fn get_inode_locks(key: &InodeKey) -> Arc<InodeLocks> {
    let mut table = LOCK_TABLE.lock();
    table
        .entry(*key)
        .or_insert_with(InodeLocks::new)
...
\end{rustcode}

\subsubsection*{序号764}
\noindent 路径：\code{os/components/wateros-vfs/vfs-impl/impl-fd-session/src/file_lock.rs}\\
\noindent 行范围：125-134\\[0.4em]

\begin{rustcode}
...
fn drop_inode_if_empty(key: &InodeKey, locks: &Arc<InodeLocks>) {
    let data = locks.data.lock();
    if data.posix.is_empty()
        && data.flock.exclusive.is_none()
        && data.flock.shared_holders.is_empty()
...
\end{rustcode}

\subsubsection*{序号765}
\noindent 路径：\code{os/components/wateros-vfs/vfs-impl/impl-fd-session/src/file_lock.rs}\\
\noindent 行范围：139-147\\[0.4em]

\begin{rustcode}
...
pub fn inode_key_from_metadata(meta: &VfsMetadata) -> Option<InodeKey> {
    if meta.node_type != VfsNodeType::File {
        return None;
    }
    Some(InodeKey {
...
\end{rustcode}

\subsubsection*{序号766}
\noindent 路径：\code{os/components/wateros-vfs/vfs-impl/impl-fd-session/src/file_lock.rs}\\
\noindent 行范围：150-156\\[0.4em]

\begin{rustcode}
...
fn lock_end(start: u64, len: u64) -> u64 {
    if len == 0 {
        u64::MAX
    } else {
        start.saturating_add(len.saturating_sub(1))
...
\end{rustcode}

\subsubsection*{序号767}
\noindent 路径：\code{os/components/wateros-vfs/vfs-impl/impl-fd-session/src/file_lock.rs}\\
\noindent 行范围：159-163\\[0.4em]

\begin{rustcode}
...
fn ranges_overlap(a_start: u64, a_len: u64, b_start: u64, b_len: u64) -> bool {
    let a_end = lock_end(a_start, a_len);
    let b_end = lock_end(b_start, b_len);
    a_start <= b_end && b_start <= a_end
}
...
\end{rustcode}

\subsubsection*{序号768}
\noindent 路径：\code{os/components/wateros-vfs/vfs-impl/impl-fd-session/src/file_lock.rs}\\
\noindent 行范围：166-172\\[0.4em]

\begin{rustcode}
...
fn flock_type_from_i16(raw: i16) -> VfsResult<LockType> {
    match raw {
        F_RDLCK => Ok(LockType::Read),
        F_WRLCK => Ok(LockType::Write),
        _ => Err(VfsError::Unsupported),
...
\end{rustcode}

\subsubsection*{序号769}
\noindent 路径：\code{os/components/wateros-vfs/vfs-impl/impl-fd-session/src/file_lock.rs}\\
\noindent 行范围：175-192\\[0.4em]

\begin{rustcode}
...
fn intersection_range(
    probe_start: u64,
    probe_len: u64,
    lock_start: u64,
    lock_len: u64,
...
\end{rustcode}

\subsubsection*{序号770}
\noindent 路径：\code{os/components/wateros-vfs/vfs-impl/impl-fd-session/src/file_lock.rs}\\
\noindent 行范围：195-200\\[0.4em]

\begin{rustcode}
...
fn getlk_conflicts_with_probe(probe_typ: LockType, held: LockType) -> bool {
    match probe_typ {
        LockType::Read => held == LockType::Write,
        LockType::Write => true,
    }
...
\end{rustcode}

\subsubsection*{序号771}
\noindent 路径：\code{os/components/wateros-vfs/vfs-impl/impl-fd-session/src/file_lock.rs}\\
\noindent 行范围：204-226\\[0.4em]

\begin{rustcode}
...
fn find_getlk_conflict(
    locks: &[PosixLock],
    pid: ProcessId,
    probe_typ: LockType,
    probe_start: u64,
...
\end{rustcode}

\subsubsection*{序号772}
\noindent 路径：\code{os/components/wateros-vfs/vfs-impl/impl-fd-session/src/file_lock.rs}\\
\noindent 行范围：229-249\\[0.4em]

\begin{rustcode}
...
fn posix_conflict(
    locks: &[PosixLock],
    pid: ProcessId,
    typ: LockType,
    start: u64,
...
\end{rustcode}

\subsubsection*{序号773}
\noindent 路径：\code{os/components/wateros-vfs/vfs-impl/impl-fd-session/src/file_lock.rs}\\
\noindent 行范围：252-257\\[0.4em]

\begin{rustcode}
...
fn flock_blocks_posix(flock: &FlockState, typ: LockType) -> bool {
    match typ {
        LockType::Read => flock.exclusive.is_some(),
        LockType::Write => flock.exclusive.is_some() || !flock.shared_holders.is_empty(),
    }
...
\end{rustcode}

\subsubsection*{序号774}
\noindent 路径：\code{os/components/wateros-vfs/vfs-impl/impl-fd-session/src/file_lock.rs}\\
\noindent 行范围：260-270\\[0.4em]

\begin{rustcode}
...
fn posix_blocks_flock(locks: &[PosixLock], pid: ProcessId, op: usize) -> bool {
    if (op & LOCK_EX) != 0 {
        return locks.iter().any(|l| l.pid != pid);
    }
    if (op & LOCK_SH) != 0 {
...
\end{rustcode}

\subsubsection*{序号775}
\noindent 路径：\code{os/components/wateros-vfs/vfs-impl/impl-fd-session/src/file_lock.rs}\\
\noindent 行范围：273-307\\[0.4em]

\begin{rustcode}
...
fn split_lock_around_region(lock: PosixLock, cut_start: u64, cut_len: u64) -> Vec<PosixLock> {
    let cut_end = lock_end(cut_start, cut_len);
    let lock_end_pos = lock_end(lock.start, lock.len);
    if cut_start > lock_end_pos || cut_end < lock.start {
        return vec![lock];
...
\end{rustcode}

\subsubsection*{序号776}
\noindent 路径：\code{os/components/wateros-vfs/vfs-impl/impl-fd-session/src/file_lock.rs}\\
\noindent 行范围：310-325\\[0.4em]

\begin{rustcode}
...
fn remove_posix_for_pid_region(
    locks: &mut Vec<PosixLock>,
    pid: ProcessId,
    start: u64,
    len: u64,
...
\end{rustcode}

\subsubsection*{序号777}
\noindent 路径：\code{os/components/wateros-vfs/vfs-impl/impl-fd-session/src/file_lock.rs}\\
\noindent 行范围：328-330\\[0.4em]

\begin{rustcode}
...
// 本方法代码由AI完成
fn apply_posix_unlock(locks: &mut Vec<PosixLock>, pid: ProcessId, start: u64, len: u64) {
    remove_posix_for_pid_region(locks, pid, start, len);
}

...
\end{rustcode}

\subsubsection*{序号778}
\noindent 路径：\code{os/components/wateros-vfs/vfs-impl/impl-fd-session/src/file_lock.rs}\\
\noindent 行范围：333-347\\[0.4em]

\begin{rustcode}
...
fn apply_posix_lock(
    locks: &mut Vec<PosixLock>,
    pid: ProcessId,
    typ: LockType,
    start: u64,
...
\end{rustcode}

\subsubsection*{序号779}
\noindent 路径：\code{os/components/wateros-vfs/vfs-impl/impl-fd-session/src/file_lock.rs}\\
\noindent 行范围：350-353\\[0.4em]

\begin{rustcode}
...
// 本方法代码由AI完成
fn posix_would_block(data: &InodeLockData, pid: ProcessId, typ: LockType, start: u64, len: u64) -> bool {
    flock_blocks_posix(&data.flock, typ)
        || posix_conflict(&data.posix, pid, typ, start, len).is_some()
}
...
\end{rustcode}

\subsubsection*{序号780}
\noindent 路径：\code{os/components/wateros-vfs/vfs-impl/impl-fd-session/src/file_lock.rs}\\
\noindent 行范围：357-402\\[0.4em]

\begin{rustcode}
...
pub fn posix_getlk(
    key: &InodeKey,
    pid: ProcessId,
    flock: &mut Flock,
) -> VfsResult<()> {
...
\end{rustcode}

\subsubsection*{序号781}
\noindent 路径：\code{os/components/wateros-vfs/vfs-impl/impl-fd-session/src/file_lock.rs}\\
\noindent 行范围：406-455\\[0.4em]

\begin{rustcode}
...
pub fn posix_setlk(
    key: &InodeKey,
    pid: ProcessId,
    flock: &Flock,
    blocking: bool,
...
\end{rustcode}

\subsubsection*{序号782}
\noindent 路径：\code{os/components/wateros-vfs/vfs-impl/impl-fd-session/src/file_lock.rs}\\
\noindent 行范围：459-540\\[0.4em]

\begin{rustcode}
...
pub fn flock_op(key: &InodeKey, pid: ProcessId, operation: usize) -> VfsResult<()> {
    let nonblocking = (operation & LOCK_NB) != 0;
    let op = operation & !LOCK_NB;

    if op != LOCK_SH && op != LOCK_EX && op != LOCK_UN {
...
\end{rustcode}

\subsubsection*{序号783}
\noindent 路径：\code{os/components/wateros-vfs/vfs-impl/impl-fd-session/src/file_lock.rs}\\
\noindent 行范围：544-555\\[0.4em]

\begin{rustcode}
...
pub fn release_process_inode_locks(pid: ProcessId, key: &InodeKey) {
    let Some(locks) = LOCK_TABLE.lock().get(key).cloned() else {
        return;
    };
    {
...
\end{rustcode}

\subsubsection*{序号784}
\noindent 路径：\code{os/components/wateros-vfs/vfs-impl/impl-fd-session/src/file_lock.rs}\\
\noindent 行范围：559-578\\[0.4em]

\begin{rustcode}
...
pub fn inherit_process_locks(parent_pid: ProcessId, child_pid: ProcessId) {
    let table = LOCK_TABLE.lock();
    for locks in table.values() {
        let mut data = locks.data.lock();
        let parent_posix: Vec<PosixLock> = data
...
\end{rustcode}

\subsubsection*{序号785}
\noindent 路径：\code{os/components/wateros-vfs/vfs-impl/impl-fd-session/src/handles.rs}\\
\noindent 行范围：1-596\\[0.4em]

\begin{rustcode}
//! 控制台与 pipe 的 [`VfsIoHandle`] 实现。
//! 本模块代码由AI完成

extern crate alloc;

...
\end{rustcode}

\subsubsection*{序号786}
\noindent 路径：\code{os/components/wateros-vfs/vfs-impl/impl-fd-session/src/handles.rs}\\
\noindent 行范围：15\\[0.4em]

\begin{rustcode}
...

// 本变量代码由AI完成
static NEXT_PIPE_INODE: AtomicU64 = AtomicU64::new(1);
// 本变量代码由AI完成
static NEXT_STREAM_PAIR_INODE: AtomicU64 = AtomicU64::new(1);
...
\end{rustcode}

\subsubsection*{序号787}
\noindent 路径：\code{os/components/wateros-vfs/vfs-impl/impl-fd-session/src/handles.rs}\\
\noindent 行范围：17\\[0.4em]

\begin{rustcode}
...
static NEXT_PIPE_INODE: AtomicU64 = AtomicU64::new(1);
// 本变量代码由AI完成
static NEXT_STREAM_PAIR_INODE: AtomicU64 = AtomicU64::new(1);
// 本变量代码由AI完成
static URANDOM_STATE: AtomicU64 = AtomicU64::new(0x6a09_e667_f3bc_c909);
...
\end{rustcode}

\subsubsection*{序号788}
\noindent 路径：\code{os/components/wateros-vfs/vfs-impl/impl-fd-session/src/handles.rs}\\
\noindent 行范围：19\\[0.4em]

\begin{rustcode}
...
static NEXT_STREAM_PAIR_INODE: AtomicU64 = AtomicU64::new(1);
// 本变量代码由AI完成
static URANDOM_STATE: AtomicU64 = AtomicU64::new(0x6a09_e667_f3bc_c909);

// 本方法代码由AI完成
...
\end{rustcode}

\subsubsection*{序号789}
\noindent 路径：\code{os/components/wateros-vfs/vfs-impl/impl-fd-session/src/handles.rs}\\
\noindent 行范围：22-24\\[0.4em]

\begin{rustcode}
...
// 本方法代码由AI完成
fn special_meta(mode: u16, inode: u64) -> VfsMetadata {
    special_dev_meta(mode, inode, 0, 0x7fff_0001)
}

...
\end{rustcode}

\subsubsection*{序号790}
\noindent 路径：\code{os/components/wateros-vfs/vfs-impl/impl-fd-session/src/handles.rs}\\
\noindent 行范围：27-40\\[0.4em]

\begin{rustcode}
...
fn special_dev_meta(mode: u16, inode: u64, device_major: u32, device_minor: u32) -> VfsMetadata {
    VfsMetadata {
        node_type: VfsNodeType::Special,
        size: 0,
        mode,
...
\end{rustcode}

\subsubsection*{序号791}
\noindent 路径：\code{os/components/wateros-vfs/vfs-impl/impl-fd-session/src/handles.rs}\\
\noindent 行范围：45\\[0.4em]

\begin{rustcode}
...
#[derive(Debug, Clone, Copy, Default)]
// 本结构代码由AI完成
pub struct ConsoleInHandle;

impl VfsIoHandle for ConsoleInHandle {
...
\end{rustcode}

\subsubsection*{序号792}
\noindent 路径：\code{os/components/wateros-vfs/vfs-impl/impl-fd-session/src/handles.rs}\\
\noindent 行范围：49-51\\[0.4em]

\begin{rustcode}
...
// 本方法代码由AI完成
    fn read(&mut self, _buf: &mut [u8]) -> VfsResult<usize> {
        Ok(0)
    }

...
\end{rustcode}

\subsubsection*{序号793}
\noindent 路径：\code{os/components/wateros-vfs/vfs-impl/impl-fd-session/src/handles.rs}\\
\noindent 行范围：54-56\\[0.4em]

\begin{rustcode}
...
// 本方法代码由AI完成
    fn seek(&mut self, _offset: i64, _whence: VfsSeekWhence) -> VfsResult<u64> {
        Err(VfsError::Unsupported)
    }

...
\end{rustcode}

\subsubsection*{序号794}
\noindent 路径：\code{os/components/wateros-vfs/vfs-impl/impl-fd-session/src/handles.rs}\\
\noindent 行范围：59-61\\[0.4em]

\begin{rustcode}
...
// 本方法代码由AI完成
    fn metadata(&self) -> VfsResult<VfsMetadata> {
        Ok(special_meta(0o20666, 1))
    }

...
\end{rustcode}

\subsubsection*{序号795}
\noindent 路径：\code{os/components/wateros-vfs/vfs-impl/impl-fd-session/src/handles.rs}\\
\noindent 行范围：64-66\\[0.4em]

\begin{rustcode}
...
// 本方法代码由AI完成
    fn duplicate(&self) -> VfsResult<Box<dyn VfsIoHandle>> {
        Ok(Box::new(*self))
    }
}
...
\end{rustcode}

\subsubsection*{序号796}
\noindent 路径：\code{os/components/wateros-vfs/vfs-impl/impl-fd-session/src/handles.rs}\\
\noindent 行范围：72\\[0.4em]

\begin{rustcode}
...
#[derive(Debug, Clone, Copy, Default)]
// 本结构代码由AI完成
pub struct ConsoleOutHandle;

impl VfsIoHandle for ConsoleOutHandle {
...
\end{rustcode}

\subsubsection*{序号797}
\noindent 路径：\code{os/components/wateros-vfs/vfs-impl/impl-fd-session/src/handles.rs}\\
\noindent 行范围：76-79\\[0.4em]

\begin{rustcode}
...
// 本方法代码由AI完成
    fn write(&mut self, buf: &[u8]) -> VfsResult<usize> {
        console::write_raw_bytes(buf);
        Ok(buf.len())
    }
...
\end{rustcode}

\subsubsection*{序号798}
\noindent 路径：\code{os/components/wateros-vfs/vfs-impl/impl-fd-session/src/handles.rs}\\
\noindent 行范围：82-90\\[0.4em]

\begin{rustcode}
...
    fn poll_revents(&mut self, events: i16) -> VfsResult<i16> {
// 本变量代码由AI完成
        const POLLOUT: i16 = 0x004;
        if events & POLLOUT != 0 {
            Ok(POLLOUT)
...
\end{rustcode}

\subsubsection*{序号799}
\noindent 路径：\code{os/components/wateros-vfs/vfs-impl/impl-fd-session/src/handles.rs}\\
\noindent 行范围：84\\[0.4em]

\begin{rustcode}
...
    fn poll_revents(&mut self, events: i16) -> VfsResult<i16> {
// 本变量代码由AI完成
        const POLLOUT: i16 = 0x004;
        if events & POLLOUT != 0 {
            Ok(POLLOUT)
...
\end{rustcode}

\subsubsection*{序号800}
\noindent 路径：\code{os/components/wateros-vfs/vfs-impl/impl-fd-session/src/handles.rs}\\
\noindent 行范围：93-95\\[0.4em]

\begin{rustcode}
...
// 本方法代码由AI完成
    fn metadata(&self) -> VfsResult<VfsMetadata> {
        Ok(special_meta(0o20666, 1))
    }

...
\end{rustcode}

\subsubsection*{序号801}
\noindent 路径：\code{os/components/wateros-vfs/vfs-impl/impl-fd-session/src/handles.rs}\\
\noindent 行范围：98-100\\[0.4em]

\begin{rustcode}
...
// 本方法代码由AI完成
    fn duplicate(&self) -> VfsResult<Box<dyn VfsIoHandle>> {
        Ok(Box::new(*self))
    }
}
...
\end{rustcode}

\subsubsection*{序号802}
\noindent 路径：\code{os/components/wateros-vfs/vfs-impl/impl-fd-session/src/handles.rs}\\
\noindent 行范围：106\\[0.4em]

\begin{rustcode}
...
#[derive(Debug, Clone, Copy, Default)]
// 本结构代码由AI完成
pub struct NullDeviceHandle;

impl VfsIoHandle for NullDeviceHandle {
...
\end{rustcode}

\subsubsection*{序号803}
\noindent 路径：\code{os/components/wateros-vfs/vfs-impl/impl-fd-session/src/handles.rs}\\
\noindent 行范围：110-112\\[0.4em]

\begin{rustcode}
...
// 本方法代码由AI完成
    fn read(&mut self, _buf: &mut [u8]) -> VfsResult<usize> {
        Ok(0)
    }

...
\end{rustcode}

\subsubsection*{序号804}
\noindent 路径：\code{os/components/wateros-vfs/vfs-impl/impl-fd-session/src/handles.rs}\\
\noindent 行范围：115-117\\[0.4em]

\begin{rustcode}
...
// 本方法代码由AI完成
    fn write(&mut self, buf: &[u8]) -> VfsResult<usize> {
        Ok(buf.len())
    }

...
\end{rustcode}

\subsubsection*{序号805}
\noindent 路径：\code{os/components/wateros-vfs/vfs-impl/impl-fd-session/src/handles.rs}\\
\noindent 行范围：120-126\\[0.4em]

\begin{rustcode}
...
    fn poll_revents(&mut self, events: i16) -> VfsResult<i16> {
// 本变量代码由AI完成
        const POLLIN: i16 = 0x001;
// 本变量代码由AI完成
        const POLLOUT: i16 = 0x004;
...
\end{rustcode}

\subsubsection*{序号806}
\noindent 路径：\code{os/components/wateros-vfs/vfs-impl/impl-fd-session/src/handles.rs}\\
\noindent 行范围：122\\[0.4em]

\begin{rustcode}
...
    fn poll_revents(&mut self, events: i16) -> VfsResult<i16> {
// 本变量代码由AI完成
        const POLLIN: i16 = 0x001;
// 本变量代码由AI完成
        const POLLOUT: i16 = 0x004;
...
\end{rustcode}

\subsubsection*{序号807}
\noindent 路径：\code{os/components/wateros-vfs/vfs-impl/impl-fd-session/src/handles.rs}\\
\noindent 行范围：124\\[0.4em]

\begin{rustcode}
...
        const POLLIN: i16 = 0x001;
// 本变量代码由AI完成
        const POLLOUT: i16 = 0x004;
        Ok(events & (POLLIN | POLLOUT))
    }
...
\end{rustcode}

\subsubsection*{序号808}
\noindent 路径：\code{os/components/wateros-vfs/vfs-impl/impl-fd-session/src/handles.rs}\\
\noindent 行范围：129-131\\[0.4em]

\begin{rustcode}
...
// 本方法代码由AI完成
    fn metadata(&self) -> VfsResult<VfsMetadata> {
        Ok(special_dev_meta(0o20666, 2, 1, 3))
    }

...
\end{rustcode}

\subsubsection*{序号809}
\noindent 路径：\code{os/components/wateros-vfs/vfs-impl/impl-fd-session/src/handles.rs}\\
\noindent 行范围：134-136\\[0.4em]

\begin{rustcode}
...
// 本方法代码由AI完成
    fn duplicate(&self) -> VfsResult<Box<dyn VfsIoHandle>> {
        Ok(Box::new(*self))
    }
}
...
\end{rustcode}

\subsubsection*{序号810}
\noindent 路径：\code{os/components/wateros-vfs/vfs-impl/impl-fd-session/src/handles.rs}\\
\noindent 行范围：142\\[0.4em]

\begin{rustcode}
...
#[derive(Debug, Clone, Copy, Default)]
// 本结构代码由AI完成
pub struct ZeroDeviceHandle;

impl VfsIoHandle for ZeroDeviceHandle {
...
\end{rustcode}

\subsubsection*{序号811}
\noindent 路径：\code{os/components/wateros-vfs/vfs-impl/impl-fd-session/src/handles.rs}\\
\noindent 行范围：146-149\\[0.4em]

\begin{rustcode}
...
// 本方法代码由AI完成
    fn read(&mut self, buf: &mut [u8]) -> VfsResult<usize> {
        buf.fill(0);
        Ok(buf.len())
    }
...
\end{rustcode}

\subsubsection*{序号812}
\noindent 路径：\code{os/components/wateros-vfs/vfs-impl/impl-fd-session/src/handles.rs}\\
\noindent 行范围：152-154\\[0.4em]

\begin{rustcode}
...
// 本方法代码由AI完成
    fn write(&mut self, buf: &[u8]) -> VfsResult<usize> {
        Ok(buf.len())
    }

...
\end{rustcode}

\subsubsection*{序号813}
\noindent 路径：\code{os/components/wateros-vfs/vfs-impl/impl-fd-session/src/handles.rs}\\
\noindent 行范围：157-163\\[0.4em]

\begin{rustcode}
...
    fn poll_revents(&mut self, events: i16) -> VfsResult<i16> {
// 本变量代码由AI完成
        const POLLIN: i16 = 0x001;
// 本变量代码由AI完成
        const POLLOUT: i16 = 0x004;
...
\end{rustcode}

\subsubsection*{序号814}
\noindent 路径：\code{os/components/wateros-vfs/vfs-impl/impl-fd-session/src/handles.rs}\\
\noindent 行范围：159\\[0.4em]

\begin{rustcode}
...
    fn poll_revents(&mut self, events: i16) -> VfsResult<i16> {
// 本变量代码由AI完成
        const POLLIN: i16 = 0x001;
// 本变量代码由AI完成
        const POLLOUT: i16 = 0x004;
...
\end{rustcode}

\subsubsection*{序号815}
\noindent 路径：\code{os/components/wateros-vfs/vfs-impl/impl-fd-session/src/handles.rs}\\
\noindent 行范围：161\\[0.4em]

\begin{rustcode}
...
        const POLLIN: i16 = 0x001;
// 本变量代码由AI完成
        const POLLOUT: i16 = 0x004;
        Ok(events & (POLLIN | POLLOUT))
    }
...
\end{rustcode}

\subsubsection*{序号816}
\noindent 路径：\code{os/components/wateros-vfs/vfs-impl/impl-fd-session/src/handles.rs}\\
\noindent 行范围：166-168\\[0.4em]

\begin{rustcode}
...
// 本方法代码由AI完成
    fn metadata(&self) -> VfsResult<VfsMetadata> {
        Ok(special_dev_meta(0o20666, 3, 1, 5))
    }

...
\end{rustcode}

\subsubsection*{序号817}
\noindent 路径：\code{os/components/wateros-vfs/vfs-impl/impl-fd-session/src/handles.rs}\\
\noindent 行范围：171-173\\[0.4em]

\begin{rustcode}
...
// 本方法代码由AI完成
    fn duplicate(&self) -> VfsResult<Box<dyn VfsIoHandle>> {
        Ok(Box::new(*self))
    }
}
...
\end{rustcode}

\subsubsection*{序号818}
\noindent 路径：\code{os/components/wateros-vfs/vfs-impl/impl-fd-session/src/handles.rs}\\
\noindent 行范围：179\\[0.4em]

\begin{rustcode}
...
#[derive(Debug, Clone, Copy, Default)]
// 本结构代码由AI完成
pub struct CpuDmaLatencyDeviceHandle;

impl VfsIoHandle for CpuDmaLatencyDeviceHandle {
...
\end{rustcode}

\subsubsection*{序号819}
\noindent 路径：\code{os/components/wateros-vfs/vfs-impl/impl-fd-session/src/handles.rs}\\
\noindent 行范围：183-185\\[0.4em]

\begin{rustcode}
...
// 本方法代码由AI完成
    fn read(&mut self, _buf: &mut [u8]) -> VfsResult<usize> {
        Ok(0)
    }

...
\end{rustcode}

\subsubsection*{序号820}
\noindent 路径：\code{os/components/wateros-vfs/vfs-impl/impl-fd-session/src/handles.rs}\\
\noindent 行范围：188-190\\[0.4em]

\begin{rustcode}
...
// 本方法代码由AI完成
    fn write(&mut self, buf: &[u8]) -> VfsResult<usize> {
        Ok(buf.len())
    }

...
\end{rustcode}

\subsubsection*{序号821}
\noindent 路径：\code{os/components/wateros-vfs/vfs-impl/impl-fd-session/src/handles.rs}\\
\noindent 行范围：193-199\\[0.4em]

\begin{rustcode}
...
    fn poll_revents(&mut self, events: i16) -> VfsResult<i16> {
// 本变量代码由AI完成
        const POLLIN: i16 = 0x001;
// 本变量代码由AI完成
        const POLLOUT: i16 = 0x004;
...
\end{rustcode}

\subsubsection*{序号822}
\noindent 路径：\code{os/components/wateros-vfs/vfs-impl/impl-fd-session/src/handles.rs}\\
\noindent 行范围：195\\[0.4em]

\begin{rustcode}
...
    fn poll_revents(&mut self, events: i16) -> VfsResult<i16> {
// 本变量代码由AI完成
        const POLLIN: i16 = 0x001;
// 本变量代码由AI完成
        const POLLOUT: i16 = 0x004;
...
\end{rustcode}

\subsubsection*{序号823}
\noindent 路径：\code{os/components/wateros-vfs/vfs-impl/impl-fd-session/src/handles.rs}\\
\noindent 行范围：197\\[0.4em]

\begin{rustcode}
...
        const POLLIN: i16 = 0x001;
// 本变量代码由AI完成
        const POLLOUT: i16 = 0x004;
        Ok(events & (POLLIN | POLLOUT))
    }
...
\end{rustcode}

\subsubsection*{序号824}
\noindent 路径：\code{os/components/wateros-vfs/vfs-impl/impl-fd-session/src/handles.rs}\\
\noindent 行范围：202-204\\[0.4em]

\begin{rustcode}
...
// 本方法代码由AI完成
    fn ioctl(&mut self, _request: usize, _arg: usize) -> VfsResult<isize> {
        Err(VfsError::Unsupported)
    }

...
\end{rustcode}

\subsubsection*{序号825}
\noindent 路径：\code{os/components/wateros-vfs/vfs-impl/impl-fd-session/src/handles.rs}\\
\noindent 行范围：207-209\\[0.4em]

\begin{rustcode}
...
// 本方法代码由AI完成
    fn metadata(&self) -> VfsResult<VfsMetadata> {
        Ok(special_dev_meta(0o20600, 5, 10, 233))
    }

...
\end{rustcode}

\subsubsection*{序号826}
\noindent 路径：\code{os/components/wateros-vfs/vfs-impl/impl-fd-session/src/handles.rs}\\
\noindent 行范围：212-214\\[0.4em]

\begin{rustcode}
...
// 本方法代码由AI完成
    fn duplicate(&self) -> VfsResult<Box<dyn VfsIoHandle>> {
        Ok(Box::new(*self))
    }
}
...
\end{rustcode}

\subsubsection*{序号827}
\noindent 路径：\code{os/components/wateros-vfs/vfs-impl/impl-fd-session/src/handles.rs}\\
\noindent 行范围：220\\[0.4em]

\begin{rustcode}
...
#[derive(Debug, Clone, Copy, Default)]
// 本结构代码由AI完成
pub struct UrandomDeviceHandle;

impl VfsIoHandle for UrandomDeviceHandle {
...
\end{rustcode}

\subsubsection*{序号828}
\noindent 路径：\code{os/components/wateros-vfs/vfs-impl/impl-fd-session/src/handles.rs}\\
\noindent 行范围：224-237\\[0.4em]

\begin{rustcode}
...
    fn read(&mut self, buf: &mut [u8]) -> VfsResult<usize> {
        let mut state = URANDOM_STATE.fetch_add(
            0x9e37_79b9_7f4a_7c15u64 ^ (buf.as_ptr() as u64).rotate_left(17),
            Ordering::Relaxed,
        );
...
\end{rustcode}

\subsubsection*{序号829}
\noindent 路径：\code{os/components/wateros-vfs/vfs-impl/impl-fd-session/src/handles.rs}\\
\noindent 行范围：240-247\\[0.4em]

\begin{rustcode}
...
    fn write(&mut self, buf: &[u8]) -> VfsResult<usize> {
        let mut mix = buf.len() as u64;
        for byte in buf.iter().take(32) {
            mix = mix.rotate_left(5) ^ (*byte as u64);
        }
...
\end{rustcode}

\subsubsection*{序号830}
\noindent 路径：\code{os/components/wateros-vfs/vfs-impl/impl-fd-session/src/handles.rs}\\
\noindent 行范围：250-256\\[0.4em]

\begin{rustcode}
...
    fn poll_revents(&mut self, events: i16) -> VfsResult<i16> {
// 本变量代码由AI完成
        const POLLIN: i16 = 0x001;
// 本变量代码由AI完成
        const POLLOUT: i16 = 0x004;
...
\end{rustcode}

\subsubsection*{序号831}
\noindent 路径：\code{os/components/wateros-vfs/vfs-impl/impl-fd-session/src/handles.rs}\\
\noindent 行范围：252\\[0.4em]

\begin{rustcode}
...
    fn poll_revents(&mut self, events: i16) -> VfsResult<i16> {
// 本变量代码由AI完成
        const POLLIN: i16 = 0x001;
// 本变量代码由AI完成
        const POLLOUT: i16 = 0x004;
...
\end{rustcode}

\subsubsection*{序号832}
\noindent 路径：\code{os/components/wateros-vfs/vfs-impl/impl-fd-session/src/handles.rs}\\
\noindent 行范围：254\\[0.4em]

\begin{rustcode}
...
        const POLLIN: i16 = 0x001;
// 本变量代码由AI完成
        const POLLOUT: i16 = 0x004;
        Ok(events & (POLLIN | POLLOUT))
    }
...
\end{rustcode}

\subsubsection*{序号833}
\noindent 路径：\code{os/components/wateros-vfs/vfs-impl/impl-fd-session/src/handles.rs}\\
\noindent 行范围：259-261\\[0.4em]

\begin{rustcode}
...
// 本方法代码由AI完成
    fn metadata(&self) -> VfsResult<VfsMetadata> {
        Ok(special_dev_meta(0o20666, 4, 1, 9))
    }

...
\end{rustcode}

\subsubsection*{序号834}
\noindent 路径：\code{os/components/wateros-vfs/vfs-impl/impl-fd-session/src/handles.rs}\\
\noindent 行范围：264-266\\[0.4em]

\begin{rustcode}
...
// 本方法代码由AI完成
    fn duplicate(&self) -> VfsResult<Box<dyn VfsIoHandle>> {
        Ok(Box::new(*self))
    }
}
...
\end{rustcode}

\subsubsection*{序号835}
\noindent 路径：\code{os/components/wateros-vfs/vfs-impl/impl-fd-session/src/handles.rs}\\
\noindent 行范围：271-275\\[0.4em]

\begin{rustcode}
...
pub struct PipeReadHandle {
    endpoint: PipeEndpoint,
    /// 合成 inode 号，供 `flock` / `stat` 区分 pipe 实例。
    inode: u64,
}
...
\end{rustcode}

\subsubsection*{序号836}
\noindent 路径：\code{os/components/wateros-vfs/vfs-impl/impl-fd-session/src/handles.rs}\\
\noindent 行范围：279-282\\[0.4em]

\begin{rustcode}
...
// 本结构代码由AI完成
pub struct PipeWriteHandle {
    endpoint: PipeEndpoint,
    inode: u64,
}
...
\end{rustcode}

\subsubsection*{序号837}
\noindent 路径：\code{os/components/wateros-vfs/vfs-impl/impl-fd-session/src/handles.rs}\\
\noindent 行范围：285-292\\[0.4em]

\begin{rustcode}
...
pub fn pipe_handle_pair(nonblocking: bool) -> (PipeReadHandle, PipeWriteHandle) {
    let (read, write) = PipeEndpoint::pair(nonblocking);
    let inode = NEXT_PIPE_INODE.fetch_add(1, Ordering::Relaxed);
    (
        PipeReadHandle { endpoint: read, inode },
...
\end{rustcode}

\subsubsection*{序号838}
\noindent 路径：\code{os/components/wateros-vfs/vfs-impl/impl-fd-session/src/handles.rs}\\
\noindent 行范围：296-298\\[0.4em]

\begin{rustcode}
...
// 本方法代码由AI完成
    fn read(&mut self, buf: &mut [u8]) -> VfsResult<usize> {
        self.endpoint.read(buf).map_err(map_pipe_err)
    }

...
\end{rustcode}

\subsubsection*{序号839}
\noindent 路径：\code{os/components/wateros-vfs/vfs-impl/impl-fd-session/src/handles.rs}\\
\noindent 行范围：301-303\\[0.4em]

\begin{rustcode}
...
// 本方法代码由AI完成
    fn poll_revents(&mut self, events: i16) -> VfsResult<i16> {
        self.endpoint.poll_revents(events).map_err(map_pipe_err)
    }

...
\end{rustcode}

\subsubsection*{序号840}
\noindent 路径：\code{os/components/wateros-vfs/vfs-impl/impl-fd-session/src/handles.rs}\\
\noindent 行范围：306-315\\[0.4em]

\begin{rustcode}
...
    fn poll_wait_for_ticks(
        &mut self,
        events: i16,
        timeout_ticks: u64,
        still_waiting: &mut dyn FnMut() -> bool,
...
\end{rustcode}

\subsubsection*{序号841}
\noindent 路径：\code{os/components/wateros-vfs/vfs-impl/impl-fd-session/src/handles.rs}\\
\noindent 行范围：318-321\\[0.4em]

\begin{rustcode}
...
// 本方法代码由AI完成
    fn close(&mut self) -> VfsResult<()> {
        self.endpoint.close();
        Ok(())
    }
...
\end{rustcode}

\subsubsection*{序号842}
\noindent 路径：\code{os/components/wateros-vfs/vfs-impl/impl-fd-session/src/handles.rs}\\
\noindent 行范围：324-326\\[0.4em]

\begin{rustcode}
...
// 本方法代码由AI完成
    fn metadata(&self) -> VfsResult<VfsMetadata> {
        Ok(special_meta(0o10600, self.inode))
    }

...
\end{rustcode}

\subsubsection*{序号843}
\noindent 路径：\code{os/components/wateros-vfs/vfs-impl/impl-fd-session/src/handles.rs}\\
\noindent 行范围：329-334\\[0.4em]

\begin{rustcode}
...
    fn duplicate(&self) -> VfsResult<Box<dyn VfsIoHandle>> {
        Ok(Box::new(Self {
            endpoint: self.endpoint.clone(),
            inode: self.inode,
        }))
...
\end{rustcode}

\subsubsection*{序号844}
\noindent 路径：\code{os/components/wateros-vfs/vfs-impl/impl-fd-session/src/handles.rs}\\
\noindent 行范围：337-343\\[0.4em]

\begin{rustcode}
...
    fn open_status_flags(&self) -> u32 {
        if self.endpoint.nonblocking() {
            0o0004000
        } else {
            0
...
\end{rustcode}

\subsubsection*{序号845}
\noindent 路径：\code{os/components/wateros-vfs/vfs-impl/impl-fd-session/src/handles.rs}\\
\noindent 行范围：346-349\\[0.4em]

\begin{rustcode}
...
// 本方法代码由AI完成
    fn set_open_status_flags(&mut self, flags: u32) -> VfsResult<()> {
        self.endpoint.set_nonblocking(flags & 0o0004000 != 0);
        Ok(())
    }
...
\end{rustcode}

\subsubsection*{序号846}
\noindent 路径：\code{os/components/wateros-vfs/vfs-impl/impl-fd-session/src/handles.rs}\\
\noindent 行范围：352-354\\[0.4em]

\begin{rustcode}
...
// 本方法代码由AI完成
    fn pipe_capacity(&self) -> Option<usize> {
        Some(self.endpoint.pipe_capacity())
    }

...
\end{rustcode}

\subsubsection*{序号847}
\noindent 路径：\code{os/components/wateros-vfs/vfs-impl/impl-fd-session/src/handles.rs}\\
\noindent 行范围：357-359\\[0.4em]

\begin{rustcode}
...
// 本方法代码由AI完成
    fn pipe_set_capacity(&mut self, capacity: usize) -> VfsResult<usize> {
        self.endpoint.set_pipe_capacity(capacity).map_err(map_pipe_err)
    }
}
...
\end{rustcode}

\subsubsection*{序号848}
\noindent 路径：\code{os/components/wateros-vfs/vfs-impl/impl-fd-session/src/handles.rs}\\
\noindent 行范围：365-367\\[0.4em]

\begin{rustcode}
...
// 本方法代码由AI完成
    fn write(&mut self, buf: &[u8]) -> VfsResult<usize> {
        self.endpoint.write(buf).map_err(map_pipe_err)
    }

...
\end{rustcode}

\subsubsection*{序号849}
\noindent 路径：\code{os/components/wateros-vfs/vfs-impl/impl-fd-session/src/handles.rs}\\
\noindent 行范围：370-372\\[0.4em]

\begin{rustcode}
...
// 本方法代码由AI完成
    fn poll_revents(&mut self, events: i16) -> VfsResult<i16> {
        self.endpoint.poll_revents(events).map_err(map_pipe_err)
    }

...
\end{rustcode}

\subsubsection*{序号850}
\noindent 路径：\code{os/components/wateros-vfs/vfs-impl/impl-fd-session/src/handles.rs}\\
\noindent 行范围：375-384\\[0.4em]

\begin{rustcode}
...
    fn poll_wait_for_ticks(
        &mut self,
        events: i16,
        timeout_ticks: u64,
        still_waiting: &mut dyn FnMut() -> bool,
...
\end{rustcode}

\subsubsection*{序号851}
\noindent 路径：\code{os/components/wateros-vfs/vfs-impl/impl-fd-session/src/handles.rs}\\
\noindent 行范围：387-390\\[0.4em]

\begin{rustcode}
...
// 本方法代码由AI完成
    fn close(&mut self) -> VfsResult<()> {
        self.endpoint.close();
        Ok(())
    }
...
\end{rustcode}

\subsubsection*{序号852}
\noindent 路径：\code{os/components/wateros-vfs/vfs-impl/impl-fd-session/src/handles.rs}\\
\noindent 行范围：393-395\\[0.4em]

\begin{rustcode}
...
// 本方法代码由AI完成
    fn metadata(&self) -> VfsResult<VfsMetadata> {
        Ok(special_meta(0o10600, self.inode))
    }

...
\end{rustcode}

\subsubsection*{序号853}
\noindent 路径：\code{os/components/wateros-vfs/vfs-impl/impl-fd-session/src/handles.rs}\\
\noindent 行范围：398-403\\[0.4em]

\begin{rustcode}
...
    fn duplicate(&self) -> VfsResult<Box<dyn VfsIoHandle>> {
        Ok(Box::new(Self {
            endpoint: self.endpoint.clone(),
            inode: self.inode,
        }))
...
\end{rustcode}

\subsubsection*{序号854}
\noindent 路径：\code{os/components/wateros-vfs/vfs-impl/impl-fd-session/src/handles.rs}\\
\noindent 行范围：406-412\\[0.4em]

\begin{rustcode}
...
    fn open_status_flags(&self) -> u32 {
        if self.endpoint.nonblocking() {
            0o0004000
        } else {
            0
...
\end{rustcode}

\subsubsection*{序号855}
\noindent 路径：\code{os/components/wateros-vfs/vfs-impl/impl-fd-session/src/handles.rs}\\
\noindent 行范围：415-418\\[0.4em]

\begin{rustcode}
...
// 本方法代码由AI完成
    fn set_open_status_flags(&mut self, flags: u32) -> VfsResult<()> {
        self.endpoint.set_nonblocking(flags & 0o0004000 != 0);
        Ok(())
    }
...
\end{rustcode}

\subsubsection*{序号856}
\noindent 路径：\code{os/components/wateros-vfs/vfs-impl/impl-fd-session/src/handles.rs}\\
\noindent 行范围：421-423\\[0.4em]

\begin{rustcode}
...
// 本方法代码由AI完成
    fn pipe_capacity(&self) -> Option<usize> {
        Some(self.endpoint.pipe_capacity())
    }

...
\end{rustcode}

\subsubsection*{序号857}
\noindent 路径：\code{os/components/wateros-vfs/vfs-impl/impl-fd-session/src/handles.rs}\\
\noindent 行范围：426-428\\[0.4em]

\begin{rustcode}
...
// 本方法代码由AI完成
    fn pipe_set_capacity(&mut self, capacity: usize) -> VfsResult<usize> {
        self.endpoint.set_pipe_capacity(capacity).map_err(map_pipe_err)
    }
}
...
\end{rustcode}

\subsubsection*{序号858}
\noindent 路径：\code{os/components/wateros-vfs/vfs-impl/impl-fd-session/src/handles.rs}\\
\noindent 行范围：433-437\\[0.4em]

\begin{rustcode}
...
pub struct UnixStreamPairEnd {
    read_end: PipeEndpoint,
    write_end: PipeEndpoint,
    inode: u64,
}
...
\end{rustcode}

\subsubsection*{序号859}
\noindent 路径：\code{os/components/wateros-vfs/vfs-impl/impl-fd-session/src/handles.rs}\\
\noindent 行范围：440-456\\[0.4em]

\begin{rustcode}
...
pub fn stream_pair_handle_pair(nonblocking: bool) -> (UnixStreamPairEnd, UnixStreamPairEnd) {
    let (read_ab, write_ab) = PipeEndpoint::pair(nonblocking);
    let (read_ba, write_ba) = PipeEndpoint::pair(nonblocking);
    let inode = NEXT_STREAM_PAIR_INODE.fetch_add(1, Ordering::Relaxed);
    (
...
\end{rustcode}

\subsubsection*{序号860}
\noindent 路径：\code{os/components/wateros-vfs/vfs-impl/impl-fd-session/src/handles.rs}\\
\noindent 行范围：460-462\\[0.4em]

\begin{rustcode}
...
// 本方法代码由AI完成
    fn read(&mut self, buf: &mut [u8]) -> VfsResult<usize> {
        self.read_end.read(buf).map_err(map_pipe_err)
    }

...
\end{rustcode}

\subsubsection*{序号861}
\noindent 路径：\code{os/components/wateros-vfs/vfs-impl/impl-fd-session/src/handles.rs}\\
\noindent 行范围：465-467\\[0.4em]

\begin{rustcode}
...
// 本方法代码由AI完成
    fn write(&mut self, buf: &[u8]) -> VfsResult<usize> {
        self.write_end.write(buf).map_err(map_pipe_err)
    }

...
\end{rustcode}

\subsubsection*{序号862}
\noindent 路径：\code{os/components/wateros-vfs/vfs-impl/impl-fd-session/src/handles.rs}\\
\noindent 行范围：470-483\\[0.4em]

\begin{rustcode}
...
    fn poll_revents(&mut self, events: i16) -> VfsResult<i16> {
// 本变量代码由AI完成
        const POLLIN: i16 = 0x001;
// 本变量代码由AI完成
        const POLLOUT: i16 = 0x004;
...
\end{rustcode}

\subsubsection*{序号863}
\noindent 路径：\code{os/components/wateros-vfs/vfs-impl/impl-fd-session/src/handles.rs}\\
\noindent 行范围：472\\[0.4em]

\begin{rustcode}
...
    fn poll_revents(&mut self, events: i16) -> VfsResult<i16> {
// 本变量代码由AI完成
        const POLLIN: i16 = 0x001;
// 本变量代码由AI完成
        const POLLOUT: i16 = 0x004;
...
\end{rustcode}

\subsubsection*{序号864}
\noindent 路径：\code{os/components/wateros-vfs/vfs-impl/impl-fd-session/src/handles.rs}\\
\noindent 行范围：474\\[0.4em]

\begin{rustcode}
...
        const POLLIN: i16 = 0x001;
// 本变量代码由AI完成
        const POLLOUT: i16 = 0x004;
        let mut revents = 0i16;
        if events & POLLIN != 0 {
...
\end{rustcode}

\subsubsection*{序号865}
\noindent 路径：\code{os/components/wateros-vfs/vfs-impl/impl-fd-session/src/handles.rs}\\
\noindent 行范围：486-507\\[0.4em]

\begin{rustcode}
...
    fn poll_wait_for_ticks(
        &mut self,
        events: i16,
        timeout_ticks: u64,
        still_waiting: &mut dyn FnMut() -> bool,
...
\end{rustcode}

\subsubsection*{序号866}
\noindent 路径：\code{os/components/wateros-vfs/vfs-impl/impl-fd-session/src/handles.rs}\\
\noindent 行范围：493\\[0.4em]

\begin{rustcode}
...
    ) -> VfsResult<()> {
// 本变量代码由AI完成
        const POLLIN: i16 = 0x001;
// 本变量代码由AI完成
        const POLLOUT: i16 = 0x004;
...
\end{rustcode}

\subsubsection*{序号867}
\noindent 路径：\code{os/components/wateros-vfs/vfs-impl/impl-fd-session/src/handles.rs}\\
\noindent 行范围：495\\[0.4em]

\begin{rustcode}
...
        const POLLIN: i16 = 0x001;
// 本变量代码由AI完成
        const POLLOUT: i16 = 0x004;
        if events & POLLIN != 0 {
            self.read_end
...
\end{rustcode}

\subsubsection*{序号868}
\noindent 路径：\code{os/components/wateros-vfs/vfs-impl/impl-fd-session/src/handles.rs}\\
\noindent 行范围：510-514\\[0.4em]

\begin{rustcode}
...
    fn close(&mut self) -> VfsResult<()> {
        self.read_end.close();
        self.write_end.close();
        Ok(())
    }
...
\end{rustcode}

\subsubsection*{序号869}
\noindent 路径：\code{os/components/wateros-vfs/vfs-impl/impl-fd-session/src/handles.rs}\\
\noindent 行范围：517-519\\[0.4em]

\begin{rustcode}
...
// 本方法代码由AI完成
    fn metadata(&self) -> VfsResult<VfsMetadata> {
        Ok(special_meta(0o140600, self.inode))
    }

...
\end{rustcode}

\subsubsection*{序号870}
\noindent 路径：\code{os/components/wateros-vfs/vfs-impl/impl-fd-session/src/handles.rs}\\
\noindent 行范围：522-528\\[0.4em]

\begin{rustcode}
...
    fn duplicate(&self) -> VfsResult<Box<dyn VfsIoHandle>> {
        Ok(Box::new(Self {
            read_end: self.read_end.clone(),
            write_end: self.write_end.clone(),
            inode: self.inode,
...
\end{rustcode}

\subsubsection*{序号871}
\noindent 路径：\code{os/components/wateros-vfs/vfs-impl/impl-fd-session/src/handles.rs}\\
\noindent 行范围：531-537\\[0.4em]

\begin{rustcode}
...
    fn open_status_flags(&self) -> u32 {
        if self.read_end.nonblocking() {
            0o0004000
        } else {
            0
...
\end{rustcode}

\subsubsection*{序号872}
\noindent 路径：\code{os/components/wateros-vfs/vfs-impl/impl-fd-session/src/handles.rs}\\
\noindent 行范围：540-545\\[0.4em]

\begin{rustcode}
...
    fn set_open_status_flags(&mut self, flags: u32) -> VfsResult<()> {
        let nonblocking = flags & 0o0004000 != 0;
        self.read_end.set_nonblocking(nonblocking);
        self.write_end.set_nonblocking(nonblocking);
        Ok(())
...
\end{rustcode}

\subsubsection*{序号873}
\noindent 路径：\code{os/components/wateros-vfs/vfs-impl/impl-fd-session/src/handles.rs}\\
\noindent 行范围：550-571\\[0.4em]

\begin{rustcode}
...
pub fn stream_pair_smoke() -> bool {
// 本变量代码由AI完成
    const POLLIN: i16 = 0x001;
    let (mut a, mut b) = stream_pair_handle_pair(false);
    if a.write(b"ab").is_err() {
...
\end{rustcode}

\subsubsection*{序号874}
\noindent 路径：\code{os/components/wateros-vfs/vfs-impl/impl-fd-session/src/handles.rs}\\
\noindent 行范围：552\\[0.4em]

\begin{rustcode}
...
pub fn stream_pair_smoke() -> bool {
// 本变量代码由AI完成
    const POLLIN: i16 = 0x001;
    let (mut a, mut b) = stream_pair_handle_pair(false);
    if a.write(b"ab").is_err() {
...
\end{rustcode}

\subsubsection*{序号875}
\noindent 路径：\code{os/components/wateros-vfs/vfs-impl/impl-fd-session/src/handles.rs}\\
\noindent 行范围：575-586\\[0.4em]

\begin{rustcode}
...
pub fn poll_pipe_smoke() -> bool {
// 本变量代码由AI完成
    const POLLIN: i16 = 0x001;
    let (mut read, mut write) = pipe_handle_pair(false);
    if read.poll_revents(POLLIN).ok() != Some(0) {
...
\end{rustcode}

\subsubsection*{序号876}
\noindent 路径：\code{os/components/wateros-vfs/vfs-impl/impl-fd-session/src/handles.rs}\\
\noindent 行范围：577\\[0.4em]

\begin{rustcode}
...
pub fn poll_pipe_smoke() -> bool {
// 本变量代码由AI完成
    const POLLIN: i16 = 0x001;
    let (mut read, mut write) = pipe_handle_pair(false);
    if read.poll_revents(POLLIN).ok() != Some(0) {
...
\end{rustcode}

\subsubsection*{序号877}
\noindent 路径：\code{os/components/wateros-vfs/vfs-impl/impl-fd-session/src/handles.rs}\\
\noindent 行范围：589-596\\[0.4em]

\begin{rustcode}
...
fn map_pipe_err(err: PipeError) -> VfsError {
    match err {
        PipeError::WouldBlock => VfsError::WouldBlock,
        PipeError::Interrupted => VfsError::Interrupted,
        PipeError::BrokenPipe => VfsError::BrokenPipe,
...
\end{rustcode}

\subsubsection*{序号878}
\noindent 路径：\code{os/components/wateros-vfs/vfs-impl/impl-fd-session/src/lib.rs}\\
\noindent 行范围：1-25\\[0.4em]

\begin{rustcode}
//! per-task fd 表与控制台/pipe [`VfsIoHandle`] 实现。
//! 本模块代码由AI完成

#![no_std]

...
\end{rustcode}

\subsubsection*{序号879}
\noindent 路径：\code{os/components/wateros-vfs/vfs-impl/impl-fd-session/src/registry.rs}\\
\noindent 行范围：1-781\\[0.4em]

\begin{rustcode}
//! 以 [`task::TaskId`] 为 key 的 per-task fd 表。
//! 本模块代码由AI完成

extern crate alloc;

...
\end{rustcode}

\subsubsection*{序号880}
\noindent 路径：\code{os/components/wateros-vfs/vfs-impl/impl-fd-session/src/registry.rs}\\
\noindent 行范围：31-40\\[0.4em]

\begin{rustcode}
...
pub struct PerTaskFdRegistry {
    tables: BTreeMap<task::TaskId, Vec<Option<Box<dyn VfsIoHandle>>>>,
    fd_flags: BTreeMap<task::TaskId, Vec<u8>>,
    owners: BTreeMap<task::TaskId, task::TaskId>,
    ref_counts: BTreeMap<task::TaskId, usize>,
...
\end{rustcode}

\subsubsection*{序号881}
\noindent 路径：\code{os/components/wateros-vfs/vfs-impl/impl-fd-session/src/registry.rs}\\
\noindent 行范围：56-71\\[0.4em]

\begin{rustcode}
...
    fn ensure_task(&mut self, task_id: task::TaskId) {
        self.owners.entry(task_id).or_insert(task_id);
        self.ref_counts.entry(task_id).or_insert(1);
        let owner = self.effective_owner(task_id);
        let table = self.tables.entry(owner).or_insert_with(Vec::new);
...
\end{rustcode}

\subsubsection*{序号882}
\noindent 路径：\code{os/components/wateros-vfs/vfs-impl/impl-fd-session/src/registry.rs}\\
\noindent 行范围：75-80\\[0.4em]

\begin{rustcode}
...
    fn effective_owner(&self, task_id: task::TaskId) -> task::TaskId {
        self.owners
            .get(&task_id)
            .copied()
            .unwrap_or(task_id)
...
\end{rustcode}

\subsubsection*{序号883}
\noindent 路径：\code{os/components/wateros-vfs/vfs-impl/impl-fd-session/src/registry.rs}\\
\noindent 行范围：83-93\\[0.4em]

\begin{rustcode}
...
    fn ensure_shared_io_state(&mut self, owner: task::TaskId) {
        let len = self.tables.get(&owner).map(Vec::len).unwrap_or(0);
        let inflight = self.io_inflight.entry(owner).or_insert_with(Vec::new);
        if inflight.len() < len {
            inflight.resize(len, 0);
...
\end{rustcode}

\subsubsection*{序号884}
\noindent 路径：\code{os/components/wateros-vfs/vfs-impl/impl-fd-session/src/registry.rs}\\
\noindent 行范围：96-108\\[0.4em]

\begin{rustcode}
...
    fn ensure_fd_not_io_busy(&self, owner: task::TaskId, fd: usize) -> VfsResult<()> {
        if self
            .io_inflight
            .get(&owner)
            .and_then(|counts| counts.get(fd))
...
\end{rustcode}

\subsubsection*{序号885}
\noindent 路径：\code{os/components/wateros-vfs/vfs-impl/impl-fd-session/src/registry.rs}\\
\noindent 行范围：111-115\\[0.4em]

\begin{rustcode}
...
    fn table_mut(&mut self, task_id: task::TaskId) -> &mut Vec<Option<Box<dyn VfsIoHandle>>> {
        self.ensure_task(task_id);
        let owner = self.effective_owner(task_id);
        self.tables.get_mut(&owner).expect("fd table owner")
    }
...
\end{rustcode}

\subsubsection*{序号886}
\noindent 路径：\code{os/components/wateros-vfs/vfs-impl/impl-fd-session/src/registry.rs}\\
\noindent 行范围：118-125\\[0.4em]

\begin{rustcode}
...
    fn ensure_flags_len(&mut self, task_id: task::TaskId, len: usize) {
        self.ensure_task(task_id);
        let owner = self.effective_owner(task_id);
        let flags = self.fd_flags.entry(owner).or_insert_with(Vec::new);
        if flags.len() < len {
...
\end{rustcode}

\subsubsection*{序号887}
\noindent 路径：\code{os/components/wateros-vfs/vfs-impl/impl-fd-session/src/registry.rs}\\
\noindent 行范围：128-140\\[0.4em]

\begin{rustcode}
...
    fn close_slot(&mut self, task_id: task::TaskId, fd: usize) -> VfsResult<()> {
        let pid = task::process_task_snapshot(task_id).map(|snap| snap.pid);
        let mut handle = self.take_fd_for_close(task_id, fd)?;
        if let Some(pid) = pid {
            if let Ok(meta) = handle.metadata() {
...
\end{rustcode}

\subsubsection*{序号888}
\noindent 路径：\code{os/components/wateros-vfs/vfs-impl/impl-fd-session/src/registry.rs}\\
\noindent 行范围：143-154\\[0.4em]

\begin{rustcode}
...
    fn take_table_handles(&mut self, owner: task::TaskId) -> Vec<Box<dyn VfsIoHandle>> {
        let mut handles = Vec::new();
        if let Some(mut table) = self.tables.remove(&owner) {
            for slot in table.iter_mut() {
                if let Some(handle) = slot.take() {
...
\end{rustcode}

\subsubsection*{序号889}
\noindent 路径：\code{os/components/wateros-vfs/vfs-impl/impl-fd-session/src/registry.rs}\\
\noindent 行范围：157-177\\[0.4em]

\begin{rustcode}
...
    pub fn take_fd_for_close(
        &mut self,
        task_id: task::TaskId,
        fd: usize,
    ) -> VfsResult<Box<dyn VfsIoHandle>> {
...
\end{rustcode}

\subsubsection*{序号890}
\noindent 路径：\code{os/components/wateros-vfs/vfs-impl/impl-fd-session/src/registry.rs}\\
\noindent 行范围：180-219\\[0.4em]

\begin{rustcode}
...
    pub fn take_fd_range_for_close(
        &mut self,
        task_id: task::TaskId,
        first: usize,
        last: usize,
...
\end{rustcode}

\subsubsection*{序号891}
\noindent 路径：\code{os/components/wateros-vfs/vfs-impl/impl-fd-session/src/registry.rs}\\
\noindent 行范围：222-240\\[0.4em]

\begin{rustcode}
...
    pub fn take_cloexec_fds_for_task(
        &mut self,
        task_id: task::TaskId,
    ) -> Vec<Box<dyn VfsIoHandle>> {
        self.ensure_task(task_id);
...
\end{rustcode}

\subsubsection*{序号892}
\noindent 路径：\code{os/components/wateros-vfs/vfs-impl/impl-fd-session/src/registry.rs}\\
\noindent 行范围：243-256\\[0.4em]

\begin{rustcode}
...
    pub fn drain_task_fd_table(&mut self, task_id: task::TaskId) -> Vec<Box<dyn VfsIoHandle>> {
        let Some(owner) = self.release_owner(task_id) else {
            return Vec::new();
        };
        let mut handles = if self.ref_counts.get(&owner).copied().unwrap_or(0) == 0 {
...
\end{rustcode}

\subsubsection*{序号893}
\noindent 路径：\code{os/components/wateros-vfs/vfs-impl/impl-fd-session/src/registry.rs}\\
\noindent 行范围：259-268\\[0.4em]

\begin{rustcode}
...
    fn release_owner(&mut self, task_id: task::TaskId) -> Option<task::TaskId> {
        let owner = self.owners.remove(&task_id)?;
        if let Some(count) = self.ref_counts.get_mut(&owner) {
            *count = count.saturating_sub(1);
            if *count == 0 {
...
\end{rustcode}

\subsubsection*{序号894}
\noindent 路径：\code{os/components/wateros-vfs/vfs-impl/impl-fd-session/src/registry.rs}\\
\noindent 行范围：271-276\\[0.4em]

\begin{rustcode}
...
    fn close_table(&mut self, owner: task::TaskId) {
        let handles = self.take_table_handles(owner);
        for mut handle in handles {
            let _ = handle.close();
        }
...
\end{rustcode}

\subsubsection*{序号895}
\noindent 路径：\code{os/components/wateros-vfs/vfs-impl/impl-fd-session/src/registry.rs}\\
\noindent 行范围：279-285\\[0.4em]

\begin{rustcode}
...
    fn open_fd_count_for_task(&self, task_id: task::TaskId) -> usize {
        let owner = self.effective_owner(task_id);
        self.tables
            .get(&owner)
            .map(|table| table.iter().filter(|slot| slot.is_some()).count())
...
\end{rustcode}

\subsubsection*{序号896}
\noindent 路径：\code{os/components/wateros-vfs/vfs-impl/impl-fd-session/src/registry.rs}\\
\noindent 行范围：288-294\\[0.4em]

\begin{rustcode}
...
    fn check_nofile_before_open(&self, task_id: task::TaskId) -> VfsResult<()> {
        let limit = task::nofile_rlimit_for_task(task_id);
        if self.open_fd_count_for_task(task_id) >= limit as usize {
            return Err(VfsError::TooManyOpenFiles);
        }
...
\end{rustcode}

\subsubsection*{序号897}
\noindent 路径：\code{os/components/wateros-vfs/vfs-impl/impl-fd-session/src/registry.rs}\\
\noindent 行范围：299-305\\[0.4em]

\begin{rustcode}
...
    fn get_io(&mut self, fd: usize) -> VfsResult<&mut (dyn VfsIoHandle + '_)> {
        let task_id = task::current_task_id().ok_or(VfsError::NoTask)?;
        match self.table_mut(task_id).get_mut(fd) {
            Some(Some(h)) => Ok(h.as_mut()),
            _ => Err(VfsError::BadFd),
...
\end{rustcode}

\subsubsection*{序号898}
\noindent 路径：\code{os/components/wateros-vfs/vfs-impl/impl-fd-session/src/registry.rs}\\
\noindent 行范围：308-326\\[0.4em]

\begin{rustcode}
...
    fn alloc_fd(&mut self, handle: Box<dyn VfsIoHandle>) -> VfsResult<usize> {
        let task_id = task::current_task_id().ok_or(VfsError::NoTask)?;
        self.check_nofile_before_open(task_id)?;
        let newfd = {
            let table = self.table_mut(task_id);
...
\end{rustcode}

\subsubsection*{序号899}
\noindent 路径：\code{os/components/wateros-vfs/vfs-impl/impl-fd-session/src/registry.rs}\\
\noindent 行范围：329-332\\[0.4em]

\begin{rustcode}
...
// 本方法代码由AI完成
    fn close_fd(&mut self, fd: usize) -> VfsResult<()> {
        let task_id = task::current_task_id().ok_or(VfsError::NoTask)?;
        self.close_fd_for_task(task_id, fd)
    }
...
\end{rustcode}

\subsubsection*{序号900}
\noindent 路径：\code{os/components/wateros-vfs/vfs-impl/impl-fd-session/src/registry.rs}\\
\noindent 行范围：338-358\\[0.4em]

\begin{rustcode}
...
    pub fn alloc_fd_for_task(
        &mut self,
        task_id: task::TaskId,
        handle: Box<dyn VfsIoHandle>,
    ) -> VfsResult<usize> {
...
\end{rustcode}

\subsubsection*{序号901}
\noindent 路径：\code{os/components/wateros-vfs/vfs-impl/impl-fd-session/src/registry.rs}\\
\noindent 行范围：362-373\\[0.4em]

\begin{rustcode}
...
    pub fn get_io_for_task(
        &mut self,
        task_id: task::TaskId,
        fd: usize,
    ) -> VfsResult<&mut (dyn VfsIoHandle + '_)> {
...
\end{rustcode}

\subsubsection*{序号902}
\noindent 路径：\code{os/components/wateros-vfs/vfs-impl/impl-fd-session/src/registry.rs}\\
\noindent 行范围：377-387\\[0.4em]

\begin{rustcode}
...
    pub fn flush_all(&mut self) -> VfsResult<()> {
        let mut first_error = None;
        for table in self.tables.values_mut() {
            for handle in table.iter_mut().flatten() {
                if let Err(err) = handle.flush() {
...
\end{rustcode}

\subsubsection*{序号903}
\noindent 路径：\code{os/components/wateros-vfs/vfs-impl/impl-fd-session/src/registry.rs}\\
\noindent 行范围：391-394\\[0.4em]

\begin{rustcode}
...
// 本方法代码由AI完成
    pub fn is_fd_table_shared(&self, task_id: task::TaskId) -> bool {
        let owner = self.effective_owner(task_id);
        self.ref_counts.get(&owner).copied().unwrap_or(0) > 1
    }
...
\end{rustcode}

\subsubsection*{序号904}
\noindent 路径：\code{os/components/wateros-vfs/vfs-impl/impl-fd-session/src/registry.rs}\\
\noindent 行范围：398-424\\[0.4em]

\begin{rustcode}
...
    pub fn begin_shared_io_for_task(
        &mut self,
        task_id: task::TaskId,
        fd: usize,
    ) -> VfsResult<(*mut dyn VfsIoHandle, Arc<Mutex<()>>)> {
...
\end{rustcode}

\subsubsection*{序号905}
\noindent 路径：\code{os/components/wateros-vfs/vfs-impl/impl-fd-session/src/registry.rs}\\
\noindent 行范围：428-435\\[0.4em]

\begin{rustcode}
...
    pub fn end_shared_io_for_task(&mut self, task_id: task::TaskId, fd: usize) {
        let owner = self.effective_owner(task_id);
        if let Some(inflight) = self.io_inflight.get_mut(&owner) {
            if fd < inflight.len() && inflight[fd] > 0 {
                inflight[fd] -= 1;
...
\end{rustcode}

\subsubsection*{序号906}
\noindent 路径：\code{os/components/wateros-vfs/vfs-impl/impl-fd-session/src/registry.rs}\\
\noindent 行范围：439-450\\[0.4em]

\begin{rustcode}
...
    pub fn take_io_for_task(
        &mut self,
        task_id: task::TaskId,
        fd: usize,
    ) -> VfsResult<Box<dyn VfsIoHandle>> {
...
\end{rustcode}

\subsubsection*{序号907}
\noindent 路径：\code{os/components/wateros-vfs/vfs-impl/impl-fd-session/src/registry.rs}\\
\noindent 行范围：454-471\\[0.4em]

\begin{rustcode}
...
    pub fn restore_io_for_task(
        &mut self,
        task_id: task::TaskId,
        fd: usize,
        handle: Box<dyn VfsIoHandle>,
...
\end{rustcode}

\subsubsection*{序号908}
\noindent 路径：\code{os/components/wateros-vfs/vfs-impl/impl-fd-session/src/registry.rs}\\
\noindent 行范围：475-477\\[0.4em]

\begin{rustcode}
...
// 本方法代码由AI完成
    pub fn close_fd_for_task(&mut self, task_id: task::TaskId, fd: usize) -> VfsResult<()> {
        self.close_slot(task_id, fd)
    }

...
\end{rustcode}

\subsubsection*{序号909}
\noindent 路径：\code{os/components/wateros-vfs/vfs-impl/impl-fd-session/src/registry.rs}\\
\noindent 行范围：481-512\\[0.4em]

\begin{rustcode}
...
    pub fn dup_fd_for_task(
        &mut self,
        task_id: task::TaskId,
        oldfd: usize,
        minfd: usize,
...
\end{rustcode}

\subsubsection*{序号910}
\noindent 路径：\code{os/components/wateros-vfs/vfs-impl/impl-fd-session/src/registry.rs}\\
\noindent 行范围：516-566\\[0.4em]

\begin{rustcode}
...
    pub fn dup3_fd_for_task(
        &mut self,
        task_id: task::TaskId,
        oldfd: usize,
        newfd: usize,
...
\end{rustcode}

\subsubsection*{序号911}
\noindent 路径：\code{os/components/wateros-vfs/vfs-impl/impl-fd-session/src/registry.rs}\\
\noindent 行范围：570-579\\[0.4em]

\begin{rustcode}
...
    pub fn get_fd_flags(&mut self, task_id: task::TaskId, fd: usize) -> VfsResult<usize> {
        self.get_io_for_task(task_id, fd)?;
        self.ensure_task(task_id);
        let owner = self.effective_owner(task_id);
        let Some(flags) = self.fd_flags.get(&owner) else {
...
\end{rustcode}

\subsubsection*{序号912}
\noindent 路径：\code{os/components/wateros-vfs/vfs-impl/impl-fd-session/src/registry.rs}\\
\noindent 行范围：583-587\\[0.4em]

\begin{rustcode}
...
    pub fn set_fd_flags(&mut self, task_id: task::TaskId, fd: usize, val: usize) -> VfsResult<()> {
        self.get_io_for_task(task_id, fd)?;
        let cloexec = (val & usize::from(FD_CLOEXEC)) != 0;
        self.set_fd_cloexec(task_id, fd, cloexec)
    }
...
\end{rustcode}

\subsubsection*{序号913}
\noindent 路径：\code{os/components/wateros-vfs/vfs-impl/impl-fd-session/src/registry.rs}\\
\noindent 行范围：591-597\\[0.4em]

\begin{rustcode}
...
    pub fn set_fd_path_only(&mut self, task_id: task::TaskId, fd: usize) -> VfsResult<()> {
        self.get_io_for_task(task_id, fd)?;
        self.ensure_flags_len(task_id, fd + 1);
        let owner = self.effective_owner(task_id);
        self.fd_flags.get_mut(&owner).expect("fd flags owner")[fd] |= FD_PATH_ONLY;
...
\end{rustcode}

\subsubsection*{序号914}
\noindent 路径：\code{os/components/wateros-vfs/vfs-impl/impl-fd-session/src/registry.rs}\\
\noindent 行范围：601-609\\[0.4em]

\begin{rustcode}
...
    pub fn is_fd_path_only(&mut self, task_id: task::TaskId, fd: usize) -> VfsResult<bool> {
        self.get_io_for_task(task_id, fd)?;
        self.ensure_task(task_id);
        let owner = self.effective_owner(task_id);
        let Some(flags) = self.fd_flags.get(&owner) else {
...
\end{rustcode}

\subsubsection*{序号915}
\noindent 路径：\code{os/components/wateros-vfs/vfs-impl/impl-fd-session/src/registry.rs}\\
\noindent 行范围：612-622\\[0.4em]

\begin{rustcode}
...
    fn set_fd_cloexec(&mut self, task_id: task::TaskId, fd: usize, cloexec: bool) -> VfsResult<()> {
        self.ensure_flags_len(task_id, fd + 1);
        let owner = self.effective_owner(task_id);
        let slot = self.fd_flags.get_mut(&owner).expect("fd flags owner");
        if cloexec {
...
\end{rustcode}

\subsubsection*{序号916}
\noindent 路径：\code{os/components/wateros-vfs/vfs-impl/impl-fd-session/src/registry.rs}\\
\noindent 行范围：625-657\\[0.4em]

\begin{rustcode}
...
    pub fn set_fd_range_cloexec(
        &mut self,
        task_id: task::TaskId,
        first: usize,
        last: usize,
...
\end{rustcode}

\subsubsection*{序号917}
\noindent 路径：\code{os/components/wateros-vfs/vfs-impl/impl-fd-session/src/registry.rs}\\
\noindent 行范围：661-663\\[0.4em]

\begin{rustcode}
...
// 本方法代码由AI完成
    pub fn init_child_fd_table(&mut self, child: task::TaskId) {
        let _ = self.table_mut(child);
    }

...
\end{rustcode}

\subsubsection*{序号918}
\noindent 路径：\code{os/components/wateros-vfs/vfs-impl/impl-fd-session/src/registry.rs}\\
\noindent 行范围：667-711\\[0.4em]

\begin{rustcode}
...
    pub fn copy_fd_table_from_parent(&mut self, child: task::TaskId, parent: task::TaskId) {
        self.ensure_task(parent);
        if self.owners.contains_key(&child) {
            self.drop_task_fd_table(child);
        }
...
\end{rustcode}

\subsubsection*{序号919}
\noindent 路径：\code{os/components/wateros-vfs/vfs-impl/impl-fd-session/src/registry.rs}\\
\noindent 行范围：715-725\\[0.4em]

\begin{rustcode}
...
    pub fn share_fd_table_from_parent(&mut self, child: task::TaskId, parent: task::TaskId) {
        self.ensure_task(parent);
        if self.owners.contains_key(&child) {
            self.drop_task_fd_table(child);
        }
...
\end{rustcode}

\subsubsection*{序号920}
\noindent 路径：\code{os/components/wateros-vfs/vfs-impl/impl-fd-session/src/registry.rs}\\
\noindent 行范围：729-740\\[0.4em]

\begin{rustcode}
...
    pub fn close_cloexec_fds_for_task(&mut self, task_id: task::TaskId) {
        self.ensure_task(task_id);
        let owner = self.effective_owner(task_id);
        let table_len = self.tables.get(&owner).map(Vec::len).unwrap_or(0);
        let flags = self.fd_flags.get(&owner).cloned().unwrap_or_default();
...
\end{rustcode}

\subsubsection*{序号921}
\noindent 路径：\code{os/components/wateros-vfs/vfs-impl/impl-fd-session/src/registry.rs}\\
\noindent 行范围：744-755\\[0.4em]

\begin{rustcode}
...
    pub fn drop_task_fd_table(&mut self, task_id: task::TaskId) {
        let Some(owner) = self.release_owner(task_id) else {
            return;
        };
        if self.ref_counts.get(&owner).copied().unwrap_or(0) == 0 {
...
\end{rustcode}

\subsubsection*{序号922}
\noindent 路径：\code{os/components/wateros-vfs/vfs-impl/impl-fd-session/src/registry.rs}\\
\noindent 行范围：759-765\\[0.4em]

\begin{rustcode}
...
fn default_stdin_handle() -> Box<dyn VfsIoHandle> {
    if let Some(dev) = default_serial_device() {
        Box::new(CharDevHandle::new_stdin(dev))
    } else {
        Box::new(ConsoleInHandle)
...
\end{rustcode}

\subsubsection*{序号923}
\noindent 路径：\code{os/components/wateros-vfs/vfs-impl/impl-fd-session/src/registry.rs}\\
\noindent 行范围：768-774\\[0.4em]

\begin{rustcode}
...
fn default_stdout_handle() -> Box<dyn VfsIoHandle> {
    if let Some(dev) = default_serial_device() {
        Box::new(CharDevHandle::new_stdout(dev))
    } else {
        Box::new(ConsoleOutHandle)
...
\end{rustcode}

\subsubsection*{序号924}
\noindent 路径：\code{os/components/wateros-vfs/vfs-impl/impl-fd-session/src/registry.rs}\\
\noindent 行范围：777-781\\[0.4em]

\begin{rustcode}
...
fn default_serial_device() -> Option<SharedCharacterDevice> {
    (0..character_device_count())
        .find(|&idx| character_device_kind_at(idx) == Some(CharacterDeviceKind::Serial))
        .and_then(character_device_at)
}
\end{rustcode}

\subsubsection*{序号925}
\noindent 路径：\code{os/components/wateros-vfs/vfs-impl/impl-fs-bridge/src/dir_handle.rs}\\
\noindent 行范围：1-166\\[0.4em]

\begin{rustcode}
//! 已打开目录句柄：供 `openat(dirfd, …)` 与 `getdents64` 使用。
//! 本模块代码由AI完成

extern crate alloc;

...
\end{rustcode}

\subsubsection*{序号926}
\noindent 路径：\code{os/components/wateros-vfs/vfs-impl/impl-fs-bridge/src/dir_handle.rs}\\
\noindent 行范围：19-24\\[0.4em]

\begin{rustcode}
...
pub struct DirectoryHandle {
    path: String,
    meta: VfsMetadata,
    dirents: Option<Vec<VfsDirEntry>>,
    next_index: usize,
...
\end{rustcode}

\subsubsection*{序号927}
\noindent 路径：\code{os/components/wateros-vfs/vfs-impl/impl-fs-bridge/src/dir_handle.rs}\\
\noindent 行范围：28-42\\[0.4em]

\begin{rustcode}
...
    pub(crate) fn open(bridge: &FsBridge, path: String) -> VfsResult<Box<dyn VfsIoHandle>> {
        if !bridge.exists(path.as_str())? {
            return Err(VfsError::NotFound);
        }
        let meta = bridge.metadata(path.as_str())?;
...
\end{rustcode}

\subsubsection*{序号928}
\noindent 路径：\code{os/components/wateros-vfs/vfs-impl/impl-fs-bridge/src/dir_handle.rs}\\
\noindent 行范围：45-52\\[0.4em]

\begin{rustcode}
...
    fn load_dirents(&mut self, bridge: &FsBridge) -> VfsResult<()> {
        if self.dirents.is_some() {
            return Ok(());
        }
        self.dirents = Some(bridge.read_dir(self.path.as_str())?);
...
\end{rustcode}

\subsubsection*{序号929}
\noindent 路径：\code{os/components/wateros-vfs/vfs-impl/impl-fs-bridge/src/dir_handle.rs}\\
\noindent 行范围：56\\[0.4em]

\begin{rustcode}
...

// 本变量代码由AI完成
const DT_REG: u8 = 8;
// 本变量代码由AI完成
const DT_DIR: u8 = 4;
...
\end{rustcode}

\subsubsection*{序号930}
\noindent 路径：\code{os/components/wateros-vfs/vfs-impl/impl-fs-bridge/src/dir_handle.rs}\\
\noindent 行范围：58\\[0.4em]

\begin{rustcode}
...
const DT_REG: u8 = 8;
// 本变量代码由AI完成
const DT_DIR: u8 = 4;
// 本变量代码由AI完成
const DT_LNK: u8 = 10;
...
\end{rustcode}

\subsubsection*{序号931}
\noindent 路径：\code{os/components/wateros-vfs/vfs-impl/impl-fs-bridge/src/dir_handle.rs}\\
\noindent 行范围：60\\[0.4em]

\begin{rustcode}
...
const DT_DIR: u8 = 4;
// 本变量代码由AI完成
const DT_LNK: u8 = 10;
// 本变量代码由AI完成
const HEADER_SIZE: usize = 19;
...
\end{rustcode}

\subsubsection*{序号932}
\noindent 路径：\code{os/components/wateros-vfs/vfs-impl/impl-fs-bridge/src/dir_handle.rs}\\
\noindent 行范围：62\\[0.4em]

\begin{rustcode}
...
const DT_LNK: u8 = 10;
// 本变量代码由AI完成
const HEADER_SIZE: usize = 19;

#[inline]
...
\end{rustcode}

\subsubsection*{序号933}
\noindent 路径：\code{os/components/wateros-vfs/vfs-impl/impl-fs-bridge/src/dir_handle.rs}\\
\noindent 行范围：66-69\\[0.4em]

\begin{rustcode}
...
// 本方法代码由AI完成
fn dirent64_reclen(name_len: usize) -> usize {
    let with_name = HEADER_SIZE + name_len + 1;
    (with_name + 7) & !7
}
...
\end{rustcode}

\subsubsection*{序号934}
\noindent 路径：\code{os/components/wateros-vfs/vfs-impl/impl-fs-bridge/src/dir_handle.rs}\\
\noindent 行范围：73-80\\[0.4em]

\begin{rustcode}
...
pub(crate) fn node_type_to_dt(t: VfsNodeType) -> u8 {
    match t {
        VfsNodeType::File => DT_REG,
        VfsNodeType::Directory => DT_DIR,
        VfsNodeType::Symlink => DT_LNK,
...
\end{rustcode}

\subsubsection*{序号935}
\noindent 路径：\code{os/components/wateros-vfs/vfs-impl/impl-fs-bridge/src/dir_handle.rs}\\
\noindent 行范围：83-91\\[0.4em]

\begin{rustcode}
...
pub(crate) fn dirent64_encode_slice(
    buf: &mut [u8],
    ino: u64,
    next_off: i64,
    name: &str,
...
\end{rustcode}

\subsubsection*{序号936}
\noindent 路径：\code{os/components/wateros-vfs/vfs-impl/impl-fs-bridge/src/dir_handle.rs}\\
\noindent 行范围：94-112\\[0.4em]

\begin{rustcode}
...
pub(crate) fn encode_one(buf: &mut [u8], ino: u64, next_off: i64, name: &str, d_type: u8) -> Option<usize> {
    let reclen = dirent64_reclen(name.len());
    if buf.len() < reclen {
        return None;
    }
...
\end{rustcode}

\subsubsection*{序号937}
\noindent 路径：\code{os/components/wateros-vfs/vfs-impl/impl-fs-bridge/src/dir_handle.rs}\\
\noindent 行范围：116-118\\[0.4em]

\begin{rustcode}
...
// 本方法代码由AI完成
    fn metadata(&self) -> VfsResult<VfsMetadata> {
        Ok(self.meta.clone())
    }

...
\end{rustcode}

\subsubsection*{序号938}
\noindent 路径：\code{os/components/wateros-vfs/vfs-impl/impl-fs-bridge/src/dir_handle.rs}\\
\noindent 行范围：121-123\\[0.4em]

\begin{rustcode}
...
// 本方法代码由AI完成
    fn directory_path(&self) -> Option<&str> {
        Some(self.path.as_str())
    }

...
\end{rustcode}

\subsubsection*{序号939}
\noindent 路径：\code{os/components/wateros-vfs/vfs-impl/impl-fs-bridge/src/dir_handle.rs}\\
\noindent 行范围：126-128\\[0.4em]

\begin{rustcode}
...
// 本方法代码由AI完成
    fn backing_path(&self) -> Option<&str> {
        Some(self.path.as_str())
    }

...
\end{rustcode}

\subsubsection*{序号940}
\noindent 路径：\code{os/components/wateros-vfs/vfs-impl/impl-fs-bridge/src/dir_handle.rs}\\
\noindent 行范围：131-133\\[0.4em]

\begin{rustcode}
...
// 本方法代码由AI完成
    fn read(&mut self, _buf: &mut [u8]) -> VfsResult<usize> {
        Err(VfsError::NotAFile)
    }

...
\end{rustcode}

\subsubsection*{序号941}
\noindent 路径：\code{os/components/wateros-vfs/vfs-impl/impl-fs-bridge/src/dir_handle.rs}\\
\noindent 行范围：136-138\\[0.4em]

\begin{rustcode}
...
// 本方法代码由AI完成
    fn write(&mut self, _buf: &[u8]) -> VfsResult<usize> {
        Err(VfsError::NotAFile)
    }

...
\end{rustcode}

\subsubsection*{序号942}
\noindent 路径：\code{os/components/wateros-vfs/vfs-impl/impl-fs-bridge/src/dir_handle.rs}\\
\noindent 行范围：141-143\\[0.4em]

\begin{rustcode}
...
// 本方法代码由AI完成
    fn duplicate(&self) -> VfsResult<Box<dyn VfsIoHandle>> {
        Ok(Box::new(self.clone()))
    }

...
\end{rustcode}

\subsubsection*{序号943}
\noindent 路径：\code{os/components/wateros-vfs/vfs-impl/impl-fs-bridge/src/dir_handle.rs}\\
\noindent 行范围：146-165\\[0.4em]

\begin{rustcode}
...
    fn fill_getdents64(&mut self, buf: &mut [u8]) -> VfsResult<usize> {
        let bridge = FsBridge;
        self.load_dirents(&bridge)?;
        let entries = self.dirents.as_ref().expect("load_dirents");
        let mut out = 0usize;
...
\end{rustcode}

\subsubsection*{序号944}
\noindent 路径：\code{os/components/wateros-vfs/vfs-impl/impl-fs-bridge/src/file_handle.rs}\\
\noindent 行范围：1-325\\[0.4em]

\begin{rustcode}
//! ext4 根卷小文件缓冲句柄；普通 `open` 主路径已改由 [`super::paged_handle`] 处理。
//! 本模块代码由AI完成

extern crate alloc;

...
\end{rustcode}

\subsubsection*{序号945}
\noindent 路径：\code{os/components/wateros-vfs/vfs-impl/impl-fs-bridge/src/file_handle.rs}\\
\noindent 行范围：20-27\\[0.4em]

\begin{rustcode}
...
pub struct BufferedFileHandle {
    path : String,
    data : Vec<u8>,
    offset : u64,
    meta : VfsMetadata,
...
\end{rustcode}

\subsubsection*{序号946}
\noindent 路径：\code{os/components/wateros-vfs/vfs-impl/impl-fs-bridge/src/file_handle.rs}\\
\noindent 行范围：34-86\\[0.4em]

\begin{rustcode}
...
    pub(crate) fn open(bridge : &FsBridge, path : String, flags : VfsOpenFlags) -> VfsResult<Self> {
        let want_read = flags.contains(VfsOpenFlags::READ) ||
                        (!flags.contains(VfsOpenFlags::WRITE) &&
                         !flags.contains(VfsOpenFlags::CREATE));
        let want_write = flags.contains(VfsOpenFlags::WRITE);
...
\end{rustcode}

\subsubsection*{序号947}
\noindent 路径：\code{os/components/wateros-vfs/vfs-impl/impl-fs-bridge/src/file_handle.rs}\\
\noindent 行范围：89-94\\[0.4em]

\begin{rustcode}
...
    pub(crate) fn open_boxed(bridge : &FsBridge,
                             path : String,
                             flags : VfsOpenFlags)
                             -> VfsResult<Box<dyn VfsIoHandle>> {
        Ok(Box::new(Self::open(bridge, path, flags)?))
...
\end{rustcode}

\subsubsection*{序号948}
\noindent 路径：\code{os/components/wateros-vfs/vfs-impl/impl-fs-bridge/src/file_handle.rs}\\
\noindent 行范围：97-109\\[0.4em]

\begin{rustcode}
...
    fn sync_dirty(&mut self) -> VfsResult<()> {
        if !self.dirty {
            return Ok(());
        }
        if !FsBridge.exists(self.path.as_str())? {
...
\end{rustcode}

\subsubsection*{序号949}
\noindent 路径：\code{os/components/wateros-vfs/vfs-impl/impl-fs-bridge/src/file_handle.rs}\\
\noindent 行范围：114-129\\[0.4em]

\begin{rustcode}
...
    fn read(&mut self, buf : &mut [u8]) -> VfsResult<usize> {
        if buf.is_empty() {
            return Ok(0);
        }
        let off = usize::try_from(self.offset).map_err(|_| VfsError::Io)?;
...
\end{rustcode}

\subsubsection*{序号950}
\noindent 路径：\code{os/components/wateros-vfs/vfs-impl/impl-fs-bridge/src/file_handle.rs}\\
\noindent 行范围：132-151\\[0.4em]

\begin{rustcode}
...
    fn write(&mut self, buf : &[u8]) -> VfsResult<usize> {
        if !self.writable {
            return Err(VfsError::Unsupported);
        }
        if buf.is_empty() {
...
\end{rustcode}

\subsubsection*{序号951}
\noindent 路径：\code{os/components/wateros-vfs/vfs-impl/impl-fs-bridge/src/file_handle.rs}\\
\noindent 行范围：154\\[0.4em]

\begin{rustcode}
...

// 本方法代码由AI完成
    fn close(&mut self) -> VfsResult<()> { self.sync_dirty() }

// 本方法代码由AI完成
...
\end{rustcode}

\subsubsection*{序号952}
\noindent 路径：\code{os/components/wateros-vfs/vfs-impl/impl-fs-bridge/src/file_handle.rs}\\
\noindent 行范围：157-161\\[0.4em]

\begin{rustcode}
...
    fn metadata(&self) -> VfsResult<VfsMetadata> {
        let mut m = self.meta.clone();
        m.size = self.data.len() as u64;
        Ok(m)
    }
...
\end{rustcode}

\subsubsection*{序号953}
\noindent 路径：\code{os/components/wateros-vfs/vfs-impl/impl-fs-bridge/src/file_handle.rs}\\
\noindent 行范围：164-166\\[0.4em]

\begin{rustcode}
...
// 本方法代码由AI完成
    fn backing_path(&self) -> Option<&str> {
        Some(self.path.as_str())
    }

...
\end{rustcode}

\subsubsection*{序号954}
\noindent 路径：\code{os/components/wateros-vfs/vfs-impl/impl-fs-bridge/src/file_handle.rs}\\
\noindent 行范围：169-181\\[0.4em]

\begin{rustcode}
...
    fn read_at(&mut self, offset : u64, buf : &mut [u8]) -> VfsResult<usize> {
        if buf.is_empty() {
            return Ok(0);
        }
        let off = usize::try_from(offset).map_err(|_| VfsError::Io)?;
...
\end{rustcode}

\subsubsection*{序号955}
\noindent 路径：\code{os/components/wateros-vfs/vfs-impl/impl-fs-bridge/src/file_handle.rs}\\
\noindent 行范围：184-202\\[0.4em]

\begin{rustcode}
...
    fn write_at(&mut self, offset : u64, buf : &[u8]) -> VfsResult<usize> {
        if !self.writable {
            return Err(VfsError::Unsupported);
        }
        if buf.is_empty() {
...
\end{rustcode}

\subsubsection*{序号956}
\noindent 路径：\code{os/components/wateros-vfs/vfs-impl/impl-fs-bridge/src/file_handle.rs}\\
\noindent 行范围：205-218\\[0.4em]

\begin{rustcode}
...
    fn truncate(&mut self, len : u64) -> VfsResult<()> {
        if !self.writable {
            return Err(VfsError::Unsupported);
        }
        let len = usize::try_from(len).map_err(|_| VfsError::InvalidPath)?;
...
\end{rustcode}

\subsubsection*{序号957}
\noindent 路径：\code{os/components/wateros-vfs/vfs-impl/impl-fs-bridge/src/file_handle.rs}\\
\noindent 行范围：221-253\\[0.4em]

\begin{rustcode}
...
    fn seek(&mut self, offset : i64, whence : VfsSeekWhence) -> VfsResult<u64> {
        let new_off = match whence {
            VfsSeekWhence::Set => {
                if offset < 0 {
                    return Err(VfsError::InvalidPath);
...
\end{rustcode}

\subsubsection*{序号958}
\noindent 路径：\code{os/components/wateros-vfs/vfs-impl/impl-fs-bridge/src/file_handle.rs}\\
\noindent 行范围：256\\[0.4em]

\begin{rustcode}
...

// 本方法代码由AI完成
    fn flush(&mut self) -> VfsResult<()> { self.sync_dirty() }

// 本方法代码由AI完成
...
\end{rustcode}

\subsubsection*{序号959}
\noindent 路径：\code{os/components/wateros-vfs/vfs-impl/impl-fs-bridge/src/file_handle.rs}\\
\noindent 行范围：259\\[0.4em]

\begin{rustcode}
...

// 本方法代码由AI完成
    fn duplicate(&self) -> VfsResult<Box<dyn VfsIoHandle>> { Ok(Box::new(self.clone())) }

// 本方法代码由AI完成
...
\end{rustcode}

\subsubsection*{序号960}
\noindent 路径：\code{os/components/wateros-vfs/vfs-impl/impl-fs-bridge/src/file_handle.rs}\\
\noindent 行范围：262-275\\[0.4em]

\begin{rustcode}
...
    fn poll_revents(&mut self, events : i16) -> VfsResult<i16> {
// 本变量代码由AI完成
        const POLLIN : i16 = 0x001;
// 本变量代码由AI完成
        const POLLOUT : i16 = 0x004;
...
\end{rustcode}

\subsubsection*{序号961}
\noindent 路径：\code{os/components/wateros-vfs/vfs-impl/impl-fs-bridge/src/file_handle.rs}\\
\noindent 行范围：264\\[0.4em]

\begin{rustcode}
...
    fn poll_revents(&mut self, events : i16) -> VfsResult<i16> {
// 本变量代码由AI完成
        const POLLIN : i16 = 0x001;
// 本变量代码由AI完成
        const POLLOUT : i16 = 0x004;
...
\end{rustcode}

\subsubsection*{序号962}
\noindent 路径：\code{os/components/wateros-vfs/vfs-impl/impl-fs-bridge/src/file_handle.rs}\\
\noindent 行范围：266\\[0.4em]

\begin{rustcode}
...
        const POLLIN : i16 = 0x001;
// 本变量代码由AI完成
        const POLLOUT : i16 = 0x004;
        let mut revents = 0i16;
        if events & POLLIN != 0 {
...
\end{rustcode}

\subsubsection*{序号963}
\noindent 路径：\code{os/components/wateros-vfs/vfs-impl/impl-fs-bridge/src/file_handle.rs}\\
\noindent 行范围：280-324\\[0.4em]

\begin{rustcode}
...
    pub(crate) fn open_path(&self,
                            path : &str,
                            flags : VfsOpenFlags)
                            -> VfsResult<Box<dyn VfsIoHandle>> {
        let abs = if path.starts_with('/') {
...
\end{rustcode}

\subsubsection*{序号964}
\noindent 路径：\code{os/components/wateros-vfs/vfs-impl/impl-fs-bridge/src/lib.rs}\\
\noindent 行范围：1-1089\\[0.4em]

\begin{rustcode}
//! 将 [`api_v0::VfsBackend`] **桥接**到 `wateros-fs` 聚合 API；不 re-export `wateros-fs` 类型。
//! 本模块代码由AI完成

#![no_std]
#![allow(static_mut_refs)]
...
\end{rustcode}

\subsubsection*{序号965}
\noindent 路径：\code{os/components/wateros-vfs/vfs-impl/impl-fs-bridge/src/lib.rs}\\
\noindent 行范围：41\\[0.4em]

\begin{rustcode}
...
#[derive(Debug, Clone, Copy, Default)]
// 本结构代码由AI完成
pub struct FsBridge;

#[inline]
...
\end{rustcode}

\subsubsection*{序号966}
\noindent 路径：\code{os/components/wateros-vfs/vfs-impl/impl-fs-bridge/src/lib.rs}\\
\noindent 行范围：45-59\\[0.4em]

\begin{rustcode}
...
pub(crate) fn map_fs_err(e : FsError) -> VfsError {
    match e {
        FsError::NotMounted => VfsError::NotMounted,
        FsError::NotFound => VfsError::NotFound,
        FsError::NotAFile => VfsError::NotAFile,
...
\end{rustcode}

\subsubsection*{序号967}
\noindent 路径：\code{os/components/wateros-vfs/vfs-impl/impl-fs-bridge/src/lib.rs}\\
\noindent 行范围：63-71\\[0.4em]

\begin{rustcode}
...
fn map_fs_kind(kind : FsKind) -> VfsFsKind {
    match kind {
        FsKind::Ext2 => VfsFsKind::Ext2,
        FsKind::Ext3 => VfsFsKind::Ext3,
        FsKind::Ext4 => VfsFsKind::Ext4,
...
\end{rustcode}

\subsubsection*{序号968}
\noindent 路径：\code{os/components/wateros-vfs/vfs-impl/impl-fs-bridge/src/lib.rs}\\
\noindent 行范围：75-82\\[0.4em]

\begin{rustcode}
...
fn map_vfs_kind(kind : VfsFsKind) -> FsKind {
    match kind {
        VfsFsKind::Ext2 => FsKind::Ext2,
        VfsFsKind::Ext3 => FsKind::Ext3,
        VfsFsKind::Ext4 => FsKind::Ext4,
...
\end{rustcode}

\subsubsection*{序号969}
\noindent 路径：\code{os/components/wateros-vfs/vfs-impl/impl-fs-bridge/src/lib.rs}\\
\noindent 行范围：86-91\\[0.4em]

\begin{rustcode}
...
fn map_access(mode : FsAccessMode) -> VfsAccessMode {
    match mode {
        FsAccessMode::ReadOnly => VfsAccessMode::ReadOnly,
        FsAccessMode::ReadWrite => VfsAccessMode::ReadWrite,
    }
...
\end{rustcode}

\subsubsection*{序号970}
\noindent 路径：\code{os/components/wateros-vfs/vfs-impl/impl-fs-bridge/src/lib.rs}\\
\noindent 行范围：94-97\\[0.4em]

\begin{rustcode}
...
// 本方法代码由AI完成
fn map_fs_cap(c : FsCapability) -> VfsCapability {
    VfsCapability::new(map_fs_kind(c.kind),
                       map_access(c.access))
}
...
\end{rustcode}

\subsubsection*{序号971}
\noindent 路径：\code{os/components/wateros-vfs/vfs-impl/impl-fs-bridge/src/lib.rs}\\
\noindent 行范围：100-111\\[0.4em]

\begin{rustcode}
...
fn map_meta(m : FsMetadata, identity : MountIdentity) -> VfsMetadata {
    VfsMetadata { node_type : map_fs_node(m.node_type),
                  size : m.size,
                  mode : m.mode,
                  device_major : identity.device_major,
...
\end{rustcode}

\subsubsection*{序号972}
\noindent 路径：\code{os/components/wateros-vfs/vfs-impl/impl-fs-bridge/src/lib.rs}\\
\noindent 行范围：114-120\\[0.4em]

\begin{rustcode}
...
fn overlay_cached_size(abs_path : &str, meta : &mut VfsMetadata) {
    if meta.node_type != VfsNodeType::File {
        return;
    }
    let cache = impl_page_cache::global_cache(fs::rootfs::active_impl::mount_generation());
...
\end{rustcode}

\subsubsection*{序号973}
\noindent 路径：\code{os/components/wateros-vfs/vfs-impl/impl-fs-bridge/src/lib.rs}\\
\noindent 行范围：124-131\\[0.4em]

\begin{rustcode}
...
pub(crate) fn map_fs_node(t : FsNodeType) -> VfsNodeType {
    match t {
        FsNodeType::File => VfsNodeType::File,
        FsNodeType::Directory => VfsNodeType::Directory,
        FsNodeType::Symlink => VfsNodeType::Symlink,
...
\end{rustcode}

\subsubsection*{序号974}
\noindent 路径：\code{os/components/wateros-vfs/vfs-impl/impl-fs-bridge/src/lib.rs}\\
\noindent 行范围：134-137\\[0.4em]

\begin{rustcode}
...
// 本方法代码由AI完成
fn map_dir_entry(e : FsDirEntry) -> VfsDirEntry {
    VfsDirEntry { name : e.name,
                  node_type : map_fs_node(e.node_type) }
}
...
\end{rustcode}

\subsubsection*{序号975}
\noindent 路径：\code{os/components/wateros-vfs/vfs-impl/impl-fs-bridge/src/lib.rs}\\
\noindent 行范围：145-157\\[0.4em]

\begin{rustcode}
...
fn fs_and_rel_rw(path : &str) -> VfsResult<(SharedRwFs, String)> {
    match resolve_route(path)? {
        FsRoute::Root { abs, .. } => Ok((root_rw()?, abs)),
        FsRoute::AuxRw { fs, rel, readonly: true, .. } => {
            let _ = (fs, rel);
...
\end{rustcode}

\subsubsection*{序号976}
\noindent 路径：\code{os/components/wateros-vfs/vfs-impl/impl-fs-bridge/src/lib.rs}\\
\noindent 行范围：160-165\\[0.4em]

\begin{rustcode}
...
fn char_dev_exists(abs : &str) -> bool {
    if is_builtin_dev_path(abs) {
        return true;
    }
    fs::devfs::active_impl::lookup_character_device(abs).is_ok()
...
\end{rustcode}

\subsubsection*{序号977}
\noindent 路径：\code{os/components/wateros-vfs/vfs-impl/impl-fs-bridge/src/lib.rs}\\
\noindent 行范围：168\\[0.4em]

\begin{rustcode}
...

// 本方法代码由AI完成
fn char_dev_metadata(abs : &str) -> VfsMetadata { impl_fd_session::metadata_for_devfs_path(abs) }

// 本方法代码由AI完成
...
\end{rustcode}

\subsubsection*{序号978}
\noindent 路径：\code{os/components/wateros-vfs/vfs-impl/impl-fs-bridge/src/lib.rs}\\
\noindent 行范围：171-174\\[0.4em]

\begin{rustcode}
...
// 本方法代码由AI完成
fn is_builtin_dev_path(abs : &str) -> bool {
    matches!(abs,
             "/dev/null" | "/dev/zero" | "/dev/random" | "/dev/urandom" | "/dev/cpu_dma_latency")
}
...
\end{rustcode}

\subsubsection*{序号979}
\noindent 路径：\code{os/components/wateros-vfs/vfs-impl/impl-fs-bridge/src/lib.rs}\\
\noindent 行范围：177-182\\[0.4em]

\begin{rustcode}
...
const UNIXBENCH_SORT_SRC : &[u8] = b"the quick brown fox jumps over the lazy dog
unixbench shell script input line alpha
the shell test sorts this text repeatedly
wateros benchmark compatibility data
the end of the small sort source
...
\end{rustcode}

\subsubsection*{序号980}
\noindent 路径：\code{os/components/wateros-vfs/vfs-impl/impl-fs-bridge/src/lib.rs}\\
\noindent 行范围：185-191\\[0.4em]

\begin{rustcode}
...
fn unixbench_virtual_file(abs : &str) -> Option<(&'static [u8], u64)> {
    match abs {
        "/glibc/sort.src" => Some((UNIXBENCH_SORT_SRC, 0x7562_0001)),
        "/musl/sort.src" => Some((UNIXBENCH_SORT_SRC, 0x7562_0002)),
        _ => None,
...
\end{rustcode}

\subsubsection*{序号981}
\noindent 路径：\code{os/components/wateros-vfs/vfs-impl/impl-fs-bridge/src/lib.rs}\\
\noindent 行范围：194-206\\[0.4em]

\begin{rustcode}
...
fn unixbench_virtual_metadata(abs : &str, identity : MountIdentity) -> Option<VfsMetadata> {
    let (data, inode) = unixbench_virtual_file(abs)?;
    Some(VfsMetadata { node_type : VfsNodeType::File,
                       size : data.len() as u64,
                       mode : 0o444,
...
\end{rustcode}

\subsubsection*{序号982}
\noindent 路径：\code{os/components/wateros-vfs/vfs-impl/impl-fs-bridge/src/lib.rs}\\
\noindent 行范围：209-217\\[0.4em]

\begin{rustcode}
...
fn copy_virtual_range(data : &[u8], offset : u64, buf : &mut [u8]) -> usize {
    let start = offset as usize;
    if start >= data.len() {
        return 0;
    }
...
\end{rustcode}

\subsubsection*{序号983}
\noindent 路径：\code{os/components/wateros-vfs/vfs-impl/impl-fs-bridge/src/lib.rs}\\
\noindent 行范围：220\\[0.4em]

\begin{rustcode}
...

// 本方法代码由AI完成
fn securityfs_exists(rel : &str) -> bool { rel == "/" }

// 本方法代码由AI完成
...
\end{rustcode}

\subsubsection*{序号984}
\noindent 路径：\code{os/components/wateros-vfs/vfs-impl/impl-fs-bridge/src/lib.rs}\\
\noindent 行范围：223-237\\[0.4em]

\begin{rustcode}
...
fn securityfs_metadata(rel : &str, identity : MountIdentity) -> VfsResult<VfsMetadata> {
    if rel != "/" {
        return Err(VfsError::NotFound);
    }
    Ok(VfsMetadata { node_type : VfsNodeType::Directory,
...
\end{rustcode}

\subsubsection*{序号985}
\noindent 路径：\code{os/components/wateros-vfs/vfs-impl/impl-fs-bridge/src/lib.rs}\\
\noindent 行范围：240-245\\[0.4em]

\begin{rustcode}
...
fn securityfs_read_dir(rel : &str) -> VfsResult<Vec<fs::FsDirEntry>> {
    if rel == "/" {
        return Ok(Vec::new());
    }
    Err(VfsError::NotFound)
...
\end{rustcode}

\subsubsection*{序号986}
\noindent 路径：\code{os/components/wateros-vfs/vfs-impl/impl-fs-bridge/src/lib.rs}\\
\noindent 行范围：249-271\\[0.4em]

\begin{rustcode}
...
    fn exists(&self, path : &str) -> VfsResult<bool> {
        let abs = normalize_absolute_path(path)?;
        if char_dev_exists(abs.as_str()) {
            return Ok(true);
        }
...
\end{rustcode}

\subsubsection*{序号987}
\noindent 路径：\code{os/components/wateros-vfs/vfs-impl/impl-fs-bridge/src/lib.rs}\\
\noindent 行范围：274-320\\[0.4em]

\begin{rustcode}
...
    fn metadata(&self, path : &str) -> VfsResult<VfsMetadata> {
        let abs = normalize_absolute_path(path)?;
        if char_dev_exists(abs.as_str()) {
            return Ok(char_dev_metadata(abs.as_str()));
        }
...
\end{rustcode}

\subsubsection*{序号988}
\noindent 路径：\code{os/components/wateros-vfs/vfs-impl/impl-fs-bridge/src/lib.rs}\\
\noindent 行范围：323-349\\[0.4em]

\begin{rustcode}
...
    fn read(&self, path : &str) -> VfsResult<Vec<u8>> {
        let abs = normalize_absolute_path(path)?;
        match resolve_route(abs.as_str())? {
            FsRoute::PseudoProc { rel, .. } => proc_view().read(rel.as_str())
                                                          .map_err(map_fs_err),
...
\end{rustcode}

\subsubsection*{序号989}
\noindent 路径：\code{os/components/wateros-vfs/vfs-impl/impl-fs-bridge/src/lib.rs}\\
\noindent 行范围：352-354\\[0.4em]

\begin{rustcode}
...
// 本方法代码由AI完成
    fn read_range(&self, path : &str, offset : u64, buf : &mut [u8]) -> VfsResult<usize> {
        FsBridge::read_range(self, path, offset, buf)
    }

...
\end{rustcode}

\subsubsection*{序号990}
\noindent 路径：\code{os/components/wateros-vfs/vfs-impl/impl-fs-bridge/src/lib.rs}\\
\noindent 行范围：357-376\\[0.4em]

\begin{rustcode}
...
    fn read_dir(&self, path : &str) -> VfsResult<Vec<VfsDirEntry>> {
        let abs = normalize_absolute_path(path)?;
        let entries = match resolve_route(abs.as_str())? {
            FsRoute::PseudoProc { rel, .. } => proc_view().read_dir(rel.as_str())
                                                          .map_err(map_fs_err)?,
...
\end{rustcode}

\subsubsection*{序号991}
\noindent 路径：\code{os/components/wateros-vfs/vfs-impl/impl-fs-bridge/src/lib.rs}\\
\noindent 行范围：379-381\\[0.4em]

\begin{rustcode}
...
// 本方法代码由AI完成
    fn boot_dump_all_paths(&self) {
        // bring-up 单 RW 根卷：启动树打印仍可由 fs 层自检触发。
    }
}
...
\end{rustcode}

\subsubsection*{序号992}
\noindent 路径：\code{os/components/wateros-vfs/vfs-impl/impl-fs-bridge/src/lib.rs}\\
\noindent 行范围：387-398\\[0.4em]

\begin{rustcode}
...
    pub(crate) fn read_dir_on_root(path : &str) -> VfsResult<Vec<VfsDirEntry>> {
        let n = normalize_absolute_path(path)?;
        let fs = root_rw()?;
        fs.lock()
          .read_dir(n.as_str())
...
\end{rustcode}

\subsubsection*{序号993}
\noindent 路径：\code{os/components/wateros-vfs/vfs-impl/impl-fs-bridge/src/lib.rs}\\
\noindent 行范围：402-440\\[0.4em]

\begin{rustcode}
...
    pub(crate) fn read_range(&self,
                             path : &str,
                             offset : u64,
                             buf : &mut [u8])
                             -> VfsResult<usize> {
...
\end{rustcode}

\subsubsection*{序号994}
\noindent 路径：\code{os/components/wateros-vfs/vfs-impl/impl-fs-bridge/src/lib.rs}\\
\noindent 行范围：445-453\\[0.4em]

\begin{rustcode}
...
pub fn mount_ext4_block_at(mount_point : &str, block_dev : &str, readonly : bool) -> VfsResult<()> {
    if readonly {
        let aux = fs::mount_aux_ro_from_block_path(block_dev).map_err(map_fs_err)?;
        mount_table::mount_aux_at_ro(mount_point, aux, block_dev)
    } else {
...
\end{rustcode}

\subsubsection*{序号995}
\noindent 路径：\code{os/components/wateros-vfs/vfs-impl/impl-fs-bridge/src/lib.rs}\\
\noindent 行范围：457-464\\[0.4em]

\begin{rustcode}
...
pub(crate) fn assert_mount_point_directory(path: &str) -> VfsResult<()> {
    let bridge = FsBridge;
    let meta = bridge.metadata(path)?;
    if meta.node_type != VfsNodeType::Directory {
        return Err(VfsError::NotAFile);
...
\end{rustcode}

\subsubsection*{序号996}
\noindent 路径：\code{os/components/wateros-vfs/vfs-impl/impl-fs-bridge/src/lib.rs}\\
\noindent 行范围：468-470\\[0.4em]

\begin{rustcode}
...
// 本方法代码由AI完成
pub fn mount_tmpfs_at(mount_point : &str) -> VfsResult<()> {
    mount_table::mount_tmpfs_at(mount_point)
}

...
\end{rustcode}

\subsubsection*{序号997}
\noindent 路径：\code{os/components/wateros-vfs/vfs-impl/impl-fs-bridge/src/lib.rs}\\
\noindent 行范围：474-476\\[0.4em]

\begin{rustcode}
...
// 本方法代码由AI完成
pub fn mount_cgroup_at(mount_point : &str, v2 : bool, options : &str) -> VfsResult<()> {
    mount_table::mount_cgroup_at(mount_point, v2, options)
}

...
\end{rustcode}

\subsubsection*{序号998}
\noindent 路径：\code{os/components/wateros-vfs/vfs-impl/impl-fs-bridge/src/lib.rs}\\
\noindent 行范围：480-482\\[0.4em]

\begin{rustcode}
...
// 本方法代码由AI完成
pub fn mount_statfs_magic(path : &str) -> Option<isize> {
    mount_table::mount_statfs_magic(path)
}

...
\end{rustcode}

\subsubsection*{序号999}
\noindent 路径：\code{os/components/wateros-vfs/vfs-impl/impl-fs-bridge/src/lib.rs}\\
\noindent 行范围：486-488\\[0.4em]

\begin{rustcode}
...
// 本方法代码由AI完成
pub fn remount_readonly_at(mount_point : &str) -> VfsResult<()> {
    mount_table::remount_aux_readonly(mount_point)
}

...
\end{rustcode}

\subsubsection*{序号1000}
\noindent 路径：\code{os/components/wateros-vfs/vfs-impl/impl-fs-bridge/src/lib.rs}\\
\noindent 行范围：492-494\\[0.4em]

\begin{rustcode}
...
// 本方法代码由AI完成
pub fn unmount_at(mount_point : &str, detach : bool) -> VfsResult<()> {
    mount_table::unmount_aux_at(mount_point, detach)
}

...
\end{rustcode}

\subsubsection*{序号1001}
\noindent 路径：\code{os/components/wateros-vfs/vfs-impl/impl-fs-bridge/src/lib.rs}\\
\noindent 行范围：498-505\\[0.4em]

\begin{rustcode}
...
pub fn ensure_proc_mount_point() -> VfsResult<()> {
    let path = "/proc";
    let bridge = FsBridge;
    if bridge.exists(path)? {
        return Ok(());
...
\end{rustcode}

\subsubsection*{序号1002}
\noindent 路径：\code{os/components/wateros-vfs/vfs-impl/impl-fs-bridge/src/lib.rs}\\
\noindent 行范围：509-511\\[0.4em]

\begin{rustcode}
...
// 本方法代码由AI完成
pub fn mount_procfs_at(mount_point : &str) -> VfsResult<()> {
    mount_table::mount_aux_proc_at(mount_point)
}

...
\end{rustcode}

\subsubsection*{序号1003}
\noindent 路径：\code{os/components/wateros-vfs/vfs-impl/impl-fs-bridge/src/lib.rs}\\
\noindent 行范围：514-516\\[0.4em]

\begin{rustcode}
...
// 本方法代码由AI完成
pub fn mount_securityfs_at(mount_point : &str) -> VfsResult<()> {
    mount_table::mount_securityfs_at(mount_point)
}

...
\end{rustcode}

\subsubsection*{序号1004}
\noindent 路径：\code{os/components/wateros-vfs/vfs-impl/impl-fs-bridge/src/lib.rs}\\
\noindent 行范围：519-521\\[0.4em]

\begin{rustcode}
...
// 本方法代码由AI完成
pub fn mount_bind_at(source : &str, target : &str, recursive : bool) -> VfsResult<()> {
    mount_table::mount_bind_at(source, target, recursive)
}

...
\end{rustcode}

\subsubsection*{序号1005}
\noindent 路径：\code{os/components/wateros-vfs/vfs-impl/impl-fs-bridge/src/lib.rs}\\
\noindent 行范围：524-526\\[0.4em]

\begin{rustcode}
...
// 本方法代码由AI完成
pub fn move_mount_at(source : &str, target : &str) -> VfsResult<()> {
    mount_table::move_mount_at(source, target)
}

...
\end{rustcode}

\subsubsection*{序号1006}
\noindent 路径：\code{os/components/wateros-vfs/vfs-impl/impl-fs-bridge/src/lib.rs}\\
\noindent 行范围：531-536\\[0.4em]

\begin{rustcode}
...
pub fn set_mount_propagation(mount_point : &str,
                             propagation : MountPropagation,
                             recursive : bool)
                             -> VfsResult<()> {
    mount_table::set_mount_propagation(mount_point, propagation, recursive)
...
\end{rustcode}

\subsubsection*{序号1007}
\noindent 路径：\code{os/components/wateros-vfs/vfs-impl/impl-fs-bridge/src/lib.rs}\\
\noindent 行范围：549-564\\[0.4em]

\begin{rustcode}
...
pub fn unlink_path(path : &str, remove_dir : bool) -> VfsResult<()> {
    mount_table::assert_path_writable(path)?;
    let (fs, rel) = fs_and_rel_rw(path)?;
    let mut sess = MountedRwSession::new(fs);
    let result = if remove_dir {
...
\end{rustcode}

\subsubsection*{序号1008}
\noindent 路径：\code{os/components/wateros-vfs/vfs-impl/impl-fs-bridge/src/lib.rs}\\
\noindent 行范围：572-579\\[0.4em]

\begin{rustcode}
...
pub fn reset_file_page_cache() -> VfsResult<()> {
    let mount_gen = fs::rootfs::active_impl::mount_generation();
    let cache = impl_page_cache::global_cache(mount_gen);
    let mut io = paged_handle::FsPageIo;
    cache.flush_all(&mut io, core::convert::identity)?;
...
\end{rustcode}

\subsubsection*{序号1009}
\noindent 路径：\code{os/components/wateros-vfs/vfs-impl/impl-fs-bridge/src/lib.rs}\\
\noindent 行范围：583-592\\[0.4em]

\begin{rustcode}
...
pub fn mkdir_path(path : &str, mode : u32) -> VfsResult<()> {
    let normalized = normalize_absolute_path(path)?;
    if normalized.as_str() == "/" {
        return Err(VfsError::Exists);
    }
...
\end{rustcode}

\subsubsection*{序号1010}
\noindent 路径：\code{os/components/wateros-vfs/vfs-impl/impl-fs-bridge/src/lib.rs}\\
\noindent 行范围：596-605\\[0.4em]

\begin{rustcode}
...
pub fn symlink_path(link_path : &str, target : &str) -> VfsResult<()> {
    let normalized = normalize_absolute_path(link_path)?;
    if normalized.as_str() == "/" {
        return Err(VfsError::InvalidPath);
    }
...
\end{rustcode}

\subsubsection*{序号1011}
\noindent 路径：\code{os/components/wateros-vfs/vfs-impl/impl-fs-bridge/src/lib.rs}\\
\noindent 行范围：609-623\\[0.4em]

\begin{rustcode}
...
pub fn read_symlink_path(path : &str) -> VfsResult<Vec<u8>> {
    let abs = normalize_absolute_path(path)?;
    match resolve_route(abs.as_str())? {
        FsRoute::PseudoProc { .. } | FsRoute::PseudoSecurity { .. } => Err(VfsError::NotAFile),
        FsRoute::Root { abs, .. } => root_rw()?.lock()
...
\end{rustcode}

\subsubsection*{序号1012}
\noindent 路径：\code{os/components/wateros-vfs/vfs-impl/impl-fs-bridge/src/lib.rs}\\
\noindent 行范围：627-644\\[0.4em]

\begin{rustcode}
...
pub fn truncate_path(path : &str, len : u64) -> VfsResult<()> {
    let normalized = normalize_absolute_path(path)?;
    if char_dev_exists(normalized.as_str()) {
        return Err(VfsError::Unsupported);
    }
...
\end{rustcode}

\subsubsection*{序号1013}
\noindent 路径：\code{os/components/wateros-vfs/vfs-impl/impl-fs-bridge/src/lib.rs}\\
\noindent 行范围：648-657\\[0.4em]

\begin{rustcode}
...
pub fn chmod_path(path : &str, mode : u32) -> VfsResult<()> {
    let normalized = normalize_absolute_path(path)?;
    if char_dev_exists(normalized.as_str()) {
        return Err(VfsError::Unsupported);
    }
...
\end{rustcode}

\subsubsection*{序号1014}
\noindent 路径：\code{os/components/wateros-vfs/vfs-impl/impl-fs-bridge/src/lib.rs}\\
\noindent 行范围：661-675\\[0.4em]

\begin{rustcode}
...
pub fn chown_path(path : &str, uid : Option<u32>, gid : Option<u32>) -> VfsResult<()> {
    let normalized = normalize_absolute_path(path)?;
    if char_dev_exists(normalized.as_str()) {
        return Err(VfsError::Unsupported);
    }
...
\end{rustcode}

\subsubsection*{序号1015}
\noindent 路径：\code{os/components/wateros-vfs/vfs-impl/impl-fs-bridge/src/lib.rs}\\
\noindent 行范围：678\\[0.4em]

\begin{rustcode}
...

// 本变量代码由AI完成
const CGROUP_SUPER_MAGIC: isize = 0x0027_e0eb;
// 本变量代码由AI完成
const CGROUP2_SUPER_MAGIC: isize = 0x6367_7270;
...
\end{rustcode}

\subsubsection*{序号1016}
\noindent 路径：\code{os/components/wateros-vfs/vfs-impl/impl-fs-bridge/src/lib.rs}\\
\noindent 行范围：680\\[0.4em]

\begin{rustcode}
...
const CGROUP_SUPER_MAGIC: isize = 0x0027_e0eb;
// 本变量代码由AI完成
const CGROUP2_SUPER_MAGIC: isize = 0x6367_7270;

// 本方法代码由AI完成
...
\end{rustcode}

\subsubsection*{序号1017}
\noindent 路径：\code{os/components/wateros-vfs/vfs-impl/impl-fs-bridge/src/lib.rs}\\
\noindent 行范围：683-688\\[0.4em]

\begin{rustcode}
...
fn path_on_cgroup_fs(path: &str) -> bool {
    matches!(
        mount_table::mount_statfs_magic(path),
        Some(CGROUP_SUPER_MAGIC) | Some(CGROUP2_SUPER_MAGIC)
    )
...
\end{rustcode}

\subsubsection*{序号1018}
\noindent 路径：\code{os/components/wateros-vfs/vfs-impl/impl-fs-bridge/src/lib.rs}\\
\noindent 行范围：692-712\\[0.4em]

\begin{rustcode}
...
pub fn validate_xattr_name(path: &str, name: &str) -> VfsResult<()> {
    if name.is_empty() || name.len() > 255 {
        return Err(VfsError::InvalidPath);
    }
    if !name.contains('.') {
...
\end{rustcode}

\subsubsection*{序号1019}
\noindent 路径：\code{os/components/wateros-vfs/vfs-impl/impl-fs-bridge/src/lib.rs}\\
\noindent 行范围：715-718\\[0.4em]

\begin{rustcode}
...
// 本方法代码由AI完成
fn xattr_route(path: &str) -> VfsResult<(SharedRwFs, String)> {
    mount_table::assert_path_writable(path)?;
    fs_and_rel_rw(path)
}
...
\end{rustcode}

\subsubsection*{序号1020}
\noindent 路径：\code{os/components/wateros-vfs/vfs-impl/impl-fs-bridge/src/lib.rs}\\
\noindent 行范围：721-726\\[0.4em]

\begin{rustcode}
...
fn map_xattr_get_err(err: VfsError) -> VfsError {
    match err {
        VfsError::InvalidPath => VfsError::NotFound,
        other => other,
    }
...
\end{rustcode}

\subsubsection*{序号1021}
\noindent 路径：\code{os/components/wateros-vfs/vfs-impl/impl-fs-bridge/src/lib.rs}\\
\noindent 行范围：730-739\\[0.4em]

\begin{rustcode}
...
pub fn setxattr_path(path: &str, name: &str, value: &[u8]) -> VfsResult<()> {
    validate_xattr_name(path, name)?;
    let normalized = normalize_absolute_path(path)?;
    if char_dev_exists(normalized.as_str()) {
        return Err(VfsError::Unsupported);
...
\end{rustcode}

\subsubsection*{序号1022}
\noindent 路径：\code{os/components/wateros-vfs/vfs-impl/impl-fs-bridge/src/lib.rs}\\
\noindent 行范围：743-753\\[0.4em]

\begin{rustcode}
...
pub fn getxattr_path(path: &str, name: &str, buf: &mut [u8]) -> VfsResult<usize> {
    validate_xattr_name(path, name)?;
    let normalized = normalize_absolute_path(path)?;
    if char_dev_exists(normalized.as_str()) {
        return Err(VfsError::Unsupported);
...
\end{rustcode}

\subsubsection*{序号1023}
\noindent 路径：\code{os/components/wateros-vfs/vfs-impl/impl-fs-bridge/src/lib.rs}\\
\noindent 行范围：757-765\\[0.4em]

\begin{rustcode}
...
pub fn listxattr_path(path: &str, buf: &mut [u8]) -> VfsResult<usize> {
    let normalized = normalize_absolute_path(path)?;
    if char_dev_exists(normalized.as_str()) {
        return Err(VfsError::Unsupported);
    }
...
\end{rustcode}

\subsubsection*{序号1024}
\noindent 路径：\code{os/components/wateros-vfs/vfs-impl/impl-fs-bridge/src/lib.rs}\\
\noindent 行范围：769-779\\[0.4em]

\begin{rustcode}
...
pub fn removexattr_path(path: &str, name: &str) -> VfsResult<()> {
    validate_xattr_name(path, name)?;
    let normalized = normalize_absolute_path(path)?;
    if char_dev_exists(normalized.as_str()) {
        return Err(VfsError::Unsupported);
...
\end{rustcode}

\subsubsection*{序号1025}
\noindent 路径：\code{os/components/wateros-vfs/vfs-impl/impl-fs-bridge/src/lib.rs}\\
\noindent 行范围：782-787\\[0.4em]

\begin{rustcode}
...
pub(crate) fn replace_file_contents(path : &str, data : &[u8]) -> VfsResult<()> {
    mount_table::assert_path_writable(path)?;
    let (fs, rel) = fs_and_rel_rw(path)?;
    let mut sess = MountedRwSession::new(fs);
    sess.write_regular_file(rel.as_str(), data)
...
\end{rustcode}

\subsubsection*{序号1026}
\noindent 路径：\code{os/components/wateros-vfs/vfs-impl/impl-fs-bridge/src/lib.rs}\\
\noindent 行范围：791-807\\[0.4em]

\begin{rustcode}
...
pub fn overwrite_file_at(path : &str, data : &[u8]) -> VfsResult<()> {
    match replace_file_contents(path, data) {
        Ok(()) => {}
        Err(VfsError::NotFound) => {
            match unlink_path(path, false) {
...
\end{rustcode}

\subsubsection*{序号1027}
\noindent 路径：\code{os/components/wateros-vfs/vfs-impl/impl-fs-bridge/src/lib.rs}\\
\noindent 行范围：811-816\\[0.4em]

\begin{rustcode}
...
pub fn mknod_socket_at(path : &str) -> VfsResult<()> {
    mount_table::assert_path_writable(path)?;
    let (fs, rel) = fs_and_rel_rw(path)?;
    let mut sess = MountedRwSession::new(fs);
    sess.mknod(rel.as_str(), 0o140_777, 0)
...
\end{rustcode}

\subsubsection*{序号1028}
\noindent 路径：\code{os/components/wateros-vfs/vfs-impl/impl-fs-bridge/src/lib.rs}\\
\noindent 行范围：820-830\\[0.4em]

\begin{rustcode}
...
pub fn rename_path(old_path : &str, new_path : &str) -> VfsResult<()> {
    mount_table::assert_path_writable(old_path)?;
    mount_table::assert_path_writable(new_path)?;
    let (fs_old, rel_old) = fs_and_rel_rw(old_path)?;
    let (fs_new, rel_new) = fs_and_rel_rw(new_path)?;
...
\end{rustcode}

\subsubsection*{序号1029}
\noindent 路径：\code{os/components/wateros-vfs/vfs-impl/impl-fs-bridge/src/lib.rs}\\
\noindent 行范围：834-836\\[0.4em]

\begin{rustcode}
...
// 本结构代码由AI完成
pub struct MountedRwSession {
    inner : SharedRwFs,
}

...
\end{rustcode}

\subsubsection*{序号1030}
\noindent 路径：\code{os/components/wateros-vfs/vfs-impl/impl-fs-bridge/src/lib.rs}\\
\noindent 行范围：840\\[0.4em]

\begin{rustcode}
...
impl MountedRwSession {
// 本方法代码由AI完成
    pub fn new(inner : SharedRwFs) -> Self { Self { inner } }
}

...
\end{rustcode}

\subsubsection*{序号1031}
\noindent 路径：\code{os/components/wateros-vfs/vfs-impl/impl-fs-bridge/src/lib.rs}\\
\noindent 行范围：845-851\\[0.4em]

\begin{rustcode}
...
    fn write_regular_file_at_root(&mut self, name : &str, data : &[u8]) -> VfsResult<()> {
        validate_root_file_name(name)?;
        self.inner
            .lock()
            .write_regular_file_at_root(name, data)
...
\end{rustcode}

\subsubsection*{序号1032}
\noindent 路径：\code{os/components/wateros-vfs/vfs-impl/impl-fs-bridge/src/lib.rs}\\
\noindent 行范围：854-860\\[0.4em]

\begin{rustcode}
...
    fn write_regular_file(&mut self, path : &str, data : &[u8]) -> VfsResult<()> {
        let n = normalize_absolute_path(path)?;
        self.inner
            .lock()
            .write_regular_file(n.as_str(), data)
...
\end{rustcode}

\subsubsection*{序号1033}
\noindent 路径：\code{os/components/wateros-vfs/vfs-impl/impl-fs-bridge/src/lib.rs}\\
\noindent 行范围：863-869\\[0.4em]

\begin{rustcode}
...
    fn unlink(&mut self, path : &str) -> VfsResult<()> {
        let n = normalize_absolute_path(path)?;
        self.inner
            .lock()
            .unlink(n.as_str())
...
\end{rustcode}

\subsubsection*{序号1034}
\noindent 路径：\code{os/components/wateros-vfs/vfs-impl/impl-fs-bridge/src/lib.rs}\\
\noindent 行范围：872-878\\[0.4em]

\begin{rustcode}
...
    fn rmdir(&mut self, path : &str) -> VfsResult<()> {
        let n = normalize_absolute_path(path)?;
        self.inner
            .lock()
            .rmdir(n.as_str())
...
\end{rustcode}

\subsubsection*{序号1035}
\noindent 路径：\code{os/components/wateros-vfs/vfs-impl/impl-fs-bridge/src/lib.rs}\\
\noindent 行范围：881-887\\[0.4em]

\begin{rustcode}
...
    fn write_range(&mut self, path : &str, offset : u64, data : &[u8]) -> VfsResult<usize> {
        let n = normalize_absolute_path(path)?;
        self.inner
            .lock()
            .write_range(n.as_str(), offset, data)
...
\end{rustcode}

\subsubsection*{序号1036}
\noindent 路径：\code{os/components/wateros-vfs/vfs-impl/impl-fs-bridge/src/lib.rs}\\
\noindent 行范围：890-896\\[0.4em]

\begin{rustcode}
...
    fn truncate(&mut self, path : &str, len : u64) -> VfsResult<()> {
        let n = normalize_absolute_path(path)?;
        self.inner
            .lock()
            .truncate(n.as_str(), len)
...
\end{rustcode}

\subsubsection*{序号1037}
\noindent 路径：\code{os/components/wateros-vfs/vfs-impl/impl-fs-bridge/src/lib.rs}\\
\noindent 行范围：899-905\\[0.4em]

\begin{rustcode}
...
    fn mkdir(&mut self, path : &str, mode : u32) -> VfsResult<()> {
        let n = normalize_absolute_path(path)?;
        self.inner
            .lock()
            .mkdir(n.as_str(), mode)
...
\end{rustcode}

\subsubsection*{序号1038}
\noindent 路径：\code{os/components/wateros-vfs/vfs-impl/impl-fs-bridge/src/lib.rs}\\
\noindent 行范围：908-914\\[0.4em]

\begin{rustcode}
...
    fn chmod(&mut self, path : &str, mode : u32) -> VfsResult<()> {
        let n = normalize_absolute_path(path)?;
        self.inner
            .lock()
            .chmod(n.as_str(), mode)
...
\end{rustcode}

\subsubsection*{序号1039}
\noindent 路径：\code{os/components/wateros-vfs/vfs-impl/impl-fs-bridge/src/lib.rs}\\
\noindent 行范围：917-923\\[0.4em]

\begin{rustcode}
...
    fn chown(&mut self, path : &str, uid : Option<u32>, gid : Option<u32>) -> VfsResult<()> {
        let n = normalize_absolute_path(path)?;
        self.inner
            .lock()
            .chown(n.as_str(), uid, gid)
...
\end{rustcode}

\subsubsection*{序号1040}
\noindent 路径：\code{os/components/wateros-vfs/vfs-impl/impl-fs-bridge/src/lib.rs}\\
\noindent 行范围：926-932\\[0.4em]

\begin{rustcode}
...
    fn setxattr(&mut self, path : &str, name : &str, value : &[u8]) -> VfsResult<()> {
        let n = normalize_absolute_path(path)?;
        self.inner
            .lock()
            .setxattr(n.as_str(), name, value)
...
\end{rustcode}

\subsubsection*{序号1041}
\noindent 路径：\code{os/components/wateros-vfs/vfs-impl/impl-fs-bridge/src/lib.rs}\\
\noindent 行范围：935-941\\[0.4em]

\begin{rustcode}
...
    fn getxattr(&self, path : &str, name : &str, buf : &mut [u8]) -> VfsResult<usize> {
        let n = normalize_absolute_path(path)?;
        self.inner
            .lock()
            .getxattr(n.as_str(), name, buf)
...
\end{rustcode}

\subsubsection*{序号1042}
\noindent 路径：\code{os/components/wateros-vfs/vfs-impl/impl-fs-bridge/src/lib.rs}\\
\noindent 行范围：944-950\\[0.4em]

\begin{rustcode}
...
    fn listxattr(&self, path : &str, buf : &mut [u8]) -> VfsResult<usize> {
        let n = normalize_absolute_path(path)?;
        self.inner
            .lock()
            .listxattr(n.as_str(), buf)
...
\end{rustcode}

\subsubsection*{序号1043}
\noindent 路径：\code{os/components/wateros-vfs/vfs-impl/impl-fs-bridge/src/lib.rs}\\
\noindent 行范围：953-959\\[0.4em]

\begin{rustcode}
...
    fn removexattr(&mut self, path : &str, name : &str) -> VfsResult<()> {
        let n = normalize_absolute_path(path)?;
        self.inner
            .lock()
            .removexattr(n.as_str(), name)
...
\end{rustcode}

\subsubsection*{序号1044}
\noindent 路径：\code{os/components/wateros-vfs/vfs-impl/impl-fs-bridge/src/lib.rs}\\
\noindent 行范围：962-969\\[0.4em]

\begin{rustcode}
...
    fn rename(&mut self, old_path : &str, new_path : &str) -> VfsResult<()> {
        let old = normalize_absolute_path(old_path)?;
        let new = normalize_absolute_path(new_path)?;
        self.inner
            .lock()
...
\end{rustcode}

\subsubsection*{序号1045}
\noindent 路径：\code{os/components/wateros-vfs/vfs-impl/impl-fs-bridge/src/lib.rs}\\
\noindent 行范围：972-979\\[0.4em]

\begin{rustcode}
...
    fn hardlink(&mut self, existing_path : &str, new_path : &str) -> VfsResult<()> {
        let existing = normalize_absolute_path(existing_path)?;
        let new = normalize_absolute_path(new_path)?;
        self.inner
            .lock()
...
\end{rustcode}

\subsubsection*{序号1046}
\noindent 路径：\code{os/components/wateros-vfs/vfs-impl/impl-fs-bridge/src/lib.rs}\\
\noindent 行范围：982-988\\[0.4em]

\begin{rustcode}
...
    fn symlink(&mut self, link_path : &str, target : &str) -> VfsResult<()> {
        let link = normalize_absolute_path(link_path)?;
        self.inner
            .lock()
            .symlink(link.as_str(), target)
...
\end{rustcode}

\subsubsection*{序号1047}
\noindent 路径：\code{os/components/wateros-vfs/vfs-impl/impl-fs-bridge/src/lib.rs}\\
\noindent 行范围：991-997\\[0.4em]

\begin{rustcode}
...
    fn mknod(&mut self, path : &str, mode : u32, rdev : u32) -> VfsResult<()> {
        let n = normalize_absolute_path(path)?;
        self.inner
            .lock()
            .mknod(n.as_str(), mode, rdev)
...
\end{rustcode}

\subsubsection*{序号1048}
\noindent 路径：\code{os/components/wateros-vfs/vfs-impl/impl-fs-bridge/src/lib.rs}\\
\noindent 行范围：1002-1006\\[0.4em]

\begin{rustcode}
...
    fn supported_capabilities(&self) -> Vec<VfsCapability> {
        fs::supported_fs_summary().into_iter()
                                  .map(map_fs_cap)
                                  .collect()
    }
...
\end{rustcode}

\subsubsection*{序号1049}
\noindent 路径：\code{os/components/wateros-vfs/vfs-impl/impl-fs-bridge/src/lib.rs}\\
\noindent 行范围：1009-1011\\[0.4em]

\begin{rustcode}
...
// 本方法代码由AI完成
    fn mount_rw_session(&self, _kind : VfsFsKind) -> VfsResult<Box<dyn RootRwSession>> {
        Ok(Box::new(MountedRwSession::new(root_rw()?)))
    }
}
...
\end{rustcode}

\subsubsection*{序号1050}
\noindent 路径：\code{os/components/wateros-vfs/vfs-impl/impl-fs-bridge/src/lib.rs}\\
\noindent 行范围：1016-1048\\[0.4em]

\begin{rustcode}
...
    fn list_dev_nodes(&self) -> Vec<VfsDevNode> {
        let mut nodes = fs::devfs::active_impl::list_nodes().into_iter()
                                                            .map(|n| {
                                                                VfsDevNode {
                path: n.path,
...
\end{rustcode}

\subsubsection*{序号1051}
\noindent 路径：\code{os/components/wateros-vfs/vfs-impl/impl-fs-bridge/src/lib.rs}\\
\noindent 行范围：1051-1053\\[0.4em]

\begin{rustcode}
...
// 本方法代码由AI完成
    fn default_root_block_path(&self) -> Option<String> {
        fs::devfs::active_impl::default_root_block_path()
    }
}
...
\end{rustcode}

\subsubsection*{序号1052}
\noindent 路径：\code{os/components/wateros-vfs/vfs-impl/impl-fs-bridge/src/lib.rs}\\
\noindent 行范围：1058-1060\\[0.4em]

\begin{rustcode}
...
// 本方法代码由AI完成
    fn open(&self, path : &str, flags : VfsOpenFlags) -> VfsResult<Box<dyn VfsIoHandle>> {
        self.open_path(path, flags)
    }
}
...
\end{rustcode}

\subsubsection*{序号1053}
\noindent 路径：\code{os/components/wateros-vfs/vfs-impl/impl-fs-bridge/src/lib.rs}\\
\noindent 行范围：1065-1068\\[0.4em]

\begin{rustcode}
...
// 本方法代码由AI完成
    fn mount_at(&mut self, mount_point : &str, _kind : VfsFsKind) -> VfsResult<()> {
        let _ = mount_point;
        Err(VfsError::Unsupported)
    }
...
\end{rustcode}

\subsubsection*{序号1054}
\noindent 路径：\code{os/components/wateros-vfs/vfs-impl/impl-fs-bridge/src/lib.rs}\\
\noindent 行范围：1071-1073\\[0.4em]

\begin{rustcode}
...
// 本方法代码由AI完成
    fn unmount_at(&mut self, mount_point : &str) -> VfsResult<()> {
        mount_table::unmount_aux_at(mount_point, false)
    }

...
\end{rustcode}

\subsubsection*{序号1055}
\noindent 路径：\code{os/components/wateros-vfs/vfs-impl/impl-fs-bridge/src/lib.rs}\\
\noindent 行范围：1076-1079\\[0.4em]

\begin{rustcode}
...
// 本方法代码由AI完成
    fn resolve_mount(&self, path : &str) -> VfsResult<&str> {
        let _ = path;
        Err(VfsError::Unsupported)
    }
...
\end{rustcode}

\subsubsection*{序号1056}
\noindent 路径：\code{os/components/wateros-vfs/vfs-impl/impl-fs-bridge/src/lib.rs}\\
\noindent 行范围：1085-1089\\[0.4em]

\begin{rustcode}
...
pub fn test() {
    api_v0::test();
    let _ = FsBridge::default();
    let _ = mount_table::mount_table_self_test();
}
\end{rustcode}

\subsubsection*{序号1057}
\noindent 路径：\code{os/components/wateros-vfs/vfs-impl/impl-fs-bridge/src/mount_ns.rs}\\
\noindent 行范围：1-141\\[0.4em]

\begin{rustcode}
//! per-task 挂载命名空间注册表（共享/复制语义对齐 cwd 表）。
//! 本模块代码由AI完成

extern crate alloc;

...
\end{rustcode}

\subsubsection*{序号1058}
\noindent 路径：\code{os/components/wateros-vfs/vfs-impl/impl-fs-bridge/src/mount_ns.rs}\\
\noindent 行范围：12-16\\[0.4em]

\begin{rustcode}
...
pub struct PerTaskMountNsRegistry {
    namespaces: BTreeMap<task::TaskId, MountNamespace>,
    owners: BTreeMap<task::TaskId, task::TaskId>,
    ref_counts: BTreeMap<task::TaskId, usize>,
}
...
\end{rustcode}

\subsubsection*{序号1059}
\noindent 路径：\code{os/components/wateros-vfs/vfs-impl/impl-fs-bridge/src/mount_ns.rs}\\
\noindent 行范围：28-32\\[0.4em]

\begin{rustcode}
...
    fn ensure_owner(&mut self, task_id: task::TaskId) -> task::TaskId {
        self.owners.entry(task_id).or_insert(task_id);
        self.ref_counts.entry(task_id).or_insert(1);
        self.effective_owner(task_id)
    }
...
\end{rustcode}

\subsubsection*{序号1060}
\noindent 路径：\code{os/components/wateros-vfs/vfs-impl/impl-fs-bridge/src/mount_ns.rs}\\
\noindent 行范围：40-49\\[0.4em]

\begin{rustcode}
...
    fn release_owner(&mut self, task_id: task::TaskId) -> Option<task::TaskId> {
        let owner = self.owners.remove(&task_id)?;
        if let Some(count) = self.ref_counts.get_mut(&owner) {
            *count = count.saturating_sub(1);
            if *count == 0 {
...
\end{rustcode}

\subsubsection*{序号1061}
\noindent 路径：\code{os/components/wateros-vfs/vfs-impl/impl-fs-bridge/src/mount_ns.rs}\\
\noindent 行范围：53-58\\[0.4em]

\begin{rustcode}
...
    pub fn init_task_mount_ns(&mut self, task_id: task::TaskId) {
        let owner = self.ensure_owner(task_id);
        self.namespaces.entry(owner).or_insert_with(|| {
            crate::mount_table::bootstrap_mount_namespace_snapshot()
        });
...
\end{rustcode}

\subsubsection*{序号1062}
\noindent 路径：\code{os/components/wateros-vfs/vfs-impl/impl-fs-bridge/src/mount_ns.rs}\\
\noindent 行范围：62-65\\[0.4em]

\begin{rustcode}
...
// 本方法代码由AI完成
    pub fn namespace_for(&self, task_id: task::TaskId) -> Option<&MountNamespace> {
        let owner = self.effective_owner(task_id);
        self.namespaces.get(&owner)
    }
...
\end{rustcode}

\subsubsection*{序号1063}
\noindent 路径：\code{os/components/wateros-vfs/vfs-impl/impl-fs-bridge/src/mount_ns.rs}\\
\noindent 行范围：69-75\\[0.4em]

\begin{rustcode}
...
    pub fn namespace_for_mut(&mut self, task_id: task::TaskId) -> &mut MountNamespace {
        self.init_task_mount_ns(task_id);
        let owner = self.effective_owner(task_id);
        self.namespaces
            .get_mut(&owner)
...
\end{rustcode}

\subsubsection*{序号1064}
\noindent 路径：\code{os/components/wateros-vfs/vfs-impl/impl-fs-bridge/src/mount_ns.rs}\\
\noindent 行范围：78-88\\[0.4em]

\begin{rustcode}
...
    pub fn drop_task(&mut self, task_id: task::TaskId) {
        let Some(owner) = self.release_owner(task_id) else {
            return;
        };
        if self.ref_counts.get(&owner).copied().unwrap_or(0) == 0 {
...
\end{rustcode}

\subsubsection*{序号1065}
\noindent 路径：\code{os/components/wateros-vfs/vfs-impl/impl-fs-bridge/src/mount_ns.rs}\\
\noindent 行范围：92-102\\[0.4em]

\begin{rustcode}
...
    pub fn copy_mount_ns_from_parent(&mut self, child: task::TaskId, parent: task::TaskId) {
        let parent_ns = self
            .namespace_for(parent)
            .cloned()
            .unwrap_or_else(|| crate::mount_table::bootstrap_mount_namespace_snapshot());
...
\end{rustcode}

\subsubsection*{序号1066}
\noindent 路径：\code{os/components/wateros-vfs/vfs-impl/impl-fs-bridge/src/mount_ns.rs}\\
\noindent 行范围：105-114\\[0.4em]

\begin{rustcode}
...
    pub fn share_mount_ns_from_parent(&mut self, child: task::TaskId, parent: task::TaskId) {
        self.init_task_mount_ns(parent);
        if self.owners.contains_key(&child) {
            self.drop_task(child);
        }
...
\end{rustcode}

\subsubsection*{序号1067}
\noindent 路径：\code{os/components/wateros-vfs/vfs-impl/impl-fs-bridge/src/mount_ns.rs}\\
\noindent 行范围：118-140\\[0.4em]

\begin{rustcode}
...
    pub fn unshare_mount_ns(&mut self, task_id: task::TaskId) {
        let owner = self.effective_owner(task_id);
        let shared = self
            .ref_counts
            .get(&owner)
...
\end{rustcode}

\subsubsection*{序号1068}
\noindent 路径：\code{os/components/wateros-vfs/vfs-impl/impl-fs-bridge/src/mount_table.rs}\\
\noindent 行范围：1-894\\[0.4em]

\begin{rustcode}
//! 辅助卷挂载表（最长前缀路由）；支持 RW、RO、procfs 伪挂载、bind 与传播类型。
//! 本模块代码由AI完成

extern crate alloc;

...
\end{rustcode}

\subsubsection*{序号1069}
\noindent 路径：\code{os/components/wateros-vfs/vfs-impl/impl-fs-bridge/src/mount_table.rs}\\
\noindent 行范围：33-42\\[0.4em]

\begin{rustcode}
...
pub enum MountPropagation {
    /// 挂载事件不向其它挂载点传播（默认）。
    Private,
    /// 挂载/卸载事件向共享组内其它挂载点传播。
    Shared,
...
\end{rustcode}

\subsubsection*{序号1070}
\noindent 路径：\code{os/components/wateros-vfs/vfs-impl/impl-fs-bridge/src/mount_table.rs}\\
\noindent 行范围：46-48\\[0.4em]

\begin{rustcode}
...
// 本方法代码由AI完成
    fn default() -> Self {
        Self::Private
    }
}
...
\end{rustcode}

\subsubsection*{序号1071}
\noindent 路径：\code{os/components/wateros-vfs/vfs-impl/impl-fs-bridge/src/mount_table.rs}\\
\noindent 行范围：53-61\\[0.4em]

\begin{rustcode}
...
struct MountEntry {
    mount_point: String,
    fs: AuxMount,
    identity: MountIdentity,
    readonly: bool,
...
\end{rustcode}

\subsubsection*{序号1072}
\noindent 路径：\code{os/components/wateros-vfs/vfs-impl/impl-fs-bridge/src/mount_table.rs}\\
\noindent 行范围：71-75\\[0.4em]

\begin{rustcode}
...
    fn clone(&self) -> Self {
        Self {
            entries: self.entries.clone(),
        }
    }
...
\end{rustcode}

\subsubsection*{序号1073}
\noindent 路径：\code{os/components/wateros-vfs/vfs-impl/impl-fs-bridge/src/mount_table.rs}\\
\noindent 行范围：79\\[0.4em]

\begin{rustcode}
...

// 本变量代码由AI完成
static DEVICE_IDS: Mutex<Vec<(String, u32)>> = Mutex::new(Vec::new());
// 本变量代码由AI完成
static NEXT_DEVICE_MINOR: AtomicU64 = AtomicU64::new(1);
...
\end{rustcode}

\subsubsection*{序号1074}
\noindent 路径：\code{os/components/wateros-vfs/vfs-impl/impl-fs-bridge/src/mount_table.rs}\\
\noindent 行范围：81\\[0.4em]

\begin{rustcode}
...
static DEVICE_IDS: Mutex<Vec<(String, u32)>> = Mutex::new(Vec::new());
// 本变量代码由AI完成
static NEXT_DEVICE_MINOR: AtomicU64 = AtomicU64::new(1);
// 本变量代码由AI完成
static NEXT_MOUNT_ID: AtomicU64 = AtomicU64::new(2);
...
\end{rustcode}

\subsubsection*{序号1075}
\noindent 路径：\code{os/components/wateros-vfs/vfs-impl/impl-fs-bridge/src/mount_table.rs}\\
\noindent 行范围：83\\[0.4em]

\begin{rustcode}
...
static NEXT_DEVICE_MINOR: AtomicU64 = AtomicU64::new(1);
// 本变量代码由AI完成
static NEXT_MOUNT_ID: AtomicU64 = AtomicU64::new(2);

/// 内核 bring-up / 无当前任务时的挂载表（`procfs` 等在 spawn 用户任务前挂载）。
...
\end{rustcode}

\subsubsection*{序号1076}
\noindent 路径：\code{os/components/wateros-vfs/vfs-impl/impl-fs-bridge/src/mount_table.rs}\\
\noindent 行范围：87\\[0.4em]

\begin{rustcode}
...
/// 内核 bring-up / 无当前任务时的挂载表（`procfs` 等在 spawn 用户任务前挂载）。
// 本变量代码由AI完成
static BOOTSTRAP_MOUNT_NS: Mutex<MountNamespace> = Mutex::new(MountNamespace { entries: Vec::new() });

// 本方法代码由AI完成
...
\end{rustcode}

\subsubsection*{序号1077}
\noindent 路径：\code{os/components/wateros-vfs/vfs-impl/impl-fs-bridge/src/mount_table.rs}\\
\noindent 行范围：90-92\\[0.4em]

\begin{rustcode}
...
// 本方法代码由AI完成
pub(crate) fn bootstrap_mount_namespace_snapshot() -> MountNamespace {
    BOOTSTRAP_MOUNT_NS.lock().clone()
}

...
\end{rustcode}

\subsubsection*{序号1078}
\noindent 路径：\code{os/components/wateros-vfs/vfs-impl/impl-fs-bridge/src/mount_table.rs}\\
\noindent 行范围：95-96\\[0.4em]

\begin{rustcode}
...

// 本变量代码由AI完成
static mut MOUNT_NS_REGISTRY: MaybeUninit<UniprocessorSafeCell<PerTaskMountNsRegistry>> =
    MaybeUninit::uninit();
// 本变量代码由AI完成
...
\end{rustcode}

\subsubsection*{序号1079}
\noindent 路径：\code{os/components/wateros-vfs/vfs-impl/impl-fs-bridge/src/mount_table.rs}\\
\noindent 行范围：98\\[0.4em]

\begin{rustcode}
...
    MaybeUninit::uninit();
// 本变量代码由AI完成
static MOUNT_NS_REGISTRY_READY: AtomicUsize = AtomicUsize::new(0);

// 本方法代码由AI完成
...
\end{rustcode}

\subsubsection*{序号1080}
\noindent 路径：\code{os/components/wateros-vfs/vfs-impl/impl-fs-bridge/src/mount_table.rs}\\
\noindent 行范围：101-109\\[0.4em]

\begin{rustcode}
...
fn registry() -> &'static UniprocessorSafeCell<PerTaskMountNsRegistry> {
    if MOUNT_NS_REGISTRY_READY.load(Ordering::Acquire) == 0 {
        unsafe {
            MOUNT_NS_REGISTRY.write(UniprocessorSafeCell::new(PerTaskMountNsRegistry::new()));
        }
...
\end{rustcode}

\subsubsection*{序号1081}
\noindent 路径：\code{os/components/wateros-vfs/vfs-impl/impl-fs-bridge/src/mount_table.rs}\\
\noindent 行范围：118-122\\[0.4em]

\begin{rustcode}
...
pub fn copy_mount_ns_from_parent(child: task::TaskId, parent: task::TaskId) {
    registry()
        .exclusive_access()
        .copy_mount_ns_from_parent(child, parent);
}
...
\end{rustcode}

\subsubsection*{序号1082}
\noindent 路径：\code{os/components/wateros-vfs/vfs-impl/impl-fs-bridge/src/mount_table.rs}\\
\noindent 行范围：126-130\\[0.4em]

\begin{rustcode}
...
pub fn share_mount_ns_from_parent(child: task::TaskId, parent: task::TaskId) {
    registry()
        .exclusive_access()
        .share_mount_ns_from_parent(child, parent);
}
...
\end{rustcode}

\subsubsection*{序号1083}
\noindent 路径：\code{os/components/wateros-vfs/vfs-impl/impl-fs-bridge/src/mount_table.rs}\\
\noindent 行范围：143-151\\[0.4em]

\begin{rustcode}
...
fn with_current_namespace<R>(f: impl FnOnce(&mut MountNamespace) -> VfsResult<R>) -> VfsResult<R> {
    if let Some(task_id) = task::current_task_id() {
        let mut reg = registry().exclusive_access();
        f(reg.namespace_for_mut(task_id))
    } else {
...
\end{rustcode}

\subsubsection*{序号1084}
\noindent 路径：\code{os/components/wateros-vfs/vfs-impl/impl-fs-bridge/src/mount_table.rs}\\
\noindent 行范围：154-156\\[0.4em]

\begin{rustcode}
...
// 本方法代码由AI完成
fn namespace_for_route(task_id: task::TaskId) -> Option<MountNamespace> {
    registry().exclusive_access().namespace_for(task_id).cloned()
}

...
\end{rustcode}

\subsubsection*{序号1085}
\noindent 路径：\code{os/components/wateros-vfs/vfs-impl/impl-fs-bridge/src/mount_table.rs}\\
\noindent 行范围：160-177\\[0.4em]

\begin{rustcode}
...
fn mount_namespace_snapshot() -> MountNamespace {
    if let Some(task_id) = task::current_task_id() {
        {
            let reg = registry().exclusive_access();
            if let Some(ns) = reg.namespace_for(task_id) {
...
\end{rustcode}

\subsubsection*{序号1086}
\noindent 路径：\code{os/components/wateros-vfs/vfs-impl/impl-fs-bridge/src/mount_table.rs}\\
\noindent 行范围：181-215\\[0.4em]

\begin{rustcode}
...
fn assert_mount_point_directory_in(ns: &MountNamespace, path: &str) -> VfsResult<()> {
    match resolve_material_route(ns, path)? {
        FsRoute::PseudoProc { .. } | FsRoute::PseudoSecurity { .. } => Err(VfsError::NotAFile),
        FsRoute::Root { abs, .. } => {
            let meta = super::root_rw()?
...
\end{rustcode}

\subsubsection*{序号1087}
\noindent 路径：\code{os/components/wateros-vfs/vfs-impl/impl-fs-bridge/src/mount_table.rs}\\
\noindent 行范围：225-234\\[0.4em]

\begin{rustcode}
...
fn device_minor_for(key: &str) -> u32 {
    let mut devices = DEVICE_IDS.lock();
    if let Some((_, minor)) = devices.iter().find(|(known, _)| known == key) {
        return *minor;
    }
...
\end{rustcode}

\subsubsection*{序号1088}
\noindent 路径：\code{os/components/wateros-vfs/vfs-impl/impl-fs-bridge/src/mount_table.rs}\\
\noindent 行范围：237-243\\[0.4em]

\begin{rustcode}
...
fn new_mount_identity(device_key: &str) -> MountIdentity {
    MountIdentity {
        device_major: 0,
        device_minor: device_minor_for(device_key),
        mount_id: NEXT_MOUNT_ID.fetch_add(1, Ordering::Relaxed),
...
\end{rustcode}

\subsubsection*{序号1089}
\noindent 路径：\code{os/components/wateros-vfs/vfs-impl/impl-fs-bridge/src/mount_table.rs}\\
\noindent 行范围：246-254\\[0.4em]

\begin{rustcode}
...
pub(crate) fn root_identity() -> MountIdentity {
    let device = fs::rootfs::active_impl::current_root_device_path()
        .unwrap_or_else(|| String::from("/dev/root"));
    MountIdentity {
        device_major: 0,
...
\end{rustcode}

\subsubsection*{序号1090}
\noindent 路径：\code{os/components/wateros-vfs/vfs-impl/impl-fs-bridge/src/mount_table.rs}\\
\noindent 行范围：271-283\\[0.4em]

\begin{rustcode}
...
fn rel_under_mount(full: &str, mount_point: &str) -> String {
    if full == mount_point {
        return String::from("/");
    }
    let rest = full.strip_prefix(mount_point).unwrap_or(full);
...
\end{rustcode}

\subsubsection*{序号1091}
\noindent 路径：\code{os/components/wateros-vfs/vfs-impl/impl-fs-bridge/src/mount_table.rs}\\
\noindent 行范围：286-292\\[0.4em]

\begin{rustcode}
...
fn join_mount_path(mount: &str, rel: &str) -> String {
    if rel == "/" {
        return String::from(mount);
    }
    let mount = mount.trim_end_matches('/');
...
\end{rustcode}

\subsubsection*{序号1092}
\noindent 路径：\code{os/components/wateros-vfs/vfs-impl/impl-fs-bridge/src/mount_table.rs}\\
\noindent 行范围：296-306\\[0.4em]

\begin{rustcode}
...
    fn clone_mount(&self) -> Self {
        match self {
            Self::Rw(fs) => Self::Rw(fs.clone()),
            Self::Ro(fs) => Self::Ro(fs.clone()),
            Self::PseudoProc => Self::PseudoProc,
...
\end{rustcode}

\subsubsection*{序号1093}
\noindent 路径：\code{os/components/wateros-vfs/vfs-impl/impl-fs-bridge/src/mount_table.rs}\\
\noindent 行范围：311-313\\[0.4em]

\begin{rustcode}
...
// 本方法代码由AI完成
    pub fn new() -> Self {
        Self { entries: Vec::new() }
    }

...
\end{rustcode}

\subsubsection*{序号1094}
\noindent 路径：\code{os/components/wateros-vfs/vfs-impl/impl-fs-bridge/src/mount_table.rs}\\
\noindent 行范围：316-331\\[0.4em]

\begin{rustcode}
...
    fn longest_match(&self, abs: &str) -> Option<(usize, &MountEntry)> {
        let mut best: Option<(usize, &MountEntry)> = None;
        for ent in self.entries.iter() {
            let mp = ent.mount_point.as_str();
            let matches =
...
\end{rustcode}

\subsubsection*{序号1095}
\noindent 路径：\code{os/components/wateros-vfs/vfs-impl/impl-fs-bridge/src/mount_table.rs}\\
\noindent 行范围：334-336\\[0.4em]

\begin{rustcode}
...
// 本方法代码由AI完成
    fn exact_mount(&self, mount_point: &str) -> Option<&MountEntry> {
        self.entries.iter().find(|e| e.mount_point == mount_point)
    }

...
\end{rustcode}

\subsubsection*{序号1096}
\noindent 路径：\code{os/components/wateros-vfs/vfs-impl/impl-fs-bridge/src/mount_table.rs}\\
\noindent 行范围：339-343\\[0.4em]

\begin{rustcode}
...
    fn exact_mount_mut(&mut self, mount_point: &str) -> Option<&mut MountEntry> {
        self.entries
            .iter_mut()
            .find(|e| e.mount_point == mount_point)
    }
...
\end{rustcode}

\subsubsection*{序号1097}
\noindent 路径：\code{os/components/wateros-vfs/vfs-impl/impl-fs-bridge/src/mount_table.rs}\\
\noindent 行范围：346-348\\[0.4em]

\begin{rustcode}
...
// 本方法代码由AI完成
    fn is_mount_point(&self, abs: &str) -> bool {
        self.exact_mount(abs).is_some()
    }

...
\end{rustcode}

\subsubsection*{序号1098}
\noindent 路径：\code{os/components/wateros-vfs/vfs-impl/impl-fs-bridge/src/mount_table.rs}\\
\noindent 行范围：351-355\\[0.4em]

\begin{rustcode}
...
    fn is_under_mount(&self, abs: &str, mount_point: &str) -> bool {
        abs == mount_point
            || abs.starts_with(mount_point)
                && abs.as_bytes().get(mount_point.len()) == Some(&b'/')
    }
...
\end{rustcode}

\subsubsection*{序号1099}
\noindent 路径：\code{os/components/wateros-vfs/vfs-impl/impl-fs-bridge/src/mount_table.rs}\\
\noindent 行范围：358-362\\[0.4em]

\begin{rustcode}
...
    fn propagation_at(&self, abs: &str) -> MountPropagation {
        self.longest_match(abs)
            .map(|(_, ent)| ent.propagation)
            .unwrap_or(MountPropagation::Private)
    }
...
\end{rustcode}

\subsubsection*{序号1100}
\noindent 路径：\code{os/components/wateros-vfs/vfs-impl/impl-fs-bridge/src/mount_table.rs}\\
\noindent 行范围：365-370\\[0.4em]

\begin{rustcode}
...
    fn bind_forbidden(&self, source: &str) -> bool {
        matches!(
            self.propagation_at(source),
            MountPropagation::Unbindable
        )
...
\end{rustcode}

\subsubsection*{序号1101}
\noindent 路径：\code{os/components/wateros-vfs/vfs-impl/impl-fs-bridge/src/mount_table.rs}\\
\noindent 行范围：374\\[0.4em]

\begin{rustcode}
...

// 本变量代码由AI完成
const BIND_CHAIN_LIMIT: usize = 32;

// 本方法代码由AI完成
...
\end{rustcode}

\subsubsection*{序号1102}
\noindent 路径：\code{os/components/wateros-vfs/vfs-impl/impl-fs-bridge/src/mount_table.rs}\\
\noindent 行范围：377-423\\[0.4em]

\begin{rustcode}
...
fn resolve_material_route(ns: &MountNamespace, abs: &str) -> VfsResult<FsRoute> {
    let abs = String::from(normalize_absolute_path(abs)?.as_str());
    let mut current = abs;
    for _ in 0..BIND_CHAIN_LIMIT {
        let Some((_, ent)) = ns.longest_match(current.as_str()) else {
...
\end{rustcode}

\subsubsection*{序号1103}
\noindent 路径：\code{os/components/wateros-vfs/vfs-impl/impl-fs-bridge/src/mount_table.rs}\\
\noindent 行范围：426-434\\[0.4em]

\begin{rustcode}
...
fn bump_mount_generation_after_cache_flush() {
    if let Err(e) = super::reset_file_page_cache() {
        log::warn!(
            "[vfs-bridge] page cache flush before mount_gen bump failed: {:?}",
            e
...
\end{rustcode}

\subsubsection*{序号1104}
\noindent 路径：\code{os/components/wateros-vfs/vfs-impl/impl-fs-bridge/src/mount_table.rs}\\
\noindent 行范围：437-465\\[0.4em]

\begin{rustcode}
...
fn mount_aux_common(
    ns: &mut MountNamespace,
    mount_point: &str,
    fs: AuxMount,
    device_key: &str,
...
\end{rustcode}

\subsubsection*{序号1105}
\noindent 路径：\code{os/components/wateros-vfs/vfs-impl/impl-fs-bridge/src/mount_table.rs}\\
\noindent 行范围：468-471\\[0.4em]

\begin{rustcode}
...
// 本方法代码由AI完成
pub(crate) fn resolve_route(path: &str) -> VfsResult<FsRoute> {
    let ns = mount_namespace_snapshot();
    resolve_material_route(&ns, path)
}
...
\end{rustcode}

\subsubsection*{序号1106}
\noindent 路径：\code{os/components/wateros-vfs/vfs-impl/impl-fs-bridge/src/mount_table.rs}\\
\noindent 行范围：475-483\\[0.4em]

\begin{rustcode}
...
pub fn assert_path_writable(path: &str) -> VfsResult<()> {
    match resolve_route(path)? {
        FsRoute::AuxRw { readonly: true, .. }
        | FsRoute::AuxRo { .. }
        | FsRoute::PseudoProc { .. }
...
\end{rustcode}

\subsubsection*{序号1107}
\noindent 路径：\code{os/components/wateros-vfs/vfs-impl/impl-fs-bridge/src/mount_table.rs}\\
\noindent 行范围：486-492\\[0.4em]

\begin{rustcode}
...
pub(crate) fn mount_aux_at_rw(mount_point: &str, fs: SharedRwFs, device_key: &str) -> VfsResult<()> {
    with_current_namespace(|ns| {
        mount_aux_common(ns, mount_point, AuxMount::Rw(fs), device_key, "ext4", false)
    })?;
    bump_mount_generation_after_cache_flush();
...
\end{rustcode}

\subsubsection*{序号1108}
\noindent 路径：\code{os/components/wateros-vfs/vfs-impl/impl-fs-bridge/src/mount_table.rs}\\
\noindent 行范围：495-501\\[0.4em]

\begin{rustcode}
...
pub(crate) fn mount_aux_at_ro(mount_point: &str, fs: SharedFs, device_key: &str) -> VfsResult<()> {
    with_current_namespace(|ns| {
        mount_aux_common(ns, mount_point, AuxMount::Ro(fs), device_key, "ext4", true)
    })?;
    bump_mount_generation_after_cache_flush();
...
\end{rustcode}

\subsubsection*{序号1109}
\noindent 路径：\code{os/components/wateros-vfs/vfs-impl/impl-fs-bridge/src/mount_table.rs}\\
\noindent 行范围：504-512\\[0.4em]

\begin{rustcode}
...
pub(crate) fn mount_tmpfs_at(mount_point: &str) -> VfsResult<()> {
    let fs: SharedRwFs =
        Arc::new(Mutex::new(LocalRwFs::new(Box::new(super::tmpfs::TmpFs::new()))));
    with_current_namespace(|ns| {
        mount_aux_common(ns, mount_point, AuxMount::Rw(fs), "tmpfs", "tmpfs", false)
...
\end{rustcode}

\subsubsection*{序号1110}
\noindent 路径：\code{os/components/wateros-vfs/vfs-impl/impl-fs-bridge/src/mount_table.rs}\\
\noindent 行范围：515-524\\[0.4em]

\begin{rustcode}
...
pub(crate) fn mount_cgroup_at(mount_point: &str, v2: bool, options: &str) -> VfsResult<()> {
    let tmp = super::tmpfs::TmpFs::new_cgroup(v2, options).map_err(super::map_fs_err)?;
    let fs: SharedRwFs = Arc::new(Mutex::new(LocalRwFs::new(Box::new(tmp))));
    let fstype = if v2 { "cgroup2" } else { "cgroup" };
    with_current_namespace(|ns| {
...
\end{rustcode}

\subsubsection*{序号1111}
\noindent 路径：\code{os/components/wateros-vfs/vfs-impl/impl-fs-bridge/src/mount_table.rs}\\
\noindent 行范围：527-540\\[0.4em]

\begin{rustcode}
...
pub(crate) fn mount_securityfs_at(mount_point: &str) -> VfsResult<()> {
    with_current_namespace(|ns| {
        mount_aux_common(
            ns,
            mount_point,
...
\end{rustcode}

\subsubsection*{序号1112}
\noindent 路径：\code{os/components/wateros-vfs/vfs-impl/impl-fs-bridge/src/mount_table.rs}\\
\noindent 行范围：543-573\\[0.4em]

\begin{rustcode}
...
pub(crate) fn mount_bind_at(source: &str, target: &str, recursive: bool) -> VfsResult<()> {
    let source = String::from(normalize_absolute_path(source)?.as_str());
    let target = String::from(normalize_absolute_path(target)?.as_str());
    if target == "/" {
        return Err(VfsError::InvalidPath);
...
\end{rustcode}

\subsubsection*{序号1113}
\noindent 路径：\code{os/components/wateros-vfs/vfs-impl/impl-fs-bridge/src/mount_table.rs}\\
\noindent 行范围：576-614\\[0.4em]

\begin{rustcode}
...
fn recursive_bind(ns: &mut MountNamespace, source_root: &str, target_root: &str) -> VfsResult<()> {
    let mut pairs: Vec<(String, String)> = Vec::new();
    pairs.push((String::from(source_root), String::from(target_root)));
    let mut mounts: Vec<String> = ns
        .entries
...
\end{rustcode}

\subsubsection*{序号1114}
\noindent 路径：\code{os/components/wateros-vfs/vfs-impl/impl-fs-bridge/src/mount_table.rs}\\
\noindent 行范围：617-649\\[0.4em]

\begin{rustcode}
...
pub(crate) fn set_mount_propagation(
    mount_point: &str,
    propagation: MountPropagation,
    recursive: bool,
) -> VfsResult<()> {
...
\end{rustcode}

\subsubsection*{序号1115}
\noindent 路径：\code{os/components/wateros-vfs/vfs-impl/impl-fs-bridge/src/mount_table.rs}\\
\noindent 行范围：652-691\\[0.4em]

\begin{rustcode}
...
pub(crate) fn move_mount_at(source: &str, target: &str) -> VfsResult<()> {
    let source = String::from(normalize_absolute_path(source)?.as_str());
    let target = String::from(normalize_absolute_path(target)?.as_str());
    if source == "/" || target == "/" {
        return Err(VfsError::InvalidPath);
...
\end{rustcode}

\subsubsection*{序号1116}
\noindent 路径：\code{os/components/wateros-vfs/vfs-impl/impl-fs-bridge/src/mount_table.rs}\\
\noindent 行范围：694-711\\[0.4em]

\begin{rustcode}
...
pub fn mount_statfs_magic(abs: &str) -> Option<isize> {
    let Ok(abs) = normalize_absolute_path(abs) else {
        return None;
    };
    let task_id = task::current_task_id()?;
...
\end{rustcode}

\subsubsection*{序号1117}
\noindent 路径：\code{os/components/wateros-vfs/vfs-impl/impl-fs-bridge/src/mount_table.rs}\\
\noindent 行范围：714-723\\[0.4em]

\begin{rustcode}
...
fn mount_statfs_magic_for_path(ns: &MountNamespace, abs: &str) -> Option<isize> {
    let route = resolve_material_route(ns, abs).ok()?;
    let path = match route {
        FsRoute::Root { abs, .. } => abs,
        FsRoute::AuxRw { .. } | FsRoute::AuxRo { .. } => return Some(0xEF53),
...
\end{rustcode}

\subsubsection*{序号1118}
\noindent 路径：\code{os/components/wateros-vfs/vfs-impl/impl-fs-bridge/src/mount_table.rs}\\
\noindent 行范围：726-740\\[0.4em]

\begin{rustcode}
...
pub(crate) fn remount_aux_readonly(mount_point: &str) -> VfsResult<()> {
    let mp = String::from(normalize_absolute_path(mount_point)?.as_str());
    with_current_namespace(|ns| {
        let ent = ns
            .exact_mount_mut(mp.as_str())
...
\end{rustcode}

\subsubsection*{序号1119}
\noindent 路径：\code{os/components/wateros-vfs/vfs-impl/impl-fs-bridge/src/mount_table.rs}\\
\noindent 行范围：743-756\\[0.4em]

\begin{rustcode}
...
pub fn mount_aux_proc_at(mount_point: &str) -> VfsResult<()> {
    with_current_namespace(|ns| {
        mount_aux_common(
            ns,
            mount_point,
...
\end{rustcode}

\subsubsection*{序号1120}
\noindent 路径：\code{os/components/wateros-vfs/vfs-impl/impl-fs-bridge/src/mount_table.rs}\\
\noindent 行范围：759-775\\[0.4em]

\begin{rustcode}
...
pub fn is_proc_mounted_at(mount_point: &str) -> bool {
    let Ok(mp) = normalize_absolute_path(mount_point) else {
        return false;
    };
    if let Some(task_id) = task::current_task_id() {
...
\end{rustcode}

\subsubsection*{序号1121}
\noindent 路径：\code{os/components/wateros-vfs/vfs-impl/impl-fs-bridge/src/mount_table.rs}\\
\noindent 行范围：778-780\\[0.4em]

\begin{rustcode}
...
// 本方法代码由AI完成
fn fstype_for(entry: &MountEntry) -> &'static str {
    entry.fstype
}

...
\end{rustcode}

\subsubsection*{序号1122}
\noindent 路径：\code{os/components/wateros-vfs/vfs-impl/impl-fs-bridge/src/mount_table.rs}\\
\noindent 行范围：783-790\\[0.4em]

\begin{rustcode}
...
fn device_for(entry: &MountEntry) -> String {
    match entry.fs {
        AuxMount::PseudoProc => String::from("proc"),
        AuxMount::PseudoSecurity => String::from("securityfs"),
        AuxMount::Bind { ref source } => source.clone(),
...
\end{rustcode}

\subsubsection*{序号1123}
\noindent 路径：\code{os/components/wateros-vfs/vfs-impl/impl-fs-bridge/src/mount_table.rs}\\
\noindent 行范围：793-795\\[0.4em]

\begin{rustcode}
...
// 本方法代码由AI完成
fn root_mount_device() -> String {
    fs::devfs::active_impl::default_root_block_path().unwrap_or_else(|| String::from("/dev/root"))
}

...
\end{rustcode}

\subsubsection*{序号1124}
\noindent 路径：\code{os/components/wateros-vfs/vfs-impl/impl-fs-bridge/src/mount_table.rs}\\
\noindent 行范围：798-828\\[0.4em]

\begin{rustcode}
...
pub fn list_proc_mount_lines() -> Vec<ProcMountLine> {
    let mut out = Vec::new();
    if fs::rootfs::active_impl::root_rw_fs().is_some() {
        out.push(ProcMountLine {
            device: root_mount_device(),
...
\end{rustcode}

\subsubsection*{序号1125}
\noindent 路径：\code{os/components/wateros-vfs/vfs-impl/impl-fs-bridge/src/mount_table.rs}\\
\noindent 行范围：831-848\\[0.4em]

\begin{rustcode}
...
pub(crate) fn unmount_aux_at(mount_point: &str, detach: bool) -> VfsResult<()> {
    let mp = String::from(normalize_absolute_path(mount_point)?.as_str());
    if mp == "/" {
        return Err(VfsError::InvalidPath);
    }
...
\end{rustcode}

\subsubsection*{序号1126}
\noindent 路径：\code{os/components/wateros-vfs/vfs-impl/impl-fs-bridge/src/mount_table.rs}\\
\noindent 行范围：851-894\\[0.4em]

\begin{rustcode}
...
pub fn mount_table_self_test() -> VfsResult<()> {
    let dev_a = new_mount_identity("/dev/__identity_test__");
    let dev_b = new_mount_identity("/dev/__identity_test__");
    assert_eq!(dev_a.device_major, dev_b.device_major);
    assert_eq!(dev_a.device_minor, dev_b.device_minor);
...
\end{rustcode}

\subsubsection*{序号1127}
\noindent 路径：\code{os/components/wateros-vfs/vfs-impl/impl-fs-bridge/src/paged_handle.rs}\\
\noindent 行范围：1-611\\[0.4em]

\begin{rustcode}
//! 大文件句柄：经全局页缓存 Direct 读写，不在 `open` 时整文件载入 RAM。
//!
//! ## Lock ordering（与 [`impl_page_cache`] 一致）
//!
//! 1. `page_cache.files`（极短）
...
\end{rustcode}

\subsubsection*{序号1128}
\noindent 路径：\code{os/components/wateros-vfs/vfs-impl/impl-fs-bridge/src/paged_handle.rs}\\
\noindent 行范围：30\\[0.4em]

\begin{rustcode}
...
/// detached 模式下单文件内核堆缓冲上限。
// 本变量代码由AI完成
const DETACHED_DATA_MAX : usize = 16 * 1024 * 1024;

// 本方法代码由AI完成
...
\end{rustcode}

\subsubsection*{序号1129}
\noindent 路径：\code{os/components/wateros-vfs/vfs-impl/impl-fs-bridge/src/paged_handle.rs}\\
\noindent 行范围：33-41\\[0.4em]

\begin{rustcode}
...
fn check_detached_len(len : usize) -> VfsResult<()> {
    if len > DETACHED_DATA_MAX {
        log::warn!("[paged_handle] detached buffer cap exceeded len={} max={}",
                   len,
                   DETACHED_DATA_MAX);
...
\end{rustcode}

\subsubsection*{序号1130}
\noindent 路径：\code{os/components/wateros-vfs/vfs-impl/impl-fs-bridge/src/paged_handle.rs}\\
\noindent 行范围：44-52\\[0.4em]

\begin{rustcode}
...
fn grow_detached_data(buf : &mut Vec<u8>, new_len : usize) -> VfsResult<()> {
    check_detached_len(new_len)?;
    if buf.len() < new_len {
        buf.try_reserve_exact(new_len - buf.len())
           .map_err(|_| VfsError::Io)?;
...
\end{rustcode}

\subsubsection*{序号1131}
\noindent 路径：\code{os/components/wateros-vfs/vfs-impl/impl-fs-bridge/src/paged_handle.rs}\\
\noindent 行范围：61-63\\[0.4em]

\begin{rustcode}
...
// 本方法代码由AI完成
    fn read_range(&self, path : &str, offset : u64, buf : &mut [u8]) -> Result<usize, VfsError> {
        FsBridge.read_range(path, offset, buf)
    }

...
\end{rustcode}

\subsubsection*{序号1132}
\noindent 路径：\code{os/components/wateros-vfs/vfs-impl/impl-fs-bridge/src/paged_handle.rs}\\
\noindent 行范围：66-93\\[0.4em]

\begin{rustcode}
...
    fn write_range(&mut self, path : &str, offset : u64, data : &[u8]) -> Result<usize, VfsError> {
        match resolve_route(path)? {
            FsRoute::Root { abs, .. } => {
                let n = normalize_absolute_path(abs.as_str())?;
                let rw = root_rw()?;
...
\end{rustcode}

\subsubsection*{序号1133}
\noindent 路径：\code{os/components/wateros-vfs/vfs-impl/impl-fs-bridge/src/paged_handle.rs}\\
\noindent 行范围：98-110\\[0.4em]

\begin{rustcode}
...
pub struct PagedFileHandle {
    path : String,
    offset : u64,
    meta : VfsMetadata,
    writable : bool,
...
\end{rustcode}

\subsubsection*{序号1134}
\noindent 路径：\code{os/components/wateros-vfs/vfs-impl/impl-fs-bridge/src/paged_handle.rs}\\
\noindent 行范围：114-129\\[0.4em]

\begin{rustcode}
...
    fn clone(&self) -> Self {
        if self.open_ref_held {
            global_cache(self.mount_gen).acquire_open_ref(self.path.as_str());
        }
        Self { path : self.path.clone(),
...
\end{rustcode}

\subsubsection*{序号1135}
\noindent 路径：\code{os/components/wateros-vfs/vfs-impl/impl-fs-bridge/src/paged_handle.rs}\\
\noindent 行范围：134-195\\[0.4em]

\begin{rustcode}
...
    pub(crate) fn open(bridge : &FsBridge,
                       path : String,
                       flags : VfsOpenFlags,
                       mut meta : VfsMetadata)
                       -> VfsResult<Self> {
...
\end{rustcode}

\subsubsection*{序号1136}
\noindent 路径：\code{os/components/wateros-vfs/vfs-impl/impl-fs-bridge/src/paged_handle.rs}\\
\noindent 行范围：167\\[0.4em]

\begin{rustcode}
...

// 本变量代码由AI完成
        const O_WRONLY : u32 = 1;
// 本变量代码由AI完成
        const O_RDWR : u32 = 2;
...
\end{rustcode}

\subsubsection*{序号1137}
\noindent 路径：\code{os/components/wateros-vfs/vfs-impl/impl-fs-bridge/src/paged_handle.rs}\\
\noindent 行范围：169\\[0.4em]

\begin{rustcode}
...
        const O_WRONLY : u32 = 1;
// 本变量代码由AI完成
        const O_RDWR : u32 = 2;
// 本变量代码由AI完成
        const O_APPEND : u32 = 0o2000;
...
\end{rustcode}

\subsubsection*{序号1138}
\noindent 路径：\code{os/components/wateros-vfs/vfs-impl/impl-fs-bridge/src/paged_handle.rs}\\
\noindent 行范围：171\\[0.4em]

\begin{rustcode}
...
        const O_RDWR : u32 = 2;
// 本变量代码由AI完成
        const O_APPEND : u32 = 0o2000;
        let accmode = if want_write && flags.contains(VfsOpenFlags::READ) {
            O_RDWR
...
\end{rustcode}

\subsubsection*{序号1139}
\noindent 路径：\code{os/components/wateros-vfs/vfs-impl/impl-fs-bridge/src/paged_handle.rs}\\
\noindent 行范围：198-203\\[0.4em]

\begin{rustcode}
...
    fn release_open_ref_if_held(&mut self) {
        if self.open_ref_held {
            global_cache(self.mount_gen).release_open_ref(self.path.as_str());
            self.open_ref_held = false;
        }
...
\end{rustcode}

\subsubsection*{序号1140}
\noindent 路径：\code{os/components/wateros-vfs/vfs-impl/impl-fs-bridge/src/paged_handle.rs}\\
\noindent 行范围：206-227\\[0.4em]

\begin{rustcode}
...
    fn sync_dirty(&mut self) -> VfsResult<()> {
        if !self.writable || self.detached {
            log::trace!("[vfs-flush] skip path={} writable={} detached={}",
                        self.path,
                        self.writable,
...
\end{rustcode}

\subsubsection*{序号1141}
\noindent 路径：\code{os/components/wateros-vfs/vfs-impl/impl-fs-bridge/src/paged_handle.rs}\\
\noindent 行范围：230-232\\[0.4em]

\begin{rustcode}
...
// 本方法代码由AI完成
    fn current_size(&self) -> u64 {
        global_cache(self.mount_gen).logical_size(self.path.as_str(), self.on_disk_size)
    }

...
\end{rustcode}

\subsubsection*{序号1142}
\noindent 路径：\code{os/components/wateros-vfs/vfs-impl/impl-fs-bridge/src/paged_handle.rs}\\
\noindent 行范围：235-249\\[0.4em]

\begin{rustcode}
...
    fn read_detached_at(&self, offset : u64, buf : &mut [u8]) -> VfsResult<usize> {
        let start = usize::try_from(offset).map_err(|_| VfsError::Io)?;
        if start >=
           self.detached_data
               .len()
...
\end{rustcode}

\subsubsection*{序号1143}
\noindent 路径：\code{os/components/wateros-vfs/vfs-impl/impl-fs-bridge/src/paged_handle.rs}\\
\noindent 行范围：252-276\\[0.4em]

\begin{rustcode}
...
    fn account_detached_write(&mut self,
                              offset : u64,
                              buf : &[u8],
                              advance_offset : bool)
                              -> VfsResult<usize> {
...
\end{rustcode}

\subsubsection*{序号1144}
\noindent 路径：\code{os/components/wateros-vfs/vfs-impl/impl-fs-bridge/src/paged_handle.rs}\\
\noindent 行范围：281-318\\[0.4em]

\begin{rustcode}
...
    fn read(&mut self, buf : &mut [u8]) -> VfsResult<usize> {
        if buf.is_empty() {
            return Ok(0);
        }
        let size = self.current_size();
...
\end{rustcode}

\subsubsection*{序号1145}
\noindent 路径：\code{os/components/wateros-vfs/vfs-impl/impl-fs-bridge/src/paged_handle.rs}\\
\noindent 行范围：321-358\\[0.4em]

\begin{rustcode}
...
    fn write(&mut self, buf : &[u8]) -> VfsResult<usize> {
        if !self.writable {
            return Err(VfsError::Unsupported);
        }
// 本变量代码由AI完成
...
\end{rustcode}

\subsubsection*{序号1146}
\noindent 路径：\code{os/components/wateros-vfs/vfs-impl/impl-fs-bridge/src/paged_handle.rs}\\
\noindent 行范围：326\\[0.4em]

\begin{rustcode}
...
        }
// 本变量代码由AI完成
        const O_APPEND : u32 = 0o2000;
        if self.status_flags & O_APPEND != 0 {
            self.offset = self.current_size();
...
\end{rustcode}

\subsubsection*{序号1147}
\noindent 路径：\code{os/components/wateros-vfs/vfs-impl/impl-fs-bridge/src/paged_handle.rs}\\
\noindent 行范围：361-381\\[0.4em]

\begin{rustcode}
...
    fn read_at(&mut self, offset : u64, buf : &mut [u8]) -> VfsResult<usize> {
        if buf.is_empty() {
            return Ok(0);
        }
        let size = self.current_size();
...
\end{rustcode}

\subsubsection*{序号1148}
\noindent 路径：\code{os/components/wateros-vfs/vfs-impl/impl-fs-bridge/src/paged_handle.rs}\\
\noindent 行范围：384-413\\[0.4em]

\begin{rustcode}
...
    fn write_at(&mut self, offset : u64, buf : &[u8]) -> VfsResult<usize> {
        if !self.writable {
            return Err(VfsError::Unsupported);
        }
        if buf.is_empty() {
...
\end{rustcode}

\subsubsection*{序号1149}
\noindent 路径：\code{os/components/wateros-vfs/vfs-impl/impl-fs-bridge/src/paged_handle.rs}\\
\noindent 行范围：416-424\\[0.4em]

\begin{rustcode}
...
    fn close(&mut self) -> VfsResult<()> {
        let sync_err = self.sync_dirty();
        self.release_open_ref_if_held();
        if let Err(e) = &sync_err {
            log::warn!("[paged_handle] close sync_dirty failed path={:?} err={e:?}",
...
\end{rustcode}

\subsubsection*{序号1150}
\noindent 路径：\code{os/components/wateros-vfs/vfs-impl/impl-fs-bridge/src/paged_handle.rs}\\
\noindent 行范围：427-431\\[0.4em]

\begin{rustcode}
...
    fn metadata(&self) -> VfsResult<VfsMetadata> {
        let mut m = self.meta.clone();
        m.size = self.current_size();
        Ok(m)
    }
...
\end{rustcode}

\subsubsection*{序号1151}
\noindent 路径：\code{os/components/wateros-vfs/vfs-impl/impl-fs-bridge/src/paged_handle.rs}\\
\noindent 行范围：434-436\\[0.4em]

\begin{rustcode}
...
// 本方法代码由AI完成
    fn backing_path(&self) -> Option<&str> {
        Some(self.path.as_str())
    }

...
\end{rustcode}

\subsubsection*{序号1152}
\noindent 路径：\code{os/components/wateros-vfs/vfs-impl/impl-fs-bridge/src/paged_handle.rs}\\
\noindent 行范围：439-471\\[0.4em]

\begin{rustcode}
...
    fn seek(&mut self, offset : i64, whence : VfsSeekWhence) -> VfsResult<u64> {
        let size = self.current_size();
        let new_off = match whence {
            VfsSeekWhence::Set => {
                if offset < 0 {
...
\end{rustcode}

\subsubsection*{序号1153}
\noindent 路径：\code{os/components/wateros-vfs/vfs-impl/impl-fs-bridge/src/paged_handle.rs}\\
\noindent 行范围：474-476\\[0.4em]

\begin{rustcode}
...
// 本方法代码由AI完成
    fn flush(&mut self) -> VfsResult<()> {
        self.sync_dirty()
    }

...
\end{rustcode}

\subsubsection*{序号1154}
\noindent 路径：\code{os/components/wateros-vfs/vfs-impl/impl-fs-bridge/src/paged_handle.rs}\\
\noindent 行范围：479-528\\[0.4em]

\begin{rustcode}
...
    fn truncate(&mut self, len : u64) -> VfsResult<()> {
        if !self.writable {
            return Err(VfsError::Unsupported);
        }
        if len > 0 {
...
\end{rustcode}

\subsubsection*{序号1155}
\noindent 路径：\code{os/components/wateros-vfs/vfs-impl/impl-fs-bridge/src/paged_handle.rs}\\
\noindent 行范围：531-533\\[0.4em]

\begin{rustcode}
...
// 本方法代码由AI完成
    fn open_status_flags(&self) -> u32 {
        self.status_flags
    }

...
\end{rustcode}

\subsubsection*{序号1156}
\noindent 路径：\code{os/components/wateros-vfs/vfs-impl/impl-fs-bridge/src/paged_handle.rs}\\
\noindent 行范围：536-538\\[0.4em]

\begin{rustcode}
...
// 本方法代码由AI完成
    fn open_accmode(&self) -> u32 {
        self.accmode
    }

...
\end{rustcode}

\subsubsection*{序号1157}
\noindent 路径：\code{os/components/wateros-vfs/vfs-impl/impl-fs-bridge/src/paged_handle.rs}\\
\noindent 行范围：541-548\\[0.4em]

\begin{rustcode}
...
    fn set_open_status_flags(&mut self, flags : u32) -> VfsResult<()> {
// 本变量代码由AI完成
        const O_APPEND : u32 = 0o2000;
// 本变量代码由AI完成
        const O_NONBLOCK : u32 = 0o4000;
...
\end{rustcode}

\subsubsection*{序号1158}
\noindent 路径：\code{os/components/wateros-vfs/vfs-impl/impl-fs-bridge/src/paged_handle.rs}\\
\noindent 行范围：543\\[0.4em]

\begin{rustcode}
...
    fn set_open_status_flags(&mut self, flags : u32) -> VfsResult<()> {
// 本变量代码由AI完成
        const O_APPEND : u32 = 0o2000;
// 本变量代码由AI完成
        const O_NONBLOCK : u32 = 0o4000;
...
\end{rustcode}

\subsubsection*{序号1159}
\noindent 路径：\code{os/components/wateros-vfs/vfs-impl/impl-fs-bridge/src/paged_handle.rs}\\
\noindent 行范围：545\\[0.4em]

\begin{rustcode}
...
        const O_APPEND : u32 = 0o2000;
// 本变量代码由AI完成
        const O_NONBLOCK : u32 = 0o4000;
        self.status_flags = flags & (O_APPEND | O_NONBLOCK);
        Ok(())
...
\end{rustcode}

\subsubsection*{序号1160}
\noindent 路径：\code{os/components/wateros-vfs/vfs-impl/impl-fs-bridge/src/paged_handle.rs}\\
\noindent 行范围：551\\[0.4em]

\begin{rustcode}
...

// 本方法代码由AI完成
    fn duplicate(&self) -> VfsResult<Box<dyn VfsIoHandle>> { Ok(Box::new(self.clone())) }

// 本方法代码由AI完成
...
\end{rustcode}

\subsubsection*{序号1161}
\noindent 路径：\code{os/components/wateros-vfs/vfs-impl/impl-fs-bridge/src/paged_handle.rs}\\
\noindent 行范围：554-567\\[0.4em]

\begin{rustcode}
...
    fn poll_revents(&mut self, events : i16) -> VfsResult<i16> {
// 本变量代码由AI完成
        const POLLIN : i16 = 0x001;
// 本变量代码由AI完成
        const POLLOUT : i16 = 0x004;
...
\end{rustcode}

\subsubsection*{序号1162}
\noindent 路径：\code{os/components/wateros-vfs/vfs-impl/impl-fs-bridge/src/paged_handle.rs}\\
\noindent 行范围：556\\[0.4em]

\begin{rustcode}
...
    fn poll_revents(&mut self, events : i16) -> VfsResult<i16> {
// 本变量代码由AI完成
        const POLLIN : i16 = 0x001;
// 本变量代码由AI完成
        const POLLOUT : i16 = 0x004;
...
\end{rustcode}

\subsubsection*{序号1163}
\noindent 路径：\code{os/components/wateros-vfs/vfs-impl/impl-fs-bridge/src/paged_handle.rs}\\
\noindent 行范围：558\\[0.4em]

\begin{rustcode}
...
        const POLLIN : i16 = 0x001;
// 本变量代码由AI完成
        const POLLOUT : i16 = 0x004;
        let mut revents = 0i16;
        if events & POLLIN != 0 {
...
\end{rustcode}

\subsubsection*{序号1164}
\noindent 路径：\code{os/components/wateros-vfs/vfs-impl/impl-fs-bridge/src/paged_handle.rs}\\
\noindent 行范围：573-611\\[0.4em]

\begin{rustcode}
...
pub(crate) fn open_file(bridge : &FsBridge,
                        path : String,
                        flags : VfsOpenFlags)
                        -> VfsResult<Box<dyn VfsIoHandle>> {
    let want_write = flags.contains(VfsOpenFlags::WRITE);
...
\end{rustcode}

\subsubsection*{序号1165}
\noindent 路径：\code{os/components/wateros-vfs/vfs-impl/impl-fs-bridge/src/proc_handle.rs}\\
\noindent 行范围：1-209\\[0.4em]

\begin{rustcode}
//! procfs 伪挂载打开句柄。
//! 本模块代码由AI完成

extern crate alloc;

...
\end{rustcode}

\subsubsection*{序号1166}
\noindent 路径：\code{os/components/wateros-vfs/vfs-impl/impl-fs-bridge/src/proc_handle.rs}\\
\noindent 行范围：19-21\\[0.4em]

\begin{rustcode}
...
// 本方法代码由AI完成
fn proc_view() -> &'static impl ProcFsView {
    fs::procfs::active_impl::view()
}

...
\end{rustcode}

\subsubsection*{序号1167}
\noindent 路径：\code{os/components/wateros-vfs/vfs-impl/impl-fs-bridge/src/proc_handle.rs}\\
\noindent 行范围：26-34\\[0.4em]

\begin{rustcode}
...
pub struct ProcDirectoryHandle {
    /// 挂载内相对路径（不含 `/proc` 前缀）。
    rel: String,
    /// 用户可见绝对路径。
    abs: String,
...
\end{rustcode}

\subsubsection*{序号1168}
\noindent 路径：\code{os/components/wateros-vfs/vfs-impl/impl-fs-bridge/src/proc_handle.rs}\\
\noindent 行范围：39-44\\[0.4em]

\begin{rustcode}
...
pub struct ProcFileHandle {
    rel: String,
    meta: VfsMetadata,
    data: Vec<u8>,
    offset: u64,
...
\end{rustcode}

\subsubsection*{序号1169}
\noindent 路径：\code{os/components/wateros-vfs/vfs-impl/impl-fs-bridge/src/proc_handle.rs}\\
\noindent 行范围：47-84\\[0.4em]

\begin{rustcode}
...
pub fn open_proc(
    rel: String,
    abs: String,
    flags: VfsOpenFlags,
    identity: MountIdentity,
...
\end{rustcode}

\subsubsection*{序号1170}
\noindent 路径：\code{os/components/wateros-vfs/vfs-impl/impl-fs-bridge/src/proc_handle.rs}\\
\noindent 行范围：88-106\\[0.4em]

\begin{rustcode}
...
    fn load_dirents(&mut self) -> VfsResult<()> {
        if self.dirents.is_some() {
            return Ok(());
        }
        let entries = proc_view()
...
\end{rustcode}

\subsubsection*{序号1171}
\noindent 路径：\code{os/components/wateros-vfs/vfs-impl/impl-fs-bridge/src/proc_handle.rs}\\
\noindent 行范围：111-113\\[0.4em]

\begin{rustcode}
...
// 本方法代码由AI完成
    fn metadata(&self) -> VfsResult<VfsMetadata> {
        Ok(self.meta.clone())
    }

...
\end{rustcode}

\subsubsection*{序号1172}
\noindent 路径：\code{os/components/wateros-vfs/vfs-impl/impl-fs-bridge/src/proc_handle.rs}\\
\noindent 行范围：116-118\\[0.4em]

\begin{rustcode}
...
// 本方法代码由AI完成
    fn directory_path(&self) -> Option<&str> {
        Some(self.abs.as_str())
    }

...
\end{rustcode}

\subsubsection*{序号1173}
\noindent 路径：\code{os/components/wateros-vfs/vfs-impl/impl-fs-bridge/src/proc_handle.rs}\\
\noindent 行范围：121-123\\[0.4em]

\begin{rustcode}
...
// 本方法代码由AI完成
    fn read(&mut self, _buf: &mut [u8]) -> VfsResult<usize> {
        Err(VfsError::NotAFile)
    }

...
\end{rustcode}

\subsubsection*{序号1174}
\noindent 路径：\code{os/components/wateros-vfs/vfs-impl/impl-fs-bridge/src/proc_handle.rs}\\
\noindent 行范围：126-128\\[0.4em]

\begin{rustcode}
...
// 本方法代码由AI完成
    fn write(&mut self, _buf: &[u8]) -> VfsResult<usize> {
        Err(VfsError::ReadOnlyFs)
    }

...
\end{rustcode}

\subsubsection*{序号1175}
\noindent 路径：\code{os/components/wateros-vfs/vfs-impl/impl-fs-bridge/src/proc_handle.rs}\\
\noindent 行范围：131-133\\[0.4em]

\begin{rustcode}
...
// 本方法代码由AI完成
    fn duplicate(&self) -> VfsResult<Box<dyn VfsIoHandle>> {
        Ok(Box::new(self.clone()))
    }

...
\end{rustcode}

\subsubsection*{序号1176}
\noindent 路径：\code{os/components/wateros-vfs/vfs-impl/impl-fs-bridge/src/proc_handle.rs}\\
\noindent 行范围：136-154\\[0.4em]

\begin{rustcode}
...
    fn fill_getdents64(&mut self, buf: &mut [u8]) -> VfsResult<usize> {
        self.load_dirents()?;
        let entries = self.dirents.as_ref().expect("load_dirents");
        let mut out = 0usize;
        let mut off = 0usize;
...
\end{rustcode}

\subsubsection*{序号1177}
\noindent 路径：\code{os/components/wateros-vfs/vfs-impl/impl-fs-bridge/src/proc_handle.rs}\\
\noindent 行范围：159-161\\[0.4em]

\begin{rustcode}
...
// 本方法代码由AI完成
    fn metadata(&self) -> VfsResult<VfsMetadata> {
        Ok(self.meta.clone())
    }

...
\end{rustcode}

\subsubsection*{序号1178}
\noindent 路径：\code{os/components/wateros-vfs/vfs-impl/impl-fs-bridge/src/proc_handle.rs}\\
\noindent 行范围：164-173\\[0.4em]

\begin{rustcode}
...
    fn read(&mut self, buf: &mut [u8]) -> VfsResult<usize> {
        let start = self.offset as usize;
        if start >= self.data.len() {
            return Ok(0);
        }
...
\end{rustcode}

\subsubsection*{序号1179}
\noindent 路径：\code{os/components/wateros-vfs/vfs-impl/impl-fs-bridge/src/proc_handle.rs}\\
\noindent 行范围：176-184\\[0.4em]

\begin{rustcode}
...
    fn read_at(&mut self, offset: u64, buf: &mut [u8]) -> VfsResult<usize> {
        let start = offset as usize;
        if start >= self.data.len() {
            return Ok(0);
        }
...
\end{rustcode}

\subsubsection*{序号1180}
\noindent 路径：\code{os/components/wateros-vfs/vfs-impl/impl-fs-bridge/src/proc_handle.rs}\\
\noindent 行范围：187-189\\[0.4em]

\begin{rustcode}
...
// 本方法代码由AI完成
    fn write(&mut self, _buf: &[u8]) -> VfsResult<usize> {
        Err(VfsError::ReadOnlyFs)
    }

...
\end{rustcode}

\subsubsection*{序号1181}
\noindent 路径：\code{os/components/wateros-vfs/vfs-impl/impl-fs-bridge/src/proc_handle.rs}\\
\noindent 行范围：192-203\\[0.4em]

\begin{rustcode}
...
    fn seek(&mut self, offset: i64, whence: VfsSeekWhence) -> VfsResult<u64> {
        let new_off = match whence {
            VfsSeekWhence::Set => offset.max(0) as u64,
            VfsSeekWhence::Cur => self.offset.saturating_add_signed(offset),
            VfsSeekWhence::End => self
...
\end{rustcode}

\subsubsection*{序号1182}
\noindent 路径：\code{os/components/wateros-vfs/vfs-impl/impl-fs-bridge/src/proc_handle.rs}\\
\noindent 行范围：206-208\\[0.4em]

\begin{rustcode}
...
// 本方法代码由AI完成
    fn duplicate(&self) -> VfsResult<Box<dyn VfsIoHandle>> {
        Ok(Box::new(self.clone()))
    }
}
\end{rustcode}

\subsubsection*{序号1183}
\noindent 路径：\code{os/components/wateros-vfs/vfs-impl/impl-fs-bridge/src/tmpfs.rs}\\
\noindent 行范围：1-702\\[0.4em]

\begin{rustcode}
//! 内存 tmpfs：供 LTP `needs_rofs` 等测例挂载可重载为只读的临时卷。
//! 本模块代码由AI完成

extern crate alloc;

...
\end{rustcode}

\subsubsection*{序号1184}
\noindent 路径：\code{os/components/wateros-vfs/vfs-impl/impl-fs-bridge/src/tmpfs.rs}\\
\noindent 行范围：52-67\\[0.4em]

\begin{rustcode}
...
    pub(crate) fn new() -> Self {
        Self {
            root: TmpNode::Dir {
                children: BTreeMap::new(),
                mode: 0o40755,
...
\end{rustcode}

\subsubsection*{序号1185}
\noindent 路径：\code{os/components/wateros-vfs/vfs-impl/impl-fs-bridge/src/tmpfs.rs}\\
\noindent 行范围：70-78\\[0.4em]

\begin{rustcode}
...
    pub(crate) fn new_cgroup(v2: bool, options: &str) -> FsResult<Self> {
        let mut fs = Self::new();
        fs.cgroup_v2 = Some(v2);
        if !v2 && !options.is_empty() {
            fs.cgroup_v1_options = Some(String::from(options));
...
\end{rustcode}

\subsubsection*{序号1186}
\noindent 路径：\code{os/components/wateros-vfs/vfs-impl/impl-fs-bridge/src/tmpfs.rs}\\
\noindent 行范围：81-86\\[0.4em]

\begin{rustcode}
...
    fn v1_has_controller(&self, name: &str) -> bool {
        self.cgroup_v1_options.as_ref().is_some_and(|opts| {
            opts.split(',')
                .any(|ctrl| ctrl.trim() == name)
        })
...
\end{rustcode}

\subsubsection*{序号1187}
\noindent 路径：\code{os/components/wateros-vfs/vfs-impl/impl-fs-bridge/src/tmpfs.rs}\\
\noindent 行范围：89-96\\[0.4em]

\begin{rustcode}
...
    fn write_control_file(&mut self, dir_path: &str, name: &str, data: &[u8]) -> FsResult<()> {
        let path = if dir_path == "/" || dir_path.is_empty() {
            alloc::format!("/{name}")
        } else {
            alloc::format!("{}/{}", dir_path.trim_end_matches('/'), name)
...
\end{rustcode}

\subsubsection*{序号1188}
\noindent 路径：\code{os/components/wateros-vfs/vfs-impl/impl-fs-bridge/src/tmpfs.rs}\\
\noindent 行范围：99-103\\[0.4em]

\begin{rustcode}
...
    fn alloc_inode(&mut self) -> u64 {
        let n = self.next_inode;
        self.next_inode += 1;
        n
    }
...
\end{rustcode}

\subsubsection*{序号1189}
\noindent 路径：\code{os/components/wateros-vfs/vfs-impl/impl-fs-bridge/src/tmpfs.rs}\\
\noindent 行范围：106-113\\[0.4em]

\begin{rustcode}
...
    fn split_path(path: &str) -> FsResult<Vec<&str>> {
        let p = path.trim();
        let p = p.strip_prefix('/').unwrap_or(p);
        if p.is_empty() {
            return Ok(Vec::new());
...
\end{rustcode}

\subsubsection*{序号1190}
\noindent 路径：\code{os/components/wateros-vfs/vfs-impl/impl-fs-bridge/src/tmpfs.rs}\\
\noindent 行范围：116-125\\[0.4em]

\begin{rustcode}
...
    fn dir_mut<'a>(root: &'a mut TmpNode, parts: &[&str]) -> FsResult<&'a mut TmpNode> {
        let mut node = root;
        for &part in parts {
            let TmpNode::Dir { children, .. } = node else {
                return Err(FsError::NotFound);
...
\end{rustcode}

\subsubsection*{序号1191}
\noindent 路径：\code{os/components/wateros-vfs/vfs-impl/impl-fs-bridge/src/tmpfs.rs}\\
\noindent 行范围：128-137\\[0.4em]

\begin{rustcode}
...
    fn dir_ref<'a>(root: &'a TmpNode, parts: &[&str]) -> FsResult<&'a TmpNode> {
        let mut node = root;
        for &part in parts {
            let TmpNode::Dir { children, .. } = node else {
                return Err(FsError::NotFound);
...
\end{rustcode}

\subsubsection*{序号1192}
\noindent 路径：\code{os/components/wateros-vfs/vfs-impl/impl-fs-bridge/src/tmpfs.rs}\\
\noindent 行范围：140-153\\[0.4em]

\begin{rustcode}
...
    fn parent_dir_mut<'a>(
        root: &'a mut TmpNode,
        parts: &[&'a str],
    ) -> FsResult<(&'a mut BTreeMap<String, TmpNode>, &'a str)> {
        if parts.is_empty() {
...
\end{rustcode}

\subsubsection*{序号1193}
\noindent 路径：\code{os/components/wateros-vfs/vfs-impl/impl-fs-bridge/src/tmpfs.rs}\\
\noindent 行范围：156-206\\[0.4em]

\begin{rustcode}
...
    fn meta_of(node: &TmpNode) -> FsMetadata {
        match node {
            TmpNode::File {
                data,
                mode,
...
\end{rustcode}

\subsubsection*{序号1194}
\noindent 路径：\code{os/components/wateros-vfs/vfs-impl/impl-fs-bridge/src/tmpfs.rs}\\
\noindent 行范围：209-222\\[0.4em]

\begin{rustcode}
...
    fn chown_node(node: &mut TmpNode, uid: Option<u32>, gid: Option<u32>) {
        match node {
            TmpNode::File { uid: u, gid: g, .. }
            | TmpNode::Dir { uid: u, gid: g, .. }
            | TmpNode::Symlink { uid: u, gid: g, .. } => {
...
\end{rustcode}

\subsubsection*{序号1195}
\noindent 路径：\code{os/components/wateros-vfs/vfs-impl/impl-fs-bridge/src/tmpfs.rs}\\
\noindent 行范围：225-231\\[0.4em]

\begin{rustcode}
...
    fn xattrs_mut(node: &mut TmpNode) -> &mut BTreeMap<String, Vec<u8>> {
        match node {
            TmpNode::File { xattrs, .. }
            | TmpNode::Dir { xattrs, .. }
            | TmpNode::Symlink { xattrs, .. } => xattrs,
...
\end{rustcode}

\subsubsection*{序号1196}
\noindent 路径：\code{os/components/wateros-vfs/vfs-impl/impl-fs-bridge/src/tmpfs.rs}\\
\noindent 行范围：234-240\\[0.4em]

\begin{rustcode}
...
    fn xattrs(node: &TmpNode) -> &BTreeMap<String, Vec<u8>> {
        match node {
            TmpNode::File { xattrs, .. }
            | TmpNode::Dir { xattrs, .. }
            | TmpNode::Symlink { xattrs, .. } => xattrs,
...
\end{rustcode}

\subsubsection*{序号1197}
\noindent 路径：\code{os/components/wateros-vfs/vfs-impl/impl-fs-bridge/src/tmpfs.rs}\\
\noindent 行范围：243-250\\[0.4em]

\begin{rustcode}
...
    fn copy_xattr_list(names: &[String]) -> Vec<u8> {
        let mut out = Vec::new();
        for name in names {
            out.extend_from_slice(name.as_bytes());
            out.push(0);
...
\end{rustcode}

\subsubsection*{序号1198}
\noindent 路径：\code{os/components/wateros-vfs/vfs-impl/impl-fs-bridge/src/tmpfs.rs}\\
\noindent 行范围：253-256\\[0.4em]

\begin{rustcode}
...
// 本方法代码由AI完成
    fn remove_leaf(root: &mut TmpNode, parts: &[&str]) -> FsResult<TmpNode> {
        let (children, name) = Self::parent_dir_mut(root, parts)?;
        children.remove(name).ok_or(FsError::NotFound)
    }
...
\end{rustcode}

\subsubsection*{序号1199}
\noindent 路径：\code{os/components/wateros-vfs/vfs-impl/impl-fs-bridge/src/tmpfs.rs}\\
\noindent 行范围：259-266\\[0.4em]

\begin{rustcode}
...
    fn insert_leaf(root: &mut TmpNode, parts: &[&str], node: TmpNode) -> FsResult<()> {
        let (children, name) = Self::parent_dir_mut(root, parts)?;
        if children.contains_key(name) {
            return Err(FsError::Exists);
        }
...
\end{rustcode}

\subsubsection*{序号1200}
\noindent 路径：\code{os/components/wateros-vfs/vfs-impl/impl-fs-bridge/src/tmpfs.rs}\\
\noindent 行范围：269-271\\[0.4em]

\begin{rustcode}
...
// 本方法代码由AI完成
    fn cgroup_v1_base_files() -> &'static [(&'static str, &'static [u8])] {
        &[("tasks", b""), ("notify_on_release", b"0\n")]
    }

...
\end{rustcode}

\subsubsection*{序号1201}
\noindent 路径：\code{os/components/wateros-vfs/vfs-impl/impl-fs-bridge/src/tmpfs.rs}\\
\noindent 行范围：274-286\\[0.4em]

\begin{rustcode}
...
    fn cgroup_v1_cpuset_root_files() -> &'static [(&'static str, &'static [u8])] {
        &[
            ("cgroup.clone_children", b"0\n"),
            ("cpuset.sched_load_balance", b"1\n"),
            ("cpuset.cpu_exclusive", b"0\n"),
...
\end{rustcode}

\subsubsection*{序号1202}
\noindent 路径：\code{os/components/wateros-vfs/vfs-impl/impl-fs-bridge/src/tmpfs.rs}\\
\noindent 行范围：289-295\\[0.4em]

\begin{rustcode}
...
    fn cgroup_v1_cpuset_child_files() -> &'static [(&'static str, &'static [u8])] {
        &[
            ("cgroup.clone_children", b"0\n"),
            ("cpuset.cpus", b""),
            ("cpuset.mems", b""),
...
\end{rustcode}

\subsubsection*{序号1203}
\noindent 路径：\code{os/components/wateros-vfs/vfs-impl/impl-fs-bridge/src/tmpfs.rs}\\
\noindent 行范围：298-307\\[0.4em]

\begin{rustcode}
...
    fn cgroup_control_files(v2: bool) -> &'static [(&'static str, &'static [u8])] {
        if v2 {
            &[("cgroup.procs", b""),
              ("cgroup.subtree_control", b""),
              ("cgroup.controllers", b"memory cpu cpuset io pids freezer\n"),
...
\end{rustcode}

\subsubsection*{序号1204}
\noindent 路径：\code{os/components/wateros-vfs/vfs-impl/impl-fs-bridge/src/tmpfs.rs}\\
\noindent 行范围：310-327\\[0.4em]

\begin{rustcode}
...
    fn seed_v1_cgroup_controls(&mut self, dir_path: &str) -> FsResult<()> {
        for (name, data) in Self::cgroup_v1_base_files() {
            self.write_control_file(dir_path, name, data)?;
        }
        if !self.v1_has_controller("cpuset") {
...
\end{rustcode}

\subsubsection*{序号1205}
\noindent 路径：\code{os/components/wateros-vfs/vfs-impl/impl-fs-bridge/src/tmpfs.rs}\\
\noindent 行范围：330-342\\[0.4em]

\begin{rustcode}
...
    fn seed_cgroup_controls(&mut self, dir_path: &str) -> FsResult<()> {
        let Some(v2) = self.cgroup_v2 else {
            return Ok(());
        };
        if v2 {
...
\end{rustcode}

\subsubsection*{序号1206}
\noindent 路径：\code{os/components/wateros-vfs/vfs-impl/impl-fs-bridge/src/tmpfs.rs}\\
\noindent 行范围：345-358\\[0.4em]

\begin{rustcode}
...
    fn reject_cpuset_tasks_write_if_needed(&self, _path: &str, _data: &[u8]) -> FsResult<()> {
        // 放行 cgroup cpuset 的 tasks 写入（不再因 cpus/mems 未配置而拒绝 attach）。
        //
        // 语义上，real Linux 在子 cpuset 未配置 cpus/mems 时拒绝 attach（write 返回
        // ENOSPC）。但本系统的 /bin/sh 是 busybox ash，其内建 echo 在“重定向输出 write
...
\end{rustcode}

\subsubsection*{序号1207}
\noindent 路径：\code{os/components/wateros-vfs/vfs-impl/impl-fs-bridge/src/tmpfs.rs}\\
\noindent 行范围：361-375\\[0.4em]

\begin{rustcode}
...
    fn is_cgroup_control_name(v2: Option<bool>, name: &str) -> bool {
        match v2 {
            Some(true) => Self::cgroup_control_files(true)
                .iter()
                .any(|(n, _)| *n == name),
...
\end{rustcode}

\subsubsection*{序号1208}
\noindent 路径：\code{os/components/wateros-vfs/vfs-impl/impl-fs-bridge/src/tmpfs.rs}\\
\noindent 行范围：380-383\\[0.4em]

\begin{rustcode}
...
// 本方法代码由AI完成
    fn mount_rw(&mut self, _device: SharedBlockDevice) -> FsResult<()> {
        self.mounted = true;
        Ok(())
    }
...
\end{rustcode}

\subsubsection*{序号1209}
\noindent 路径：\code{os/components/wateros-vfs/vfs-impl/impl-fs-bridge/src/tmpfs.rs}\\
\noindent 行范围：386-388\\[0.4em]

\begin{rustcode}
...
// 本方法代码由AI完成
    fn is_mounted(&self) -> bool {
        self.mounted
    }

...
\end{rustcode}

\subsubsection*{序号1210}
\noindent 路径：\code{os/components/wateros-vfs/vfs-impl/impl-fs-bridge/src/tmpfs.rs}\\
\noindent 行范围：391-393\\[0.4em]

\begin{rustcode}
...
// 本方法代码由AI完成
    fn write_regular_file_at_root(&mut self, name: &str, data: &[u8]) -> FsResult<()> {
        self.write_regular_file(alloc::format!("/{name}").as_str(), data)
    }

...
\end{rustcode}

\subsubsection*{序号1211}
\noindent 路径：\code{os/components/wateros-vfs/vfs-impl/impl-fs-bridge/src/tmpfs.rs}\\
\noindent 行范围：396-413\\[0.4em]

\begin{rustcode}
...
    fn write_regular_file(&mut self, path: &str, data: &[u8]) -> FsResult<()> {
        self.reject_cpuset_tasks_write_if_needed(path, data)?;
        let parts = Self::split_path(path)?;
        if parts.is_empty() {
            return Err(FsError::InvalidPath);
...
\end{rustcode}

\subsubsection*{序号1212}
\noindent 路径：\code{os/components/wateros-vfs/vfs-impl/impl-fs-bridge/src/tmpfs.rs}\\
\noindent 行范围：416-425\\[0.4em]

\begin{rustcode}
...
    fn unlink(&mut self, path: &str) -> FsResult<()> {
        let parts = Self::split_path(path)?;
        let (children, name) = Self::parent_dir_mut(&mut self.root, &parts)?;
        let node = children.get(name).ok_or(FsError::NotFound)?;
        if matches!(node, TmpNode::Dir { .. }) {
...
\end{rustcode}

\subsubsection*{序号1213}
\noindent 路径：\code{os/components/wateros-vfs/vfs-impl/impl-fs-bridge/src/tmpfs.rs}\\
\noindent 行范围：428-443\\[0.4em]

\begin{rustcode}
...
    fn symlink(&mut self, link_path: &str, target: &str) -> FsResult<()> {
        let parts = Self::split_path(link_path)?;
        if parts.is_empty() {
            return Err(FsError::InvalidPath);
        }
...
\end{rustcode}

\subsubsection*{序号1214}
\noindent 路径：\code{os/components/wateros-vfs/vfs-impl/impl-fs-bridge/src/tmpfs.rs}\\
\noindent 行范围：446-463\\[0.4em]

\begin{rustcode}
...
    fn rmdir(&mut self, path: &str) -> FsResult<()> {
        let parts = Self::split_path(path)?;
        let cgroup_v2 = self.cgroup_v2;
        let (children, name) = Self::parent_dir_mut(&mut self.root, &parts)?;
        let node = children.get(name).ok_or(FsError::NotFound)?;
...
\end{rustcode}

\subsubsection*{序号1215}
\noindent 路径：\code{os/components/wateros-vfs/vfs-impl/impl-fs-bridge/src/tmpfs.rs}\\
\noindent 行范围：466-485\\[0.4em]

\begin{rustcode}
...
    fn write_range(&mut self, path: &str, offset: u64, data: &[u8]) -> FsResult<usize> {
        if data.is_empty() {
            return Ok(0);
        }
        self.reject_cpuset_tasks_write_if_needed(path, data)?;
...
\end{rustcode}

\subsubsection*{序号1216}
\noindent 路径：\code{os/components/wateros-vfs/vfs-impl/impl-fs-bridge/src/tmpfs.rs}\\
\noindent 行范围：488-497\\[0.4em]

\begin{rustcode}
...
    fn truncate(&mut self, path: &str, len: u64) -> FsResult<()> {
        let parts = Self::split_path(path)?;
        let node = Self::dir_mut(&mut self.root, &parts)?;
        let TmpNode::File { data, .. } = node else {
            return Err(FsError::NotAFile);
...
\end{rustcode}

\subsubsection*{序号1217}
\noindent 路径：\code{os/components/wateros-vfs/vfs-impl/impl-fs-bridge/src/tmpfs.rs}\\
\noindent 行范围：500-516\\[0.4em]

\begin{rustcode}
...
    fn mkdir(&mut self, path: &str, mode: u32) -> FsResult<()> {
        let parts = Self::split_path(path)?;
        if parts.is_empty() {
            return Err(FsError::Exists);
        }
...
\end{rustcode}

\subsubsection*{序号1218}
\noindent 路径：\code{os/components/wateros-vfs/vfs-impl/impl-fs-bridge/src/tmpfs.rs}\\
\noindent 行范围：519-534\\[0.4em]

\begin{rustcode}
...
    fn chmod(&mut self, path: &str, mode: u32) -> FsResult<()> {
        let parts = Self::split_path(path)?;
        let node = Self::dir_mut(&mut self.root, &parts)?;
        match node {
            TmpNode::File { mode: m, .. } => {
...
\end{rustcode}

\subsubsection*{序号1219}
\noindent 路径：\code{os/components/wateros-vfs/vfs-impl/impl-fs-bridge/src/tmpfs.rs}\\
\noindent 行范围：537-545\\[0.4em]

\begin{rustcode}
...
    fn chown(&mut self, path: &str, uid: Option<u32>, gid: Option<u32>) -> FsResult<()> {
        if uid.is_none() && gid.is_none() {
            return self.metadata(path).map(|_| ());
        }
        let parts = Self::split_path(path)?;
...
\end{rustcode}

\subsubsection*{序号1220}
\noindent 路径：\code{os/components/wateros-vfs/vfs-impl/impl-fs-bridge/src/tmpfs.rs}\\
\noindent 行范围：548-556\\[0.4em]

\begin{rustcode}
...
    fn setxattr(&mut self, path: &str, name: &str, value: &[u8]) -> FsResult<()> {
        if name.is_empty() || name.contains('\0') {
            return Err(FsError::InvalidPath);
        }
        let parts = Self::split_path(path)?;
...
\end{rustcode}

\subsubsection*{序号1221}
\noindent 路径：\code{os/components/wateros-vfs/vfs-impl/impl-fs-bridge/src/tmpfs.rs}\\
\noindent 行范围：559-576\\[0.4em]

\begin{rustcode}
...
    fn getxattr(&self, path: &str, name: &str, buf: &mut [u8]) -> FsResult<usize> {
        if name.is_empty() || name.contains('\0') {
            return Err(FsError::InvalidPath);
        }
        let parts = Self::split_path(path)?;
...
\end{rustcode}

\subsubsection*{序号1222}
\noindent 路径：\code{os/components/wateros-vfs/vfs-impl/impl-fs-bridge/src/tmpfs.rs}\\
\noindent 行范围：579-593\\[0.4em]

\begin{rustcode}
...
    fn listxattr(&self, path: &str, buf: &mut [u8]) -> FsResult<usize> {
        let parts = Self::split_path(path)?;
        let node = Self::dir_ref(&self.root, &parts)?;
        let mut names: alloc::vec::Vec<String> = Self::xattrs(node).keys().cloned().collect();
        names.sort();
...
\end{rustcode}

\subsubsection*{序号1223}
\noindent 路径：\code{os/components/wateros-vfs/vfs-impl/impl-fs-bridge/src/tmpfs.rs}\\
\noindent 行范围：596-606\\[0.4em]

\begin{rustcode}
...
    fn removexattr(&mut self, path: &str, name: &str) -> FsResult<()> {
        if name.is_empty() || name.contains('\0') {
            return Err(FsError::InvalidPath);
        }
        let parts = Self::split_path(path)?;
...
\end{rustcode}

\subsubsection*{序号1224}
\noindent 路径：\code{os/components/wateros-vfs/vfs-impl/impl-fs-bridge/src/tmpfs.rs}\\
\noindent 行范围：609-618\\[0.4em]

\begin{rustcode}
...
    fn rename(&mut self, old_path: &str, new_path: &str) -> FsResult<()> {
        let old_parts = Self::split_path(old_path)?;
        let new_parts = Self::split_path(new_path)?;
        let node = Self::remove_leaf(&mut self.root, &old_parts)?;
        if Self::dir_ref(&self.root, &new_parts).is_ok() {
...
\end{rustcode}

\subsubsection*{序号1225}
\noindent 路径：\code{os/components/wateros-vfs/vfs-impl/impl-fs-bridge/src/tmpfs.rs}\\
\noindent 行范围：621-627\\[0.4em]

\begin{rustcode}
...
    fn exists(&self, path: &str) -> FsResult<bool> {
        let parts = Self::split_path(path)?;
        if parts.is_empty() {
            return Ok(true);
        }
...
\end{rustcode}

\subsubsection*{序号1226}
\noindent 路径：\code{os/components/wateros-vfs/vfs-impl/impl-fs-bridge/src/tmpfs.rs}\\
\noindent 行范围：630-636\\[0.4em]

\begin{rustcode}
...
    fn metadata(&self, path: &str) -> FsResult<FsMetadata> {
        let parts = Self::split_path(path)?;
        if parts.is_empty() {
            return Ok(Self::meta_of(&self.root));
        }
...
\end{rustcode}

\subsubsection*{序号1227}
\noindent 路径：\code{os/components/wateros-vfs/vfs-impl/impl-fs-bridge/src/tmpfs.rs}\\
\noindent 行范围：639-646\\[0.4em]

\begin{rustcode}
...
    fn read(&self, path: &str) -> FsResult<Vec<u8>> {
        let parts = Self::split_path(path)?;
        let node = Self::dir_ref(&self.root, &parts)?;
        let TmpNode::File { data, .. } = node else {
            return Err(FsError::NotAFile);
...
\end{rustcode}

\subsubsection*{序号1228}
\noindent 路径：\code{os/components/wateros-vfs/vfs-impl/impl-fs-bridge/src/tmpfs.rs}\\
\noindent 行范围：649-656\\[0.4em]

\begin{rustcode}
...
    fn read_symlink(&self, path: &str) -> FsResult<Vec<u8>> {
        let parts = Self::split_path(path)?;
        let node = Self::dir_ref(&self.root, &parts)?;
        let TmpNode::Symlink { target, .. } = node else {
            return Err(FsError::NotAFile);
...
\end{rustcode}

\subsubsection*{序号1229}
\noindent 路径：\code{os/components/wateros-vfs/vfs-impl/impl-fs-bridge/src/tmpfs.rs}\\
\noindent 行范围：659-675\\[0.4em]

\begin{rustcode}
...
    fn read_range(&self, path: &str, offset: u64, buf: &mut [u8]) -> FsResult<usize> {
        if buf.is_empty() {
            return Ok(0);
        }
        let parts = Self::split_path(path)?;
...
\end{rustcode}

\subsubsection*{序号1230}
\noindent 路径：\code{os/components/wateros-vfs/vfs-impl/impl-fs-bridge/src/tmpfs.rs}\\
\noindent 行范围：678-701\\[0.4em]

\begin{rustcode}
...
    fn read_dir(&self, path: &str) -> FsResult<Vec<FsDirEntry>> {
        let parts = Self::split_path(path)?;
        let node = if parts.is_empty() {
            &self.root
        } else {
...
\end{rustcode}

\subsubsection*{序号1231}
\noindent 路径：\code{os/components/wateros-vfs/vfs-impl/impl-page-cache/src/lib.rs}\\
\noindent 行范围：1-1115\\[0.4em]

\begin{rustcode}
//! 全局共享文件页缓存（Direct 模式）：按 `(mount_gen, path, page_index)` 键 LRU 缓存页帧。
//!
//! 每个文件条目使用 [`Arc<RwLock<FileEntryInner>>`]：允许多个读者并发，写/刷盘独占。
//!
//! ## Lock ordering
...
\end{rustcode}

\subsubsection*{序号1232}
\noindent 路径：\code{os/components/wateros-vfs/vfs-impl/impl-page-cache/src/lib.rs}\\
\noindent 行范围：31\\[0.4em]

\begin{rustcode}
...

// 本变量代码由AI完成
const FLUSH_RUN_MAX_PAGES : usize = 64;

/// 区间读写下层（通常由 `FsBridge` 委托 `ReadOnlyFs` / `ReadWriteFs`）。
...
\end{rustcode}

\subsubsection*{序号1233}
\noindent 路径：\code{os/components/wateros-vfs/vfs-impl/impl-page-cache/src/lib.rs}\\
\noindent 行范围：37\\[0.4em]

\begin{rustcode}
...
    type Error;
// 本方法代码由AI完成
    fn read_range(&self, path : &str, offset : u64, buf : &mut [u8]) -> Result<usize, Self::Error>;
// 本方法代码由AI完成
    fn write_range(&mut self,
...
\end{rustcode}

\subsubsection*{序号1234}
\noindent 路径：\code{os/components/wateros-vfs/vfs-impl/impl-page-cache/src/lib.rs}\\
\noindent 行范围：39-43\\[0.4em]

\begin{rustcode}
...
    fn write_range(&mut self,
                   path : &str,
                   offset : u64,
                   data : &[u8])
                   -> Result<usize, Self::Error>;
...
\end{rustcode}

\subsubsection*{序号1235}
\noindent 路径：\code{os/components/wateros-vfs/vfs-impl/impl-page-cache/src/lib.rs}\\
\noindent 行范围：49-52\\[0.4em]

\begin{rustcode}
...
// 本结构代码由AI完成
pub struct FileCacheKey {
    pub mount_gen : u64,
    pub path : Arc<str>,
}
...
\end{rustcode}

\subsubsection*{序号1236}
\noindent 路径：\code{os/components/wateros-vfs/vfs-impl/impl-page-cache/src/lib.rs}\\
\noindent 行范围：55-59\\[0.4em]

\begin{rustcode}
...
struct PageFrame {
    key : Option<(FileCacheKey, u64)>,
    data : Vec<u8>,
    dirty : bool,
}
...
\end{rustcode}

\subsubsection*{序号1237}
\noindent 路径：\code{os/components/wateros-vfs/vfs-impl/impl-page-cache/src/lib.rs}\\
\noindent 行范围：62-68\\[0.4em]

\begin{rustcode}
...
struct GlobalCacheState {
    capacity : usize,
    frames : Vec<PageFrame>,
    index : BTreeMap<(FileCacheKey, u64), usize>,
    lru : VecDeque<usize>,
...
\end{rustcode}

\subsubsection*{序号1238}
\noindent 路径：\code{os/components/wateros-vfs/vfs-impl/impl-page-cache/src/lib.rs}\\
\noindent 行范围：72-90\\[0.4em]

\begin{rustcode}
...
    fn new() -> Self {
        let cap = FILE_PAGE_CACHE_CAPACITY;
        let mut frames = Vec::new();
        let mut free = Vec::new();
        if cap > 0 {
...
\end{rustcode}

\subsubsection*{序号1239}
\noindent 路径：\code{os/components/wateros-vfs/vfs-impl/impl-page-cache/src/lib.rs}\\
\noindent 行范围：93-102\\[0.4em]

\begin{rustcode}
...
    fn touch_lru(&mut self, idx : usize) {
        if let Some(p) = self.lru
                             .iter()
                             .position(|&x| x == idx)
        {
...
\end{rustcode}

\subsubsection*{序号1240}
\noindent 路径：\code{os/components/wateros-vfs/vfs-impl/impl-page-cache/src/lib.rs}\\
\noindent 行范围：105-124\\[0.4em]

\begin{rustcode}
...
    fn pop_free_or_lru_index(&mut self) -> usize {
        if let Some(idx) = self.free.pop() {
            return idx;
        }
        if let Some(idx) = self.lru.pop_front() {
...
\end{rustcode}

\subsubsection*{序号1241}
\noindent 路径：\code{os/components/wateros-vfs/vfs-impl/impl-page-cache/src/lib.rs}\\
\noindent 行范围：127-148\\[0.4em]

\begin{rustcode}
...
    fn detach_slot_for_reuse(&mut self,
                             idx : usize)
                             -> Option<((FileCacheKey, u64), Vec<u8>)> {
        let old = self.frames[idx].key.take();
        if let Some(ref key) = old {
...
\end{rustcode}

\subsubsection*{序号1242}
\noindent 路径：\code{os/components/wateros-vfs/vfs-impl/impl-page-cache/src/lib.rs}\\
\noindent 行范围：151-159\\[0.4em]

\begin{rustcode}
...
    fn return_detached_slot(&mut self, idx : usize) {
        if self.frames[idx].key.is_none() &&
           !self.free
                .iter()
                .any(|&free_idx| free_idx == idx)
...
\end{rustcode}

\subsubsection*{序号1243}
\noindent 路径：\code{os/components/wateros-vfs/vfs-impl/impl-page-cache/src/lib.rs}\\
\noindent 行范围：164-179\\[0.4em]

\begin{rustcode}
...
    fn clear_in_place(&mut self) {
        for frame in self.frames
                         .iter_mut()
        {
            frame.key = None;
...
\end{rustcode}

\subsubsection*{序号1244}
\noindent 路径：\code{os/components/wateros-vfs/vfs-impl/impl-page-cache/src/lib.rs}\\
\noindent 行范围：184-189\\[0.4em]

\begin{rustcode}
...
struct FileEntryInner {
    logical_size : u64,
    dirty_pages : BTreeMap<u64, ()>,
    /// 上次 read 结束页号；用于顺序读检测后再预取（F-14）。
    last_read_end_page : Option<u64>,
...
\end{rustcode}

\subsubsection*{序号1245}
\noindent 路径：\code{os/components/wateros-vfs/vfs-impl/impl-page-cache/src/lib.rs}\\
\noindent 行范围：193-199\\[0.4em]

\begin{rustcode}
...
pub struct GlobalFilePageCache {
    mount_gen : AtomicU64,
    state : Mutex<GlobalCacheState>,
    files : Mutex<BTreeMap<FileCacheKey, Arc<RwLock<FileEntryInner>>>>,
    /// 仍被 [`PagedFileHandle`] 持有的路径数；归零时在 `close` 后回收该路径缓存条目。
...
\end{rustcode}

\subsubsection*{序号1246}
\noindent 路径：\code{os/components/wateros-vfs/vfs-impl/impl-page-cache/src/lib.rs}\\
\noindent 行范围：204-209\\[0.4em]

\begin{rustcode}
...
    pub fn new(mount_gen : u64) -> Self {
        Self { mount_gen : AtomicU64::new(mount_gen),
               state : Mutex::new(GlobalCacheState::new()),
               files : Mutex::new(BTreeMap::new()),
               open_refs : Mutex::new(BTreeMap::new()) }
...
\end{rustcode}

\subsubsection*{序号1247}
\noindent 路径：\code{os/components/wateros-vfs/vfs-impl/impl-page-cache/src/lib.rs}\\
\noindent 行范围：219-231\\[0.4em]

\begin{rustcode}
...
    pub fn reset_to_gen(&self, new_gen : u64) {
        self.state
            .lock()
            .clear_in_place();
        self.files
...
\end{rustcode}

\subsubsection*{序号1248}
\noindent 路径：\code{os/components/wateros-vfs/vfs-impl/impl-page-cache/src/lib.rs}\\
\noindent 行范围：235-238\\[0.4em]

\begin{rustcode}
...
// 本方法代码由AI完成
    fn file_key(&self, path : &str) -> FileCacheKey {
        FileCacheKey { mount_gen : self.mount_gen(),
                       path : Arc::from(path) }
    }
...
\end{rustcode}

\subsubsection*{序号1249}
\noindent 路径：\code{os/components/wateros-vfs/vfs-impl/impl-page-cache/src/lib.rs}\\
\noindent 行范围：241-256\\[0.4em]

\begin{rustcode}
...
    fn get_file_entry(&self, path : &str, initial_size : u64) -> Arc<RwLock<FileEntryInner>> {
        let key = self.file_key(path);
        let mut files = self.files.lock();
        if let Some(e) = files.get(&key) {
            let entry = e.clone();
...
\end{rustcode}

\subsubsection*{序号1250}
\noindent 路径：\code{os/components/wateros-vfs/vfs-impl/impl-page-cache/src/lib.rs}\\
\noindent 行范围：260-264\\[0.4em]

\begin{rustcode}
...
    pub fn acquire_open_ref(&self, path : &str) {
        let key = self.file_key(path);
        let mut refs = self.open_refs.lock();
        *refs.entry(key).or_insert(0) += 1;
    }
...
\end{rustcode}

\subsubsection*{序号1251}
\noindent 路径：\code{os/components/wateros-vfs/vfs-impl/impl-page-cache/src/lib.rs}\\
\noindent 行范围：268-286\\[0.4em]

\begin{rustcode}
...
    pub fn release_open_ref(&self, path : &str) {
        let key = self.file_key(path);
        let should_purge = {
            let mut refs = self.open_refs.lock();
            let Some(count) = refs.get_mut(&key) else {
...
\end{rustcode}

\subsubsection*{序号1252}
\noindent 路径：\code{os/components/wateros-vfs/vfs-impl/impl-page-cache/src/lib.rs}\\
\noindent 行范围：289-297\\[0.4em]

\begin{rustcode}
...
    fn note_page_written_back(&self, path : &str, page_idx : u64) {
        let key = self.file_key(path);
        let files = self.files.lock();
        if let Some(entry) = files.get(&key) {
            entry.write()
...
\end{rustcode}

\subsubsection*{序号1253}
\noindent 路径：\code{os/components/wateros-vfs/vfs-impl/impl-page-cache/src/lib.rs}\\
\noindent 行范围：300-319\\[0.4em]

\begin{rustcode}
...
    fn writeback_evicted_page<Io, E>(&self,
                                       io : &mut Io,
                                       key : &FileCacheKey,
                                       page_idx : u64,
                                       saved_data : &[u8],
...
\end{rustcode}

\subsubsection*{序号1254}
\noindent 路径：\code{os/components/wateros-vfs/vfs-impl/impl-page-cache/src/lib.rs}\\
\noindent 行范围：322-330\\[0.4em]

\begin{rustcode}
...
    fn logical_size_for_key(&self, key : &FileCacheKey, fallback : u64) -> u64 {
        let files = self.files.lock();
        files.get(key)
             .map(|entry| {
                 entry.read()
...
\end{rustcode}

\subsubsection*{序号1255}
\noindent 路径：\code{os/components/wateros-vfs/vfs-impl/impl-page-cache/src/lib.rs}\\
\noindent 行范围：334-342\\[0.4em]

\begin{rustcode}
...
    fn queue_flush_batch(batches : &mut Vec<(u64, Vec<u8>)>,
                         batch_start : &mut Option<u64>,
                         batch_data : &mut Vec<u8>) {
        if let Some(start_page) = batch_start.take() {
            if !batch_data.is_empty() {
...
\end{rustcode}

\subsubsection*{序号1256}
\noindent 路径：\code{os/components/wateros-vfs/vfs-impl/impl-page-cache/src/lib.rs}\\
\noindent 行范围：345-425\\[0.4em]

\begin{rustcode}
...
    fn flush_dirty_run<Io, E>(&self,
                              io : &mut Io,
                              key : &FileCacheKey,
                              pages : &[u64],
                              logical_size : u64,
...
\end{rustcode}

\subsubsection*{序号1257}
\noindent 路径：\code{os/components/wateros-vfs/vfs-impl/impl-page-cache/src/lib.rs}\\
\noindent 行范围：428-502\\[0.4em]

\begin{rustcode}
...
    fn install_page<Io, E>(&self,
                           io : &mut Io,
                           path : &str,
                           page_idx : u64,
                           file_size : u64,
...
\end{rustcode}

\subsubsection*{序号1258}
\noindent 路径：\code{os/components/wateros-vfs/vfs-impl/impl-page-cache/src/lib.rs}\\
\noindent 行范围：505-562\\[0.4em]

\begin{rustcode}
...
    fn install_zero_page<Io, E>(&self,
                                io : &mut Io,
                                path : &str,
                                page_idx : u64,
                                map_err : fn(Io::Error) -> E)
...
\end{rustcode}

\subsubsection*{序号1259}
\noindent 路径：\code{os/components/wateros-vfs/vfs-impl/impl-page-cache/src/lib.rs}\\
\noindent 行范围：566-619\\[0.4em]

\begin{rustcode}
...
    pub fn read<Io, E>(&self,
                       io : &mut Io,
                       path : &str,
                       file_size : u64,
                       offset : u64,
...
\end{rustcode}

\subsubsection*{序号1260}
\noindent 路径：\code{os/components/wateros-vfs/vfs-impl/impl-page-cache/src/lib.rs}\\
\noindent 行范围：623-681\\[0.4em]

\begin{rustcode}
...
    pub fn write<Io, E>(&self,
                        io : &mut Io,
                        path : &str,
                        file_size : u64,
                        offset : u64,
...
\end{rustcode}

\subsubsection*{序号1261}
\noindent 路径：\code{os/components/wateros-vfs/vfs-impl/impl-page-cache/src/lib.rs}\\
\noindent 行范围：685-694\\[0.4em]

\begin{rustcode}
...
    pub fn logical_size(&self, path : &str, fallback : u64) -> u64 {
        let key = self.file_key(path);
        let files = self.files.lock();
        files.get(&key)
             .map(|e| {
...
\end{rustcode}

\subsubsection*{序号1262}
\noindent 路径：\code{os/components/wateros-vfs/vfs-impl/impl-page-cache/src/lib.rs}\\
\noindent 行范围：698-702\\[0.4em]

\begin{rustcode}
...
    pub fn set_logical_size(&self, path : &str, size : u64) {
        let entry = self.get_file_entry(path, size);
        entry.write()
             .logical_size = size;
    }
...
\end{rustcode}

\subsubsection*{序号1263}
\noindent 路径：\code{os/components/wateros-vfs/vfs-impl/impl-page-cache/src/lib.rs}\\
\noindent 行范围：706-716\\[0.4em]

\begin{rustcode}
...
    pub fn dirty_page_count(&self, path : &str) -> usize {
        let key = self.file_key(path);
        let files = self.files.lock();
        files.get(&key)
             .map(|entry| {
...
\end{rustcode}

\subsubsection*{序号1264}
\noindent 路径：\code{os/components/wateros-vfs/vfs-impl/impl-page-cache/src/lib.rs}\\
\noindent 行范围：720-781\\[0.4em]

\begin{rustcode}
...
    pub fn flush<Io, E>(&self,
                        io : &mut Io,
                        path : &str,
                        map_err : fn(Io::Error) -> E)
                        -> Result<(), E>
...
\end{rustcode}

\subsubsection*{序号1265}
\noindent 路径：\code{os/components/wateros-vfs/vfs-impl/impl-page-cache/src/lib.rs}\\
\noindent 行范围：786-799\\[0.4em]

\begin{rustcode}
...
    pub fn flush_all<Io, E>(&self, io : &mut Io, map_err : fn(Io::Error) -> E) -> Result<(), E>
        where Io : PageCacheIo
    {
        let paths : Vec<Arc<str>> = {
            let files = self.files.lock();
...
\end{rustcode}

\subsubsection*{序号1266}
\noindent 路径：\code{os/components/wateros-vfs/vfs-impl/impl-page-cache/src/lib.rs}\\
\noindent 行范围：804-840\\[0.4em]

\begin{rustcode}
...
    pub fn purge_closed_file(&self, path : &str) {
        let key = self.file_key(path);
        self.open_refs
            .lock()
            .remove(&key);
...
\end{rustcode}

\subsubsection*{序号1267}
\noindent 路径：\code{os/components/wateros-vfs/vfs-impl/impl-page-cache/src/lib.rs}\\
\noindent 行范围：844-906\\[0.4em]

\begin{rustcode}
...
    pub fn truncate(&self, path : &str, len : u64) {
        let key = self.file_key(path);
        let entry = self.get_file_entry(path, len);
        {
            let mut guard = entry.write();
...
\end{rustcode}

\subsubsection*{序号1268}
\noindent 路径：\code{os/components/wateros-vfs/vfs-impl/impl-page-cache/src/lib.rs}\\
\noindent 行范围：910\\[0.4em]

\begin{rustcode}
...

// 本变量代码由AI完成
static GLOBAL_CACHE : Mutex<Option<Arc<GlobalFilePageCache>>> = Mutex::new(None);

/// 根卷重挂载或辅助卷挂载/卸载后调用：原地清空缓存并绑定新 `mount_gen`，
...
\end{rustcode}

\subsubsection*{序号1269}
\noindent 路径：\code{os/components/wateros-vfs/vfs-impl/impl-page-cache/src/lib.rs}\\
\noindent 行范围：915-922\\[0.4em]

\begin{rustcode}
...
pub fn reset_global_cache(mount_gen : u64) {
    let mut g = GLOBAL_CACHE.lock();
    if let Some(ref existing) = *g {
        existing.reset_to_gen(mount_gen);
    } else {
...
\end{rustcode}

\subsubsection*{序号1270}
\noindent 路径：\code{os/components/wateros-vfs/vfs-impl/impl-page-cache/src/lib.rs}\\
\noindent 行范围：926-945\\[0.4em]

\begin{rustcode}
...
pub fn global_cache(mount_gen : u64) -> Arc<GlobalFilePageCache> {
    let mut g = GLOBAL_CACHE.lock();
    if let Some(ref existing) = *g {
        let current = existing.mount_gen();
        if current != mount_gen {
...
\end{rustcode}

